% AA Classe Del Documento
\documentclass[11pt, a4paper]{article}

% AA Tipografia E Lingua
% BB Lingua E Codifica Font
\usepackage[italian]{babel}
\usepackage[T1]{fontenc}
\usepackage{lmodern}

% BB Geometria Pagina E Microtipografia
\usepackage[margin=2cm, top=1.5cm, bottom=1.5cm]{geometry}
\usepackage{parskip}
\usepackage{microtype}

% AA Gestione Dei Colori
% BB Pacchetto Colori
\usepackage[table]{xcolor}

% BB Palette Colori Personalizzata
\definecolor{darkgreen}{RGB}{0,100,0}
\definecolor{theoryblue}{RGB}{0,84,147}
\definecolor{backcolour}{RGB}{245,245,245}
\definecolor{codegray}{rgb}{0.5,0.5,0.5}

\definecolor{codebg}{HTML}{F8F9FA}
\definecolor{codeframe}{HTML}{E2E8F0}
\definecolor{codekeywords}{HTML}{7C3AED}
\definecolor{codedocker}{HTML}{0284C7}
\definecolor{codecomments}{HTML}{64748B}
\definecolor{codestrings}{HTML}{0D9488}
\definecolor{codenumbers}{HTML}{94A3B8}
\definecolor{codetext}{HTML}{1E293B}

% AA Matematica E Notazione Scientifica
\usepackage{amsmath, amssymb, amsfonts}
\usepackage{mathtools}
\usepackage{siunitx}

% BB OTTIMIZZAZIONE MATRICI LUNGHE PER LA STAMPA
\setlength{\arraycolsep}{10pt}
\usepackage{bm}

% AA Tabelle Ed Allineamenti
\usepackage{booktabs}
\usepackage{array}
\usepackage{multirow}
\usepackage{tabularx}
\usepackage{arydshln}

% AA Grafica E Disegno Vettoriale
% BB Gestione Figure E Posizionamento
\usepackage{graphicx}
\graphicspath{{./Images/CloudComputing/}}
\usepackage{subcaption}
\usepackage{float}

% BB TikZ E Pgfplots
\usepackage{tikz}
\usetikzlibrary{arrows.meta, patterns, angles, quotes, calc}
\usetikzlibrary{positioning}
\usepackage{pgfplots}
\pgfplotsset{compat=1.18}

% AA Listati Di Codice (Listings)
% BB Configurazione Generale
\usepackage{listings}
\usepackage{upquote}

% BB Configurazione Generale
\lstset{
    basicstyle=\ttfamily\footnotesize\color{codetext},
    backgroundcolor=\color{codebg},
    frame=single,
    rulecolor=\color{codeframe},
    framesep=6pt,
    xleftmargin=8pt,
    xrightmargin=8pt,
    breaklines=true,
    breakatwhitespace=true,
    showstringspaces=false,
    keepspaces=true,
    columns=flexible,
    upquote=true,
    tabsize=2,
    captionpos=b,
    % Stili e Colori Predefiniti Globali
    keywordstyle=\bfseries\color{codekeywords},
    commentstyle=\itshape\color{codecomments},
    stringstyle=\color{codestrings}
}

% BB Definizione Dei Linguaggi

% CC YAML
\lstdefinelanguage{yaml}{
    keywords={true,false,null,yes,no,on,off},
    keywordstyle=\bfseries\color{codekeywords},
    comment=[l]{\#},
    commentstyle=\itshape\color{codecomments},
    stringstyle=\color{codestrings},
    morestring=[b]',
    morestring=[b]",
    sensitive=false
}

% CC Dockerfile
\lstdefinelanguage{dockerfile}{
    keywords={FROM, RUN, CMD, LABEL, EXPOSE, ENV, ADD, COPY, ENTRYPOINT,
              VOLUME, USER, WORKDIR, ARG, ONBUILD, STOPSIGNAL, HEALTHCHECK,
              SHELL, AS},
    keywordstyle=\bfseries\color{codedocker},
    comment=[l]{\#},
    commentstyle=\itshape\color{codecomments},
    stringstyle=\color{codestrings},
    morestring=[b]",
    morestring=[b]',
    sensitive=true
}

% CC INI
\lstdefinelanguage{ini}{
    comment=[l]{\#},
    morecomment=[l]{;},
    commentstyle=\itshape\color{codecomments},
    stringstyle=\color{codestrings},
    morestring=[b]",
    morestring=[b]',
    keywordstyle=\bfseries\color{codekeywords},
    sensitive=false
}

% CC JSON
\lstdefinelanguage{json}{
    keywords={true,false,null},
    keywordstyle=\bfseries\color{codekeywords},
    comment=[l]{//},
    morecomment=[s]{/*}{*/},
    commentstyle=\itshape\color{codecomments},
    stringstyle=\color{codestrings},
    morestring=[b]",
    sensitive=true
}

% CC PYTHON (Dialetto nativo pulito)
\lstset{defaultdialect=[3]Python}
\lstdefinelanguage{python}[]{Python}{
    morekeywords={as, async, await, True, False, None},
    sensitive=true
}

% CC BASH / SHELL
\lstdefinelanguage{bash}[]{sh}{
    morekeywords={docker, cd, run, build, stop, rm, rmi, push, pull, network, compose, openstack, rq, worker, keystone-manage, time, sleep, echo, exit, sudo, export, systemctl},
    sensitive=true
}

% CC HTML (Autonomo, previene il loop infinito del parser XML/HTML)
\lstdefinelanguage{html}{
    keywords={img, src, alt, class, id, div, span, table, tr, td, th, href, a, p, h1, h2, h3, form, input, button},
    keywordstyle=\bfseries\color{codekeywords},
    comment=[s]{<!--}{-->},
    commentstyle=\itshape\color{codecomments},
    stringstyle=\color{codestrings},
    morestring=[b]",
    morestring=[b]',
    sensitive=false
}

% CC SQL / CYPHER (Autonomo con supporto per // di Cypher)
\lstdefinelanguage{sql}{
    keywords={MATCH, RETURN, DETACH, DELETE, WHERE, CREATE, SET, WITH, ORDER, BY, ASC, DESC, LIMIT, SELECT, INSERT, INTO, VALUES, FROM, UPDATE, DROP, TABLE, JOIN, ON, AND, OR, NOT},
    keywordstyle=\bfseries\color{codekeywords},
    comment=[l]{//},
    morecomment=[l]{--},
    commentstyle=\itshape\color{codecomments},
    stringstyle=\color{codestrings},
    morestring=[b]',
    morestring=[b]",
    sensitive=false
}

% AA Algoritmi E Pseudocodice
\usepackage[ruled,vlined,linesnumbered]{algorithm2e}
\SetKwData{Dati}{Input}
\SetKwData{Risultato}{Output}
\SetKw{Inizializza}{Inizializza}
\SetKwFor{Per}{per}{fare}{fine per}
\SetKwComment{Commento}{// }{}

% BB Costrutti Di Controllo Aggiuntivi
\SetKwIF{Se}{AltrSe}{Altrimenti}{Se}{Allora}{Altrimenti Se}{Altrimenti}{Fine Se}
\SetKwFor{Mentre}{mentre}{Fare}{Fine Mentre}

% AA Struttura, Titoli Ed Ambienti
% BB Personalizzazione Titoli Sezioni
\usepackage{titlesec}
\titleformat{\section}{\normalfont\LARGE\bfseries}{\thesection}{1em}{}
\titleformat{\subsection}{\normalfont\Large\bfseries}{\thesubsection}{1em}{}
\titleformat{\subsubsection}{\normalfont\large\bfseries}{\thesubsubsection}{1em}{}

% BB Teoremi Ed Ambienti Ausiliari
\usepackage{comment}
\newtheorem{theorem}{Definizione}

% BB Liste Ed Elenchi Puntati
\usepackage{enumitem}
\setlist{noitemsep}
\usepackage{pifont}

% AA Riferimenti, Link Ipertestuali E Segnalibri
\usepackage[unicode]{hyperref}
\hypersetup{
    colorlinks=true,
    linkcolor=blue,
    urlcolor=blue,
    citecolor=blue,
    pdftitle={2) Cloud Computing},
    pdfauthor={Francesco Lo Verde},
    pdfpagemode=FullScreen,
}
\usepackage{bookmark}

% AA Metadati Del Documento
\title{Cloud Computing}

\date{\today}

\author{Francesco Lo Verde}
\begin{document}


\section{Token JSON Web Signature (JWS): Lo Standard Asimmetrico E Scenari Avanzati}

A partire dalla release OpenStack Stein, è stato introdotto il supporto per il provider di token basato su **JSON Web Signature (JWS)**, standardizzato a livello internazionale tramite la specifica **RFC 7519** (\textbf{JSON Web Token} o \textbf{JWT}).

\subsection{Fondamenti Del Formato JWS}

Un token JWT/JWS definisce un metodo compatto, autonomo (\emph{self-contained}) e sicuro per trasmettere attestazioni informative tra entità di rete sotto forma di oggetto JSON.

Le sue peculiarità strutturali includono:
\begin{itemize}[noitemsep]
    \item \textbf{Firma asimmetrica a curve ellittiche}: l'implementazione del provider JWS in Keystone adotta l'algoritmo crittografico standardizzato **ES256** (\emph{ECDSA con curva P-256 e digest SHA-256}). Keystone firma il token impiegando la propria chiave privata, mentre qualsiasi servizio o terza parte può verificarne l'integrità e l'autenticità utilizzando la corrispondente chiave pubblica.
    \item \textbf{Token trasparente (non opaco)}: a differenza di Fernet, in cui il payload è cifrato e inaccessibile a chiunque non possieda la chiave simmetrica, nei token JWS il payload è unicamente firmato digitalmente e codificato in Base64URL. Di conseguenza, le informazioni descrittive dell'utente, del progetto e dei ruoli sono **in chiaro e leggibili da chiunque** sia in possesso della stringa del token.
    \item \textbf{Carattere effimero}: al pari di Fernet, i token JWS non necessitano di persistenza sul database di Keystone né di repliche sincrone di tabelle tra i nodi di calcolo.
\end{itemize}

\subsection{Analisi Comparativa: JWS Verso Fernet}

L'adozione di JWS è stata proposta per risolvere scenari architetturali complessi:
\begin{itemize}[noitemsep]
    \item \textbf{Federazione cross-cloud e Identity Provider eterogenei}: l'uso di standard RFC aperti (JWT) semplifica grandemente l'integrazione di OpenStack con tecnologie come \textbf{OAuth2} e \textbf{OpenID Connect (OIDC)}.
    \item \textbf{Eliminazione della condivisione di segreti simmetrici}: nei token Fernet, tutti i nodi autorizzati a decifrare o emettere token devono condividere le medesime chiavi simmetriche AES/HMAC. Se un nodo dell'infrastruttura venisse compromesso, le chiavi simmetriche verrebbero esposte, mettendo a rischio l'intera cloud. In JWS, la chiave privata di firma rimane rigidamente confinata sul server centrale di Keystone; gli altri servizi ricevono unicamente la chiave pubblica di validazione.
\end{itemize}

Tuttavia, esistono precise controindicazioni ingegneristiche:
\begin{itemize}
    \item \textbf{Assenza di riservatezza}: poiché il payload JWS è decodificabile da chiunque, eventuali informazioni sensibili inserite nei metadati risulterebbero visibili durante il transito o nel browser del client.
    \item \textbf{Invarianza del payload rispetto al contratto API}: la comunità di sviluppo di OpenStack ha chiarito che il payload interno di un token JWS non costituisce un contratto API immutabile; Keystone si riserva il diritto di variare la formattazione interna del token tra versioni diverse senza garantire retrocompatibilità. Pertanto, la raccomandazione operativa per i servizi è di validare comunque il token richiamando le API di Keystone qualora si necessiti di una lettura garantita dei metadati associati.
    \item \textbf{Complessità di gestione crittografica}: per installazioni cloud standard, la gestione e distribuzione delle chiavi simmetriche Fernet risulta operativamente più snella e matura rispetto al ciclo di vita delle coppie di chiavi asimmetriche JWS, motivo per cui **Fernet rimane tuttora il provider predefinito e raccomandato** nella maggioranza dei deploy OpenStack.
\end{itemize}

\newpage

\section{Quadro Sinottico Comparativo Tra I Formati Di Token In Keystone}

A conclusione dell'analisi, la seguente tabella riassume e confronta le proprietà architetturali e prestazionali dei quattro formati di token supportati nel corso dell'evoluzione di OpenStack Keystone.

\begin{table}[H]
\centering
\small
\begin{tabularx}{\textwidth}{p{2.3cm} p{2.8cm} p{3.2cm} p{2.8cm} X}
\toprule
\textbf{Parametro} & \textbf{UUID} & \textbf{PKI / PKIZ} & \textbf{Fernet} & \textbf{JWS} \\
\midrule
\textbf{Natura Stato} & Stateful (persistente su DB) & Stateful / Caching certificati & Ephemeral (non-persistent) & Ephemeral (non-persistent) \\
\midrule
\textbf{Dimensione Media} & 36 caratteri (128 bit / 16 byte netti) & Tipicamente $> \SI{1.7}{\kilo\byte}$, fino a oltre $\SI{8}{\kilo\byte}$ & Rigorosamente $< \SI{250}{\byte}$ & Compatta (leggermente superiore a Fernet) \\
\midrule
\textbf{Crittografia Adottata} & Nessuna (identificatore casuale opaco) & Asimmetrica X.509 e CMS (compressione zlib in PKIZ) & Simmetrica: AES-128-CBC + HMAC-SHA256 & Asimmetrica: ES256 (ECDSA P-256 + SHA-256) \\
\midrule
\textbf{Opacità del Payload} & Completamente opaco & Non opaco (documento CMS ispezionabile) & Opaco (il payload è ciphertext cifrato) & Trasparente (decodificabile liberamente in JSON) \\
\midrule
\textbf{Impatto sul Database} & Elevato (DB bloat, query lente, pruning periodico) & Elevato per revoche e stato iniziale & Nullo (nessun salvataggio su DB) & Nullo (nessun salvataggio su DB) \\
\midrule
\textbf{Meccanismo Validazione} & Query sincrona a Keystone per ogni chiamata & Locale su ogni nodo tramite certificato pubblico & Verifica simmetrica locale o centrale via chiavi & Convalida asimmetrica tramite chiave pubblica \\
\midrule
\textbf{Stato Operativo} & Deprecato & Deprecato / Rimosso & \textbf{Standard Corrente Predefinito} & Supportato (per scenari federati avanzati) \\
\bottomrule
\end{tabularx}
\caption{Confronto sinottico esaustivo tra i meccanismi di emissione e validazione token in OpenStack Keystone.}
\end{table}

% CC ----------------------------------------------------------------------------------

% ==============================================================================
% CAPITOLO: ARCHITETTURA OPENSTACK: IDENTITÀ, AUTORIZZAZIONE, GESTIONE IMMAGINI E CALCOLO DISTRIBUITO
% ==============================================================================

\section{Autenticazione E Autorizzazione In OpenStack: Il Ruolo Di Keystone E Il Flusso JWS}

All'interno di un'architettura di calcolo distribuito multi-tenant quale \textbf{OpenStack}, la sicurezza dell'intero sistema si fonda sulla rigorosa separazione tra la verifica dell'identità (\emph{Authentication}, comunemente abbreviata in \textbf{Authn}) e la verifica dei privilegi operativi (\emph{Authorization}, abbreviata in \textbf{Authz}).
\begin{itemize}[noitemsep]
    \item L'\textbf{autenticazione} risponde alla domanda: \emph{«L'utente o il servizio è effettivamente chi dichiara di essere?»}. Essa consiste nella validazione delle credenziali primarie (coppie nome utente/password, chiavi API o certificati crittografici).
    \item L'\textbf{autorizzazione} risponde alla domanda: \emph{«Questo soggetto, operando in un determinato contesto o progetto, dispone dei permessi necessari per invocare una specifica operazione su una determinata risorsa?»}.
\end{itemize}

Il servizio deputato all'identità in OpenStack è \textbf{Keystone}. Nelle prime iterazioni dell'architettura OpenStack, l'autenticazione avveniva mediante \emph{token} opachi in formato \textbf{UUID} (Universal Unique Identifier). Sebbene estremamente semplici, i token UUID presentavano un grave limite architetturale: erano intrinsecamente con stato (\emph{stateful}). Ogni volta che un servizio periferico (come il modulo di calcolo Nova o il gestore dei volumi Cinder) riceveva una richiesta HTTP contenente un token UUID, era costretto a effettuare una chiamata sincrona verso Keystone per verificarne la validità, estrarre il contesto utente e recuperare i ruoli associati. In installazioni su larga scala, questa dinamica generava un insostenibile collo di bottiglia su Keystone, degradando drasticamente il tempo di risposta globale delle API.

Per ovviare a questo problema di scalabilità orizzontale, OpenStack ha introdotto l'impiego di token basati sullo standard aperto \textbf{JSON Web Signature} (\textbf{JWS}), derivato dalla famiglia JOSE (\emph{Javascript Object Signing and Encryption}) e strutturato analogamente ai JSON Web Token (JWT). I token JWS contengono al proprio interno l'intero carico informativo (\emph{payload}) necessario per la presa di decisione locale, crittograficamente firmato dall'emittente.

% [IMMAGINE DEL FILE DI SLIDE CON NOME "slide_openstack.pdf"] Flusso sequenziale dell'autenticazione iniziale e generazione del token JWS firmato da Keystone — Slide 60 / Pagina 2 del PDF

\subsection{Flusso Di Creazione Del Token JWS}

Il ciclo di vita di un token crittografico inizia con la fase di accreditamento iniziale (\emph{Initial Authentication}):
\begin{enumerate}
    \item \textbf{Inoltro della richiesta di autenticazione}: Il client (ad esempio l'interfaccia da riga di comando \texttt{openstack CLI} o la dashboard web \textbf{Horizon}) invia una richiesta HTTP con metodo \texttt{POST} all'endpoint di Keystone dedicato alla gestione dei token, convenzionalmente attestato sul percorso:
    \[
    \texttt{POST /v3/auth/tokens}
    \]
    Tale richiesta incapsula le credenziali primarie (nome utente, password, identificativo del dominio o progetto) oppure, nel caso di sessioni preesistenti, un token in scadenza da rinnovare.

    \item \textbf{Verifica delle credenziali}: Keystone convalida l'identità del richiedente interrogando il proprio database locale o federando la richiesta verso un sistema di identità aziendale preesistente (ad esempio LDAP o Active Directory).

    \item \textbf{Costruzione del payload JSON}: Una volta validata l'identità, Keystone predispone un documento strutturato in formato JSON che serializza in chiaro l'intero contesto operativo del client, includendo:
    \begin{itemize}[noitemsep]
        \item \emph{User Identity}: identificativo univoco (\texttt{user\_id}) e nome dell'utente.
        \item \emph{Tenants and Associated Roles}: l'elenco dei progetti (\emph{tenant}) a cui l'utente appartiene e la lista dei ruoli specifici (es. \texttt{admin}, \texttt{member}, \texttt{reader}) assegnati all'utente all'interno di ciascun progetto.
        \item \emph{Scope}: l'ambito operativo a cui il token è vincolato, che può essere limitato a uno specifico progetto (\emph{project-scoped token}) o a un intero dominio amministrativo (\emph{domain-scoped token}).
        \item \emph{Expiration Date}: marcatura temporale (\emph{timestamp}) che definisce con precisione il termine di validità del token, prevenendo attacchi di tipo \emph{replay} protratti nel tempo.
    \end{itemize}

    \item \textbf{Firma asimmetrica}: Keystone serializza il payload e ne genera la firma digitale mediante la propria \textbf{chiave privata crittografica} (usando algoritmi di firma a chiave asimmetrica come RSA o ECDSA). La firma garantisce matematicamente l'integrità del payload e l'autenticità della provenienza.

    \item \textbf{Risposta al client}: Il token JWS risultante viene restituito al client all'interno dell'header HTTP della risposta (tipicamente nell'header \texttt{X-Subject-Token}).
\end{enumerate}

\subsection{Utilizzo Del Token E Validazione Distribuita}

La proprietà architetturale fondamentale introdotta dai token JWS è la \textbf{decentralizzazione della verifica}, descritta in dettaglio nella Slide 61.

% [IMMAGINE DEL FILE DI SLIDE CON NOME "slide_openstack.pdf"] Flusso di utilizzo del token JWS verso i servizi periferici e verifica locale tramite chiave pubblica — Slide 61 / Pagina 3 del PDF

Quando il client necessita di interagire con un servizio periferico dell'ecosistema OpenStack (ad esempio il demone \textbf{Nova} per il calcolo o \textbf{Swift} per l'archiviazione a oggetti), inserisce il token JWS nell'header della richiesta HTTP:
\[
\texttt{X-Auth-Token: <JWS\_Token>}
\]

Il servizio ricevente gestisce la richiesta secondo una logica interamente locale:
\begin{enumerate}
    \item \textbf{Ricezione ed estrazione}: Il middleware del servizio ricevente intercetta la richiesta HTTP ed estrae il token dal relativo header.
    \item \textbf{Verifica crittografica senza interpellare Keystone}: Il servizio ricevente dispone localmente della \textbf{chiave pubblica} di Keystone. Utilizzando tale chiave, il servizio verifica matematicamente la firma del JWS. Questa operazione si svolge in modo totalmente autonomo: non viene effettuata alcuna chiamata di rete verso Keystone, azzerando la latenza di validazione e scongiurando l'interruzione del servizio per congestione dell'Identity Provider.
    \item \textbf{Ispezione del payload}: Avendo accertato l'integrità del token, il servizio decodifica il payload JSON ed estrae l'identità dell'utente, i ruoli associati e l'ambito di validità (\emph{scope}).
    \item \textbf{Controlli di conformità e concessione dell'accesso}: Come illustrato nella Slide 62, ogni servizio esegue tassativamente due controlli:
    \begin{itemize}[noitemsep]
        \item \textbf{Controllo di integrità}: La firma digitale deve corrispondere al payload; se un attaccante tenta di manomettere il token (ad esempio modificando il ruolo da \texttt{member} ad \texttt{admin}), la verifica crittografica fallisce istantaneamente.
        \item \textbf{Controllo temporale di scadenza}: Il timestamp corrente viene confrontato con il campo \emph{expiration date}. Se la data limite è stata superata, il token viene rifiutato.
    \end{itemize}
    Se uno qualsiasi dei controlli fallisce, il servizio respinge la richiesta HTTP sollevando un codice di errore \texttt{401 Unauthorized}. In caso contrario, l'accesso alla risorsa richiesta viene autorizzato.
\end{enumerate}

La distribuzione delle chiavi pubbliche di Keystone ai vari servizi dell'infrastruttura avviene periodicamente mediante download programmato oppure in tempo reale attraverso il sistema di notifiche asincrone su bus di messaggi. Quando un token scade, il client non può più rinnovarlo per via locale, ma è tenuto a ripetere integralmente la procedura di autenticazione iniziale.

\subsection{Il Servizio Catalogo (\emph{Catalog Service}) Ed Endpoint Discovery}

Affinché un client o un servizio interno possa consumare le risorse distribuite nel cloud, deve conoscere la locazione di rete (URL, porte, versioni API) di ciascun componente. Il modulo di Keystone deputato a tale scopo è il \textbf{Catalog Service}.

Il Catalogo dei Servizi implementa un vero e proprio registro degli endpoint (\emph{endpoint registry}) utilizzato per l'\emph{endpoint discovery} dinamico. Esso garantisce che i client non debbano memorizzare staticamente la topologia dell'infrastruttura: interrogando Keystone, un client ottiene l'elenco aggiornato dei servizi disponibili. Tale catalogo non è globale e monolitico, bensì \textbf{dinamico e contestualizzato}: può differire da deployment a deployment, da utente a utente e da progetto a progetto in base alle autorizzazioni applicate.

Per ogni servizio registrato (ad esempio \texttt{glance}, \texttt{keystone}, \texttt{heat}, \texttt{neutron}), OpenStack espone tre tipologie di endpoint:
\begin{itemize}[noitemsep]
    \item \textbf{Public}: Accessibile dall'esterno del cloud, attraverso reti instradabili pubblicamente; è l'interfaccia utilizzata da utenti finali, CLI remote e dashboard web.
    \item \textbf{Internal}: Riservato alle comunicazioni intra-cloud tra i servizi stessi (ad esempio Nova che comunica con Glance); è instradato su una rete di gestione interna protetta e isolata.
    \item \textbf{Admin}: Destinato a chiamate con privilegi amministrativi elevati (configurazioni globali, diagnostica avanzata o allocazione di quote).
\end{itemize}

% [IMMAGINE DEL FILE DI SLIDE CON NOME "slide_openstack.pdf"] Screenshot del terminale raffigurante l'output del comando openstack catalog list — Slide 63 / Pagina 5 del PDF

Il comando amministrativo standard per ispezionare il registro degli endpoint è \texttt{openstack catalog list}. Come mostrato nell'output di laboratorio della Slide 63, ciascun servizio è censito per nome, tipo, regione geografica (\texttt{RegionOne}) e tripla di URL:

\begin{lstlisting}[language=bash, caption={Esempio di output del comando \texttt{openstack catalog list}}]
$ openstack catalog list
+-----------+---------------+----------------------------------------------------------------+
| Name      | Type          | Endpoints                                                      |
+-----------+---------------+----------------------------------------------------------------+
| glance    | image         | RegionOne                                                      |
|           |               |   public: http://controller:9292                               |
|           |               | RegionOne                                                      |
|           |               |   internal: http://controller:9292                             |
|           |               | RegionOne                                                      |
|           |               |   admin: http://controller:9292                                |
| keystone  | identity      | RegionOne                                                      |
|           |               |   public: http://controller:5000/v3/                           |
|           |               |   admin: http://controller:5000/v3/                            |
|           |               |   internal: http://controller:5000/v3/                         |
| heat      | orchestration | RegionOne                                                      |
|           |               |   internal: http://controller:8004/v1/0bf1e07c065b4dc9a909a300a5a8497c |
|           |               |   public: http://controller:8004/v1/0bf1e07c065b4dc9a909a300a5a8497c   |
|           |               |   admin: http://controller:8004/v1/0bf1e07c065b4dc9a909a300a5a8497c    |
| aodh      | alarming      | RegionOne (public, internal, admin: port 8042)                 |
| placement | placement     | RegionOne (public, internal, admin: port 8778)                 |
| neutron   | network       | RegionOne (public: port 9696)                                  |
+-----------+---------------+----------------------------------------------------------------+
\end{lstlisting}

\section{Il Motore Di Autorizzazione: Il Servizio Policy E Il Controllo Basato Su Ruoli}

Se Keystone definisce chi sia l'utente e quali ruoli possieda, spetta a ciascun servizio stabilire se quel determinato ruolo consenta di invocare una specifica operazione. Questa separazione delle responsabilità è gestita tramite il \textbf{Policy Service}.

Il Policy Service fornisce un motore di autorizzazione basato su regole (\emph{rule-based authorization engine}) e la relativa interfaccia di configurazione. Ogni singolo servizio di OpenStack (Nova per il calcolo, Cinder per i volumi, Glance per le immagini, Neutron per la rete) dichiara esplicitamente le regole di accesso alle proprie risorse all'interno di un file di policy dedicato, originariamente denominato \texttt{policy.json} (progressivamente convertito in \texttt{policy.yaml} nelle release recenti).

Una risorsa protetta può essere un endpoint API, la capacità di connettere un volume di archiviazione a un'istanza, o il permesso di istanziare nuove macchine virtuali. Ogni qualvolta giunge una chiamata API al servizio, il motore di policy intercetta il contesto della richiesta ed esegue la valutazione della regola associata.

\subsection{Sintassi E Semantica Delle Regole Di Policy}

I file di policy associano a ogni azione (\emph{target}) una regola booleana (\emph{rule}) secondo la seguente notazione:
\[
\texttt{"<target>" : "<rule>"}
\]
\begin{itemize}[noitemsep]
    \item Il \textbf{target} rappresenta l'operazione API o l'azione controllata (ad esempio \texttt{"compute:create"} per l'avvio di un'istanza o \texttt{"volume:attach"} per il collegamento di un volume).
    \item La \textbf{rule} è un'espressione logica che stabilisce le condizioni affinché l'azione sia lecita. Può combinare ruoli (ad esempio \texttt{"role:admin"}), controlli di proprietà sul progetto (\texttt{"project\_id:\%(project\_id)s"}), regole generiche o espressioni booleane complesse tramite operatori \texttt{and}, \texttt{or}, \texttt{not}.
\end{itemize}

Una caratteristica fondamentale del meccanismo di policy in OpenStack è la sua \textbf{reattività a caldo}: le modifiche apportate ai file di policy hanno efficacia immediata. Il demone del servizio rileva l'aggiornamento del file sul filesystem (o lo ricarica all'invio di un segnale di reload) senza richiedere il riavvio del processo. Ciò consente agli amministratori di rimodellare le politiche di accesso su un'infrastruttura di produzione attiva senza causare alcuna interruzione di servizio.

\subsection{Mappatura Dei File Di Policy Nel Sistema}

Come evidenziato dalle slide del corso, i file di configurazione delle policy risiedono nei percorsi di installazione e configurazione dei rispettivi moduli Python. L'interrogazione tramite il comando \texttt{locate policy.json} su un nodo controller rivela come ciascun componente possieda il proprio file autonomo:

\begin{lstlisting}[language=bash, caption={Individuazione dei file di policy nel filesystem del controller}]
$ locate policy.json
/usr/lib/python2.7/dist-packages/cinder/tests/unit/policy.json
/usr/lib/python2.7/dist-packages/glance/tests/etc/policy.json
/usr/lib/python2.7/dist-packages/gnocchi/rest/policy.json
/usr/lib/python2.7/dist-packages/neutron/tests/etc/policy.json
/usr/lib/python2.7/dist-packages/openstack_auth/tests/conf/keystone_policy.json
/usr/lib/python2.7/dist-packages/openstack_auth/tests/conf/nova_policy.json
/etc/openstack-dashboard/cinder_policy.json
/etc/openstack-dashboard/glance_policy.json
/etc/openstack-dashboard/keystone_policy.json
/etc/openstack-dashboard/neutron_policy.json
/etc/openstack-dashboard/nova_policy.json
\end{lstlisting}

La presenza di file di policy replicati anche per \texttt{openstack\_dashboard} (l'interfaccia grafica Horizon) garantisce che la GUI possa disabilitare preventivamente pulsanti o form di configurazione se il token dell'utente non possiede i privilegi necessari per eseguire le rispettive chiamate API sottostanti.

\newpage

\section{OpenStack Glance: Il Servizio Di Gestione Delle Immagini Di Macchine Virtuali}

\textbf{OpenStack Glance} è il componente IaaS responsabile della scoperta, registrazione, catalogazione, recupero e gestione del ciclo di vita delle immagini di macchine virtuali e dei relativi \emph{snapshot}. Implementa la specifica formale \textbf{Images API v2}.

Per comprendere la funzione di Glance nel cloud, è opportuno chiarire la nozione ingegneristica di \textbf{immagine}: un'immagine di macchina virtuale è un singolo file contenente una copia conforme di un disco fisso virtuale, comprensivo di bootloader, sistema operativo installato, configurazioni di base del kernel e pacchetti software preinstallati. Glance consente agli utenti di istanziare decine o centinaia di VM identiche partendo da un unico \emph{template} immutabile.

% [IMMAGINE DEL FILE DI SLIDE CON NOME "slide_openstack.pdf"] Diagramma dell'architettura interna di OpenStack Glance (Client, API, Domain Controller, Registry Layer, DAL e Glance Store) — Slide 66 / Pagina 8 del PDF

\subsection{Architettura Interna E Componenti Di Glance}

L'architettura di Glance adotta un modello client-server modulare altamente disaccoppiato, strutturato nei seguenti componenti:
\begin{itemize}
    \item \textbf{Glance Client}: Qualsiasi entità software (la CLI di OpenStack, il demone nova-compute o un tool di orchestrazione come Terraform) che interagisca via rete con il server Glance.
    \item \textbf{REST API}: Il punto di ingresso unificato che espone le funzionalità del servizio mediante verbi HTTP e formati JSON standard.
    \item \textbf{Glance Domain Controller}: Il cuore dell'elaborazione di Glance. È un middleware a strati concentrici che orchestra l'intera logica di business prima che i dati o i metadati tocchino i livelli di memorizzazione sottostanti.
    \item \textbf{Glance Store}: Libreria specializzata (\texttt{glance\_store}) responsabile dell'interazione fisica con i backend di archiviazione esterni e con il filesystem locale. Essa astrae le peculiarità dei diversi sistemi di storage sottostanti, esponendo un'interfaccia di lettura/scrittura omogenea.
    \item \textbf{Database Abstraction Layer (DAL)}: Un'interfaccia di programmazione interna che standardizza e unifica la comunicazione tra la logica di dominio e il database relazionale, astraendo le query SQL.
    \item \textbf{Registry Layer}: Uno strato software opzionale progettato per garantire comunicazioni crittografate e rigidamente isolate tra il controllore di dominio e la DAL, impiegando un servizio di rete separato per ambienti ad alto isolamento di sicurezza.
    \item \textbf{Glance Database (Glance DB)}: Un database relazionale centralizzato (tipicamente istanziato su MySQL, MariaDB o SQLite per test) condiviso tra i processi di Glance, deputato a registrare esclusivamente i \textbf{metadati} dell'immagine (nome, ID univoco, dimensione, visibilità, hash crittografico, formati, proprietario). I dati pesanti del disco virtuale (\emph{image payload}) non transitano né risiedono mai all'interno del DB relazionale.
\end{itemize}

\subsection{I Sette Strati Del Glance Domain Controller}

La logica applicativa del Glance Domain Controller è rigorosamente suddivisa in una pipeline sequenziale a sette strati funzionali (\emph{domain model layers}), analizzati in dettaglio nelle slide 66:

\begin{enumerate}
    \item \textbf{Authorization Layer}: È il primo strato operativo. Verifica se l'immagine o le sue proprietà possano essere modificate dall'utente che ha inviato la richiesta. Estrapola il contesto dell'utente dalla richiesta HTTP autenticata e lo confronta con l'identificativo del proprietario dell'immagine (\emph{image owner}). Le modifiche sono consentite unicamente all'amministratore di sistema o al proprietario legittimo dell'immagine; in caso contrario, il flusso si arresta con errore immediato.

    \item \textbf{Property Protection Layer}: Strato opzionale, attivato qualora nel file di configurazione principale di Glance sia valorizzato il parametro \texttt{property\_protection\_file}. In Glance, le proprietà si distinguono in:
    \begin{itemize}[noitemsep]
        \item \emph{Core Properties}: Attributi formali e standard definiti dallo schema ufficiale dell'immagine (ad esempio \texttt{disk\_format}, \texttt{container\_format}, \texttt{size}).
        \item \emph{Meta Properties}: Coppie chiave/valore arbitrarie definite dall'utente per descrivere peculiarità del sistema operativo (es. architettura, driver specifici, flag di ottimizzazione).
    \end{itemize}
    Questo strato governa e limita l'accesso in lettura e scrittura a tali metadati avanzati tramite le chiamate API pubbliche.

    \item \textbf{Notifier Layer}: Rappresenta il terzo anello della catena. Inoltra al bus di messaggi di OpenStack (la Message Queue AMQP) gli eventi salienti del ciclo di vita dell'immagine. Pubblica notifiche relative alle modifiche di stato, nonché avvisi diagnostici (\emph{warnings}) ed eccezioni operative verificatesi durante l'interazione con l'immagine.

    \item \textbf{Policy Layer}: Il quarto strato valida l'azione richiesta a fronte delle regole formali di autorizzazione RBAC codificate nel file \texttt{policy.yaml}. Verifica la liceità dell'azione richiesta e monitora l'applicazione coerente delle regole.

    \item \textbf{Quota Layer}: Il quinto strato applica i vincoli di allocazione delle risorse stabiliti dall'amministratore per ciascun tenant. Durante la fase di caricamento (\emph{upload}) e salvataggio di un'immagine, calcola lo spazio totale occupato dall'utente e verifica se la nuova immagine ecceda il limite di quota disco consentito. Se la quota viene superata, la richiesta fallisce e l'upload viene abortito.

    \item \textbf{Location Layer}: Il sesto strato gestisce l'interfacciamento con la libreria \texttt{glance\_store}. È deputato a:
    \begin{itemize}[noitemsep]
        \item Validare la correttezza formale e la raggiungibilità dell'URI di memorizzazione (\emph{Location URI}) quando viene aggiunta una nuova sorgente.
        \item Provvedere all'effettiva cancellazione dei blocchi dati dal backend fisico quando una locazione viene rimossa o modificata.
        \item Verificare e scongiurare duplicazioni incoerenti della posizione dell'immagine nel backend.
    \end{itemize}

    \item \textbf{Database Layer}: Il settimo e ultimo strato implementa le interazioni dirette con l'API del database. Effettua il marshalling e unmarshalling dei dati: converte i dati grezzi ricevuti dalle tabelle SQL in istanze strutturate di oggetti software \texttt{Image} e viceversa, garantendo la persistenza delle informazioni di stato.
\end{enumerate}

\subsection{Backend Di Memorizzazione Supportati}

Glance offre un disaccoppiamento totale tra la gestione logica dei metadati e il repository fisico dei bit. I driver di \texttt{glance\_store} consentono di memorizzare il payload su un ventaglio eterogeneo di tecnologie:
\begin{itemize}[noitemsep]
    \item \textbf{Filesystem locale}: Soluzione basilare in cui le immagini sono salvate come file ordinari in una cartella montata localmente sul nodo (convenzionalmente posizionata in \texttt{/var/lib/glance/images}).
    \item \textbf{OpenStack Swift}: Il servizio di archiviazione a oggetti nativo di OpenStack, altamente distribuito, scalabile ed esente da singoli punti di vulnerabilità (\emph{no single point of failure}).
    \item \textbf{Ceph RBD}: Storage unificato a blocchi distribuito, lo standard de facto nelle infrastrutture industriali per garantire replicazione sincrona e clonazione istantanea tramite \emph{Copy-on-Write} (CoW).
    \item \textbf{Sheepdog}: Sistema di archiviazione distribuito a oggetti progettato specificamente per le macchine virtuali basate su QEMU/KVM.
\end{itemize}

Ogni immagine salvata viene univocamente referenziata e raggiunta tramite un URI avente la seguente struttura:
\[
\texttt{<Glance Server Location>/v2/images/<Image\_UUID>}
\]

\subsection{Tassonomia Dei Formati: Formati Di Disco E Contenitore}

Al momento della registrazione di un'immagine, l'amministratore o l'utente devono specificare due attributi fondamentali che ne definiscono la compatibilità con l'hypervisor sottostante: il formato del disco (\textbf{Disk Format}) e il formato del contenitore (\textbf{Container Format}).

\subsubsection{Formati Di Disco (\emph{Disk Formats})}
Il formato di disco definisce la struttura interna dei blocchi logici del disco fisso virtualizzato:
\begin{itemize}
    \item \textbf{\texttt{raw}}: È il formato binario non strutturato più semplice in assoluto. Corrisponde alla copia conforme bit-a-bit del dispositivo a blocchi, concettualmente equivalente all'output generato dal comando Unix \texttt{dd if=/dev/sda of=image.raw}. Offre le migliori prestazioni di I/O poiché l'hypervisor non deve applicare alcun livello di traduzione, ma occupa l'intero spazio nominale del disco senza compressione nativa. Supportato nativamente sia da KVM che da Xen.
    \item \textbf{\texttt{qcow2}}: (\emph{QEMU Copy-On-Write v2}). È il formato preferenziale nell'ecosistema Linux KVM/QEMU. Supporta la crescita dinamica dello spazio su disco (allocazione \emph{thin provisioning}), la compressione dei blocchi, gli snapshot multipli e la tecnica \emph{Copy-on-Write}, permettendo a decine di VM di condividere una singola immagine di base read-only scrivendo unicamente i blocchi modificati.
    \item \textbf{\texttt{vhd}}: Formato VHD (\emph{Virtual Hard Disk}) impiegato nativamente da hypervisor quali Microsoft Hyper-V, Citrix XenServer e VirtualBox.
    \item \textbf{\texttt{vhdx}}: Evoluzione avanzata del formato VHD, supporta dischi di capacità superiore (fino a \SI{64}{\tera\byte}) e garantisce resilienza contro il danneggiamento dei metadati in caso di blackout o caduta improvvisa di tensione.
    \item \textbf{\texttt{vmdk}}: (\emph{Virtual Machine Disk}). Formato standard proprietario e industriale introdotto da VMware per i propri hypervisor ESXi e Workstation, ampiamente supportato da tool di conversione.
    \item \textbf{\texttt{vdi}}: Formato proprietario supportato dall'emulatore open source Oracle VirtualBox e da QEMU.
    \item \textbf{\texttt{iso}}: File archivio conforme allo standard per supporti ottici (CD-ROM/DVD-ROM), tipicamente impiegato per le immagini di boot e installazione iniziale del sistema operativo.
    \item \textbf{\texttt{ploop}}: Formato di dischi virtuali a blocchi con allocazione dinamica supportato e utilizzato dal sistema di containerizzazione a livello di sistema operativo Virtuozzo.
    \item \textbf{\texttt{aki}, \texttt{ari}, \texttt{ami}}: Terzina di formati proprietari dell'ecosistema Amazon Web Services: rispettivamente \emph{Amazon Kernel Image}, \emph{Amazon Ramdisk Image} e \emph{Amazon Machine Image}.
\end{itemize}

\subsubsection{Formati Di Contenitore (\emph{Container Formats})}
Il formato contenitore indica se l'immagine è impacchettata in un involucro che racchiude, oltre al payload del disco virtuale, metadati aggiuntivi specifici per l'esecuzione della VM.
\begin{itemize}[noitemsep]
    \item \textbf{\texttt{bare}}: Indica l'assenza totale di qualsiasi contenitore o busta di metadati. È il formato più comune quando si caricano dischi crudi (\texttt{raw}) o file \texttt{qcow2}. Nel dubbio, è la scelta predefinita raccomandata.
    \item \textbf{\texttt{ovf}}: (\emph{Open Virtualization Format}). Standard aperto per la distribuzione di \emph{virtual appliance}, contenente file di configurazione XML e riferimenti alle immagini disco.
    \item \textbf{\texttt{ova}}: Archivio tar non compresso contenente un intero package OVF.
    \item \textbf{\texttt{docker}}: Archivio tar rappresentativo del filesystem a strati per container gestiti dal demone Docker.
    \item \textbf{\texttt{aki}, \texttt{ari}, \texttt{ami}}: Specificano che l'involucro dell'immagine adotta il contenitore previsto da Amazon EC2.
\end{itemize}

\subsection{Ciclo Di Vita E Macchina A Stati Dell'Immagine}

% [IMMAGINE DEL FILE DI SLIDE CON NOME "slide_openstack.pdf"] Macchina a stati finiti con le transizioni tra gli stati di un'immagine Glance — Slide 71 / Pagina 16 del PDF

Durante la sua esistenza nel registro, un'immagine Glance evolve attraverso una serie formale di stati operativi (\emph{Image Statuses}), riassunti in Tabella~\ref{tab:glance_statuses}:

\begin{table}[H]
\centering
\small
\begin{tabularx}{\textwidth}{l p{3.2cm} X}
\toprule
\textbf{Stato} & \textbf{Nome Tecnico} & \textbf{Descrizione Operativa} \\
\midrule
\texttt{queued} & Accodata & L'identificatore dell'immagine è stato riservato nel catalogo di Glance (tramite chiamata \texttt{POST /v2/images}), ma nessun dato di payload è stato ancora caricato e la dimensione non è stata fissata a zero. \\
\addlinespace
\texttt{saving} & Salvataggio in corso & I blocchi dati grezzi dell'immagine sono in fase di caricamento verso il backend di storage designato. (Nota: se l'immagine è creata passando esplicitamente l'header \texttt{x-image-meta-location}, non transita mai per lo stato \texttt{saving}, poiché il payload risiede già su una locazione esterna nota). \\
\addlinespace
\texttt{uploading} & Caricamento & I dati sono stati allocati (\emph{staged}) all'interno del processo standardizzato di importazione interoperabile (\emph{interoperable image import process}). È stato introdotto a partire dalle API Glance v2.6. La risorsa non è ancora pronta. \\
\addlinespace
\texttt{importing} & Importazione & Il payload è in fase di elaborazione attiva e validazione all'interno del flusso di importazione interoperabile, ma non è ancora pronto per l'uso. \\
\addlinespace
\texttt{active} & Attiva & L'immagine è completata, integra e disponibile per il consumo operativo. Le istanze Nova possono essere avviate a partire da questo stato. Si raggiunge anche quando la dimensione è fissata esplicitamente a 0 in fase di creazione. \\
\addlinespace
\texttt{deactivated} & Disattivata & L'accesso ai blocchi dati dell'immagine è interdetto a qualunque utente non amministratore. L'inibizione del download impedisce operazioni derivate come il cloning o l'export verso formati esterni. \\
\addlinespace
\texttt{killed} & Fallita / Abortita & Si è verificato un errore critico durante la fase di upload del payload; i dati risultano parziali o illeggibili, rendendo l'immagine inservibile. \\
\addlinespace
\texttt{deleted} & Cancellata & L'immagine non è più disponibile all'uso; Glance conserva temporaneamente i metadati all'interno del DB per audit log. Verrà rimossa fisicamente in un secondo momento. \\
\addlinespace
\texttt{pending\_delete} & Cancellazione pendente & Stato analogo a \texttt{deleted}, ma il backend di archiviazione fisico non ha ancora eliminato i blocchi disco associati. La risorsa non è in alcun modo recuperabile. \\
\bottomrule
\end{tabularx}
\caption{Stati del ciclo di vita di un'immagine in OpenStack Glance.}
\label{tab:glance_statuses}
\end{table}

\subsection{Modello Di Visibilità Multi-Tenant}

In un'architettura cloud condivisa tra progetti distinti (\emph{tenants}), la regolamentazione dell'accesso alle immagini è disciplinata dall'attributo di \textbf{visibilità} (\emph{Visibility}). Un principio cardine stabilito dalle slide è che:
\begin{center}
\emph{La visibilità è un attributo ortogonale rispetto alla proprietà (\textbf{ownership}).}
\end{center}
La proprietà definisce chi sia il tenant proprietario dell'immagine (chi può modificarla o distruggerla), mentre la visibilità disciplina chi possa vederla in elenco ed eseguire il boot di nuove macchine virtuali da essa. I valori ammessi per la visibilità sono:

\begin{itemize}
    \item \textbf{\texttt{public}}: Qualsiasi utente autenticato nel cloud può leggere l'immagine, ispezionarne i metadati e usarla per avviare istanze. L'immagine compare automaticamente nella lista predefinita (\emph{default image-list}) di tutti i tenant.
    \item \textbf{\texttt{private}}: Esclusivamente gli utenti appartenenti al medesimo tenant proprietario ($\text{tenantId} == \text{tenantId}(\text{owner})$) possono visualizzare i dettagli dell'immagine, elencarla di default e avviare VM da essa.
    \item \textbf{\texttt{community}}: Qualsiasi utente del cloud può accedere ai metadati ed effettuare il boot da tale immagine, ma essa non compare nell'elenco predefinito di utenti diversi dal proprietario a meno che non venga esplicitamente interrogata tramite il suo identificatore univoco. Se un utente esterno viene inserito nella \emph{member-list} dell'immagine e imposta il proprio stato a \texttt{'accepted'}, l'immagine comparirà nella sua lista di default.
    \item \textbf{\texttt{shared}}: L'accesso è esteso selettivamente tramite una lista di appartenenza (\emph{image members list}). Oltre al proprietario, un utente esterno può usufruire dell'immagine solo se il suo \texttt{tenantId} è presente nella lista dei membri. Se il membro imposta il flag \texttt{member\_status == 'accepted'}, l'immagine comparirà nella sua lista di default. Se non viene specificata alcuna visibilità all'atto della creazione, l'immagine assume questa modalità per default (ma con \emph{member-list} inizialmente vuota).
\end{itemize}

\newpage

\section{OpenStack Nova: Il Cuore Del Calcolo Distribuito}

\textbf{OpenStack Nova} è il modulo fondamentale responsabile della gestione delle risorse computazionali del cloud IaaS. Il suo compito principale risiede nel provisionare, configurare, monitorare e terminare istanze di macchine virtuali su una flotta distribuita di server fisici (\emph{compute nodes}). Nova espone l'\textbf{OpenStack Compute API} ed è progettato per interfacciarsi in modo trasparente con molteplici hypervisor eterogenei: principalmente \textbf{KVM} (Kernel-based Virtual Machine) orchestrato via \textbf{libvirt}, ma anche \textbf{Xen}, \textbf{VMware ESXi} o macchine bare-metal (tramite Ironic).

% [IMMAGINE DEL FILE DI SLIDE CON NOME "slide_openstack.pdf"] Schema architetturale complessivo di OpenStack Nova con interconnessione tra demoni, DB, Message Queue e nodi di calcolo — Slide 74 / Pagina 21 del PDF

\subsection{Architettura A Microservizi E Componenti Principali}

Nova adotta un'architettura completamente modulare basata su microservizi cooperanti, mostrata nella Slide 75:

% [IMMAGINE DEL FILE DI SLIDE CON NOME "slide_openstack.pdf"] Flussi di comunicazione tra i componenti di Nova (nova-api, scheduler, conductor, compute, placement, DB e oslo.messaging) — Slide 75 / Pagina 22 del PDF

\begin{enumerate}
    \item \textbf{Nova API (\texttt{nova-api})}: È il punto di accesso unificato per i client esterni, la CLI e la dashboard. Riceve le richieste HTTP, valida i token JWS, verifica i file di policy ed esegue il parsing dei comandi inviandoli al resto del sistema tramite code di messaggi.
    \item \textbf{Nova API Metadata (\texttt{nova-api-metadata})}: Servizio specializzato che risponde alle richieste di metadati provenienti dall'interno delle istanze VM in esecuzione (ad esempio la chiave pubblica SSH da iniettare tramite \texttt{cloud-init} al primo boot).
    \item \textbf{Nova Scheduler (\texttt{nova-scheduler})}: Il cervello decisionale del calcolo. Estrae una richiesta di creazione istanza dalla coda e determina su quale specifico host fisico debba essere allocata la VM, massimizzando l'efficienza delle risorse.
    \item \textbf{Nova Compute (\texttt{nova-compute})}: Demone operaio (\emph{worker daemon}) in esecuzione direttamente su ogni nodo di calcolo fisico. Riceve i comandi di creazione, spegnimento, riavvio o migrazione e interagisce con l'hypervisor locale attraverso le rispettive librerie applicative (principalmente \texttt{libvirt}).
    \item \textbf{Nova Conductor (\texttt{nova-conductor})}: Componente di mediazione di cruciale rilevanza architetturale, concepito come proxy applicativo tra il demone \texttt{nova-compute} e il database centrale.
    \item \textbf{Nova SQL Database}: Repository centrale basato su un database relazionale (MySQL, MariaDB o PostgreSQL gestiti mediante l'ORM Python \textbf{SQLAlchemy}) che memorizza lo stato transazionale dell'infrastruttura: istanze attive, tipi di istanza (\emph{flavors}), associazioni di rete e mapping sui progetti.
    \item \textbf{Message Queue}: Bus di comunicazione orientato ai messaggi (Message-Oriented Middleware, MOM) basato sul protocollo \textbf{AMQP} (\emph{Advanced Message Queuing Protocol}). L'implementazione predefinita è \textbf{RabbitMQ}, sebbene sia supportato anche Apache Qpid.
\end{enumerate}

\subsection{Disaccoppiamento Tramite Message Queue: AMQP E RPC}

Le comunicazioni tra i componenti interni di Nova non avvengono tramite chiamate HTTP sincrone punto-a-punto, bensì mediante l'interposizione della \textbf{Message Queue}, implementando il paradigma \textbf{Publish/Subscribe}:
\begin{itemize}[noitemsep]
    \item I componenti operano in modalità debolmente accoppiata (\emph{loosely coupled}): i nodi produttori generano messaggi instradandoli su specifici canali tematici (\emph{topics}), e uno o più consumatori in ascolto prelevano ed elaborano la richiesta.
    \item L'astrazione software utilizzata da OpenStack per gestire questo scambio è la libreria unificata \textbf{Oslo Messaging}, che modella chiamate di procedura remota (\textbf{RPC}) sul bus AMQP.
\end{itemize}

Si distinguono due primitive fondamentali di interscambio:
\begin{itemize}
    \item \textbf{\texttt{rpc.call()}}: Chiamata sincrona con attesa di risposta. L'invocatore deposita un messaggio sulla coda del destinatario e si pone in attesa bloccante (con un meccanismo di timeout) su una coda privata di risposta (\emph{reply queue}) fino alla restituzione del valore calcolato.
    \item \textbf{\texttt{rpc.cast()}}: Chiamata asincrona \emph{fire-and-forget}. Il componente emette il messaggio verso la coda del broker e prosegue immediatamente la propria esecuzione, senza attendere riscontro operativo.
\end{itemize}

\subsection{Il Ruolo Architetturale Del Nova Conductor}

Nella progettazione originaria di OpenStack, ciascun demone \texttt{nova-compute} installato sui nodi fisici possedeva le credenziali di rete per connettersi direttamente al database SQL centrale. Questo approccio presentava due criticità insormontabili:
\begin{enumerate}
    \item \textbf{Vulnerabilità di sicurezza catastrofica}: Se un attaccante riusciva a evadere da una macchina virtuale (\emph{hypervisor breakout}) e a compromettere l'ambiente di root del nodo di calcolo, otteneva l'accesso diretto e non filtrato al database principale di OpenStack, potendo corrompere l'intera infrastruttura multi-tenant del cloud.
    \item \textbf{Complessità di aggiornamento e accoppiamento}: Il codice di \texttt{nova-compute} conteneva query SQL dirette. L'aggiornamento a una versione successiva con modifiche allo schema relazionale imponeva l'arresto e la migrazione simultanea di tutti i nodi di calcolo dell'infrastruttura.
\end{enumerate}

Per sanare queste problematiche è stato introdotto il modulo \textbf{nova-conductor}, illustrato nelle Slide 78 e 79.

% [IMMAGINE DEL FILE DI SLIDE CON NOME "slide_openstack.pdf"] Schema di interazione tra controller node e compute node tramite nova-conductor ed exchange RPC rpc.call() — Slide 79 / Pagina 27 del PDF

Il Nova Conductor opera come un server RPC stateless con i seguenti compiti:
\begin{itemize}
    \item \textbf{Proxy verso il database}: I demoni \texttt{nova-compute} non detengono alcuna credenziale SQL; ogni necessità di lettura o scrittura dello stato dell'istanza viene convertita in un messaggio RPC instradato via RabbitMQ verso \texttt{nova-conductor}, che valida l'operazione e interagisce fisicamente con il database per conto del compute node.
    \item \textbf{Isolamento e sicurezza}: Se un nodo di calcolo viene violato, l'attaccante può unicamente inviare messaggi RPC limitati, senza avere la possibilità di manipolare direttamente tabelle arbitrarie sul DB centrale. \textbf{Regola tassativa}: \texttt{nova-conductor} non deve mai essere installato sui nodi di calcolo, bensì segregato nei nodi di controllo (\emph{controller nodes}).
    \item \textbf{Scalabilità orizzontale}: Essendo un demone stateless, è possibile istanziare multipli worker di \texttt{nova-conductor} distribuiti su macchine diverse, scalando le prestazioni di gestione delle transazioni.
    \item \textbf{Orchestrazione di task a lunga durata}: Coordina operazioni complesse che coinvolgono molteplici nodi di calcolo, quali la migrazione a caldo (\emph{live migration}), il ridimensionamento di una VM (\emph{resize}) e la ricostruzione dell'istanza (\emph{rebuild}).
    \item \textbf{Sincronizzazione degli oggetti Nova (\emph{Nova Objects})}: Nova adotta strutture dati orientate agli oggetti (es. \texttt{Instance}, \texttt{Flavor}). Il Conductor garantisce la serializzazione, deserializzazione e compatibilità di versione di tali oggetti durante la transizione tra release software eterogenee.
\end{itemize}

\newpage

\section{Il Servizio Placement E La Gestione Delle Risorse Eterogenee}

Con l'evoluzione del cloud, le risorse hardware non si limitano più a semplici quote uniformi di CPU e RAM, ma includono pool di storage condiviso (Ceph, NFS), pool di indirizzi IP, schede di accelerazione hardware (GPU, FPGA) e schede di rete con supporto SR-IOV. Per tracciare in modo atomico, granulare ed efficiente l'inventario e il consumo di queste risorse eterogenee, OpenStack ha introdotto il servizio \textbf{Placement}.

\subsection{Evoluzione Storica E Motivazione Di Sistema}

La genesi e l'estrazione architetturale di Placement seguono due tappe fondamentali:
\begin{itemize}[noitemsep]
    \item \textbf{Release 14.0.0 Newton}: Il servizio Placement API nasce come codice integrato all'interno del repository software di Nova.
    \item \textbf{Release 19.0.0 Stein}: Placement viene estratto da Nova e promosso a progetto indipendente con un proprio repository e una propria identità architetturale.
\end{itemize}

La separazione è scaturita dalla necessità di consentire anche ad altri sottosistemi di OpenStack (non solo Nova, ma ad esempio Neutron per le allocazioni di banda di rete, o Cyborg per le periferiche FPGA) di registrare e consumare risorse senza dipendere monoliticamente dal demone del calcolo. Oggi, Placement espone una propria REST API autonoma e deve essere censito nel catalogo di Keystone con il tipo di servizio formale:
\[
\texttt{service\_type = placement}
\]

\subsection{Il Modello Dati Di Placement}

L'architettura di Placement poggia su cinque entità logiche fondamentali, descritte nelle slide 81:

\begin{itemize}
    \item \textbf{Resource Provider (RP)}: Qualsiasi entità fisica o logica capace di erogare risorse nell'infrastruttura. Esempi tipici sono: un nodo di calcolo fisico (che fornisce vCPU e memoria RAM), un pool di storage condiviso esterno (che fornisce spazio su disco in \SI{}{\giga\byte}) o un pool di allocazione di indirizzi IP gestito dal servizio di networking.
    \item \textbf{Resource Class}: La tipologia astratta della risorsa gestita. Placement possiede un elenco normalizzato di classi standard (es. \texttt{VCPU}, \texttt{MEMORY\_MB}, \texttt{DISK\_GB}) e consente la definizione di classi personalizzate (\emph{custom resource classes}) come \texttt{CUSTOM\_GPU\_A100}.
    \item \textbf{Inventory}: L'insieme quantitativo delle \emph{resource classes} messe a disposizione da un determinato \emph{resource provider}. Un compute node, ad esempio, dichiarerà al servizio Placement un inventario pari a 64 unità di \texttt{VCPU}, 262144 unità di \texttt{MEMORY\_MB} e 1000 unità di \texttt{DISK\_GB}, specificando anche eventuali tassi di sovrascrizione (\emph{allocation ratio}).
    \item \textbf{Usage Data}: Le metriche mantenute da Placement che tengono traccia della quota di inventario correntemente consumata per ciascuna classe di risorsa su ogni provider.
    \item \textbf{Consumer}: L'entità fruitrice della risorsa, identificata tipicamente dall'UUID della specifica istanza di macchina virtuale creata nel cloud.
    \item \textbf{Allocations}: Il record logico che formalizza il legame tra un \emph{Consumer} e uno o più \emph{Resource Providers}. L'allocazione rappresenta la riserva vincolante e atomica di risorse; ad esempio, un'istanza VM (consumer) vanterà un'allocazione di 4 \texttt{VCPU} e 8192 \texttt{MEMORY\_MB} presso il compute node, e contemporaneamente 50 \texttt{DISK\_GB} presso un pool di storage esterno.
\end{itemize}

\subsection{Ciclo Di Interazione Congiunto Tra Nova E Placement}

La cooperazione tra i nodi di calcolo, lo schedulatore di Nova e il servizio Placement durante il ciclo di vita dell'infrastruttura segue un flusso standardizzato:
\begin{enumerate}
    \item \textbf{Registrazione dell'inventario}: All'avvio del demone \texttt{nova-compute} su un nodo fisico, il componente interno denominato \textbf{Resource Tracker} interroga Keystone per individuare l'endpoint di Placement. Successivamente, invia una richiesta HTTP a Placement per registrare il nodo come \emph{Resource Provider} e dichiararne l'inventario hardware complessivo disponibile.
    \item \textbf{Fase di Scheduling}: Quando giunge una richiesta di creazione di una nuova istanza, il \textbf{nova-scheduler} deve individuare il nodo fisico idoneo a ospitarla. Invece di dover scorrere esaustivamente l'intero database locale delle istanze, lo scheduler invia una singola richiesta filtrata a Placement, specificando i vincoli di capacità del \emph{flavor} desiderato (ad esempio, host con almeno 8 vCPU e 16 GB di RAM liberi).
    \item \textbf{Risposta ed elezione del candidato}: Placement esegue una query efficiente sul proprio data model e restituisce a \texttt{nova-scheduler} l'insieme dei soli \emph{Resource Providers} che dispongono di capacità residua sufficiente.
    \item \textbf{Creazione dell'allocazione}: Una volta eletto il nodo target, lo scheduler (o il conductor per suo conto) crea contestualmente su Placement un record di \textbf{Allocation} registrando l'UUID della nuova VM come consumatore delle rispettive quote di risorse.
    \item \textbf{Garanzia di consistenza atomica}: Questo meccanismo previene in modo deterministico condizioni di \emph{race condition} o situazioni di sovra-allocazione non pianificata (\emph{over-commitment}) che potrebbero scaturire da richieste concorrenti di avvio macchine virtuali su nodi condivisi.
\end{enumerate}

% CC -------------------------------------------------------------------------------

\section{Gestione Delle Risorse E Schedulazione In OpenStack Nova}

Nei sistemi di calcolo distribuito per infrastrutture \emph{Cloud IaaS} (\emph{Infrastructure as a Service}), uno dei compiti più complessi e gravosi dal punto di vista architetturale consiste nell'allocazione ottima delle risorse fisiche eterogenee (quali cicli di clock della CPU, banchi di memoria RAM e blocchi di storage persistente o effimero) a fronte di richieste asincrone generate dagli utenti. In \textbf{OpenStack}, il componente preposto alla gestione del ciclo di vita delle istanze di calcolo è \textbf{Nova}.

Nelle prime versioni del framework, il sottosistema di scheduling di Nova doveva mantenere simultaneamente in memoria e nel database relazionale sia lo stato globale dei server fisici (\emph{hypervisor}), sia il monitoraggio dettagliato di quali risorse fossero riservate a ciascuna macchina virtuale. Questa commistione di responsabilità portava a problemi endemici di \emph{race condition} durante le fasi di schedulazione concorrente ad alto carico e a disallineamenti tra l'effettiva capacità residua dei nodi e i vincoli espressi dai \emph{flavor} (profili dimensionali della macchina virtuale).

Per risolvere strutturalmente questo collo di bottiglia, l'architettura di OpenStack ha introdotto un servizio dedicato, completamente disaccoppiato dal punto di vista logico e di persistenza: il servizio di \textbf{Placement}.

\subsection{Il Servizio Di Placement: Modello Relazionale Delle Risorse}

Il servizio di \textbf{Placement} traccia l'inventario hardware e governa i vincoli quantitativi del consumo di risorse dell'intero data center. Esso espone un'interfaccia RESTful autonoma interrogata costantemente sia dal modulo di pianificazione (\textbf{nova-scheduler}), sia dai demoni di calcolo eseguiti sugli hypervisor (\textbf{nova-compute}).

% [IMMAGINE DEL FILE DI SLIDE CON NOME "openstack_nova_neutron.pdf"] Schema relazionale del Placement Service: Consumer, Allocation, Resource Provider, Inventory e Resource Class — Slide 82 / Pagina 2 del PDF

Il modello logico-relazionale su cui si fonda Placement si articola su cinque costrutti fondamentali strettamente interconnessi:

\begin{enumerate}[leftmargin=*]
    \item \textbf{Resource Provider (RP)}: rappresenta una qualsiasi entità fisica o logica in grado di fornire capacità computazionale o di archiviazione. Nel caso canonico di calcolo, ogni nodo fisico (\emph{compute node}) su cui è attivo un hypervisor (es. KVM, QEMU) viene registrato presso il Placement API come un distinto \emph{Resource Provider}. Ciascun provider possiede un identificativo univoco (UUID) e un insieme di attributi descrittivi.

    \item \textbf{Resource Class}: formalizza un tipo standardizzato e non fungibile di risorsa hardware quantificabile gestita dalla piattaforma. Placement definisce un catalogo rigido di classi standard conformi allo standard POSIX e all'ingegneria del silicio, tra cui:
    \begin{itemize}
        \item \texttt{VCPU}: processori virtuali assegnabili;
        \item \texttt{MEMORY\_MB}: quantità di memoria volatile espressa in megabyte;
        \item \texttt{DISK\_GB}: volume di archiviazione su disco locale espresso in gigabyte.
    \end{itemize}
    Tali classi consentono di uniformare la semantica delle metriche evitando ambiguità tra hypervisor differenti.

    \item \textbf{Inventory}: stabilisce la relazione quantitativa tra un determinato \emph{Resource Provider} e una specifica \emph{Resource Class}. L'\emph{Inventory} memorizza la capacità totale teorica disponibile sul nodo, la granularità minima e massima di allocazione, l'eventuale rapporto di \emph{overcommit} (rapporto di sovrasottoscrizione consentito, ad esempio permettendo un rapporto di $16:1$ per i core virtuali rispetto a quelli fisici) e una quota fissa di riserva riservata al sistema operativo ospitante.

    \item \textbf{Consumer}: identifica qualsiasi entità che richiede ed occupa risorse nell'infrastruttura. Nel contesto di Nova Compute, il \emph{Consumer} coincide ordinariamente con l'istanza di macchina virtuale (rappresentata anch'essa dal suo UUID di server OpenStack).

    \item \textbf{Allocation}: costituisce la transazione di consumo effettiva. Un'\emph{Allocation} modella il record contabile bidirezionale che collega un \emph{Consumer} a uno o più \emph{Resource Provider}, specificando per ciascuno l'esatta frazione numerica di \emph{Resource Class} consumata. Ad esempio, un record di allocazione registra formalmente che l'istanza $U_1$ (Consumer) impegna $2$ unità di \texttt{VCPU} e $4096$ megabyte di \texttt{MEMORY\_MB} a carico del nodo fisico \texttt{ComputeNode1} (Resource Provider).
\end{enumerate}

\subsubsection{Dinamica Operativa Del Flusso Di Allocazione}

Quando un utente invia una richiesta per istanziare una nuova macchina virtuale specificando una determinata immagine e un certo \emph{flavor}, il processo di negoziazione e mutua esclusione delle risorse si articola nei seguenti passaggi sequenziali:

\begin{enumerate}
    \item La richiesta di creazione perviene a \textbf{nova-api}, che la valida e delega la fase di individuazione del nodo a \textbf{nova-scheduler}.
    \item Il componente \textbf{nova-scheduler} interroga l'API REST di Placement inviando il profilo delle risorse necessarie (estratte dai requisiti di vCPU, RAM e disco definiti nel \emph{flavor}).
    \item Placement esegue una query interna incrociando gli inventari (\emph{Inventory}) di tutti i provider registrati al netto delle allocazioni già attive (\emph{Allocation}). Restituisce a Nova una lista di candidati idonei, ossia nodi che possiedono una capacità residua strettamente sufficiente a soddisfare la richiesta.
    \item Nova seleziona l'host definitivo attraverso i propri algoritmi interni di filtro e pesatura (\emph{Filter Scheduler}).
    \item Prima che l'istanza venga effettivamente costruita sull'hypervisor di destinazione, viene inviata una chiamata atomica a Placement per creare l'\textbf{Allocation}: la risorsa viene riservata formalmente a nome dell'UUID dell'istanza. Se l'operazione ha successo, il demone \textbf{nova-compute} locale procede con il provisioning della VM; qualora si verifichi un conflitto di allocazione concorrente, la transazione fallisce a livello di Placement prevenendo qualsiasi fenomeno di \emph{overbooking} accidentale.
\end{enumerate}

\newpage

\subsection{Il Modello Di Pianificazione: Il Filter Scheduler Di Nova}

Una volta ottenuta la platea dei nodi computazionali dotati di inventario sufficiente tramite Placement, la decisione vincolante su quale specifico elaboratore debba ospitare la macchina virtuale viene demandata al modulo \textbf{nova-scheduler}. L'algoritmo standard e maggiormente diffuso in ambiente di produzione è il \textbf{Filter Scheduler}.

% [IMMAGINE DEL FILE DI SLIDE CON NOME "openstack_nova_neutron.pdf"] Diagramma di funzionamento del Filter Scheduler: fase di filtraggio dei nodi candidati e successiva pesatura tramite weighers — Slide 89 / Pagina 11 del PDF

Il Filter Scheduler opera secondo una pipeline concettuale a due stadi successivi e rigidamente separati:
\[
    \text{Insieme di Tutti i Nodi} \xrightarrow{\quad\text{Filtraggio (Filtri)}\quad} \text{Nodi Idonei} \xrightarrow{\quad\text{Pesatura (Weighers)}\quad} \text{Nodo Selezionato}
\]

\subsubsection{Fase 1: Il Meccanismo Di Filtraggio (\emph{Filtering})}

Nel primo stadio, lo scheduler itera sull'intero insieme dei nodi di calcolo disponibili nel sistema, applicando in cascata una catena di regole booleane denominate \textbf{Filtri} (\emph{Filters}). Ciascun filtro analizza un requisito specifico del carico di lavoro e lo confronta con lo stato dell'host: se l'host non soddisfa il criterio, viene scartato in modo permanente per quella determinata richiesta.

I principali filtri standard integrati nel modulo \texttt{nova.scheduler.filters} includono:
\begin{itemize}[leftmargin=*]
    \item \textbf{AllHostsFilter}: un filtro passante di test che non applica alcuna esclusione, restituendo la totalità degli host censiti.
    \item \textbf{AvailabilityZoneFilter}: analizza la zona di disponibilità richiesta dal client e scarta tutti gli elaboratori non appartenenti alla medesima partizione logico-fisica. Supporta l'immissione di elenchi separati da virgole per garantire la compatibilità con zone multiple.
    \item \textbf{ImagePropertiesFilter}: ispeziona i metadati iniettati nell'immagine del sistema operativo (gestita da Glance), verificando ad esempio se l'hypervisor del nodo supporti una specifica architettura (es. \texttt{x86\_64} vs \texttt{aarch64}), una modalità di virtualizzazione o determinate estensioni di virtualizzazione annidata.
    \item \textbf{ComputeCapabilitiesFilter}: confronta le capacità computazionali dichiarate dal servizio di calcolo dell'host con le specifiche addizionali (\emph{extra specs}) definite nel \emph{flavor} dell'istanza.
    \item \textbf{ComputeFilter}: ammette esclusivamente i nodi fisici il cui servizio \texttt{nova-compute} sia attivo, operativamente integro e non posto esplicitamente in stato disabilitato dall'amministratore.
    \item \textbf{CoreFilter}, \textbf{RamFilter}, \textbf{DiskFilter}: controllano rispettivamente che l'host disponga di un quantitativo di core virtuali, memoria RAM e spazio disco sufficiente a soddisfare la richiesta, tenendo in debita considerazione gli indici di overcommit impostati nei file di configurazione locali.
\end{itemize}

\subsubsection{Fase 2: Il Meccanismo Di Pesatura (\emph{Weighting})}

Gli host che superano con successo la totalità dei filtri configurati vengono inoltrati al secondo stadio: la \textbf{Pesatura} (\emph{Weighting}). In questa fase, a ogni nodo idoneo viene assegnato un punteggio numerico normalizzato, calcolato come combinazione lineare di diverse metriche di carico attraverso moduli detti \emph{Weighers}.

Il weigher predefinito di Nova è il \textbf{RAM Weigher}. Il suo principio operativo standard premia gli host dotati del maggior quantitativo assoluto di memoria RAM libera. Tale comportamento realizza per default una politica di tipo \emph{spread} (distribuzione bilanciata): le macchine virtuali vengono distribuite uniformemente tra tutti gli elaboratori fisici disponibili, massimizzando le risorse residue su ciascuno ed evitando precoci saturazioni locali.

Il comportamento del motore di pesatura è interamente modificabile intervenendo sui parametri del file di configurazione \texttt{nova.conf}:
\begin{itemize}[leftmargin=*]
    \item \texttt{ram\_weight\_multiplier}: rappresenta il fattore moltiplicativo applicato alla quantità di RAM libera. Impostando questo valore su un numero positivo (comportamento standard), il sistema massimizza la RAM libera distribuendo i nodi. Qualora l'amministratore configuri invece un valore \textbf{negativo}, la logica di calcolo si inverte completamente: lo scheduler preferirà gli host con la \emph{minore} quantità di memoria disponibile residua, attuando una politica di tipo \emph{stacking} (consolidamento). In questo modo, le nuove macchine virtuali vengono ammassate sul medesimo server fino a completa saturazione, consentendo ai nodi non utilizzati di rimanere inattivi o di essere posti in modalità a basso consumo energetico.
    \item \texttt{scheduler\_host\_subset\_size}: governa il grado di casualità nella scelta finale. Al termine della fase di pesatura, gli host vengono ordinati in modo decrescente in base al punteggio finale ottenuto. Se \texttt{scheduler\_host\_subset\_size = 1}, lo scheduler sceglierà deterministicamente il nodo al primo posto (il migliore in assoluto). Se invece il parametro è impostato a un intero $N > 1$, lo scheduler estrae casualmente l'host di destinazione all'interno dell'insieme dei \textbf{primi $N$ migliori candidati}. Questo meccanismo previene contese simultanee qualora più istanze vengano schedulate parallelamente da thread concorrenti.
    \item Disabilitazione del RAM Weigher: qualora il componente venga rimosso dai pesatori attivi, la scelta tra gli host che hanno superato il filtraggio avviene in maniera puramente pseudo-casuale equiprobabile.
\end{itemize}

\subsection{Sequenza Di Bootstrap: Lancio Di Un'Istanza Base}

Per comprendere come interagiscono i moduli fondamentali di OpenStack in assenza di servizi ausiliari complessi (quali reti avanzate con Neutron o storage a blocchi persistente con Cinder), è utile esaminare il diagramma temporale che descrive l'evasione di una richiesta di creazione VM mediante storage locale ed elaborazione pura.

% [IMMAGINE DEL FILE DI SLIDE CON NOME "openstack_nova_neutron.pdf"] Diagramma temporale di lancio istanza: interazione tra User, Keystone, Nova e Glance — Slide 94 / Pagina 16 del PDF

La sequenza temporale si sviluppa attraverso i seguenti scambi funzionali:
\begin{enumerate}
    \item \textbf{Autenticazione iniziale}: l'utente invia le proprie credenziali (username, password o token applicativo) al servizio di identità \textbf{Keystone}, specificando il dominio o il tenant di riferimento.
    \item \textbf{Emissione del token}: Keystone valida le credenziali all'interno del proprio backend e restituisce al client un token di autenticazione crittografato (tipicamente in formato Fernet), associato ai permessi RBAC e all'identificativo del progetto (\emph{project\_id}).
    \item \textbf{Invocazione dell'API di calcolo}: il client trasmette una richiesta HTTP \texttt{POST} all'endpoint di Nova sull'URI standard:
    \begin{center}
        \texttt{POST /v2.1/\{project\_id\}/servers}
    \end{center}
    allegando nell'header \texttt{X-Auth-Token} il token emesso da Keystone e nel payload JSON i dettagli dell'istanza (nome del server, identificativo del \emph{flavor} e identificativo dell'immagine del sistema operativo).
    \item \textbf{Validazione del token da parte di Nova}: il servizio \textbf{nova-api} intercetta la richiesta e si interfaccia in modalità interattiva con \textbf{Keystone} per verificare la firma, l'integrità crittografica e la mancata scadenza del token ricevuto.
    \item \textbf{Risoluzione dell'immagine con Glance}: accertata la validità del richiedente, Nova contatta il servizio di gestione delle immagini \textbf{Glance}, trasmettendo il token autorizzativo dell'utente ed effettuando una richiesta per ottenere i metadati descrittivi e il puntatore al binario dell'immagine disco (\emph{VM Image}) richiesta.
    \item \textbf{Esito positivo e restituzione al client}: una volta verificata l'esistenza dell'immagine e generato il record dell'istanza a livello del database di controllo, Nova risponde tempestivamente al client con un codice di stato \texttt{202 Accepted} (\emph{Success}), restituendo l'UUID univoco della neonata macchina virtuale. La successiva trafila di download dell'immagine, esecuzione del provisioning sull'hypervisor e allocazione delle periferiche prosegue in modalità completamente asincrona.
\end{enumerate}

\newpage

\section{Scalabilità E Resilienza: L'Architettura OpenStack Cells}

\subsection{Limiti Del Modello Monolitico A Database E Coda Unica}

Nelle installazioni iniziali di modeste dimensioni, tutti i nodi di elaborazione e i microservizi di Nova dipendono da una singola istanza logica di database relazionale (es. MariaDB Galera Cluster) e da un unico broker di messaggistica asincrona (es. RabbitMQ). L'intero scambio di comunicazioni RPC (\emph{Remote Procedure Call}) tra i servizi di controllo (\texttt{nova-api}, \texttt{nova-conductor}, \texttt{nova-scheduler}) e le centinaia di demoni \texttt{nova-compute} periferici transita attraverso il medesimo bus di messaggi.

Questo approccio presenta gravi limiti fisici e architetturali:
\begin{itemize}[leftmargin=*]
    \item \textbf{Saturazione della coda di messaggistica}: al crescere del numero di hypervisor (oltre 500--1000 nodi), la frequenza dei messaggi periodici di stato (\emph{heartbeat}) e delle chiamate di stato satura le connessioni TCP aperte verso RabbitMQ, innescando degradazioni prestazionali irreversibili.
    \item \textbf{Collo di bottiglia sul Database}: le transazioni di lettura e scrittura concorrenti sulle tabelle delle istanze generano frequenti fenomeni di contesa sui blocchi (\emph{lock contention}), riducendo drasticamente il throughput di provisioning.
    \item \textbf{Rischio di guasto sistemico (\emph{Single Point of Failure})}: un guasto o un rallentamento sul cluster di database centrale blocca istantaneamente l'intera infrastruttura cloud a livello globale.
\end{itemize}

\subsection{Il Modello Ad Albero Di OpenStack Cells}

Per superare tali barriere di scalabilità, OpenStack ha introdotto l'architettura a partizioni scalabili denominata \textbf{Cells} (evoluta stabilmente nel modello \textbf{Cells v2}). L'idea fondamentale consiste nello \emph{sharding} orizzontale dei nodi computazionali, suddividendo l'infrastruttura in raggruppamenti logici autonomi detti \textbf{celle}, disposti secondo una gerarchia ad albero.

Nel modello contemporaneo, ogni distribuzione di OpenStack Nova è intrinsecamente strutturata come un'architettura a celle: anche un'installazione minima a singolo server rappresenta, per definizione formale del codice sorgente, un sistema a cella singola. Di conseguenza, l'espansione verso cluster multi-cella con decine di migliaia di nodi non richiede un cambio di paradigma applicativo, ma esclusivamente l'aggiunta modulare di nuovi segmenti di calcolo.

% [IMMAGINE DEL FILE DI SLIDE CON NOME "openstack_nova_neutron.pdf"] Architettura di OpenStack Cells: confronto topologico tra installazione a cella singola (Single cell) e installazione distribuita multi-cella (Multiple cells) — Slide 84 / Pagina 5 del PDF

\subsubsection{Topologia Dei Servizi E Partizionamento Dei Dati}

La struttura organizzativa separa rigorosamente il livello apicale globale dai livelli periferici:
\begin{itemize}[leftmargin=*]
    \item \textbf{Livello Apicale (\emph{Top-Level Cell})}: ospita il servizio globale di ricezione delle richieste \textbf{nova-api} e il componente di pianificazione \textbf{nova-scheduler}, ma non contiene alcun nodo computazionale né istanze in esecuzione.
    \item \textbf{Livello Periferico (\emph{Child Cells})}: ogni singola cella racchiude al proprio interno una porzione dei nodi di calcolo (\texttt{nova-compute}) e dispone in via esclusiva di un proprio server database dedicato e di un proprio broker RabbitMQ isolato. In tali celle sono attivi tutti i servizi ordinari di Nova a eccezione di \texttt{nova-api}.
\end{itemize}

La persistenza relazionale dell'infrastruttura viene conseguentemente scissa in due database con compiti ben distinti:
\begin{itemize}
    \item \textbf{API Database (Globale)}: memorizza le entità a validità trasversale per l'intero cloud, quali la mappatura delle celle (\emph{cell mappings}), le regole di sicurezza, i profili dimensionali (\emph{flavors}) e le tabelle che associano ciascun UUID di istanza alla specifica cella in cui risiede fisicamente.
    \item \textbf{Cell Database (Locale)}: ciascuna cella detiene un proprio database locale completamente disconnesso da quello delle altre celle. Questo database locale memorizza la quasi totalità delle informazioni dettagliate relative alle istanze appartenenti a quella cella specifica (stato operativo, allocazione disco, interfacce di rete collegate, metadati).
\end{itemize}

Quando un nuovo compute node viene configurato, i parametri di connessione definiti nel file locale puntano univocamente al database di cella e al RabbitMQ di cella. Il nodo si registra nel database locale, e un'apposita utility a linea di comando, \texttt{nova-manage}, deve essere eseguita per censire l'host all'interno dell'API database globale, rendendolo individuabile dallo scheduler.

\subsubsection{Il Ruolo Speciale Di Cell0}

All'interno dell'architettura a celle risiede una partizione del tutto peculiare denominata \textbf{cell0}.
\begin{itemize}[leftmargin=*]
    \item Il \textbf{cell0 database} ricalca perfettamente lo schema tabelle di un database di cella convenzionale, ma non dispone di alcun nodo computazionale reale né di broker RabbitMQ associato.
    \item La sua funzione è quella di ospitare il record di tutte quelle istanze la cui schedulazione è fallita in modo irreversibile (ad esempio a causa della totale assenza di nodi fisici provvisti delle risorse richieste dai filtri).
    \item Registrando l'istanza fallita in cell0, il sistema garantisce che l'utente possa comunque interrogare l'API REST globale per verificare lo stato di errore del server ed eventualmente procedere alla sua cancellazione, senza che sia mai stata impegnata alcuna risorsa su una cella operativa reale.
\end{itemize}

\subsubsection{Il Super Conductor}

La scissione dei canali di comunicazione RPC ha imposto l'introduzione di una figura specializzata del servizio di coordinamento: il \textbf{Super Conductor}. Mentre all'interno di ciascuna cella opera un'istanza locale di \textbf{nova-conductor} deputata a interagire esclusivamente con il database locale e i nodi computazionali della cella stessa, a livello globale viene istanziato il Super Conductor.

Questo componente ha accesso contemporaneo sia all'API Database sia a tutti i database di cella: coordina l'attività di \texttt{nova-scheduler}, gestisce la transizione dell'istanza tra il livello API e la cella prescelta e convoglia i record in cell0 qualora la fase di pianificazione termini con esito negativo.

\newpage

\section{Macchina A Stati E Ciclo Di Vita Delle Istanze}

La gestione delle macchine virtuali all'interno di un'architettura cloud richiede un modello formale rigoroso che governi l'evoluzione del ciclo di vita dei carichi di lavoro. Nova implementa una macchina a stati finiti strutturata su due livelli ortogonali: lo stato macroscopico stabile dell'istanza (\textbf{vm\_state}) e lo stato transitorio o puntuale del comando in esecuzione (\textbf{task\_state}).

\subsection{Stati Macroscopici Dell'Istanza (\texttt{vm\_state})}

Lo stato identificato dalla proprietà \texttt{vm\_state} descrive la condizione operativa permanente della macchina virtuale dal punto di vista dell'infrastruttura e del database di persistenza.

% [IMMAGINE DEL FILE DI SLIDE CON NOME "openstack_nova_neutron.pdf"] Diagramma a grafo degli stati e delle transizioni della macchina virtuale (vm_state) in OpenStack — Slide 85 / Pagina 6 del PDF

Le principali condizioni operative previste dal motore di calcolo sono le seguenti:

\begin{itemize}[leftmargin=*]
    \item \textbf{INITIALIZED} (precedentemente denominato \textbf{BUILDING}): la macchina virtuale è stata formalmente registrata a livello di database centrale ma non è ancora stata creata né istanziata fisicamente sui nodi hypervisor.
    \item \textbf{ACTIVE}: l'istanza è pienamente operativa, in esecuzione sull'hypervisor di destinazione e sta eseguendo il sistema operativo caricato dall'immagine di avvio.
    \item \textbf{STOPPED}: la macchina virtuale si trova in stato spento; i cicli di calcolo della CPU e la memoria volatile risultano deallocati, mentre l'immagine e i file del disco persistono intatti sullo storage dell'host.
    \item \textbf{PAUSED}: l'esecuzione dell'istanza viene congelata temporaneamente sul nodo. Lo stato della CPU viene bloccato e i consumi di calcolo azzerati, ma l'intero footprint di memoria rimane allocato all'interno della RAM fisica dell'hypervisor.
    \item \textbf{SUSPENDED}: l'istanza viene posta in sospensione profonda (\emph{hibernation}). L'hypervisor effettua un dump completo dello stato della memoria volatile su un file di snapshot su disco; la memoria fisica viene integralmente rilasciata per altre elaborazioni. Al risveglio, lo stato di esecuzione viene ripristinato esattamente dal file salvato.
    \item \textbf{RESCUED}: l'istanza viene avviata collegandovi un'immagine kernel/disco di emergenza (\emph{rescue image}), mantenendo il disco originario montato come periferica secondaria. Questa modalità viene tipicamente attivata per consentire all'amministratore operazioni straordinarie di recupero da filesystem danneggiati o problemi di boot.
    \item \textbf{RESIZED}: stato caratteristico che segue un'operazione di ridimensionamento del \emph{flavor} o migrazione a freddo. L'istanza è arrestata sul nodo di origine e avviata sul nodo di destinazione con le nuove caratteristiche dimensionali; i dischi sussistono temporaneamente in entrambe le posizioni in attesa che l'utente confermi formalmente l'operazione (\emph{confirm}) o ne richieda il ripristino allo stato precedente (\emph{revert}).
    \item \textbf{SOFT\_DELETED}: l'istanza cessa la propria esecuzione e viene scaricata dai nodi di calcolo, ma la sua immagine su disco viene preservata per un intervallo di conservazione configurabile. In tale finestra temporale, un utente autorizzato ha la possibilità di ripristinare il server (\emph{restore}) al suo normale stato operativo. Solo allo scadere del timeout il server viene eliminato in modo irrevocabile.
    \item \textbf{HARD\_DELETED}: l'istanza cessa di esistere formalmente sia per il registro di contabilità (\emph{billing}) sia per i massimali di quota utente. I file su disco associati e le strutture di memoria residue sull'hypervisor vengono distrutti.
    \item \textbf{ERROR}: indica il verificarsi di un'anomalia irreversibile durante un'operazione di provisioning, migrazione o esecuzione. Da questa condizione l'unica azione permessa è la cancellazione dell'istanza (\emph{delete}).
\end{itemize}

\subsection{Meccanismi Di Shelving Ed Evacuazione Delle Risorse}

Nelle infrastrutture cloud è frequente la necessità di conservare un'istanza complessa configurata dall'utente senza tuttavia consumare inutilmente preziose quote di RAM e vCPU sugli hypervisor (ad esempio durante i periodi di fermo programmato nel fine settimana). A questo scopo Nova implementa la procedura di \textbf{Shelving}.

\begin{figure}[H]
    \centering
    \begin{tabularx}{\textwidth}{p{4.2cm}X}
        \toprule
        \textbf{Stato / Comando} & \textbf{Descrizione e Implicazioni Operative} \\
        \midrule
        \texttt{SHELVED}\newline(\texttt{nova shelve}) & Arresta in modo controllato l'istanza; se la VM poggia su disco effimero (non basata su volumi Cinder), genera un'istantanea (\emph{snapshot}) dell'immagine e la invia al registro Glance. I dati in memoria volatile vengono persi. Le associazioni logiche dell'istanza rimangono inalterate. \\
        \addlinespace
        \texttt{SHELVED\_OFFLOADED}\newline(\texttt{nova shelve-offload}) & Rimuove completamente l'istanza dall'hypervisor (\emph{offloading}), liberando integralmente memoria, core CPU e spazio disco locale sull'host fisico. Per essere riutilizzata, richiede l'invocazione di un'operazione di risveglio (\texttt{unshelve}). \\
        \addlinespace
        \texttt{nova unshelve} & Ripristina l'istanza precedentemente archiviata: risottomette la richiesta di pianificazione a Nova Scheduler e Placement per individuare un nodo fisico con capacità disponibile, ricrea l'ambiente e re-inietta lo snapshot. \\
        \bottomrule
    \end{tabularx}
    \caption{Tassonomia e comandi per la gestione dello shelving in OpenStack Nova.}
\end{figure}

Il comando \texttt{nova shelve} tenta preventivamente di impartire uno shutdown controllato (\emph{graceful shutdown}) al sistema operativo guest inviando un segnale ACPI. La durata di questa finestra di attesa è vincolata da precisi parametri gerarchici:
\begin{enumerate}
    \item \textbf{Parametro globale}: definito nel file di configurazione \texttt{nova.conf} sotto la direttiva:
    \begin{center}
        \texttt{shutdown\_timeout = 60}
    \end{center}
    Esprime il valore predefinito pari a $60$ secondi concessi all'OS ospite prima che Nova forzi lo spegnimento brutale dell'alimentazione virtuale.
    \item \textbf{Override per-immagine}: il valore globale può essere scavalcato impostando direttamente nei metadati dell'immagine Glance la proprietà:
    \begin{center}
        \texttt{os\_shutdown\_timeout}
    \end{center}
    Tale meccanismo consente di accomodare sistemi operativi lenti a completare la disattivazione dei servizi di sistema (ad esempio server Windows o appliance complesse) garantendo una chiusura integra del filesystem.
\end{enumerate}

\subsection{Gli Stati Di Transizione Operativa (\texttt{task\_state})}

Mentre \texttt{vm\_state} delinea la macro-condizione strutturale, la proprietà \textbf{task\_state} modella l'esatta fase transitoria associata a una specifica chiamata di Compute API in esecuzione. Durante il normale funzionamento asincrono, il valore di default è impostato su \texttt{None}, a indicare che nessun processo straordinario è in corso.

La tabella seguente riassume la catalogazione esaustiva degli stati transitori supportati:

\begin{table}[H]
    \centering
    \begin{tabularx}{\textwidth}{p{4.5cm}X}
        \toprule
        \textbf{Ambito Funzionale} & \textbf{Elenco Valori \texttt{task\_state}} \\
        \midrule
        \textbf{Creazione e Provisioning} & \texttt{SCHEDULING}, \texttt{BLOCK\_DEVICE\_MAPPING}, \texttt{NETWORKING}, \texttt{SPAWNING}, \texttt{BUILDING} \\
        \addlinespace
        \textbf{Copia e Manutenzione} & \texttt{IMAGE\_SNAPSHOTTING}, \texttt{IMAGE\_BACKINGUP}, \texttt{UPDATING\_PASSWORD}, \texttt{REBUILDING} \\
        \addlinespace
        \textbf{Controllo Esecuzione} & \texttt{PAUSING}, \texttt{UNPAUSING}, \texttt{SUSPENDING}, \texttt{RESUMING}, \texttt{STOPPING}, \texttt{STARTING}, \texttt{POWERING\_OFF}, \texttt{POWERING\_ON}, \texttt{REBOOTING} \\
        \addlinespace
        \textbf{Ripristino e Rimozione} & \texttt{RESCUING}, \texttt{UNRESCUING}, \texttt{DELETING} \\
        \addlinespace
        \textbf{Ridimensionamento / Migrazione} & \texttt{RESIZE\_PREP}, \texttt{RESIZE\_MIGRATING}, \texttt{RESIZE\_MIGRATED}, \texttt{RESIZE\_FINISH}, \texttt{RESIZING}, \texttt{RESIZE\_CONFIRMING}, \texttt{RESIZE\_REVERTING} \\
        \bottomrule
    \end{tabularx}
    \caption{Quadro esaustivo degli stati transitori del task engine (\texttt{task\_state}).}
\end{table}

\newpage

\section{Architettura Di Rete Con OpenStack Neutron}

Nel paradigma IaaS, la fornitura isolata di pura capacità di calcolo risulta priva di utilità se non accompagnata da un'infrastruttura programmabile capace di collegare tra loro le istanze in topologie arbitrarie. In OpenStack, questo ruolo è assolto da \textbf{Neutron}, il quale eroga connettività di rete su richiesta (\emph{Network-as-a-Service}) per i dispositivi di interfaccia virtuali (\textbf{vNIC}) istanziati da Nova.

\subsection{Fondamenti Di Virtual Networking E Primitive Di Neutron}

Per comprendere il funzionamento di Neutron è opportuno richiamare i meccanismi con cui i sistemi operativi moderni collegano le macchine virtuali allo stack di rete:
\begin{itemize}[leftmargin=*]
    \item Un hypervisor genera per ciascuna macchina virtuale un'interfaccia virtuale nota come \textbf{VIF} (\emph{Virtual Network Interface}).
    \item Nel kernel dell'elaboratore host, tale interfaccia si manifesta solitamente sotto forma di dispositivo \textbf{TAP}, ovvero un'estremità virtuale in grado di iniettare ed estrarre trame Ethernet direttamente a Livello 2 (Data Link).
    \item Per consentire la comunicazione tra più interfacce TAP sul medesimo server fisico, il kernel Linux ricorre a ponti software (\textbf{bridge}), realizzati tramite il modulo canonico \textbf{Linux Bridge} o mediante soluzioni avanzate a commutazione di flusso quale \textbf{Open vSwitch (OVS)}.
\end{itemize}

% [IMMAGINE DEL FILE DI SLIDE CON NOME "openstack_nova_neutron.pdf"] Schema della connettività su singolo compute node: macchine virtuali collegate all'Integration Bridge (br-int) — Slide 97 / Pagina 19 del PDF

Neutron fornisce un'astrazione dichiarativa di questa tecnologia mediante quattro primitive concettuali cardine:
\begin{enumerate}[leftmargin=*]
    \item \textbf{Network}: modella un segmento virtuale di Livello 2 isolato. Dal punto di vista del funzionamento logico, una rete Neutron equivale a uno switch virtuale dotato di porte create e rimosse in modo dinamico. Garantisce la diffusione del traffico broadcast e unicast a livello di trama Ethernet tra tutte le interfacce collegate.
    \item \textbf{Port}: rappresenta la porta logica di connessione a un determinato switch virtuale (Network). A ogni porta è associato un indirizzo MAC univoco e una tupla di parametri di sicurezza (es. filtri anti-spoofing). Un router virtuale o la vNIC di una VM si connettono alla rete impegnando una porta.
    \item \textbf{Subnet}: specifica un pool coerente di indirizzi IPv4 o IPv6 appartenenti a un blocco CIDR (es. \texttt{192.168.10.0/24}). Assegna le coordinate di Livello 3 alle porte, governando l'allocazione tramite server DHCP virtuale o meccanismi stateless. Due diverse subnet collegate alla medesima o a differenti reti non possono comunicare direttamente senza l'interposizione di un router.
    \item \textbf{Router}: dispositivo di instradamento virtuale di Livello 3 deputato alla selezione dell'instradamento e all'inoltro dei pacchetti IP tra differenti subnet e verso l'esterno (\emph{gateway} e SNAT/DNAT verso la rete pubblica tramite \emph{Floating IP}).
\end{enumerate}

\subsection{Isolamento Locale E Connettività Distribuita}

Il soddisfacimento dei requisiti di connettività in ambiente multi-inquilino (\emph{multi-tenant}) richiede che tenant distinti possano istanziare topologie aventi i medesimi piani di indirizzamento privato (es. entrambi utilizzanti la subnet \texttt{10.0.0.0/24}) in totale assenza di collisioni di pacchetti e senza compromissione della sicurezza.

\subsubsection{L'Integration Bridge E Le VLAN Locali}

Sul singolo nodo computazionale, Neutron convoglia ordinariamente il traffico proveniente dalle interfacce VIF delle varie macchine virtuali su un unico commutatore centrale Open vSwitch denominato \textbf{Integration Bridge} (\texttt{br-int}).

Tuttavia, poiché su tale bridge transitano simultaneamente flussi generati da tenant reciprocamente isolati, l'apparato non può operare come un semplice hub condiviso. Per separare i domini di broadcast a livello locale senza sprecare etichette globali, Neutron applica il principio delle \textbf{VLAN Locali}:
\begin{itemize}
    \item Alle porte virtuali che collegano le VIF delle macchine virtuali al bridge \texttt{br-int} viene assegnato un tag VLAN interno (compreso tra $1$ e $4094$), valido tassativamente ed esclusivamente all'interno del singolo host fisico.
    \item A ogni rete Neutron attiva su quell'elaboratore viene associato un tag VLAN locale differente. Se due macchine virtuali risiedono sul medesimo nodo ma appartengono a tenant diversi, esse vengono taggate con identificatori numerici disgiunti: l'interruttore OVS impedisce intrinsecamente qualsiasi transito di frame tra i due flussi.
\end{itemize}

\subsubsection{Connettività Inter-Host: Modelli Flat, VLAN E Reti Overlay}

Nel momento in cui due macchine virtuali appartenenti alla medesima rete logica vengono allocate su server computazionali fisicamente distanti, l'infrastruttura deve estendere il dominio di Livello 2 attraverso l'infrastruttura di rete fisica (\emph{underlay network}).

% [IMMAGINE DEL FILE DI SLIDE CON NOME "openstack_nova_neutron.pdf"] Schema di interconnessione delle macchine virtuali e degli Integration Bridge distribuiti su nodi computazionali distinti — Slide 99 / Pagina 21 del PDF

A tale scopo Neutron supporta tre paradigmi implementativi principali:
\begin{enumerate}[leftmargin=*]
    \item \textbf{Reti Flat}: l'integrazione bridge viene collegato direttamente, senza incapsulamento né tag, a una scheda di rete fisica (\emph{physical network interface}) dell'host. Tutte le istanze collegate appartengono al medesimo dominio di broadcast fisico e condividono il segmento Ethernet con i dispositivi reali. In questo scenario vi può essere al più \textbf{un'unica rete flat per deployment} (o per interfaccia fisica dedicata), annullando ogni capacità di multi-tenancy avanzata.

    \item \textbf{Reti VLAN (Trunking Globale)}: il traffico in uscita dall'integration bridge viene etichettato con un tag IEEE 802.1Q conforme a un ID VLAN globale assegnato da Neutron. Tale ID globale è comune a tutti gli host per quella specifica rete, ma è logicamente disaccoppiato dall'ID VLAN locale impiegato su \texttt{br-int}. Sebbene robusto e nativamente supportato da tutti gli switch fisici aziendali, questo approccio soffre del limite strutturale del protocollo 802.1Q, che impone un tetto massimo invalicabile di 4094 VLAN totali sull'intera infrastruttura del data center.

    \item \textbf{Reti Tunneling (Overlay Networks)}: per superare la restrizione delle 4094 VLAN e rendere la topologia logica del cloud del tutto indipendente dalla configurazione degli switch fisici, Neutron impiega protocolli di incapsulamento \emph{overlay}. La trama Ethernet di Livello 2 generata dalla VM viene incapsulata all'interno di un pacchetto di Livello 3/4 (tipicamente UDP o IP grezzo) e instradata attraverso una normale rete IP routed. I protocolli supportati includono:
    \begin{itemize}
        \item \textbf{VXLAN} (\emph{Virtual Extensible LAN}): incapsula i frame Ethernet in datagrammi UDP sulla porta convenzionale 4789. Introduce un identificatore a 24 bit (\textbf{VNI}, \emph{Virtual Network Identifier}), espandendo il limite teorico fino a oltre $16$ milioni di reti virtuali distinte ($2^{24}$).
        \item \textbf{GRE} (\emph{Generic Routing Encapsulation}): protocollo di incapsulamento IP-in-IP che adotta chiavi a 32 bit per isolare i canali logici.
        \item \textbf{Geneve} (\emph{Generic Network Virtualization Encapsulation}): protocollo moderno concepito per superare i limiti di estendibilità di VXLAN e GRE, consentendo il trasporto di metadati con opzioni di lunghezza variabile TLV (\emph{Type-Length-Value}).
    \end{itemize}
\end{enumerate}

\newpage

\subsection{Architettura Interna: Demone, Moduli E Agenti Distribuiti}

L'architettura del sottosistema di networking è costruita su un modello modulare orientato ai servizi, disaccoppiando l'elaborazione centrale delle chiamate API dall'attuazione concreta dei comandi sugli hypervisor e sui nodi di rete.

% [IMMAGINE DEL FILE DI SLIDE CON NOME "openstack_nova_neutron.pdf"] Architettura di dettaglio di OpenStack Neutron: interazione tra neutron-server, database relazionale, bus AMQP ed agenti periferici — Slide 107 / Pagina 31 del PDF

La cooperazione tra i vari sottosistemi si struttura secondo la gerarchia illustrata di seguito:

\subsubsection{Il Controller Node E Il Demone \texttt{neutron-server}}

Il nodo di controllo esegue il demone principale \textbf{neutron-server}. Questo componente:
\begin{itemize}
    \item Espone gli endpoint REST per la ricezione delle richieste API provenienti dagli utenti o dagli altri servizi di OpenStack (in primis Nova).
    \item Integra le \textbf{API Extensions}: moduli software ausiliari in linguaggio Python caricati all'avvio da una cartella dedicata, che introducono funzionalità avanzate non previste dalle API primitive (es. \emph{Quality of Service} QoS, \emph{Port Security} per il filtraggio degli indirizzi MAC/IP abusivi, allocazioni avanzate e \emph{Firewall-as-a-Service} FWaaS). Le estensioni possono essere di azione (\emph{action}), di risorsa (\emph{resource}) o di arricchimento campi (\emph{request}).
    \item Inizializza il motore RPC per l'invio asincrono delle direttive operative agli agenti distribuiti tramite il message broker (AMQP / RabbitMQ).
    \item Esegue il codice dei \textbf{Plugin}, memorizzando la persistenza delle entità (reti, subnet, porte, regole di firewall) nel database relazionale di Neutron.
\end{itemize}

\subsubsection{La Tassonomia Dei Plugin: Core E Service Plugins}

I plugin di Neutron sono classi Python plug-in eseguite all'interno del processo \texttt{neutron-server} sul nodo controller. Essi hanno il compito di tradurre le direttive astratte ricevute via API in sequenze operative destinate alle piattaforme di rete sottostanti (interruttori fisici, SDN controller o agenti host):
\begin{itemize}[leftmargin=*]
    \item \textbf{Core Plugin}: governa le primitive essenziali di Livello 2 (reti, subnet, porte logiche) e la gestione delle allocazioni IP. Nelle distribuzioni storiche era implementato da plugin monolitici specifici per tecnologia (es. il vecchio \emph{Open vSwitch Plugin}); nelle architetture odierne è integralmente sostituito dal modulo modulare \textbf{ML2}.
    \item \textbf{Service Plugin}: introduce e gestisce funzionalità di rete avanzate e servizi di livello superiore, quali:
    \begin{itemize}
        \item \textbf{L3 Router Plugin}: governa la creazione di router virtuali, la traslazione degli indirizzi di rete (NAT) e l'assegnazione dei Floating IP;
        \item \textbf{FWaaS} (\emph{Firewall-as-a-Service}): attua regole di filtraggio perimetrale del traffico a monte dei router di tenant;
        \item \textbf{LBaaS} (\emph{Load-Balancer-as-a-Service}): distribuisce il traffico in ingresso tra più istanze di backend.
    \end{itemize}
\end{itemize}

\subsubsection{Gli Agenti Distribuiti (\emph{Neutron Agents})}

Poiché il demone \texttt{neutron-server} non manipola direttamente gli stack di rete locali dei nodi di calcolo, tale compito è affidato a una schiera di agenti specializzati eseguiti in background sui nodi computazionali e di rete:
\begin{itemize}[leftmargin=*]
    \item \textbf{L2 Agent (es. OVS Agent o LinuxBridge Agent)}: riceve chiamate RPC dal controller e manipola i bridge locali (\texttt{br-int}, \texttt{br-tun}), programmando le regole di flusso e applicando le etichette VLAN necessarie all'isolamento delle porte.
    \item \textbf{DHCP Agent}: implementa server DHCP isolati (allocati ordinariamente all'interno di \emph{network namespaces} dedicati) per erogare in modo dinamico gli indirizzi IP e i parametri di default gateway alle VM.
    \item \textbf{L3 Agent}: presidia i nodi di rete (\emph{Network Nodes}), istanziando router virtuali mediante iptables/nftables e gestendo i flussi da e verso le reti esterne.
    \item \textbf{Metadata Agent}: consente alle macchine virtuali prive di indirizzo pubblico di interrogare il servizio di metadati di Nova (ad esempio per ottenere le chiavi SSH pubbliche o i file di configurazione \emph{cloud-init}) effettuando chiamate HTTP su un indirizzo IP link-local riservato (\texttt{169.254.169.254}).
\end{itemize}

\newpage

\subsection{Il Modulo Modulare ML2 (Modular Layer 2)}

Il plugin \textbf{ML2} (\emph{Modular Layer 2}) rappresenta il pilastro fondamentale della virtualizzazione di rete in OpenStack. Concepito per superare l'incompatibilità dei vecchi plugin monolitici (i quali costringevano l'intero cloud a basarsi su un'unica ed esclusiva implementazione tecnologica), ML2 opera come un framework altamente scalabile che consente la coesistenza simultanea di differenti metodologie di incapsulamento e apparati multivendor eterogenei.

% [IMMAGINE DEL FILE DI SLIDE CON NOME "openstack_nova_neutron.pdf"] Struttura interna del plugin ML2: correlazione e cooperazione tra Type Drivers e Mechanism Drivers — Slide 105 / Pagina 28 del PDF

L'architettura di ML2 separa in modo netto la nozione astratta di come una rete viene strutturata logicamente dal modo concreto in cui i pacchetti vengono commutati sui nodi. Tale disaccoppiamento è attuato tramite due categorie complementari di driver:

\subsubsection{1. Type Drivers}

I \textbf{Type Drivers} definiscono e mantengono lo stato logico associato a una determinata tecnologia di rete di Livello 2.
\begin{itemize}
    \item Ciascun tipo di rete supportato da Neutron (quali \texttt{local}, \texttt{flat}, \texttt{vlan}, \texttt{vxlan}, \texttt{gre}, \texttt{geneve}) è governato da un rispettivo Type Driver.
    \item I Type Driver convalidano i parametri immessi dall'amministratore per le reti di provider (\emph{provider networks}) e gestiscono l'allocazione autonoma di un segmento libero (ad esempio estraendo un tag VLAN o un VNI VXLAN da un pool preconfigurato) quando un tenant genera una nuova rete privata.
\end{itemize}

\subsubsection{2. Mechanism Drivers}

I \textbf{Mechanism Drivers} hanno il compito di recepire le informazioni di segmentazione formalizzate dal Type Driver e applicarle concretamente sul backend tecnologico abilitato.
\begin{itemize}
    \item Possono pilotare agenti residenti sui nodi computazionali mediante comunicazioni RPC (es. agendo su Open vSwitch o sul sottosistema Linux Bridge) oppure interagire in via diretta tramite protocolli REST/NETCONF con apparati hardware esterni, smartNIC o SDN controller (es. Cisco Nexus, Cisco ACI, Arista, driver SR-IOV per commutazione diretta in hardware ad altissime prestazioni).
    \item Uno dei vantaggi più significativi di ML2 è la capacità di abilitare simultaneamente più Mechanism Driver: ad esempio, una medesima rete di Livello 2 può avere alcune porte instanziate mediante Open vSwitch standard per macchine virtuali convenzionali e altre porte attestate tramite SR-IOV per elaborazioni scientifiche a bassa latenza.
\end{itemize}

\subsubsection{Profili Di Configurazione Di Rete}

L'amministratore istruisce il framework ML2 modificando i file di configurazione standard residenti sui nodi controller e computazionali. In particolare, il file centrale di configurazione di ML2 è situato nel percorso:
\begin{center}
    \texttt{/etc/neutron/plugins/ml2/ml2\_conf.ini}
\end{center}
All'interno di questo file, una tipica configurazione specifica l'elenco dei Type Driver abilitati, i Mechanism Driver ammessi e i range di allocazione automatica delle reti tenant:

\begin{lstlisting}[language=bash]
[ml2]
# Definizione dei tipi di rete supportati contemporaneamente
type_drivers = flat,vlan,vxlan

# Meccanismi di attuazione fisica abilitati
mechanism_drivers = openvswitch,linuxbridge

# Tipi di rete assegnabili automaticamente ai progetti tenant
tenant_network_types = vxlan,vlan

[ml2_type_vlan]
# Intervallo dei tag VLAN fisici disponibili per i progetti
network_vlan_ranges = physnet1:100:200

[ml2_type_vxlan]
# Range identificatori VXLAN (VNI) riservati ai tenant
vni_ranges = 10001:20000
\end{lstlisting}

Sui nodi computazionali in cui l'attuazione è delegata a Linux Bridge, il relativo agente periferico viene configurato nel file complementare:
\begin{center}
    \texttt{/etc/neutron/plugins/ml2/linuxbridge\_agent.ini}
\end{center}
in cui vengono esplicitate le corrispondenze puntuali tra le etichette logiche delle reti fisiche e le interfacce Ethernet reali dell'hypervisor (es. \texttt{physnet1:eth1}), garantendo che ogni pacchetto iniettato dalle interfacce VIF trovi l'instradamento corretto verso l'infrastruttura di commutazione del data center.

% CC ---------------------------------------------------------------------------

\section{Estensioni Architetturali Di OpenStack Neutron: Driver E Servizi Avanzati}

Il framework di networking \textbf{OpenStack Neutron} è strutturato secondo un paradigma modulare a plug-in che consente di disaccoppiare la logica di controllo astratta dalle effettive tecnologie di commutazione e instradamento fisico o virtuale. Per comprendere appieno la flessibilità di Neutron, è opportuno richiamare brevemente la separazione concettuale fondamentale introdotta dal modello \emph{Software-Defined Networking} (SDN): il piano di controllo (\emph{control plane}), centralizzato logicamente e demandato all'orchestrazione delle politiche di connettività, e il piano dati (\emph{data plane} o \emph{forwarding plane}), distribuito sui singoli nodi di calcolo (\emph{compute nodes}) ed execution engine preposti all'instradamento effettivo dei frame e dei pacchetti.

All'interno di tale ecosistema, il modulo centrale \textbf{ML2} (\emph{Modular Layer 2}) adotta un'architettura a due livelli: i \emph{Type Driver}, incaricati di definire il tipo di isolamento di rete (quali \texttt{flat}, \texttt{vlan}, \texttt{vxlan}, \texttt{gre}), e i \textbf{Mechanism Driver}, responsabili di tradurre i requisiti ad alto livello in configurazioni concrete per specifici switch virtuali o apparati fisici di rete.

\subsection{Meccanismi E Driver Specifici Dei Vendor}

Il panorama dei driver supportati all'interno della distribuzione principale e manutenuti direttamente dalla comunità di Neutron include sia implementazioni software puramente open-source sia moduli di integrazione proprietari:
\begin{itemize}[noitemsep]
    \item \textbf{Open vSwitch (OVS)} e \textbf{Linux Bridge}: costituiscono i due capisaldi software per la commutazione di rete in ambiente host Linux, supportando incapsulamento VLAN e overlay di tipo VXLAN e GRE.
    \item \textbf{Cisco UCS / Nexus} e \textbf{Cisco Nexus 1000v}: moduli di integrazione hardware e software per apparati industriali Cisco.
    \item \textbf{Modular Layer 2 (ML2)}: il meta-plugin di riferimento che accoglie la cooperazione eterogenea di driver multipli.
    \item \textbf{Nicira Network Virtualization Platform (NVP)}: piattaforma pionieristica di virtualizzazione di rete SDN basata su tunnel logici.
    \item \textbf{Ryu OpenFlow Controller} e \textbf{NEC OpenFlow}: implementazioni che dialogano direttamente con controller basati sullo standard aperto OpenFlow per la programmazione reattiva e proattiva delle tabelle di flusso.
    \item \textbf{Big Switch Controller Plugin}: integrazione di tessiture di commutazione centralizzate commerciali.
    \item \textbf{Cloudbase Hyper-V}: supporto nativo all'hypervisor Microsoft per ambienti eterogenei.
    \item \textbf{MidoNet}: architettura SDN distribuita per la virtualizzazione di rete senza hop centralizzati.
    \item \textbf{Brocade Neutron Plugin}, \textbf{PLUMgrid}, \textbf{Mellanox Neutron Plugin} ed \textbf{Embrane Neutron Plugin}: driver per piattaforme ad alte prestazioni, con particolare attenzione all'offloading hardware (es. SR-IOV per Mellanox) e all'orchestrazione distribuita di servizi di livello superiore.
    \item Ulteriori plugin del catalogo comunitario comprendono: \textbf{IBM SDN-VE}, \textbf{CPLANE NETWORKS}, \textbf{Nuage Networks}, \textbf{OpenContrail}, \textbf{Lenovo Networking} e \textbf{Avaya Neutron Plugin}.
\end{itemize}

Oltre all'albero di codice principale integrato nella distribuzione, sono disponibili soluzioni di terze parti (\emph{out-of-tree plugins}) sviluppate da vendor e consorzi indipendenti, tra cui: \textbf{Extreme Networks Plugin}, \textbf{Ruijie Networks Plugin}, \textbf{Juniper Networks Neutron Plugin}, \textbf{Calico Neutron Plugin} (orientato a un instradamento puramente Layer 3 basato su BGP per scenari ad altissima densità di container e microservizi) e il modulo \textbf{BNC Plugin}.

\subsection{Plugin Di Servizio E Servizi Di Rete Avanzati}

Mentre il plugin di commutazione core garantisce la connettività di base Layer 2, i requisiti applicativi aziendali impongono l'erogazione trasparente di servizi di rete a valore aggiunto. In Neutron, tali funzionalità sono formalizzate come \textbf{Service Plugins} (\emph{plugin di servizio}):
\begin{itemize}
    \item \textbf{Firewall-as-a-Service (FWaaS)}: fornisce un meccanismo di filtraggio delle connessioni perimetrali. Storicamente, l'implementazione di FWaaS risiede ed è eseguita all'interno dell'agente Layer 3 (\texttt{neutron-l3-agent}), intercettando il traffico prima che venga instradato verso le reti interne o le interfacce esterne.
    \item \textbf{VPN-as-a-Service (VPNaaS)}: deriva direttamente e amplia le funzionalità dell'agente L3, permettendo la configurazione di tunnel IPsec crittografati \emph{site-to-site} per collegare in sicurezza reti di progetto private a sedi aziendali remote o data center secondari.
    \item \textbf{Load Balancing-as-a-Service (LBaaS)}: consente la ripartizione del carico applicativo su insiemi eterogenei di istanze compute, implementato storicamente mediante stack proxy o daemon dedicati (come HAProxy).
    \item \textbf{Metering Agent}: monitora e quantifica il volume di traffico IP in transito per finalità di \emph{accounting}, \emph{auditing} e fatturazione delle risorse per tenant.
    \item \textbf{Trunk Plugin}: consente il passaggio di traffico appartenente a VLAN multiple attraverso una singola porta logica (\emph{trunk port}) associata a una VM, funzionalità critica per l'esecuzione di Network Functions Virtualization (NFV) o per l'orchestrazione interna di cluster di container annidati.
\end{itemize}

\subsection{Agenti Di Neutron: Architettura E Cooperazione Distribuita}

La scalabilità e la robustezza di Neutron sono assicurate da una netta suddivisione funzionale: il server di controllo (\textbf{Neutron Server}) espone l'interfaccia API REST centralizzata, autentica le chiamate, valida lo schema logico e memorizza lo stato delle risorse sul database relazionale persistente. Tuttavia, le modifiche effettive alla topologia di rete non vengono attuate dal server centrale, bensì distribuite ai singoli nodi di calcolo (\emph{Compute Node}) e ai nodi di rete dedicati (\emph{Network Node}) attraverso appositi daemon locali definiti \textbf{Agenti di Neutron} (\emph{Neutron Agents}).

Gli agenti comunicano con il \texttt{neutron-server} in modo asincrono sfruttando il \textbf{Message Bus} (basato sullo standard AMQP, ad esempio implementato da RabbitMQ). Gli agenti ascoltano le code e traducono i messaggi in comandi operativi sul kernel host, avvalendosi di tecnologie native quali \emph{network namespace}, tabelle di routing di Linux, regole \texttt{iptables}/\texttt{nftables}, interfacce veth e ponti di commutazione logici (Open vSwitch o bridge Linux).

Gli agenti architetturali primari includono:
\begin{enumerate}
    \item \textbf{Plugin Agent} (\texttt{neutron-*-agent}, es. \texttt{neutron-openvswitch-agent} o \texttt{neutron-linuxbridge-agent}): attivo su tutti i nodi di calcolo in cui risiedono macchine virtuali. Riceve le istruzioni relative a porte logiche e VLAN/VNI assegnate, instaura le interfacce virtuali (TAP device), applica le regole dei gruppi di sicurezza (\emph{Security Groups}) e coordina l'elaborazione dei frame di livello 2.
    \item \textbf{DHCP Agent} (\texttt{neutron-dhcp-agent}): preposto all'assegnazione automatica degli indirizzi IP, gateway, maschere di sottorete e record DNS alle istanze dei tenant. L'agente genera tipicamente processi \texttt{dnsmasq} isolati all'interno di specifici \emph{network namespace} (identificati dal prefisso \texttt{qdhcp-<network-id>}), garantendo che le allocazioni IP e i broadcast di sottoreti isolate e sovrapposte non interferiscano tra loro.
    \item \textbf{L3 Agent} (\texttt{neutron-l3-agent}): implementa la logica di instradamento di Livello 3 e la traduzione degli indirizzi di rete (\textbf{NAT}). È il componente responsabile della gestione dei router virtuali di OpenStack, permettendo la comunicazione tra sottoreti distinte appartenenti al medesimo tenant (\emph{inter-tenant communication}) e l'accesso verso la rete esterna e la rete Internet. L'agente esegue le funzioni di instradamento confinandole in namespace di rete distinti (\texttt{qrouter-<router-id>}).
\end{enumerate}

\newpage

\section{Modelli Di Configurazione Di Rete E Strategie Di Connettività}

La virtualizzazione di rete in OpenStack prevede l'astrazione di elementi convenzionali: \textbf{Networks} (domini di broadcast virtuali di livello 2), \textbf{Subnets} (intervalli di indirizzamento CIDR IPv4 o IPv6 associati a un blocco logico) e \textbf{Ports} (punti di connessione virtuali che associano parametri quali MAC address, IP address e vincoli di sicurezza a un'interfaccia di rete di una VM istanziata da Nova). OpenStack Compute (\textbf{Nova}) opera come consumatore primario di tali risorse.

La progettazione dell'infrastruttura di rete si suddivide in due macro-modelli: le \textbf{Provider Networks} e le \textbf{Self-Service Networks} (dette anche \emph{Project Networks}).

\begin{table}[H]
\centering
\small
\caption{Confronto architetturale tra Provider Networks e Self-Service Networks}
\label{tab:provider_vs_selfservice}
\begin{tabularx}{\textwidth}{l p{6.5cm} p{6.5cm}}
\toprule
\textbf{Caratteristica} & \textbf{Provider Networks} & \textbf{Self-Service (Project) Networks} \\
\midrule
\textbf{Privilegi di Gestione} & Creabili e modificabili unicamente dall'amministratore di sistema (\emph{admin}). & Autonomia di gestione concessa ai singoli tenant/progetti non privilegiati. \\
\addlinespace
\textbf{Mappatura Fisica} & Mappatura 1:1 diretta su segmenti L2 fisici del data center (tipicamente VLAN 802.1Q). & Reti virtuali sovrapposte (\emph{overlay}), incapsulate in tunnel (VXLAN, GRE, GENEVE). \\
\addlinespace
\textbf{Instradamento L3} & L2 diretto verso lo switch fisico dell'infrastruttura; routing L3 gestito all'esterno. & Richiedono router virtuali interni orchestrati da \texttt{neutron-l3-agent}. \\
\addlinespace
\textbf{Spazio d'Indirizzamento} & Condiviso con la rete fisica aziendale; vincolato a range fisici prestabiliti. & Indirizzi privati RFC 1918 arbitrari; possibilità di IP sovrapposti (\emph{overlapping}). \\
\addlinespace
\textbf{Vantaggi Operativi} & Elevata semplicità architetturale, massime prestazioni, basso overhead, affidabilità. & Massima flessibilità, automazione self-service, isolamento multi-tenant su larga scala. \\
\bottomrule
\end{tabularx}
\end{table}

\subsection{Provider Networks}

Le \textbf{Provider Networks} realizzano un'estensione diretta dell'infrastruttura fisica di rete esistente all'interno del data center fino all'interno delle istanze virtuali. Offrono connettività puramente di Livello 2, con supporto opzionale ai servizi locali di DHCP e al servizio di metadati (\emph{metadata service}).
\begin{itemize}[noitemsep]
    \item La segmentazione tra reti diverse è tradizionalmente ottenuta impiegando il tagging \textbf{VLAN IEEE 802.1Q}. Ogni rete logica corrisponde univocamente a un tag VLAN instradato sui trunk switch fisici.
    \item Tale approccio garantisce semplicità e massime prestazioni poiché elimina la complessità dell'incapsulamento logico aggiuntivo, riducendo l'\emph{overhead} sui pacchetti.
    \item Il rovescio della medaglia risiede nella ridotta flessibilità: la creazione di tali reti richiede l'intervento degli amministratori di cloud e di rete, dovendo coordinare le allocazioni VLAN e la configurazione degli apparati fisici sottostanti.
\end{itemize}

\subsection{Self-Service (Project) Networks}

Le \textbf{Self-Service Networks} consentono ai singoli tenant di definire topologie di rete arbitrarie in piena autonomia, senza richiedere privilegi amministrativi né riconfigurazioni dell'hardware fisico.
\begin{itemize}
    \item \textbf{Protocolli Overlay}: Poiché lo standard VLAN 802.1Q limita la scalabilità a soli 4096 identificativi di rete (dei quali solo circa 4000 effettivamente utilizzabili), le reti self-service sfruttano protocolli di incapsulamento logico \emph{overlay} su IP, quali \textbf{VXLAN} (\emph{Virtual Extensible LAN}) o \textbf{GRE} (\emph{Generic Routing Encapsulation}). Il protocollo VXLAN adotta un identificatore a 24 bit (\emph{VNI - VXLAN Network Identifier}), estendendo lo spazio di partizionamento a oltre 16 milioni di domini L2 distinti.
    \item \textbf{Instradamento e Attraversamento L3}: Le reti di progetto sono confinate logicamente e necessitano di router virtuali software per dialogare sia con le reti di altri tenant sia con le reti provider o la rete Internet pubblica. A differenza delle provider networks (dove le VM si attestano a livello L2 direttamente sulla rete fisica), i pacchetti generati all'interno di una rete self-service \textbf{devono necessariamente attraversare un agente Layer 3} (\texttt{neutron-l3-agent}) per uscire dal dominio di isolamento.
    \item \textbf{Gestione degli Indirizzi (IPv4 vs IPv6)}:
    \begin{itemize}[noitemsep]
        \item In \textbf{IPv4}, le reti di progetto impiegano quasi universalmente intervalli di indirizzi privati (conformi a RFC 1918, ad es. \texttt{10.0.0.0/24}). L'accesso in uscita verso la rete esterna/Internet viene mediato da un meccanismo di \textbf{Source NAT (SNAT)} operato dal router virtuale Neutron. Per consentire l'accesso bidirezionale o l'ingresso dall'esterno verso le VM interne, si adotta il paradigma dei \textbf{Floating IP}.
        \item In \textbf{IPv6}, in virtù dell'immensità dello spazio di indirizzamento globale, le reti self-service impiegano intervalli di indirizzi pubblici routabili. Di conseguenza, l'interazione con le reti esterne non richiede meccanismi di traduzione NAT, ma viene realizzata dai router virtuali mediante tabelle e politiche di \textbf{static routing}.
    \end{itemize}
\end{itemize}

\subsection{Meccanismo E Flusso Dei Floating IP}

Un'istanza istanziata all'interno di una rete self-service riceve un indirizzo IP privato (\emph{Fixed IP}) assegnato dal servizio DHCP. Tale indirizzo è valido unicamente all'interno della rete logica del tenant e, trovandosi a monte di un router NAT, non può essere instradato né indirizzato direttamente dall'esterno o da Internet.

Per rendere l'istanza accessibile pubblicamente, OpenStack adotta l'astrazione dei \textbf{Floating IP}. Un Floating IP è un indirizzo IP instradabile pubblicamente, allocato da un pool associato alla rete esterna (provider network), che viene associato dinamicamente alla porta logica privata di una specifica istanza.

% [IMMAGINE DEL FILE DI SLIDE CON NOME "10_OpenStack.pdf"] Schema architetturale di rete con Virtual Router, External Network, Floating IPs e Private Network con due istanze — Slide 115 / Pagina 8 del PDF

L'analisi del flusso di pacchetti evidenzia il funzionamento congiunto delle tecniche di NAT implementate dal router virtuale gestito dal \texttt{neutron-l3-agent}:
\begin{enumerate}
    \item \textbf{Traffico In Entrata (DNAT)}:
    I pacchetti provenienti dall'esterno e destinati all'istanza hanno come indirizzo di destinazione il Floating IP (ad esempio \texttt{172.24.4.2}). Quando il frame raggiunge l'interfaccia esterna del router SDN (la quale funge da terminazione ARP per gli indirizzi del blocco pubblico), il kernel applica una regola di \textbf{Destination NAT (DNAT)}:
    \[
        \text{IP Destinazione: } 172.24.4.2 \;\xrightarrow{\text{DNAT}}\; 10.0.0.2
    \]
    Il pacchetto, con l'header IP riscritto, viene quindi instradato sulla porta interna del router (\texttt{10.0.0.1}) e inoltrato attraverso la rete privata fino alla scheda di rete virtuale (\emph{vNIC}) dell'Istanza 1. In modo del tutto analogo, il traffico diretto all'indirizzo \texttt{172.24.4.3} subisce la traslazione verso l'IP privato \texttt{10.0.0.3} dell'Istanza 2.
    \item \textbf{Traffico In Uscita (SNAT)}:
    Quando l'Istanza invia pacchetti verso un host esterno, se essa possiede un Floating IP associato, il router virtuale provvede a sostituire l'indirizzo sorgente privato con il Floating IP corrispondente mediante \textbf{Source NAT (SNAT)} 1-a-1:
    \[
        \text{IP Sorgente: } 10.0.0.2 \;\xrightarrow{\text{SNAT}}\; 172.24.4.2
    \]
    Qualora l'istanza non disponga di un Floating IP, il router applicherà un meccanismo di mascheramento NAPT (\emph{PAT/Masquerade}) condividendo l'IP assegnato al gateway esterno del router virtuale (\texttt{172.24.4.1}).
\end{enumerate}

Le proprietà architetturali fondamentali dei Floating IP comprendono:
\begin{itemize}[noitemsep]
    \item \textbf{Coesistenza sull'interfaccia}: La scheda di rete virtuale (NIC) dell'istanza non ha contezza diretta del Floating IP; la VM continua ad avere configurato localmente il solo indirizzo IP privato (\emph{Fixed IP}). L'associazione biunivoca tra Fixed IP e Floating IP è gestita integralmente dal router SDN virtuale.
    \item \textbf{Assegnazione manuale e disaccoppiata}: I Floating IP non vengono allocati o agganciati automaticamente all'avvio dell'istanza; l'utente deve allocarli esplicitamente dal pool esterno e legarli alla porta di rete dell'istanza (\emph{association}). Possono inoltre essere sganciati e riassociati a caldo (\emph{hot-swapping}) a un'altra VM per scenari di \emph{failover} e alta affidabilità.
\end{itemize}

\newpage

\section{Architettura Dello Storage In OpenStack}

La persistenza e la gestione dei dati in OpenStack si articolano su tre modelli complementari, ciascuno indirizzato a soddisfare vincoli distinti in termini di latenza, ciclo di vita, semantica d'accesso e scalabilità:

\begin{figure}[H]
\centering
\begin{tikzpicture}[
    box/.style={rectangle, draw=theoryblue, fill=theoryblue!10, rounded corners, thick, text width=4.5cm, minimum height=3.5cm, align=center, inner sep=6pt},
    header/.style={font=\bfseries\color{theoryblue}, depth=0pt, text height=1.5ex}
]
\node[box] (ephemeral) at (0,0) {
    \textbf{\large Ephemeral Storage}\\[2mm]
    \small
    \textbf{Gestito da:} Nova\\
    \textbf{Ciclo di vita:} Legato alla VM\\
    \textbf{File System:} Locale / Root Disk\\
    \textbf{Uso tipico:} OS Image, swap, file temporanei di computazione
};

\node[box] (block) at (5.5,0) {
    \textbf{\large Block Storage}\\[2mm]
    \small
    \textbf{Gestito da:} Cinder\\
    \textbf{Ciclo di vita:} Persistente (indipendente)\\
    \textbf{File System:} Formattabile dal guest\\
    \textbf{Uso tipico:} Database, volumi dati, partizioni raw persistenti
};

\node[box] (object) at (11,0) {
    \textbf{\large Object Storage}\\[2mm]
    \small
    \textbf{Gestito da:} Swift\\
    \textbf{Ciclo di vita:} Altamente duraturo\\
    \textbf{Accesso:} API REST (HTTP GET/PUT)\\
    \textbf{Uso tipico:} File non strutturati, immagini VM (Glance), backup
};
\end{tikzpicture}
\caption{Tassonomia delle soluzioni di archiviazione dati in OpenStack}
\label{fig:openstack_storage_models}
\end{figure}

\begin{enumerate}
    \item \textbf{Ephemeral Storage} (Storage Effimero):
    Rappresenta lo spazio disco primario assegnato a ciascuna VM per alloggiare il sistema operativo dell'immagine avviata e i dati temporanei di lavoro. La caratteristica dirimente risiede nel fatto che \textbf{il ciclo di vita dello storage coincide con quello dell'istanza}: qualora la macchina virtuale venga distrutta o terminata, tutti i dati residenti sullo storage effimero vengono eliminati irrevocabilmente.
    \item \textbf{Block Storage} (Cinder):
    Fornisce volumi a blocchi persistenti collegabili a caldo (\emph{attach/detach}) come dischi secondari (o come dischi di boot persistenti) alle macchine virtuali gestite da Nova. Il volume a blocchi viene presentato al sistema operativo guest come un dispositivo a blocchi raw (es. \texttt{/dev/vdb}) su cui la VM può inizializzare una tabella delle partizioni e creare un proprio file system locale (es. \texttt{ext4}, \texttt{xfs}). I dati persistono alla distruzione della VM.
    \item \textbf{Object Storage} (Swift):
    Sistema di archiviazione distribuito e massivamente scalabile concepito per memorizzare dati non strutturati sotto forma di oggetti binari immutabili. L'accesso prescinde completamente da logiche di partizionamento, montaggio di dischi o gerarchie POSIX: le applicazioni interagiscono con Swift direttamente tramite chiamate HTTP RESTful (\texttt{GET}, \texttt{PUT}, \texttt{DELETE}, \texttt{HEAD}). Si presta in modo ideale alla conservazione di media, documenti, archivi di log e immagini di Glance.
\end{enumerate}

\newpage

\section{OpenStack Cinder: Il Servizio Di Block Storage}

OpenStack \textbf{Cinder} astrae e virtualizza la gestione dei dispositivi fisici di memorizzazione a blocchi. L'obiettivo primario di Cinder è fornire all'utente un'interfaccia API \emph{self-service} per allocare, dimensionare, creare istantanee (\emph{snapshot}) e collegare volumi di storage senza dover conoscere i dettagli implementativi del backend sottostante né la topologia dell'infrastruttura di storage fisica.

\subsection{Architettura E Componenti Del Control Plane Di Cinder}

L'architettura di Cinder rispetta il paradigma comune a tutti i servizi OpenStack: netta separazione tra i moduli di interfaccia (\emph{API}), il coordinatore decisionale (\emph{Scheduler}) e i moduli esecutivi di backend (\emph{Volume Drivers}), il tutto integrato per mezzo di un broker di messaggistica (\emph{Message Bus}) e di un database relazionale di stato.

% [IMMAGINE DEL FILE DI SLIDE CON NOME "10_OpenStack.pdf"] Diagramma architetturale di OpenStack Cinder: Client REST, Cinder API, Cinder Scheduler, Message Bus, Cinder Volume e Volume Drivers verso Storage Backend A/B/C — Slide 118 / Pagina 11 del PDF

La cooperazione tra i componenti interni di Cinder si sviluppa come segue:
\begin{enumerate}
    \item \textbf{Client / SDK / Horizon}: L'utente o i moduli interni (come Nova durante il provisioning) avviano una richiesta di allocazione volume inviando una chiamata HTTP conforme alle specifiche Cinder REST API.
    \item \textbf{cinder-api}: È il demone endpoint che riceve, valida e autentica (tramite Keystone) le richieste REST in ingresso. Non esegue operazioni di storage dirette; inserisce la richiesta sul message bus indirizzandola allo scheduler o ai worker dei volumi.
    \item \textbf{cinder-scheduler daemon}: Valuta i requisiti della richiesta (ad esempio capacità richiesta, tipo di volume, zona di disponibilità, parametri di qualità del servizio) e interroga lo stato delle risorse disponibili. Il suo compito è \textbf{selezionare il nodo o il pool di storage ottimale} sul quale istanziare il volume.
    \begin{itemize}[noitemsep]
        \item Cinder supporta algoritmi di schedulazione elementari (es. \emph{Round-Robin}), ma adotta di default il sofisticato \textbf{Filter Scheduler}.
        \item Il \textbf{Filter Scheduler} applica una serie di filtri sequenziali escludendo i nodi non idonei (in base a vincoli quali capacità residua, \emph{Availability Zone}, compatibilità del \emph{Volume Type} e capacità architetturali esportate dal backend), ordinando poi i candidati rimanenti tramite funzioni di pesatura (\emph{weighing}).
    \end{itemize}
    \item \textbf{Message Bus}: Canale di comunicazione asincrono e disaccoppiato su cui transitano i comandi RPC (Remote Procedure Call) che collegano le API, lo scheduler e i nodi esecutivi di volume.
    \item \textbf{cinder-volume}: Demone esecutivo che controlla e manipola direttamente i dispositivi di backend. Risiede sui nodi che hanno connettività diretta con l'infrastruttura di archiviazione (siano essi nodi controller, compute o nodi di storage dedicati). Interagisce con le tecnologie fisiche attraverso una vasta collezione di \textbf{Volume Drivers}.
    \item \textbf{cinder-backup daemon} (opzionale): Servizio dedicato alla creazione di copie di backup complete o incrementali di volumi Cinder verso storage target esterni (quali cluster OpenStack Swift, sistemi Ceph o NFS).
\end{enumerate}

\subsection{Backend Di Memorizzazione E Modello Di Accesso Multi-Tenant}

Il design di Cinder supporta ambienti con \textbf{storage backends eterogenei}: all'interno del medesimo deployment cloud, differenti classi di volume (es. \emph{Gold}, basato su array SAN ad altissime prestazioni; \emph{Silver}, basato su cluster distribuiti Ceph; \emph{Bronze}, basato su array locali LVM) possono coesistere ed essere selezionate dinamicamente tramite metadati \texttt{volume-type}.

\begin{itemize}
    \item \textbf{Default Backend}: L'implementazione di riferimento predefinita nei laboratori e nelle installazioni base si basa su \textbf{LVM} (\emph{Logical Volume Manager}) configurato su un volume group locale denominato \texttt{cinder-volumes}. Cinder gestisce l'allocazione creando volumi logici LVM ed esportandoli all'hypervisor (es. KVM/QEMU) tramite protocolli di trasporto a blocchi di rete come \textbf{iSCSI}.
    \item \textbf{Dispositivi Enterprise e Software-Defined}: Cinder integra driver per storage appliance commerciali e sistemi RAID esterni (NetApp, Dell EMC, HPE) e tecnologie SDS ampiamente diffuse quali \textbf{Ceph RBD} (\emph{Rados Block Device}). Con KVM/QEMU come hypervisor, tali backend permettono la configurazione di dimensioni dei blocchi (\emph{block sizes}) personalizzate per ottimizzare le prestazioni I/O.
    \item \textbf{Controllo degli Accessi (RBAC)}: Cinder garantisce l'isolamento multi-tenant attraverso un sistema di autorizzazioni basato su ruoli (\textbf{RBAC}). Nella configurazione predefinita, gran parte delle operazioni di creazione, attach e snapshot di volumi sono accessibili a qualunque utente autenticato all'interno del proprio progetto. Politiche restrittive più granulari (es. quote di spazio massimo per tenant o permessi di visibilità dei volumi condivisi) sono configurabili dagli amministratori modificando il file di policy di Cinder (\texttt{policy.yaml} o \texttt{policy.json}).
\end{itemize}

\newpage

\section{OpenStack Swift: Il Servizio Di Object Storage}

\textbf{OpenStack Swift} è un sistema di archiviazione a oggetti distribuito, ridondato e privo di punti singoli di fallimento (\emph{no single point of failure}). Il suo obiettivo architetturale è garantire la massima disponibilità dei dati e scalabilità orizzontale (\emph{scale-out}), consentendo l'archiviazione a lungo termine di moli di dati dell'ordine dei petabyte o degli exabyte con costi infrastrutturali ridotti, basandosi su hardware di tipo \emph{commodity}.

A differenza dei sistemi a blocchi o dei tradizionali file system gerarchici, in Swift non esistono cartelle nidificate a livello di kernel né tabelle di allocazione dei blocchi; i file vengono salvati come entità binarie complete, identificate univocamente e arricchite da metadati.

\subsection{Architettura Dei Server Di Swift E Gerarchia Dello Storage}

L'architettura logica di Swift articola l'intero spazio di memorizzazione in tre livelli di entità nidificate:
\begin{enumerate}[noitemsep]
    \item \textbf{Account}: Rappresenta il namespace radice, associato a un'area di storage di progetto/utente. \emph{Nota terminologica}: in Swift, il termine ``account'' definisce un'area di memorizzazione logica (\emph{storage area}) e non l'identità autenticata dell'utente (la quale fa capo a Keystone).
    \item \textbf{Container}: È l'analogo dei \emph{bucket} di Amazon S3; funge da contenitore logico (namespace) all'interno di un account per raggruppare un insieme di oggetti. Non possono esistere container annidati dentro altri container.
    \item \textbf{Object}: È il file binario effettivo caricato dall'utente, corredato da metadati descrittivi.
\end{enumerate}

Le richieste esterne e la persistenza interna sono delegate a due macro-famiglie di nodi: i \textbf{Proxy Server} e gli \textbf{Storage Server}.

% [IMMAGINE DEL FILE DI SLIDE CON NOME "10_OpenStack.pdf"] Schema ad albero di Swift: swift-proxy connesso a account, container, object server e rispettivi DB/store — Slide 121 / Pagina 15 del PDF

\subsubsection{Swift Proxy Server}
Il \textbf{Proxy Server} (\texttt{swift-proxy}) è l'unico componente del cluster che dialoga direttamente con i client esterni. Esso implementa l'interfaccia HTTP RESTful di Swift. Per ogni richiesta in transito:
\begin{itemize}[noitemsep]
    \item Autentica il token Keystone o le credenziali fornite.
    \item Ispeziona gli hash logici e consulta l'anello di mappatura (\textbf{Ring}) per determinare quali nodi fisici di backend ospitano le repliche del dato richiesto.
    \item Coordina le chiamate verso gli storage server di backend, raccogliendo le risposte e restituendo il codice di stato HTTP e l'eventuale payload al client.
\end{itemize}

\subsubsection{Storage Server Di Backend}
I server di memorizzazione erogano a loro volta un'API REST interna e si articolano in tre processi distinti a seconda dello strato logico di competenza:
\begin{itemize}
    \item \textbf{Account Server}: Gestisce i metadati relativi ai singoli account e mantiene l'elenco dei container associati a ciascuno di essi. Tali informazioni non sono archiviate come file di testo, bensì memorizzate in leggeri database \textbf{SQLite} residenti su disco.
    \item \textbf{Container Server}: Tiene traccia dei metadati del container (es. parametri di auto-scadenza \texttt{X-Container-Meta-TTL}) e mantiene l'indice degli oggetti in esso contenuti. \`E cruciale notare che \textbf{il container server non memorizza il dato binario dell'oggetto né la sua posizione fisica sui dischi}, ma funge esclusivamente da indice logico. Anche in questo caso le voci sono memorizzate mediante database \textbf{SQLite}.
    \item \textbf{Object Server}: Ha la responsabilità della persistenza fisica dei dati binari sui dischi dei nodi di storage. Gli oggetti vengono salvati come file binari grezzi all'interno del file system host (tipicamente formattato con supporto esteso ad \texttt{xfs}). Il percorso su disco è ricavato combinando l'identificativo di partizione dell'anello con il \emph{timestamp} dell'operazione (fondamentale per tracciare il versioning o sovrascritture). I metadati associati al dato (sia di sistema, come il \texttt{Content-Type}, sia definiti dall'utente) sono memorizzati direttamente negli \textbf{attributi estesi del file system (\emph{xattrs})}: ciò assicura che il dato binario e i suoi descrittori viaggino e vengano duplicati sempre come un blocco atomico.
\end{itemize}

\subsection{Anatomia Di Una Richiesta RESTful In Swift}

Ogni interazione con Swift rispetta i crismi del protocollo HTTP su API RESTful e si compone di quattro sezioni fondamentali:
\begin{enumerate}
    \item \textbf{Metodo/Verbo HTTP}: Specifica l'azione desiderata:
    \begin{itemize}[noitemsep]
        \item \texttt{PUT}: caricamento o creazione di un nuovo account, container o oggetto.
        \item \texttt{GET}: lettura e download del payload o recupero delle liste/elenchi.
        \item \texttt{DELETE}: eliminazione logico-fisica della risorsa.
        \item \texttt{HEAD}: interrogazione dei soli metadati presenti negli header HTTP.
    \end{itemize}
    \item \textbf{Informazioni di Autenticazione}: Header contenente il token di sessione autenticato rilasciato da Keystone (es. \texttt{X-Auth-Token}).
    \item \textbf{Storage URL}: L'indirizzo completo della risorsa, strutturato come segue:
    \[
        \underbrace{\texttt{https://swift.example.com/v1/}}_{\text{Cluster Location (Endpoint REST)}} \quad \underbrace{\texttt{/account/container/object}}_{\text{Storage Location}}
    \]
    La porzione di \emph{Storage Location} modella gerarchicamente la profondità della richiesta:
    \begin{itemize}[noitemsep]
        \item \texttt{/account}: opera sull'account, restituendo la lista dei container o i metadati di storage globale.
        \item \texttt{/account/container}: opera sul container, permettendo l'interrogazione dell'elenco degli oggetti contenuti.
        \item \texttt{/account/container/object}: indirizza l'oggetto binario specifico per il trasferimento o la cancellazione.
    \end{itemize}
    \item \textbf{Corpo/Payload Opzionale}: Dati binari o header personalizzati inviati a supporto della richiesta di scrittura.
\end{enumerate}

\subsection{Gestione Della Consistenza Nel Teorema CAP}

Nel campo dei sistemi distribuiti, il \textbf{Teorema CAP} (formulato da Eric Brewer) stabilisce l'impossibilità teorica per un archivio dati distribuito di soddisfare contemporaneamente le tre seguenti proprietà:
\begin{itemize}
    \item \textbf{Consistency (C)}: Ogni lettura riceve la scrittura più recente o, in alternativa, solleva un errore bloccante.
    \item \textbf{Availability (A)}: Ogni richiesta non-errata riceve una risposta priva di errore, indipendentemente dal fatto che contenga l'aggiornamento più recente.
    \item \textbf{Partition Tolerance (P)}: Il sistema continua a operare e soddisfare le richieste anche a fronte di partizionamenti di rete o guasti di comunicazione arbitrari tra i nodi.
\end{itemize}

\begin{figure}[H]
\centering
\begin{tikzpicture}
    \draw[thick, fill=blue!10] (-2,0) circle (1.8cm);
    \draw[thick, fill=red!10] (2,0) circle (1.8cm);
    \draw[thick, fill=green!10] (0,2.5) circle (1.8cm);
    \node at (-2,0) {\textbf{Consistency (C)}};
    \node at (2,0) {\textbf{Partition Tol. (P)}};
    \node at (0,2.5) {\textbf{Availability (A)}};
    \node[draw, fill=white, rounded corners, thick] at (1.2,1.2) {\textbf{Swift (AP)}};
\end{tikzpicture}
\caption{Collocazione di OpenStack Swift nel quadro del Teorema CAP come sistema AP}
\label{fig:cap_swift}
\end{figure}

Data l'inevitabilità di guasti fisici su scala estesa (condizione per cui la \emph{Partition Tolerance} è un vincolo non negoziabile nell'ingegneria del cloud), \textbf{OpenStack Swift opta esplicitamente per la coppia Availability e Partition Tolerance (sistema AP), sacrificando la consistenza forte immediata}.

Di conseguenza, Swift persegue un modello di \textbf{Eventual Consistency} (consistenza a tendere). Il sistema garantisce che una richiesta HTTP riceva sempre risposta positiva e che i dati siano dispersi ridondantemente; tuttavia, in condizioni anomale (quali sovraccarico improvviso di nodi, partizionamento del cluster o failure di rete), un client potrebbe temporaneamente leggere una versione precedente di un oggetto prima che il disallineamento venga risanato.

\subsection{Servizi Di Background Per La Consistenza: Auditor, Updater E Replicator}

Per ridurre al minimo la finestra temporale di inconsistenza (\emph{window of inconsistency}) ed eliminare anomalie dovute a guasti hardware o corruzione silente (\emph{bit rot}), Swift esegue in background tre categorie fondamentali di demoni di consistenza:

\begin{enumerate}
    \item \textbf{Auditor}:
    Esegue una scansione ciclica e continua dei dischi fisici locali, verificando l'integrità strutturale dei database di account/container e ricalcolando il checksum crittografico (MD5) degli oggetti per confrontarlo con i metadati memorizzati.
    \begin{itemize}[noitemsep]
        \item Se rileva un oggetto corrotto o alterato a livello di bit, l'Auditor lo sposta tempestivamente in una directory di \textbf{quarantena} (\emph{quarantine area}) per impedirne l'erogazione ai client.
        \item Il recupero della versione integra viene poi demandato automaticamente al \textbf{Replicator}.
    \end{itemize}
    \item \textbf{Updater}:
    Interviene per recuperare scritture asincrone non andate a buon fine. In condizioni normali, durante l'upload di un nuovo oggetto (\texttt{PUT}), l'Object Server tenta di aggiornare contestualmente l'indice sul Container Server. Qualora il Container Server sia irraggiungibile o sovraccarico, l'operazione di aggiornamento fallisce ma l'oggetto viene comunque salvato.
    \begin{itemize}[noitemsep]
        \item L'\textbf{Object Updater} accoda la modifica pendente e scansiona periodicamente i dischi per inviare l'aggiornamento al Container Server non appena torna raggiungibile.
        \item Analogamente, il \textbf{Container Updater} colma le discrepanze verso l'Account Server aggiornando il conteggio aggregato di container e spazio consumato.
        \item \emph{Dettaglio implementativo}: Esistono solo Container Updater e Object Updater; \textbf{non esiste un Account Updater}, in quanto non sussiste un'entità gerarchica superiore a cui notificare variazioni.
    \end{itemize}
    \item \textbf{Replicator}:
    Assicura che nel cluster persista costantemente il corretto numero di copie repliche e che esse siano collocate esattamente sui nodi designati dall'anello (\emph{Ring}).
    \begin{itemize}[noitemsep]
        \item Ripara lo stato dei dati trasferendo copie conformi in caso di guasti a dischi o nodi.
        \item Gestisce la propagazione delle cancellazioni (tramite \emph{tombstone marker}).
        \item Svolge le migrazioni fisiche dei dati durante le fasi di ribilanciamento (\emph{rebalance}) del Ring.
        \item Si divide internamente in \textbf{Database Replicator} (per i DB SQLite di account e container, basato su sincronizzazione via \texttt{rsync}) e \textbf{Object Replicator} (per i payload binari dei file).
    \end{itemize}
\end{enumerate}

\newpage

\section{Consistent Hashing E Meccanismo Del Ring}

Il cuore algoritmico su cui si fondano la dispersione e l'indicizzazione dei dati in Swift è il \textbf{Consistent Hashing} (hashing coerente). In un'architettura distribuita tradizionale, mappare un oggetto su un insieme di $M$ nodi di storage tramite la formula aritmetica convenzionale:
\[
    \text{Nodo} = \text{hash}(\text{chiave}) \pmod M
\]
presenta un limite critico: se un nodo si guasta o viene aggiunto un nuovo nodo (variando $M$ in $M+1$), quasi la totalità degli hash ricalcolati cambia valore, provocando un massiccio spostamento forzato di dati tra tutti i server del cluster con conseguente saturazione della rete.

Per risolvere questo problema, Swift modella l'intero spazio di calcolo degli hash (tipicamente su valori a 128 o 160 bit generati tramite hash crittografico MD5) come uno spazio circolare chiuso continuo, denominato \textbf{Ring} (Anello).

% [IMMAGINE DEL FILE DI SLIDE CON NOME "10_OpenStack.pdf"] Diagramma del Consistent Hashing Ring con partizioni, intervallo hash circolare da 0 a 2^160 e assegnazione nodi — Slide 129 / Pagina 25 del PDF

\subsection{Partizioni, Potenza Di Partizione Ed Esempio Di Ribilanciamento}

Per consentire l'espansione e la redistribuzione efficiente dei dati, lo spazio continuo dell'anello viene suddiviso a priori in un numero discreto e fisso di blocchi logici equivalenti chiamati \textbf{Partizioni} (\emph{Partitions}):
\begin{itemize}
    \item Il numero totale di partizioni è una potenza intera di due, fissata in fase di configurazione tramite la \textbf{Partition Power} ($P$):
    \[
        \text{Totale Partizioni} = 2^P
    \]
    (Ad esempio, per $P = 16$, il ring comprende $2^{16} = 65{,}536$ partizioni).
    \item Ogni partizione è mappata staticamente su dispositivi fisici distinti appartenenti a specifici nodi. La partizione costituisce l'unità logica minima di spostamento durante il ribilanciamento.
\end{itemize}

\paragraph{Esempio di Ribilanciamento Analitico:}
Si ipotizzi un'infrastruttura con uno spazio logico diviso in $2^{10} = 1024$ partizioni, distribuite inizialmente su un cluster di $N = 4$ nodi omogenei.
La quota di partizioni assegnata a ciascun singolo nodo corrisponde a:
\[
    \text{Partizioni per nodo} = \frac{2^{10}}{4} = \frac{1024}{4} = 256
\]
Nel momento in cui si decide di scalare l'infrastruttura introducendo un 5° nodo ($N=5$), non si ricalcolano gli hash di tutti gli oggetti archiviati: l'algoritmo effettua una ricalibrazione (\emph{rebalance}) prelevando una frazione bilanciata pari a $\frac{1}{5}$ delle partizioni da ciascuno dei nodi originari per migrarle verso il nuovo nodo entrante. Al termine, la distribuzione risulterà uniforme con:
\[
    \text{Nuova quota per nodo} = \frac{2^{10}}{5} = \frac{1024}{5} \approx 204 \text{ partizioni per nodo}
\]
Ciascun dispositivo può essere inoltre dotato di un \textbf{peso} (\emph{weight}) arbitrario per assegnare proporzionalmente più partizioni ai server equipaggiati con dischi a maggiore capienza. Il limite massimo teorico di scalabilità garantito da una partizione ad esponente $P = 10$ corrisponde a supportare fino a $2^{10}$ nodi di memorizzazione fisici indipendenti.

\subsection{Replicazione, Domini Di Guasto E Zone Di Isolamento}

Per assicurare la persistenza dei dati e la tolleranza ai guasti senza ricorrere a RAID hardware centralizzati, Swift introduce la duplicazione geografica e logica:
\begin{itemize}
    \item \textbf{Fattore di Replica ($R$)}: Parametro configurato in \texttt{swift.conf} (tramite la direttiva \texttt{default\_replicas}, convenzionalmente pari a $R = 3$). Il parametro $R$ copre un duplice ruolo logico:
    \begin{enumerate}
        \item Ogni oggetto caricato viene duplicato esattamente in $R$ copie conformi distinte.
        \item Ogni partizione logica dell'anello possiede $R$ repliche fisiche residenti su dispositivi distinti.
    \end{enumerate}
    \item \textbf{Failure Domains e Zone}: I dati replicati non devono risiedere su apparati soggetti al medesimo evento di guasto. Swift introduce l'astrazione di \textbf{Zone}: una zona è definita manualmente dall'amministratore e modella un dominio di guasto isolato (un singolo disco, un server fisico, un gruppo di continuità comune, o un intero rack/armadio di rete).
    \item \textbf{Vincolo di Allocazione}: L'algoritmo di composizione del Ring impone rigidamente che \textbf{le $R$ repliche di una medesima partizione debbano essere posizionate in $R$ zone completamente distinte}.
\end{itemize}

\paragraph{Calcolo Analitico delle Partizioni Medie per Dispositivo:}
Dati i seguenti parametri architetturali del cluster:
\begin{itemize}[noitemsep]
    \item $P$: Partition Power (che definisce il totale delle partizioni uniche $2^P$);
    \item $R$: Numero di repliche per partizione;
    \item $N$: Numero di nodi di storage attivi nel cluster;
    \item $D$: Numero medio di dispositivi (dischi fisici) equipaggiati per nodo;
\end{itemize}
il numero complessivo di partizioni duplicate da memorizzare sull'intero parco dischi equivale a $2^P \cdot R$, mentre il numero totale di drive presenti nel sistema è pari a $N \cdot D$. Ne consegue che il numero medio di repliche di partizione residenti su ciascun singolo disco fisico è espresso analiticamente da:
\begin{equation}
    \boxed{\text{Partizioni Medie per Disco} = \frac{2^P \cdot R}{N \cdot D}}
\end{equation}

\subsection{Architettura Multi-Ring}

Poiché i requisiti operativi, i carichi di I/O e le tecnologie di archiviazione differiscono drasticamente a seconda che si gestiscano metadati o file binari, Swift non impiega un unico anello globale:
\begin{itemize}[noitemsep]
    \item Vengono istanziati e mantenuti \textbf{tre anelli operativi distinti e indipendenti}:
    \begin{enumerate}
        \item \textbf{Account Ring}: governa la dislocazione delle informazioni relative agli account;
        \item \textbf{Container Ring}: governa l'indicizzazione e la collocazione dei container;
        \item \textbf{Object Ring}: governa il posizionamento dei file binari veri e propri.
    \end{enumerate}
    \item I primi due strati (Account e Container) memorizzano lo stato in forma di leggeri database relazionali \textbf{SQLite}, con forte enfasi su frequenti aggiornamenti atomici dei metadati. Il terzo strato (Object) opera direttamente memorizzando \textbf{file binari grezzi} sul file system locale (\texttt{xfs}).
    \item La separazione logica permette l'isolamento fisico: i tre ring possono essere ospitati e serviti da insiemi hardware completamente differenti e dedicati (ad esempio, posizionando i database degli account e container su dischi SSD ad altissima velocità di IOPS, relegando i payload dell'Object Ring su cluster composti da dischi magnetici convenzionali ad alta capacità).
\end{itemize}

\newpage

\section{Interfacce Di Gestione E Autenticazione In OpenStack}

L'ecosistema di OpenStack è concepito secondo una filosofia interamente \emph{API-First}: ogni azione, dalla creazione di un volume alla configurazione di un'interfaccia di rete, è formalizzata come una chiamata stateless al servizio RESTful competente.

\subsection{Modalità Di Accesso Ai Servizi OpenStack}

Un operatore o un'applicazione cloud-native può interagire con le REST API di OpenStack attraverso cinque modalità primarie:
\begin{enumerate}
    \item \textbf{OpenStack Dashboard (Horizon)}: Interfaccia grafica web-based orientata all'amministrazione visuale e al provisioning rapido per utenti non specialisti.
    \item \textbf{Client REST Generici}: Qualsiasi client in grado di emettere richieste HTTP/HTTPS (ad es. comandi CLI via \texttt{cURL}, strumenti come Postman o librerie software integrate) inviando payload JSON formattati.
    \item \textbf{Librerie e SDK Ufficiali (OpenStack Python SDK)}: Moduli software che incapsulano le chiamate REST all'interno di classi, metodi ed eccezioni native (ad es. \texttt{keystoneauth1}, \texttt{openstacksdk}).
    \item \textbf{Command-Line Clients (OpenStack CLI)}: Il tool unificato a riga di comando \texttt{openstack}, che astrae le chiamate API dietro una sintassi coerente e scriptabile per amministratori di sistema.
    \item \textbf{HOT Scripts (Heat Orchestration Template)}: File dichiarativi scritti in formato YAML interpretati da \textbf{OpenStack Heat} per l'orchestrazione automatizzata dell'infrastruttura secondo il paradigma \emph{Infrastructure-as-Code} (IaC).
\end{enumerate}

\subsection{Autenticazione Tramite OpenStack Keystone: Esempio cURL}

Il servizio di identità \textbf{Keystone} opera da snodo centrale per l'autenticazione, la validazione dei tenant e il rilascio dei token di accesso necessari per dialogare con gli altri servizi dell'infrastruttura.

L'interazione a basso livello tramite \texttt{cURL} per l'ottenimento di un token di autenticazione v3 prevede l'invio di una richiesta HTTP \texttt{POST} verso l'endpoint delle API di Keystone (convenzionalmente in ascolto sulla porta TCP 5000):

\begin{lstlisting}[language=bash]
# Esecuzione della chiamata di autenticazione verso Keystone v3
curl -i -H "Content-Type: application/json" \
  -d @my_credentials \
  http://localhost:5000/v3/auth/tokens
\end{lstlisting}

Il payload JSON, racchiuso all'interno del file locale descrittore \texttt{my\_credentials}, esplicita il metodo di convalida (password) e il contesto del dominio e dell'utente:

\begin{lstlisting}
{
  "auth": {
    "identity": {
      "methods": [
        "password"
      ],
      "password": {
        "user": {
          "name": "admin",
          "domain": {
            "id": "default"
          },
          "password": "mypassword"
        }
      }
    }
  }
}
\end{lstlisting}

Nel modello di Keystone v3, in risposta a una richiesta valida con codice \texttt{HTTP/1.1 201 Created}, il token di autorizzazione generato non viene restituito all'interno del body JSON, bensì all'interno dell'header di risposta HTTP \texttt{X-Subject-Token}. Tale token dovrà essere successivamente incluso in ogni successiva invocazione REST verso Nova, Neutron, Cinder o Swift valorizzando l'header \texttt{X-Auth-Token}.

\subsection{Autenticazione Programmatica Con Python SDK: Gestione Delle Sessioni}

Nello sviluppo applicativo, l'emissione manuale di chiamate cURL e la gestione diretta del ciclo di vita dei token risultano inefficienti e proni a errori di scadenza temporale. L'approccio raccomandato da OpenStack consiste nell'uso delle \textbf{Sessioni} fornite dalla libreria \texttt{keystoneauth1}.

L'oggetto \textbf{Session} astrae e gestisce in modo trasparente l'intero ciclo di vita dell'autenticazione:
\begin{itemize}[noitemsep]
    \item Memorizza e incapsula le credenziali di accesso fornite dal plugin di autenticazione.
    \item Riconosce quando un token di autenticazione è scaduto, richiedendo automaticamente un nuovo token a Keystone senza interruzioni operative per l'applicazione chiamante.
    \item Permette a client di servizi differenti (Nova, Glance, Neutron, Cinder) di riutilizzare la medesima connessione logica TCP/TLS autenticata.
\end{itemize}

Lo snippet Python seguente illustra il ciclo completo di autenticazione tramite la creazione di un plugin di credenziali \texttt{v3.Password}, l'instanziazione di una \texttt{Session} e l'aggancio del client Keystone:

\begin{lstlisting}[language=python]
from keystoneauth1.identity import v3
from keystoneauth1 import session
from keystoneclient.v3 import client

# Definizione del plugin di autenticazione basato su credenziali utente e dominio
auth = v3.Password(
    user_domain_name='default',
    username='admin',
    password='mypassword',
    project_domain_name='default',
    project_name='admin',
    auth_url='http://localhost:5000/v3'
)

# Creazione dell'oggetto Session per la gestione trasparente e persistente dei token
mysess = session.Session(auth=auth)

# Istanziazione del client Keystone associato alla sessione autenticata
keystone = client.Client(session=mysess)
\end{lstlisting}

Il modello a plugin di autenticazione (\emph{authentication plugin}) disaccoppia il codice applicativo dal metodo concreto di verifica dell'identità: qualora l'infrastruttura aziendale decida di sostituire la convalida via password con certificati X.509, SAML o OpenID Connect federato, sarà sufficiente istanziare il relativo plugin di autenticazione in \texttt{keystoneauth1} lasciando invariato il resto del codice che opera sulla sessione.

% CC -------------------------------------------------------------------------------------------------------------------

\section{La Convergenza Tra Intelligenza Artificiale E Cloud Computing}

Il paradigma del \emph{Cloud Computing} e i modelli di \textbf{Intelligenza Artificiale} (AI) rappresentano due pilastri tecnologici la cui intersezione genera una profonda sinergia bidirezionale. Da un lato, il cloud computing fornisce l'infrastruttura computazionale elastica, le capacità di archiviazione scalabili su vasta scala e i framework distribuiti indispensabili per addestrare ed eseguire modelli di \emph{Machine Learning} (ML) e \emph{Deep Learning} (DL) caratterizzati da miliardi di parametri. Dall'altro lato, le metodologie dell'intelligenza artificiale vengono integrate nativamente all'interno dei meccanismi operativi del cloud, automatizzando e ottimizzando funzioni architetturali critiche: dal provisioning dinamico delle risorse all'\emph{autoscaling}, dal bilanciamento energetico alla rilevazione tempestiva di violazioni di sicurezza.

% [IMMAGINE DEL FILE DI SLIDE CON NOME "CloudComputing_4"] Frontespizio del corso: "AI in Cloud Computing" a cura di Gianluca Reali — Slide 1 / Pagina 1 del PDF

Tale unificazione trasforma il cloud da mero fornitore passivo di risorse virtualizzate (\emph{compute}, \emph{storage}, \emph{networking}) a una piattaforma attiva, adattiva e intelligente. Nelle architetture moderne, l'elaborazione dei carichi di lavoro non segue più rigide e statiche configurazioni prefissate, ma viene rimodellata in tempo reale in base a previsioni analitiche, ottimizzazioni multi-obiettivo e reazioni guidate da agenti intelligenti.

% [IMMAGINE DEL FILE DI SLIDE CON NOME "CloudComputing_4"] Schema concettuale della sinergia tra AI e Cloud Computing e convergenza nei servizi digitali — Slide 2 / Pagina 2 del PDF

\subsection{La Sinergia Sistemica Ed Esempi Applicativi}

L'esempio paradigmatico della cooperazione simbiotica tra cloud e AI è rappresentato dagli assistenti vocali su larga scala, quali Amazon Alexa, Apple Siri o Google Assistant. Il ciclo operativo di tali piattaforme evidenzia la ripartizione dei compiti computazionali:
\begin{enumerate}
    \item \textbf{Cattura e trasmissione a bassa latenza}: il dispositivo periferico (\emph{edge/IoT device}) acquisisce il segnale acustico ed esegue un pre-filtraggio minimale del rumore, aprendo una connessione di trasporto sicura verso il cluster cloud.
    \item \textbf{Elaborazione elastica su infrastruttura cloud}: il cluster cloud riceve il flusso audio e alloca istantaneamente istanze computazionali scalabili per convertire la forma d'onda in testo (\emph{Automatic Speech Recognition}, ASR).
    \item \textbf{Comprensione semantica e inferenza}: algoritmi di \emph{Natural Language Understanding} (NLU) estraggono l'intento dell'utente e gli argomenti d'interesse (\emph{entities}), interrogando contemporaneamente basi di conoscenza e microservizi distribuiti.
    \item \textbf{Sintesi e restituzione}: viene formulata la risposta, sintetizzata in formato audio (\emph{Text-to-Speech}, TTS) e trasmessa al dispositivo periferico con latenze nell'ordine delle centinaia di millisecondi.
\end{enumerate}

Senza l'elasticità computazionale del cloud, i dispositivi finali non disporrebbero dell'accelerazione hardware necessaria (GPU, TPU, NPU) né della memoria richiesta per mantenere in locale modelli linguistici complessi. Parallelamente, senza l'automazione intelligente, il data center non potrebbe smistare e bilanciare milioni di flussi audio concorrenti garantendo vincoli stringenti di qualità del servizio (\emph{Quality of Service}, QoS).

Un secondo impatto fondamentale riguarda la democratizzazione dell'adozione tecnologica per le imprese. Storicamente, l'integrazione di pipeline analitiche imponeva un compromesso binario penalizzante: allestire costosi reparti di ricerca e sviluppo interni supportati da data center proprietari on-premises, oppure esternalizzare l'intero stack a fornitori verticali proprietari. Il cloud computing istituisce una via intermedia accessibile (\emph{cloud-hosted AI services}), abilitando l'accesso \emph{on-demand} sia a modelli pre-addestrati consumabili via API REST o gRPC, sia a framework MLOps nativi su Kubernetes per orchestrare modelli proprietari su nodi virtuali noleggiati al secondo.

\subsection{Rilevanza Strategica E Trasformazione Dello Stack Cloud}

L'adozione dell'AI nell'ecosistema cloud non costituisce un semplice livello applicativo accessorio, ma ristruttura dalle fondamenta l'intero \emph{stack} architetturale. Le aziende che considerano l'AI come componente opzionale incorrono in gravi inefficienze operative; al contrario, le infrastrutture di ultima generazione sono esplicitamente ingegnerizzate attorno ai principi dell'\emph{AIOps} (\emph{Artificial Intelligence for IT Operations}).

% [IMMAGINE DEL FILE DI SLIDE CON NOME "CloudComputing_4"] I quattro pilastri dell'importanza strategica dell'AI nel Cloud Computing — Slide 3 / Pagina 3 del PDF

L'impatto architetturale si articola lungo quattro direttrici primarie:
\begin{itemize}
    \item \textbf{Efficienza Operativa su Vasta Scala}: l'approccio tradizionale alla gestione dell'infrastruttura si basa su soglie reattive (ad esempio: \emph{``se la CPU media del cluster supera l'80\% per più di 5 minuti, avvia due nuove repliche''}). Questo meccanismo soffre intrinsecamente di isteresi e ritardo temporale d'innesco, causando frequenti degradazioni di throughput durante transitori rapidi. L'AI permette il passaggio a operazioni \emph{predittive}: le anomalie vengono rilevate prima che si traducano in disservizi e il dimensionamento dei carichi anticipa i picchi di traffico.
    \item \textbf{Allocazione Ottimizzata Delle Risorse}: motori basati su algoritmi euristici e di apprendimento rinforzato calibrano dinamicamente la ripartizione di capacità di calcolo, memoria volatile, larghezza di banda e storage a blocchi, massimizzando il consolidamento dei carichi di lavoro (\emph{workload packing}) ed eliminando il fenomeno dell'\emph{over-provisioning} in configurazioni ibride e multi-cloud.
    \item \textbf{Sicurezza Continua e Rilevamento in Tempo Reale}: le minacce informatiche moderne evolvono rapidamente eludendo i controlli a firma statica tipici dei firewall tradizionali. L'integrazione di modelli di apprendimento automatico consente di stabilire basi di comportamento nominale (\emph{behavioral baselining}) a livello di chiamate di sistema del kernel, traffico di rete e log applicativi, isolando tempestivamente exploit \emph{zero-day} e attacchi complessi \emph{multi-stage}.
    \item \textbf{Accelerazione Del Ciclo Di Vita Del Software}: l'offerta di piattaforme integrate per il tracciamento, versionamento e rilascio di modelli (\emph{MLOps}) e l'accesso a microservizi di inferenza per modelli generativi riducono drasticamente l'intervallo temporale che intercorre tra l'ideazione dell'algoritmo e la sua messa in produzione.
\end{itemize}

% [IMMAGINE DEL FILE DI SLIDE CON NOME "CloudComputing_4"] Mappa concettuale delle 13 macro-aree di impatto tra AI e Cloud Computing — Slide 4 / Pagina 4 del PDF

L'intersezione tra AI e cloud investe una pluralità di dimensioni tecniche ed economiche. Nelle sezioni successive verranno analizzate in dettaglio le macro-aree fondamentali: dalla gestione automatizzata dell'infrastruttura all'abbattimento dei costi, dalla scalabilità elastica all'analisi predittiva, fino alla sicurezza distribuita e all'accessibilità dei dati tramite interfacce semantiche.

\newpage

% ==============================================================================
% SEZIONE: GESTIONE DEL CLOUD TRAMITE INTELLIGENZA ARTIFICIALE
% ==============================================================================

\section{Gestione Del Cloud Tramite Intelligenza Artificiale}

La gestione convenzionale delle infrastrutture distribuite impone agli operatori sistemistici compiti complessi di configurazione manuale, allocazione statica e diagnostica post-guasto. L'integrazione di algoritmi intelligenti abilita l'automazione completa dell'intero ciclo di vita delle risorse: dal provisioning iniziale dell'hardware virtualizzato al monitoraggio continuo, fino all'auto-riparazione (\emph{self-healing}) e all'ottimizzazione energetica.

% [IMMAGINE DEL FILE DI SLIDE CON NOME "CloudComputing_4"] Panoramica architetturale sulla gestione del cloud con AI e riduzione del TCO — Slide 5 / Pagina 5 del PDF

Un driver essenziale nell'ingegneria del cloud è la minimizzazione del \textbf{Total Cost of Ownership} (TCO). Il TCO è un parametro quantitativo che valuta l'impegno finanziario complessivo generato da un'infrastruttura tecnologica lungo l'intero arco temporale del suo utilizzo, integrando non solo i costi di acquisizione iniziale (\emph{Capital Expenditure}, CapEx), ma soprattutto le spese correnti di alimentazione, raffreddamento, canoni di licenza, manutenzione, personale operativo e perdite dovute a fermi improvvisi di servizio (\emph{Operational Expenditure}, OpEx).

\begin{equation}
    \text{TCO} = \text{CapEx}_{\text{acquisizione}} + \sum_{t=1}^{T} \frac{\text{OpEx}_{\text{infrastruttura}}(t) + \text{OpEx}_{\text{gestione}}(t) + C_{\text{downtime}}(t)}{(1 + r)^t}
\end{equation}
dove $r$ rappresenta il tasso di attualizzazione finanziaria e $T$ l'orizzonte temporale considerato. L'applicazione dell'AI incide direttamente riducendo $\text{OpEx}_{\text{gestione}}$ attraverso l'automazione dei processi e minimizzando $C_{\text{downtime}}$ mediante l'identificazione precoce dei guasti.

% [IMMAGINE DEL FILE DI SLIDE CON NOME "CloudComputing_4"] Tabella riassuntiva: algoritmi e tecnologie per Provisioning e Orchestrazione — Slide 6 / Pagina 6 del PDF

\subsection{Provisioning Ed Orchestrazione Predittiva}

Nel provisioning e nell'orchestrazione di macchine virtuali e container, gli obiettivi primari sono tre:
\begin{enumerate}
    \item Prevedere accuratamente la richiesta di risorse hardware future (\emph{resource demand});
    \item Anticipare le azioni di scalabilità orizzontale e verticale prima che si verifichi la saturazione dei nodi;
    \item Ottimizzare il piazzamento dei carichi computazionali (\emph{VM/container placement}) minimizzando la frammentazione e il consumo energetico.
\end{enumerate}

Per soddisfare questi requisiti si ricorre a un'ampia classe di modelli matematici:
\begin{itemize}
    \item \textbf{Modelli di Serie Temporali Classici}: approcci autoregressivi integrati a media mobile, quali \textbf{ARIMA} (\emph{Autoregressive Integrated Moving Average}) e la sua controparte stagionale \textbf{SARIMA}, espressi formalmente come:
    \begin{equation}
        \Phi(B) (1 - B)^d X_t = \Theta(B) \epsilon_t
    \end{equation}
    dove $B$ è l'operatore di ritardo ($B^k X_t = X_{t-k}$), $d$ è l'ordine di differenziazione necessario a indurre stazionarietà, e $\Phi, \Theta$ sono polinomi autoregressivi e di media mobile. A questi si affianca \textbf{Prophet}, un modello decomponibile additivo basato su regressione a spezzate con termini di stagionalità fittizi per catturare tendenze settimanali e festività:
    \begin{equation}
        y(t) = g(t) + s(t) + h(t) + \epsilon_t
    \end{equation}
    dove $g(t)$ modella la tendenza non periodica, $s(t)$ i cambiamenti periodici (stagionalità) e $h(t)$ gli effetti di eventi discreti non ricorrenti.
    \item \textbf{Deep Learning per Serie Temporali}: architetture a memoria ricorrente di tipo \textbf{LSTM} (\emph{Long Short-Term Memory}), in grado di mantenere traccia di dipendenze temporali estese superando il problema dell'attenuazione del gradiente (\emph{vanishing gradient}), e \textbf{Temporal Convolutional Networks} (TCN), le quali impiegano convoluzioni dilatate unidimensionali causalmente coerenti per estrarre caratteristiche su orizzonti ricettivi ampi con parallelizzazione completa durante l'addestramento.
    \item \textbf{Apprendimento per Rinforzo} (\emph{Reinforcement Learning}, RL): formalizzato come un Processo Decisionale di Markov (\emph{Markov Decision Process}, MDP) descritto dalla quadrupla $(\mathcal{S}, \mathcal{A}, \mathcal{P}, \mathcal{R})$. Lo stato $s_t \in \mathcal{S}$ descrive l'impegno di CPU, memoria e traffico di rete del cluster; l'azione $a_t \in \mathcal{A}$ consiste nel variare il numero di pod o migrare una VM; la funzione di ricompensa $r_t \in \mathcal{R}$ penalizza la violazione dei \emph{Service Level Agreements} (SLA) bilanciandola con il costo orario delle istanze attive. Vengono impiegati algoritmi a controllo di valore come il \textbf{Q-Learning} tabellare o reti \textbf{DQN} (\emph{Deep Q-Networks}) con esperienza rigiocata (\emph{experience replay}), e algoritmi a gradiente di politica come \textbf{PPO} (\emph{Proximal Policy Optimization}) che assicurano aggiornamenti stabili limitando la divergenza di Kullback-Leibler tra la vecchia e la nuova politica.
    \item \textbf{Euristiche di Ottimizzazione Combinatoria}: \textbf{Algoritmi Genetici}, \textbf{Simulated Annealing} e solutori a vincoli (\emph{Constraint Solvers}) che risolvono varianti del problema del \emph{Bin Packing} multi-dimensionale, minimizzando il numero totale di server fisici attivi.
\end{itemize}

Questi modelli operano a supporto degli strumenti industriali di orchestrazione e \emph{Infrastructure as Code} (IaC):
\begin{itemize}
    \item \textbf{Terraform}: definisce in modo dichiarativo lo stato desiderato delle risorse IaaS multi-provider;
    \item \textbf{Pulumi}: permette di esprimere le logiche infrastrutturali tramite linguaggi imperativi ad alto livello (TypeScript, Python, Go);
    \item \textbf{Cloud Custodian}: motore di regole per applicare politiche automatizzate di conformità e spegnimento delle risorse orfane;
    \item \textbf{Kubernetes}: il control plane nativo estendibile tramite operatori personalizzati per pilotare lo scheduling e il ridimensionamento dei nodi in base agli output generati dai modelli analitici.
\end{itemize}

% [IMMAGINE DEL FILE DI SLIDE CON NOME "CloudComputing_4"] Tassonomia di algoritmi per monitoraggio e anomaly detection e librerie ML — Slide 7 / Pagina 7 del PDF

\subsection{Monitoraggio Ed Anomaly Detection}

Il monitoraggio avanzato di una flotta cloud deve rilevare pattern anomali di consumo di clock, saturazione della coda I/O sui dischi ed esaurimento dei socket di rete, individuando sia derive lente del comportamento nominale (\emph{baseline drift}) sia picchi istantanei isolati.

I modelli di apprendimento non supervisionato (\emph{Unsupervised ML}) risultano fondamentali data la quasi totale assenza di dataset annotati contenenti etichette di guasto pregresse:
\begin{itemize}
    \item \textbf{Isolation Forest}: algoritmo basato sulla costruzione ricorsiva di alberi decisionali binari casuali (\emph{iTrees}). Il principio teorico sfrutta il fatto che i dati anomali sono rari e presentano attributi molto distanti dalla norma; di conseguenza, all'interno dello spazio delle caratteristiche, vengono isolati con un numero di partizioni binarie decisamente inferiore rispetto ai punti normali, traducendosi in una lunghezza media del percorso radice-foglia $h(x)$ sensibilmente ridotta. Il punteggio di anomalia $s(x, n)$ per un campione $x$ su una collezione di $n$ istanze è espresso da:
    \begin{equation}
        s(x, n) = 2^{-\frac{\mathbb{E}[h(x)]}{c(n)}}
    \end{equation}
    dove $c(n) = 2 \ln(n - 1) + 0.5772156649 - \frac{2(n - 1)}{n}$ rappresenta la lunghezza media di ricerca di un albero binario non riuscito.
    \item \textbf{One-Class Support Vector Machine} (OC-SVM): mappa i vettori delle metriche in uno spazio a dimensionalità infinita mediante kernel non lineari (es. RBF Gaussiano) e identifica l'iperpiano a massimo margine che separa l'origine dalla regione in cui si concentra la densità dei dati nominali.
    \item \textbf{DBSCAN} (\emph{Density-Based Spatial Clustering of Applications with Noise}): raggruppa campioni vicini in base a parametri di raggio spaziale $\varepsilon$ e numero minimo di punti vicini \emph{MinPts}. I campioni situati in aree a densità insufficiente vengono classificati come rumore o anomalie infrastrutturali.
    \item \textbf{Local Outlier Factor} (\textbf{LOF}): quantifica la densità locale di un punto rispetto a quella dei suoi $k$-vicini prossimi. Un punto immerso in un'area a densità marcatamente inferiore rispetto a quella dell'intorno circostante riceve uno score di anomalia elevato, consentendo di identificare anomalie contestuali al variare delle condizioni d'uso del cluster.
\end{itemize}

% [IMMAGINE DEL FILE DI SLIDE CON NOME "CloudComputing_4"] Modelli predittivi multivariati e Change-Point Detection per il monitoraggio cloud — Slide 8 / Pagina 8 del PDF

Parallelamente, per analizzare l'interazione contemporanea tra centinaia di flussi di telemetria, si impiegano modelli generativi di tipo \textbf{Multivariate Forecasting}:
\begin{itemize}
    \item \textbf{LSTM Autoencoder}: una rete neurale ricorrente comprime la sequenza multivariata di telemetria in uno spazio latente a bassa dimensionalità (\emph{encoder}) per poi tentare di ricostruire la sequenza temporale di partenza (\emph{decoder}). La funzione di perdita è governata dall'errore quadratico medio di ricostruzione:
    \begin{equation}
        \mathcal{L}_{\text{rec}} = \frac{1}{T} \sum_{t=1}^T \lVert \bm{x}_t - \hat{\bm{x}}_t \rVert_2^2
    \end{equation}
    In presenza di un degrado non nominale o di un'interferenza tra container, la correlazione multivariata salta e il modello produce un errore $\mathcal{L}_{\text{rec}}$ marcatamente superiore alla soglia di guardia calibrata.
    \item \textbf{Variational Autoencoder} (VAE): a differenza dell'autoencoder tradizionale, il VAE modella la distribuzione dei dati all'interno dello spazio latente in termini statistici, imponendo che il codificatore produca i parametri di una distribuzione gaussiana (media $\bm{\mu}$ e varianza $\bm{\sigma}^2$). L'ottimizzazione massimizza il limite inferiore dell'evidenza (\emph{Evidence Lower Bound}, ELBO):
    \begin{equation}
        \text{ELBO} = \mathbb{E}_{q_\phi(\bm{z}|\bm{x})}[\ln p_\theta(\bm{x}|\bm{z})] - D_{\text{KL}}(q_\phi(\bm{z}|\bm{x}) \,||\, p(\bm{z}))
    \end{equation}
    consentendo un'identificazione robusta di stati anomali su base strettamente probabilistica.
\end{itemize}

Nei casi in cui sia critico distinguere tra fluttuazioni transitorie e variazioni strutturali permanenti (es. guasto imminente di un modulo di memoria o degrado del cavo di rete), intervengono gli algoritmi di \textbf{Change-Point Detection} (CPD):
\begin{itemize}
    \item \textbf{Ruptures}: libreria per l'analisi \emph{offline} delle serie storiche, orientata a segmentare i segnali temporali identificando le discontinuità strutturali pregresse;
    \item \textbf{Bayesian Online Change Detection} (\textbf{BOCD}): approccio ricorsivo \emph{online} che calcola in tempo reale la distribuzione di probabilità a posteriori del parametro \emph{run-length} $r_t$ (tempo trascorso dall'ultimo punto di rottura sistemica), garantendo un allineamento immediato delle soglie operative.
\end{itemize}

Lo stack software industriale che raccoglie e alimenta questi algoritmi si fonda su:
\begin{itemize}
    \item \textbf{Prometheus}: time-series database orientato alla raccolta pull di metriche scalari da endpoint HTTP;
    \item \textbf{Grafana, Mimir e Loki}: piattaforme di visualizzazione, storage distribuito di metriche a lungo termine e aggregazione centralizzata dei log di sistema;
    \item \textbf{OpenTelemetry}: framework aperto e unificato per la standardizzazione della raccolta di tracce distribuite, parametri e log;
    \item \textbf{PyOD}, \textbf{GluonTS} e \textbf{ThirdEye} (sviluppato da LinkedIn): librerie specializzate per l'inferenza rapida di anomalie e per l'analisi guidata della causa radice (\emph{Root Cause Analysis}, RCA).
\end{itemize}

% [IMMAGINE DEL FILE DI SLIDE CON NOME "CloudComputing_4"] Tabella riassuntiva: algoritmi e tecnologie per Workload Optimization — Slide 9 / Pagina 9 del PDF

\subsection{Ottimizzazione Del Carico Di Lavoro}

L'ottimizzazione del carico di lavoro (\emph{Workload Optimization}) mira ad auto-sintonizzare i parametri di esecuzione dei processi distribuiti, pianificare le scalature e definire politiche di scheduling guidate dal risparmio energetico o dalla tariffazione oraria del fornitore di infrastruttura.

L'auto-tuning dei parametri di esecuzione (come la dimensione dell'\emph{heap} della JVM, i buffer di thread o i flag di compilazione JIT) viene trattato come l'ottimizzazione di una funzione scatola nera (\emph{black-box}) costosa da campionare. In questo contesto eccelle l'\textbf{Ottimizzazione Bayesiana} basata su \textbf{Processi Gaussiani} ($GP$):
\begin{equation}
    f(\bm{x}) \sim \mathcal{GP}\left(m(\bm{x}), k(\bm{x}, \bm{x}')\right)
\end{equation}
La configurazione successiva $\bm{x}_{t+1}$ da validare a runtime viene selezionata massimizzando una funzione di acquisizione, come l'\emph{Expected Improvement} ($EI$), minimizzando il numero di iterazioni degradate richieste per convergere alla configurazione ottimale.

Parallelamente, lo scheduling su nodi fisici computazionalmente eterogenei integra formulazioni matematiche esatte di \textbf{Programmazione Lineare Intera} (\emph{Integer Linear Programming}, ILP), euristiche basate su \textbf{Particle Swarm Optimization} (PSO) e algoritmi di apprendimento per rinforzo continuo.

Sul fronte applicativo, queste metodologie potenziano il \textbf{Vertical Pod Autoscaler} (VPA) di Kubernetes, permettendo al sistema di variare dinamicamente le clausole di \texttt{resources.requests} e \texttt{resources.limits} in base ai requisiti stimati, e si integrano all'interno del \texttt{kube-scheduler} tramite plugin personalizzati che tengono conto di metriche avanzate di affinità hardware, consumi termici e costi di rete intra-zona.

% [IMMAGINE DEL FILE DI SLIDE CON NOME "CloudComputing_4"] Tabella riassuntiva: algoritmi e tecnologie per Cost Optimization e riduzione del TCO — Slide 10 / Pagina 10 del PDF

\subsection{Ottimizzazione Dei Costi E Riduzione Del TCO}

L'ultimo pilastro della gestione automatizzata concerne il controllo analitico del bilancio economico cloud. L'obiettivo consiste nel prevedere con elevata accuratezza l'esborso finanziario generato dall'infrastruttura ed eseguire il dimensionamento ottimale (\emph{right-sizing}) delle istanze attive, eliminando le allocazioni sovradimensionate.

I carichi vengono segmentati tramite algoritmi di clustering non supervisionato (\textbf{K-Means}) per identificare categorie di utilizzo omogenee (es. istanze a ciclo continuo versus carichi \emph{batch} intermittenti), dopodiché modelli di regressione robusta (\textbf{XGBoost}, \textbf{Random Forest}) o previsori neurali (\textbf{LSTM}) quantificano l'evoluzione del consumo orario. Il dimensionamento ottimale viene quindi calcolato risolvendo un problema di ottimizzazione vincolata:
\begin{equation}
    \min_{\bm{w}} \sum_{i=1}^N c_i w_i \quad \text{soggetto a} \quad \sum_{i=1}^N R_{ij} w_i \ge D_j, \quad \forall j \in \{\text{CPU}, \text{RAM}, \text{IOPS}\}
\end{equation}
dove $c_i$ indica la tariffa dell'istanza $i$-esima, $w_i \in \{0, 1\}$ l'assegnazione dell'istanza, $R_{ij}$ la capacità della risorsa $j$, e $D_j$ il fabbisogno minimo calcolato per evitare colli di bottiglia computazionali.

Strumenti a supporto di questo processo includono:
\begin{itemize}
    \item \textbf{Kubecost} e la sua declinazione aperta sotto egida CNCF, \textbf{OpenCost}: strumenti specializzati nel monitoraggio granulare delle spese generate da namespace, pod e servizi su cluster Kubernetes, integrando previsori di spesa automatizzati;
    \item \textbf{Kubeflow} e \textbf{Ray}: piattaforme scalabili di livello enterprise deputate all'esecuzione e al deployment distribuito dei modelli di apprendimento per rinforzo dedicati al controllo delle risorse.
\end{itemize}

\newpage

% ==============================================================================
% SEZIONE: OTTIMIZZAZIONE DEI COSTI NEL CALCOLO PARALLELO ED ELASTICO
% ==============================================================================

\section{Ottimizzazione Dei Costi Nel Calcolo Parallelo Ed Elastico}

L'adozione del modello di fatturazione \emph{pay-as-you-go} caratteristico del cloud computing solleva le organizzazioni dalla necessità di allocare ingenti investimenti di capitale iniziale per l'allestimento di cluster di calcolo dedicati. È possibile attivare al bisogno pool eterogenei di acceleratori grafici ad alte prestazioni (\textbf{GPU}), memorie a banda larga (\textbf{HBM}) e storage a oggetti, saldando unicamente l'effettivo consumo orario o al millisecondo.

% [IMMAGINE DEL FILE DI SLIDE CON NOME "CloudComputing_4"] Principi base della convenienza economica del cloud per l'elaborazione di modelli AI — Slide 11 / Pagina 11 del PDF

Tuttavia, l'addestramento distribuito di reti neurali profonde con centinaia di miliardi di parametri (\emph{Large Language Models} e modelli di visione computazionale) richiede centinaia di migliaia di ore-GPU. In tale contesto, l'efficienza algoritmica di calcolo e la gestione della memoria hardware assumono un ruolo economico determinante: una riduzione del 30\% dell'impegno di memoria permette di dimezzare il numero di istanze noleggiate, incidendo direttamente sui costi operativi.

% [IMMAGINE DEL FILE DI SLIDE CON NOME "CloudComputing_4"] Tabella riassuntiva: algoritmi per efficient training e training and memory optimization — Slide 12 / Pagina 12 del PDF

\subsection{Tecniche Di Ottimizzazione Della Memoria E Addestramento Distribuito}

I requisiti di memoria per l'addestramento di una rete neurale derivano da quattro contributi primari: i \textbf{pesi del modello}, i \textbf{gradienti}, gli \textbf{stati dell'ottimizzatore} (es. vettori del primo e secondo momento in Adam) e le \textbf{attivazioni intermedie} generate durante la fase di propagazione in avanti (\emph{forward pass}) e mantenute in memoria per calcolare i gradienti tramite la regola della catena nella retro-propagazione (\emph{backward pass}).

Per abbattere l'impronta computazionale e velocizzare l'elaborazione su GPU si impiegano le seguenti metodologie:
\begin{itemize}
    \item \textbf{Mixed-Precision Training}: anziché mantenere l'intero processo su valori in virgola mobile a singola precisione a 32 bit (\texttt{FP32}), si introducono formati a precisione ridotta a 16 bit, quali \texttt{FP16} o \textbf{BF16} (\emph{Brain Floating Point}, che preserva i medesimi 8 bit di esponente di \texttt{FP32} riducendo la mantissa a 7 bit). L'utilizzo di registri a 16 bit dimezza l'occupazione di memoria e attiva le unità di calcolo tensoriale dedicate sulle GPU (\emph{Tensor Cores}). Per prevenire il fenomeno dell'\emph{underflow} numerico dei gradienti a valori prossimi allo zero, la tecnica implementa un fattore di scala dinamico della perdita (\emph{dynamic loss scaling}): la perdita calcolata viene moltiplicata per un fattore $S$ prima del backward pass e i gradienti risultanti vengono riscalati di $1/S$ prima dell'aggiornamento dei pesi.
    \item \textbf{Gradient Checkpointing}: memorizzare le attivazioni intermedie di ogni strato satura la memoria HBM della GPU. Questa tecnica memorizza selettivamente soltanto un sottoinsieme ridotto di attivazioni di confine (i \emph{checkpoint}); durante la fase di retropropagazione, le attivazioni interne mancanti vengono ricalcolate al volo eseguendo una breve propagazione in avanti locale. Si introduce un modesto sovraccarico computazionale aggiuntivo (tipicamente attorno al 20-30\% di tempo di calcolo in più), ma si ottiene una riduzione drastica dell'occupazione di picco della memoria, che scala da un ordine $\mathcal{O}(L)$ (dove $L$ è la profondità della rete) a $\mathcal{O}(\sqrt{L})$.
    \item \textbf{Tensor Parallelism}: quando un singolo strato o blocco di attenzione eccede la capacità di memoria fisica di una singola GPU, lo strato viene frammentato su più dispositivi. Nelle moltiplicazioni matrice-vettore $Y = XW$, la matrice dei pesi $W$ viene suddivisa per colonne o per righe tra le GPU cooperanti, sincronizzando i risultati intermedi ad alta velocità tramite primitive di comunicazione collettiva \texttt{All-Reduce} sul bus hardware \textbf{NVLink}.
    \item \textbf{Pipeline Parallelism}: la sequenza verticale degli strati della rete neurale viene ripartita su dispositivi distinti organizzati a catena. Per evitare il sotto-utilizzo computazionale (la cosiddetta \emph{pipeline bubble}), il lotto di dati (\emph{batch}) viene suddiviso in \emph{micro-batch} elaborati in parallelo secondo schemi a \emph{pipelining} temporale (es. scheduling 1F1B: \emph{One Forward, One Backward}).
    \item \textbf{ZeRO} (\emph{Zero Redundancy Optimizer}): algoritmo introdotto nel framework DeepSpeed per abbattere le ridondanze tipiche del puro parallelismo dati. Nel parallelismo convenzionale, ciascun nodo GPU mantiene una copia completa di stati dell'ottimizzatore, gradienti e parametri. ZeRO elimina tale duplicazione suddividendo gli stati della memoria tra le GPU:
    \begin{itemize}
        \item \emph{ZeRO-Stage 1}: partiziona gli stati dell'ottimizzatore tra i nodi, riducendo l'impronta di memoria di circa 4 volte senza generare traffico di comunicazione aggiuntivo;
        \item \emph{ZeRO-Stage 2}: partiziona sia gli stati dell'ottimizzatore sia i gradienti;
        \item \emph{ZeRO-Stage 3}: partiziona tutti e tre i componenti (inclusi i parametri del modello), richiedendo la raccolta temporanea dei pesi tramite chiamate \texttt{All-Gather} broadcast unicamente nel momento in cui il singolo strato deve eseguire il calcolo.
    \end{itemize}
\end{itemize}

% [IMMAGINE DEL FILE DI SLIDE CON NOME "CloudComputing_4"] Dettaglio algoritmi di ottimizzazione, compressione del modello e hyperparameter tuning — Slide 13 / Pagina 13 del PDF

\subsection{Algoritmi Di Ottimizzazione Ed Efficienza Algoritmica}

La convergenza del processo di addestramento distribuito viene governata da ottimizzatori specificamente progettati per mitigare l'impatto della latenza di rete tra nodi:
\begin{itemize}
    \item \textbf{Distributed SGD}: parallelizza la discesa stocastica del gradiente calcolando i gradienti locali $\nabla_{\bm{\theta}} \mathcal{L}_k$ su $K$ nodi indipendenti, calcolando la media aggregata dei contributi prima dell'aggiornamento:
    \begin{equation}
        \bm{\theta}_{t+1} = \bm{\theta}_t - \eta \cdot \frac{1}{K} \sum_{k=1}^K \nabla_{\bm{\theta}} \mathcal{L}_k(\bm{\theta}_t)
    \end{equation}
    \item \textbf{AdamW}: variante fondamentale del classico algoritmo Adam, in cui la regolarizzazione mediante decadimento dei pesi (\emph{weight decay}) viene esplicitamente disaccoppiata dall'aggiornamento adattivo dei momenti dei gradienti:
    \begin{align}
        \bm{m}_t &= \beta_1 \bm{m}_{t-1} + (1 - \beta_1) \bm{g}_t \\
        \bm{v}_t &= \beta_2 \bm{v}_{t-1} + (1 - \beta_2) \bm{g}_t^2 \\
        \hat{\bm{m}}_t &= \frac{\bm{m}_t}{1 - \beta_1^t}, \quad \hat{\bm{v}}_t = \frac{\bm{v}_t}{1 - \beta_2^t} \\
        \bm{\theta}_{t+1} &= \bm{\theta}_t - \eta_t \left( \frac{\hat{\bm{m}}_t}{\sqrt{\hat{\bm{v}}_t} + \epsilon} + \lambda \bm{\theta}_t \right)
    \end{align}
    Questo disaccoppiamento ripristina la corretta interpretazione geometrica della norma $L_2$ e garantisce una generalizzazione superiore nell'addestramento di architetture Transformer.
    \item \textbf{LAMB} (\emph{Layer-wise Adaptive Moments optimizer for Batch training}): specificamente ingegnerizzato per scalare l'addestramento con batch estremamente grandi (superiori a $32\,000$ campioni per passo). Adatta il tasso di apprendimento in modo indipendente per ciascuno strato basandosi sul rapporto tra la norma euclidea dei parametri dello strato e la norma dell'aggiornamento adattivo:
    \begin{equation}
        r_t^{(l)} = \frac{\lVert \bm{\theta}_t^{(l)} \rVert_2}{\lVert \bm{u}_t^{(l)} \rVert_2}
    \end{equation}
    evitando instabilità numeriche e divergenze del gradiente che altrimenti impedirebbero la convergenza su ampi cluster di GPU.
\end{itemize}

\subsection{Compressione Del Modello Per L'Inferenza Economica}

Una volta concluso l'addestramento, l'esecuzione continua dell'inferenza in cloud rappresenta la voce di costo più rilevante. La compressione dei modelli consente di impiegare hardware meno costoso o di ridurre la latenza:
\begin{itemize}
    \item \textbf{Quantizzazione} (a 8 bit o 4 bit): converte i tensori dei pesi e delle attivazioni dal formato floating point continuo (\texttt{FP16}) a rappresentazioni a numeri interi discreti (\texttt{INT8}, \texttt{INT4}). Mappando linearmente l'intervallo continuo $[x_{\min}, x_{\max}]$ nell'intervallo intero $[-128, 127]$ tramite un fattore di scala $S$ e un punto zero $Z$:
    \begin{equation}
        q = \text{round}\left( \frac{x}{S} \right) + Z
    \end{equation}
    Questa procedura abbatte l'occupazione di memoria fino al 75\% e incrementa la velocità di throughput sfruttando le istruzioni vettoriali intere senza decadimento apprezzabile dell'accuratezza.
    \item \textbf{Pruning}: identifica ed azzera i parametri o le connessioni neuronali che apportano un contributo trascurabile alla funzione obiettivo. Può essere \emph{strutturato} (rimozione di interi filtri convoluzionali, teste di attenzione o canali) per tradursi immediatamente in un guadagno prestazionale sulle architetture hardware standard, oppure \emph{non strutturato} (imposizione di matrici sparse casuali).
    \item \textbf{Distillazione Della Conoscenza} (\emph{Knowledge Distillation}): addestra una rete neurale compatta (\emph{student}) inducendola ad approssimare la distribuzione di probabilità d'uscita (\emph{soft targets}) prodotta da un modello ad alta complessità (\emph{teacher}) precedentemente addestrato. La funzione di costo viene calibrata tramite un parametro di temperatura $T$:
    \begin{equation}
        \mathcal{L}_{\text{KD}} = (1 - \alpha) \mathcal{L}_{\text{CE}}(y, \hat{y}_{\text{student}}) + \alpha T^2 \, D_{\text{KL}}\!\left( \sigma\left(\frac{\bm{z}_{\text{teacher}}}{T}\right) \Bigg|\Bigg| \, \sigma\left(\frac{\bm{z}_{\text{student}}}{T}\right) \right)
    \end{equation}
    \item \textbf{Low-Rank Adaptation} (\textbf{LoRA}): tecnica di sintonizzazione fine (\emph{fine-tuning}) parametro-efficiente. Anziché ricalcolare tutti i pesi $W_0 \in \mathbb{R}^{d \times k}$ di un modello pre-addestrato, $W_0$ viene congelato e la variazione $\Delta W$ viene vincolata a una decomposizione a basso rango:
    \begin{equation}
        W = W_0 + \Delta W = W_0 + \frac{\alpha}{r} (B \cdot A), \quad B \in \mathbb{R}^{d \times r}, \; A \in \mathbb{R}^{r \times k}
    \end{equation}
    con rango intrinseco $r \ll \min(d, k)$. Ciò riduce di ordini di grandezza il numero di parametri da addestrare e memorizzare per singolo tenant o cliente cloud.
\end{itemize}

% [IMMAGINE DEL FILE DI SLIDE CON NOME "CloudComputing_4"] Tecniche di Hyperparameter Optimization: Bayesian, Hyperband, PBT — Slide 14 / Pagina 14 del PDF

\subsection{Ottimizzazione Degli Iperparametri}

La ricerca manuale delle configurazioni ottimali comporta un immenso spreco di ore-GPU. I framework cloud integrano metodologie di \emph{Hyperparameter Optimization} (HPO) automatizzate:
\begin{itemize}
    \item \textbf{Ottimizzazione Bayesiana}: campiona iterativamente lo spazio degli iperparametri costruendo un modello probabilistico surrogato della funzione di perdita e massimizzando l'utilità marginale attesa.
    \item \textbf{Hyperband}: algoritmo guidato dall'allocazione delle risorse che estende la ricerca causale (\emph{random search}). Formula l'esplorazione come un problema di \emph{bandit} multi-braccio, sfruttando la tecnica di \textbf{Successive Halving}: avvia un gran numero di configurazioni assegnando loro una quantità ridotta di epoche computazionali, valuta le prestazioni intermedie, interrompe anticipatamente il 50\% peggiore delle configurazioni e raddoppia le risorse allocate per quelle promettenti, concentrando le ore-macchina sui candidati migliori.
    \item \textbf{Population-Based Training} (\textbf{PBT}): combina l'ottimizzazione degli iperparametri e l'addestramento del modello in un unico processo evolutivo continuo. Una popolazione di modelli viene addestrata contemporaneamente; periodicamente, i modelli che evidenziano metriche inferiori ereditano i pesi e clonano gli iperparametri delle istanze migliori, applicando mutazioni casuali per continuare la ricerca nello spazio di esplorazione senza dover riavviare l'addestramento dall'inizio.
\end{itemize}

% [IMMAGINE DEL FILE DI SLIDE CON NOME "CloudComputing_4"] Tabella riassuntiva: Data Handling e Large-Scale Storage — Slide 15 / Pagina 15 del PDF

\subsection{Gestione Dei Dati E Storage Su Larga Scala}

La gestione efficiente dei dati di addestramento e dei vettori latenti su larga scala richiede stack di archiviazione specializzati:
\begin{itemize}
    \item \textbf{Algoritmi di Ricerca Vettoriale}: per confrontare rapidamente embedding ad altissima dimensionalità senza scansioni esaustive, si usano strutture ad indice basate su grafi approssimati come \textbf{HNSW} (\emph{Hierarchical Navigable Small World}), partizioni cellulari Voronoi con indici invertiti \textbf{IVF} (\emph{Inverted File Index}), quantizzazioni vettoriali \textbf{PQ} (\emph{Product Quantization}) e \textbf{ScaNN}.
    \item \textbf{Deduplicazione dei Dati}: per ripulire dataset di centinaia di terabyte da documenti duplicati o ridondanti, si utilizzano funzioni hash ad alta efficienza statistica: \textbf{MinHash} per approssimare l'indice di somiglianza di Jaccard tra insiemi lessicali e \textbf{SimHash} per valutare la distanza angolare di coseno su vettori pesati.
    \item \textbf{Preprocessing Distribuito}: basato su architetture a flussi con calcolo parallelo, shuffling distribuito (\emph{MapReduce}) e protocolli di streaming unificati.
\end{itemize}

Gli storage tecnologici impiegati includono sistemi distribuiti come \textbf{Apache Hadoop / HDFS}, formati aperti di \emph{lakehouse} che garantiscono proprietà ACID e versionamento temporale delle tabelle come \textbf{Apache Iceberg} e \textbf{Delta Lake}, e piattaforme di \emph{Object Storage} distribuito compatibili S3 come \textbf{Ceph} e \textbf{MinIO}.

% [IMMAGINE DEL FILE DI SLIDE CON NOME "CloudComputing_4"] Tabella riassuntiva: algoritmi per automazione ed efficienza operativa — Slide 16 / Pagina 16 del PDF

\subsection{Automazione Ed Efficienza Operativa}

A livello orizzontale, l'integrazione di tecniche di \textbf{Robotic Process Automation} (RPA), accoppiate con modelli di elaborazione del linguaggio naturale (NLP) e agenti basati su modelli linguistici (LLM), permette di automatizzare workflow operativi sistemistici complessi. La convergenza con sistemi di \emph{Infrastructure as Code} come **Terraform**, **Pulumi**, **Ansible** e **Crossplane** (per il provisioning nativo su Kubernetes) crea architetture a ciclo chiuso, dove le raccomandazioni analitiche si traducono direttamente in riconfigurazioni automatiche del cluster senza intervento umano.

\newpage

% ==============================================================================
% SEZIONE: SCALABILITÀ ED AGILITÀ DELLE RISORSE CLOUD
% ==============================================================================

\section{Scalabilità Ed Agilità Delle Risorse Cloud}

La combinazione tra la disponibilità di risorse \emph{on-demand} propria delle piattaforme cloud e le capacità di auto-apprendimento dell'AI offre un livello di flessibilità e scalabilità senza precedenti. Le organizzazioni possono scalare verticalmente (\emph{scale-up}) o orizzontalmente (\emph{scale-out}) i nodi computazionali senza dover effettuare pianificazioni statiche di acquisto hardware.

% [IMMAGINE DEL FILE DI SLIDE CON NOME "CloudComputing_4"] Concetti chiave di Scalabilità & Agilità e aggiornamento continuo dei modelli — Slide 17 / Pagina 17 del PDF

L'agilità architetturale garantisce che le pipeline di elaborazione possano venire aggiornate o riaddestrate continuamente sul campo (\emph{continuous training}) senza generare interruzioni di servizio (\emph{zero-downtime deployments}), adattando l'infrastruttura sottostante ai carichi dinamici.

% [IMMAGINE DEL FILE DI SLIDE CON NOME "CloudComputing_4"] Tabella: algoritmi di Self-Learning e tecnologie Kubernetes per autoscaling avanzato — Slide 18 / Pagina 18 del PDF

\subsection{Algoritmi Di Ottimizzazione Ad Auto-Apprendimento}

Nelle infrastrutture cloud native basate su container, l'autoscaling convenzionale (come l'\textbf{Horizontal Pod Autoscaler} di Kubernetes standard) valuta metriche istantanee mediante soglie lineari reattive:
\begin{equation}
    N_{\text{desiderato}} = \left\lceil N_{\text{corrente}} \times \frac{\text{Metrica}_{\text{corrente}}}{\text{Metrica}_{\text{target}}} \right\rceil
\end{equation}
Tale formula è priva di memoria e ignora le latenze di avvio dei container (\emph{cold start}), l'inerzia termica dei nodi e i vincoli incrociati di rete e I/O.

Per ovviare a questi limiti, la scalabilità elastica viene affidata a controllori basati su \textbf{Reinforcement Learning} (RL) e algoritmi di \textbf{Multi-Armed Bandit}:
\begin{itemize}
    \item \textbf{Formulazione MDP per l'Autoscaling}: lo stato sistemico $s_t$ raccoglie il tasso di arrivo delle richieste $\lambda_t$, il tempo medio di risposta $W_t$, la lunghezza della coda $L_{q, t}$ e il numero attuale di nodi attivi. L'azione $a_t$ determina la variazione del numero di istanze $\Delta N$. La ricompensa $R_t$ penalizza sia la violazione dei vincoli di latenza sia il costo associato all'allocazione delle risorse:
    \begin{equation}
        R_t = - \left( \alpha \cdot \max(0, W_t - \text{SLA}_{\text{lat}}) + \beta \cdot \text{Costo}(N_t) \right)
    \end{equation}
    \item \textbf{Algoritmi RL Utilizzati}: per spazi di stato e azione discreti si impiegano varianti potenziate di **DQN** con reti di stima separate (\emph{Target Network}) per evitare divergenze; per spazi di controllo continui (come la modulazione fine della frequenza della CPU o delle quote di allocazione dei millicore) si applicano algoritmi actor-critic quali **PPO** (\emph{Proximal Policy Optimization}).
    \item \textbf{Multi-Armed Bandit}: modelli come LinUCB governano scenari in cui la decisione di bilanciamento o allocazione della specifica taglia di istanza (es. rapporto ottimale tra GPU di tipo T4 versus A100 per un dato tipo di carico di inferenza) deve essere effettuata rapidamente con un bilanciamento continuo tra esplorazione e sfruttamento.
\end{itemize}

Lo stack applicativo sfrutta le estensioni dell'ecosistema Kubernetes:
\begin{itemize}
    \item \textbf{StableBaselines3}: libreria Python contenente implementazioni affidabili e testate di algoritmi RL utilizzata per addestrare i controllori;
    \item \textbf{KEDA} (\emph{Kubernetes Event-Driven Autoscaling}): modulo avanzato che permette l'attivazione e la scalatura dei pod in funzione del volume di code di messaggi asincrone (es. Kafka lag, RabbitMQ queues);
    \item \textbf{Vertical Pod Autoscaler} (\textbf{VPA}): demone di riconfigurazione automatica delle richieste di memoria e CPU dei singoli container;
    \item \textbf{Karpenter}: modulo open-source ad altissime prestazioni per Kubernetes progettato per monitorare i pod in stato di attesa (\emph{pending}), calcolando in millisecondi la combinazione ideale di istanze EC2/bare-metal da lanciare ed eliminando la latenza introdotta dai gruppi di nodi gestiti convenzionali.
\end{itemize}

% [IMMAGINE DEL FILE DI SLIDE CON NOME "CloudComputing_4"] Tabella: algoritmi di piazzamento carichi e architettura telemetrica di monitoraggio — Slide 19 / Pagina 19 del PDF

\subsection{Piazzamento Dei Carichi Di Lavoro E Schedulazione Delle Risorse}

Una volta scalato il numero di container o VM, è necessario decidere su quale nodo fisico allocare ciascuna entità (\emph{Workload Placement}). Questo problema ricade nella categoria dei problemi NP-difficili di \textbf{Multi-Dimensional Bin Packing}: ciascun nodo fisico rappresenta un contenitore (\emph{bin}) avente capacità finite su $d$ dimensioni (CPU, RAM, larghezza di banda PCIe, porte di rete), e ciascun task ha una richiesta vettoriale $\bm{r} = (r_1, r_2, \dots, r_d)$.

Gli scheduler moderni integrano:
\begin{itemize}
    \item \textbf{Algoritmi di Bin Packing euristici}: strategie come il \emph{Best-Fit Decreasing} multivariato, che cerca di minimizzare la varianza delle risorse residue per nodo al fine di prevenire colli di bottiglia;
    \item \textbf{Algoritmi Genetici e Simulated Annealing}: usati per calcolare periodicamente piani ottimi di deframmentazione e consolidamento del cluster (\emph{live migration});
    \item \textbf{Programmazione Lineare/Intera} (MILP): garantisce l'ottimalità dell'assegnazione considerando vincoli hard (es. anti-affinità per alta affidabilità);
    \item \textbf{Heuristics for GPU/CPU Affinity}: politiche che tengono conto della topologia hardware NUMA (\emph{Non-Uniform Memory Access}) e della gerarchia delle interconnessioni PCIe e switch NVLink, allocando i pod che cooperano sullo stesso nodo o sullo stesso socket per abbattere le latenze di trasferimento dei tensori.
\end{itemize}

La tempestività di reazione di questi controllori poggia su un'architettura telemetrica unificata che rileva le variazioni di carico:
\begin{itemize}
    \item \textbf{Prometheus}: estrazione e calcolo di serie temporali metriche (es. saturazione CPU, caduta di pacchetti);
    \item \textbf{Grafana}: visualizzazione e allarmistica su pannelli operativi;
    \item \textbf{Loki}: tracciamento log indicizzato tramite etichette, compatibile nativamente con la sintassi di Prometheus;
    \item \textbf{OpenTelemetry}: standard di propagazione di contesti e tracce per rilevare latenze anomale tra microservizi e pipeline ML.
\end{itemize}

\newpage

% ==============================================================================
% SEZIONE: ANALISI PREDITTIVA PER DECISIONI PROATTIVE
% ==============================================================================

\section{Analisi Predittiva Per Decisioni Proattive}

I modelli di analisi predittiva in esecuzione sull'infrastruttura cloud analizzano flussi di dati massivi in tempo reale per estrarre pattern, trend stocastici e deviazioni funzionali ben prima che gli operatori umani possano individuarli. Ciò consente di anticipare colli di bottiglia nelle catene logistiche, mitigare i picchi di consumo e isolare guasti tecnici prima che si verifichino interruzioni di servizio.

% [IMMAGINE DEL FILE DI SLIDE CON NOME "CloudComputing_4"] Obiettivi strategici dell'analisi predittiva e proattività operativa — Slide 20 / Pagina 20 del PDF

% [IMMAGINE DEL FILE DI SLIDE CON NOME "CloudComputing_4"] Tabella riassuntiva: algoritmi di Time-Series Forecasting e architetture Real-Time Data Pipelines — Slide 21 / Pagina 21 del PDF

\subsection{Algoritmi Di Previsione Delle Serie Temporali}

La modellazione matematica dell'evoluzione temporale delle metriche operative (richiesta di banda, volumi di transazioni, prezzi spot delle istanze cloud) si avvale di una vasta gerarchia di algoritmi:
\begin{itemize}
    \item \textbf{ARIMA e SARIMA}: formulazioni stazionarie per segnali con componenti autoregressive e di media mobile;
    \item \textbf{Livellamento Esponenziale di Holt-Winters}: tecnica statistica che cattura livello, trend lineare e stagionalità moltiplicativa o additiva tramite tre equazioni ricorsive aggiornate con costanti di attenuazione $(\alpha, \beta, \gamma)$:
    \begin{align}
        \ell_t &= \alpha (y_t - s_{t-m}) + (1 - \alpha)(\ell_{t-1} + b_{t-1}) \\
        b_t &= \beta (\ell_t - \ell_{t-1}) + (1 - \beta) b_{t-1} \\
        s_t &= \gamma (y_t - \ell_{t-1} - b_{t-1}) + (1 - \gamma) s_{t-m}
    \end{align}
    \item \textbf{Facebook Prophet}: modello di regressione additiva scalabile per dati con forti pattern stagionali e festività;
    \item \textbf{Reti Ricorrenti Avanzate} (\textbf{LSTM / GRU}): modelli neurali basati su porte di controllo (\emph{input, forget, output gates} nelle LSTM; \emph{reset, update gates} nelle GRU) progettati per memorizzare dipendenze a lungo termine superando i limiti numerici delle reti ricorrenti canoniche;
    \item \textbf{Temporal Convolutional Networks} (\textbf{TCN}): reti convoluzionali 1D con convoluzioni causali dilatate, che garantiscono invarianza temporale ed efficienza parallela;
    \item \textbf{Architetture Transformer per Serie Temporali}: modelli di ultima generazione specificamente adattati per superare la complessità computazionale quadratica $\mathcal{O}(L^2)$ del meccanismo di attenzione su sequenze temporali lunghe. Tra questi spiccano:
    \begin{itemize}
        \item \textbf{Informer}: introduce il meccanismo di \emph{ProbSparse Self-Attention} riducendo la complessità a $\mathcal{O}(L \ln L)$ selezionando le query dominanti;
        \item \textbf{Autoformer}: integra meccanismi di auto-correlazione a livello di sottosequenza al posto dell'attenzione puntiforme, decomponendo intrinsecamente il trend e le componenti stagionali;
        \item \textbf{TimesNet}: rimappa le serie 1D in tensori 2D per catturare variazioni multi-periodiche intra-periodo e inter-periodo tramite blocchi convoluzionali 2D;
        \item \textbf{PatchTST}: raggruppa le serie temporali in porzioni (\emph{patches}) di campioni contigui, riducendo la lunghezza effettiva della sequenza e preservando le relazioni semantiche locali con calcolo parallelo.
    \end{itemize}
\end{itemize}

\subsection{Pipeline Di Ingestion E Calcolo In Tempo Reale}

Questi modelli predittivi richiedono un'infrastruttura di alimentazione ad altissimo throughput e minima latenza:
\begin{itemize}
    \item \textbf{Apache Kafka}: spina dorsale per lo streaming distribuito di eventi basata su log commit immutabili partizionati, garantendo un disaccoppiamento robusto e tollerante ai guasti tra produttori di dati e consumatori;
    \item \textbf{Apache Flink}: motore di elaborazione di flussi nativo (\emph{true stream processing}) a latenze nell'ordine del millisecondo, dotato di gestione interna dello stato (\emph{stateful computing}) e supporto avanzato per \emph{Complex Event Processing} (CEP) e finestre temporali guidate dall'\emph{event-time};
    \item \textbf{Apache Spark Structured Streaming}: motore di elaborazione a micro-lotti (\emph{micro-batching}) che unifica le API di calcolo batch e streaming, facilitando l'integrazione di pipeline analitiche su vasta scala;
    \item \textbf{Apache Spark MLlib}: libreria di apprendimento automatico distribuito per clustering, classificazione e modelli previsionali su volumi di dati dell'ordine dei terabyte;
    \item \textbf{Ray}: framework distribuito orientato all'addestramento e al serving di modelli deep learning e RL a bassa latenza;
    \item \textbf{Dask}: libreria nativa Python per il parallelismo analitico flessibile su grafi computazionali diretti aciclici (DAG).
\end{itemize}

% [IMMAGINE DEL FILE DI SLIDE CON NOME "CloudComputing_4"] Tabella riassuntiva: algoritmi di Manutenzione Predittiva e librerie open-source per serie temporali — Slide 22 / Pagina 22 del PDF

\subsection{Applicazioni Di Manutenzione Predittiva}

Nel contesto dei data center e delle applicazioni industriali connesse al cloud, la manutenzione predittiva mira a calcolare il guasto imminente dei componenti hardware (dischi NVMe, alimentatori, ventole, banchi RAM):
\begin{itemize}
    \item \textbf{Alberi di Decisione Potenziati da Gradiente} (\textbf{XGBoost}, \textbf{LightGBM}, \textbf{Random Forest}): modelli supervisionati basati sull'aggregazione di alberi di regressione deboli, che elaborano metriche hardware discrete (\emph{S.M.A.R.T.} logs) per classificare la probabilità di rottura;
    \item \textbf{Stima della Remaining Useful Life} (\textbf{RUL}): regressione non lineare operata tramite modelli ibridi \textbf{CNN-LSTM}, in cui la componente convoluzionale estrae le correlazioni spaziali tra sensori multipli e la rete ricorrente modella il degrado temporale progressivo;
    \item \textbf{Modelli di Analisi di Sopravvivenza} (\emph{Survival Analysis}):
    \begin{itemize}
        \item \emph{Cox Proportional Hazards Model}: modella il tasso di rischio istantaneo di guasto $h(t | \bm{x})$ condizionato al vettore delle covariate $\bm{x}$:
        \begin{equation}
            h(t | \bm{x}) = h_0(t) \exp(\bm{\beta}^T \bm{x})
        \end{equation}
        dove $h_0(t)$ rappresenta il tasso di rischio di base;
        \item \emph{Random Survival Forests}: estensione non parametrica che integra foreste casuali per trattare dati censurati a destra, stimando curve di sopravvivenza affidabili per ciascun nodo hardware.
    \end{itemize}
\end{itemize}

Lo sviluppo di tali soluzioni poggia su un ecosistema unificato di librerie open-source: **GluonTS**, **Kats** (sviluppata da Meta), **Merlion** (da Salesforce per rilevamento anomalie e forecasting integrato), **PyTorch Forecasting**, **sktime** e **Darts**.

% [IMMAGINE DEL FILE DI SLIDE CON NOME "CloudComputing_4"] Tabella: algoritmi di ottimizzazione per Supply Chain e tecnologie di programmazione matematica — Slide 23 / Pagina 23 del PDF

\subsection{Ottimizzazione Della Supply Chain E Logistica Cloud}

L'allocazione delle scorte, la gestione del traffico di rete tra data center distribuiti geograficamente e l'instradamento logistico vengono formalizzati come problemi di programmazione matematica:
\begin{itemize}
    \item \textbf{MILP} (\emph{Mixed-Integer Linear Programming}): risolti tramite motori di computazione scientifica:
    \begin{itemize}
        \item \textbf{COIN-OR CBC} e \textbf{GLPK}: solutori aperti ampiamente diffusi per problemi di programmazione lineare ed intera;
        \item \textbf{SCIP}: solutore avanzato per problemi complessi di tipo MILP e MINLP (\emph{Mixed-Integer Non-Linear Programming});
        \item \textbf{Pyomo}: framework di modellazione algebrica avanzato in Python che interfaccia l'applicazione con molteplici solutori computazionali;
    \end{itemize}
    \item \textbf{Apprendimento per Rinforzo e Algoritmi Evolutivi}: ambienti complessi di simulazione logistica vengono risolti integrando ambienti standardizzati come **OpenAI Gymnasium**, supportati da framework di RL quali **Stable-Baselines3** e **RLlib** (in Ray), affiancati da librerie di calcolo genetico come **DEAP**, **PyGAD** e **Inspyred**.
\end{itemize}

\newpage

% ==============================================================================
% SEZIONE: ESPERIENZE UTENTE IPER-PERSONALIZZATE
% ==============================================================================

\section{Esperienze Utente Iper-Personalizzate}

L'esecuzione nel cloud consente di distribuire su scala globale piattaforme capaci di erogare contenuti, interfacce applicative e suggerimenti commerciali calibrati in tempo reale sul profilo e sulle interazioni di ogni singolo utente, senza introdurre ritardi o sovraccarichi architetturali sul backend.

% [IMMAGINE DEL FILE DI SLIDE CON NOME "CloudComputing_4"] Introduzione alla personalizzazione scalabile tramite modelli AI su cloud — Slide 24 / Pagina 24 del PDF

% [IMMAGINE DEL FILE DI SLIDE CON NOME "CloudComputing_4"] Tabella: Algoritmi di Recommendation System e framework specializzati — Slide 25 / Pagina 25 del PDF

\subsection{Framework E Reti Neurali Di Raccomandazione}

I moderni sistemi di raccomandazione hanno superato i modelli statici tabellari basati su scomposizione matriciale semplice, migrando verso architetture neurali ad alta complessità computazionale:
\begin{itemize}
    \item \textbf{Neural Collaborative Filtering} (\textbf{NCF}): sostituisce il prodotto scalare tra vettori latenti di utente e oggetto con una rete neurale \emph{feed-forward} multistrato, apprendendo funzioni di interazione non lineari complesse;
    \item \textbf{Deep Matrix Factorization}: integra rappresentazioni non lineari sia per descrivere le matrici di valutazione (\emph{rating}) sia per catturare informazioni testuali o visive associate agli elementi;
    \item \textbf{Wide \& Deep Networks}: architettura sviluppata da Google che combina una componente lineare (\emph{Wide}) deputata a memorizzare correlazioni storiche frequenti con una componente a strati profondi (\emph{Deep}) deputata a generalizzare verso nuove associazioni latenti;
    \item \textbf{DeepFM}: integra la potenza descrittiva delle \emph{Factorization Machines} per catturare interazioni di secondo ordine tra feature con una rete profonda senza richiedere feature engineering manuale;
    \item \textbf{Transformers per Raccomandazioni Sequenziali}:
    \begin{itemize}
        \item \textbf{SASRec} (\emph{Self-Attention based Sequential Recommendation}): impiega l'auto-attenzione unidirezionale per modellare la cronologia delle azioni dell'utente, dando il corretto peso sia a preferenze a lungo termine sia ad azioni recenti;
        \item \textbf{BERT4Rec}: adotta blocchi Transformer bidirezionali con meccanismi di mascheramento dei token (\emph{Cloze task}) per modellare il contesto complessivo della sessione.
    \end{itemize}
    \item \textbf{Raccomandazioni Contestuali}: algoritmi come \textbf{LinUCB} e \textbf{Thompson Sampling} formulano la selezione del contenuto come un problema di bandito contestuale, massimizzando il click-through rate bilanciando l'esplorazione di nuovi item con lo sfruttamento di quelli noti.
\end{itemize}

Tra i framework software dedicati si annoverano **TensorFlow Recommenders (TFRS)**, **Microsoft Recommenders**, **Surprise**, **Implicit** (specializzato nell'algoritmo di minimi quadrati alternati, \emph{Alternating Least Squares}, ALS, per matrici sparse) e **RecBole** (suite integrata con oltre 90 algoritmi allo stato dell'arte).

% [IMMAGINE DEL FILE DI SLIDE CON NOME "CloudComputing_4"] Tecnologie di inferenza a bassa latenza e segmentazione in tempo reale — Slide 26 / Pagina 26 del PDF

\subsection{Personalizzazione In Tempo Reale Ed Inferenza Su Microservizi}

La reattività dell'interfaccia utente richiede che le azioni eseguite nel corso della sessione corrente (click, scroll, visualizzazioni) modifichino istantaneamente i contenuti presentati. L'architettura prevede:
\begin{itemize}
    \item \textbf{Streaming Event Hubs}: pipeline basate su **Apache Kafka**, **Kafka Streams** e **Apache Flink** che intercettano gli eventi di navigazione e calcolano metriche di sessione in finestre temporali di pochi millisecondi;
    \item \textbf{Materialize}: motore SQL reattivo che mantiene viste materializzate continuamente aggiornate a ogni nuovo evento sul bus di streaming;
    \item \textbf{Piattaforme di Model Serving Distribuito}: motori specializzati nell'ospitare i modelli di inferenza come microservizi scalabili all'interno di Kubernetes:
    \begin{itemize}
        \item \textbf{Seldon Core} e \textbf{KServe}: orchestratori cloud-native per il ciclo di vita dell'inferenza, con supporto integrato per bilanciamento del traffico canarino (\emph{canary deployment}), trasformazione dei payload e rilevamento del \emph{drift} dei dati;
        \item \textbf{BentoML}: framework standardizzato per il packaging containerizzato di modelli ML ad alte prestazioni;
        \item \textbf{Ray Serve}: layer di erogazione a bassa latenza programmabile in Python nativo, progettato per gestire pipeline complesse di composizione tra più modelli.
    \end{itemize}
\end{itemize}

Parallelamente, per segmentare gli utenti in coorti omogenee, vengono eseguiti algoritmi di clustering non supervisionato: **K-Means / MiniBatch K-Means**, **DBSCAN / HDBSCAN**, modelli a mistura gaussiana (**Gaussian Mixture Models**, GMM), mappe auto-organizzanti (**Self-Organizing Maps**, SOM) e **Spectral Clustering**.

% [IMMAGINE DEL FILE DI SLIDE CON NOME "CloudComputing_4"] Tabella riassuntiva: Sequence modeling e banche dati vettoriali (Vector DBs) — Slide 27 / Pagina 27 del PDF

\subsection{Modellazione Del Comportamento Dell'Utente E Database Vettoriali}

La modellazione del comportamento temporale dell'utente sfrutta modelli sequenziali quali \textbf{LSTM}, \textbf{GRU}, \textbf{TCN}, \textbf{Hidden Markov Models} (HMM) e formulazioni bayesiane dell'utente, integrate con approcci canonici di \emph{Collaborative Filtering} user-user, item-item, matrici fattorizzate a valori singolari (**SVD**, **ALS**) e fattorizzazioni probabilistiche (**PMF**).

Per estrarre raccomandazioni e similarità semantica in pochi millisecondi su cataloghi di centinaia di milioni di prodotti, l'architettura si affida a **Vector Databases** e librerie di ricerca sui vicini prossimi approssimati (\emph{Approximate Nearest Neighbors}, ANN):
\begin{itemize}
    \item \textbf{FAISS} (sviluppato da Meta): libreria altamente ottimizzata per calcolo vettoriale parallelo massivo su GPU;
    \item \textbf{Annoy} (sviluppato da Spotify): libreria efficiente basata su alberi casuali di partizione dello spazio;
    \item \textbf{HNSWlib}: implementazione rapida basata su grafi navigabili a piccolo mondo gerarchici;
    \item \textbf{Database Vettoriali Distribuiti}: piattaforme cloud-native complete con supporto per persistenza, interrogazioni scalabili su cloud e filtri metadati avanzati, quali \textbf{Milvus}, \textbf{Qdrant} e \textbf{Weaviate}.
\end{itemize}

\newpage

% ==============================================================================
% SEZIONE: SICUREZZA AVANZATA E RILEVAMENTO DELLE MINACCE
% ==============================================================================

\section{Sicurezza Avanzata E Rilevamento Delle Minacce}

La protezione delle infrastrutture cloud richiede un cambio di paradigma rispetto alla sicurezza perimetrale tradizionale basata su firewall statici e ispezione a firme note. La natura dinamica dei container, delle comunicazioni est-ovest intra-cluster e l'emergere di attacchi condotti con tecniche evasive impongono l'adozione di modelli di intelligenza artificiale per il monitoraggio comportamentale in tempo reale e il rilevamento delle anomalie.

% [IMMAGINE DEL FILE DI SLIDE CON NOME "CloudComputing_4"] Principi base di sicurezza avanzata con modelli AI e rilevamento zero-day — Slide 28 / Pagina 28 del PDF

% [IMMAGINE DEL FILE DI SLIDE CON NOME "CloudComputing_4"] Tabella: Classificazione minacce, SIEM, IDS/IPS e strumenti cloud-native come Falco — Slide 29 / Pagina 29 del PDF

\subsection{Classificazione Delle Minacce E Sistemi Di Analisi}

L'identificazione di comportamenti ostili noti e sconosciuti sfrutta una combinazione di modelli supervisionati e di analisi sequenziale dei log di sistema:
\begin{itemize}
    \item \textbf{Modelli Ensemble ad Albero}: algoritmi di gradient boosting come \textbf{XGBoost}, \textbf{LightGBM} e \textbf{CatBoost}, uniti a foreste casuali (\textbf{Random Forests}), per classificare eventi di telemetria estraendo correlazioni complesse da tabelle con centinaia di indicatori;
    \item \textbf{Riconoscimento Convoluzionale dei Pattern nei Log}: tecniche \emph{log-as-image}, in cui sequenze di eventi di testo e matrici binarie di esecuzione vengono convertite in matrici bidimensionali processate tramite \textbf{Convolutional Neural Networks} (CNN) per isolare anomalie strutturali visive;
    \item \textbf{Modelli Sequenziali Ricorrenti}: reti \textbf{RNN} ed \textbf{LSTM} impiegate per tracciare la successione cronologica delle chiamate di sistema (\emph{syscalls}) effettuate dai processi, identificando alterazioni rispetto al flusso di esecuzione normale;
    \item \textbf{Architetture Transformer per la Correlazione dei Log}: modelli specializzati quali \textbf{LogBERT} e \textbf{AuditBERT}, che trattano gli eventi di log come sequenze lessicali e sfruttano meccanismi di attenzione per cogliere correlazioni a lungo raggio tra eventi registrati su nodi differenti.
\end{itemize}

Questi modelli operano all'interno di architetture di sicurezza aziendali stratificate:
\begin{itemize}
    \item \textbf{Piattaforme SIEM} (\emph{Security Information and Event Management}) e XDR: sistemi per l'ingestion, la conservazione e l'analisi forense centralizzata, quali **Wazuh**, **OpenSearch Security Analytics**, **AlienVault OSSIM** e **OSSEC**;
    \item \textbf{Orchestrazione della Risposta}: framework quali **TheHive** (gestione dei ticket di sicurezza) e **Cortex** (automazione dei workflow di indagine ed esecuzione di script di contenimento);
    \item \textbf{Sistemi IDS/IPS di Rete}: sonde di ispezione profonda dei pacchetti (\emph{Deep Packet Inspection}) ad altissima efficienza, quali **Suricata**, **Zeek** (precedentemente Bro) e **Snort**;
    \item \textbf{Sicurezza Cloud-Native a Livello Kernel}:
    \begin{itemize}
        \item \textbf{Falco} (progetto graduato CNCF): sfrutta la tecnologia del kernel Linux \textbf{eBPF} (\emph{Extended Berkeley Packet Filter}) per intercettare direttamente a livello kernel ogni chiamata di sistema generata all'interno dei container, correlandola con i metadati del pod in tempo reale;
        \item \textbf{Kubescape} e \textbf{KubeArmor}: agenti specializzati nell'applicare politiche di minimizzazione dei privilegi e profilazione dei profili AppArmor/SELinux;
        \item \textbf{ClamAV}: scansione e protezione continua contro payload virali e malware depositati sullo storage.
    \end{itemize}
\end{itemize}

% [IMMAGINE DEL FILE DI SLIDE CON NOME "CloudComputing_4"] Tabella: Rilevamento attacchi Zero-Day, Graph Analytics e ricerca semantica minacce — Slide 30 / Pagina 30 del PDF

\subsection{Rilevamento Degli Attacchi Zero-Day E Analisi Su Grafi}

Gli attacchi di tipo \emph{Zero-Day} sfruttano vulnerabilità non note, per le quali non esiste una firma predefinita. Il loro rilevamento si basa interamente su modelli non supervisionati:
\begin{itemize}
    \item \textbf{Autoencoders, LOF ed Isolation Forest}: calibrati unicamente sul traffico e sull'attività dei processi considerata nominale;
    \item \textbf{Apprendimento Auto-supervisionato}: modelli addestrati a ricostruire token mascherati all'interno di normali tracce di audit (\emph{masked token prediction}); quando si presenta un evento di attacco, la discrepanza tra l'output predetto dal modello e il reale flusso di eventi genera picchi di entropia immediatamente rilevabili;
    \item \textbf{Graph Neural Networks} (\textbf{GNN}) per la Lateral Movement Detection: l'infrastruttura di rete, i processi e le identità di accesso IAM vengono formalizzati come nodi di un grafo eterogeneo, e le connessioni, comunicazioni e chiamate RPC come archi pesati. Tramite architetture a passaggio di messaggi (\emph{Message Passing Neural Networks}) si estraggono embedding dei cammini di attacco (\emph{attack paths}), identificando pattern tipici del movimento laterale operato dagli attaccanti.
\end{itemize}

Le tecnologie utilizzate per l'elaborazione di grafi di sicurezza ad alta intensità computazionale includono **Neo4j** per l'interrogazione delle relazioni, **Graphistry** per l'ispezione visiva accelerata su GPU, **NetworkX** per analisi analitica algoritmica, e **RAPIDS cuGraph** per eseguire algoritmi su grafi direttamente in memoria GPU. Parallelamente, database vettoriali quali **FAISS**, **Milvus**, **Qdrant** e **Weaviate** consentono di confrontare semanticamente gli embedding degli eventi registrati con basi di conoscenza globali sulle minacce informatiche (\emph{Threat Intelligence}).

% [IMMAGINE DEL FILE DI SLIDE CON NOME "CloudComputing_4"] Tabella: Reinforcement Learning per risposta a intrusioni e modelli open-source di cybersecurity — Slide 31 / Pagina 31 del PDF

\subsection{Risposta Automatizzata Tramite Reinforcement Learning}

Una volta confermata la presenza di un'intrusione attiva, la latenza necessaria all'intervento umano per bloccare l'attaccante può risultare fatale. L'integrazione di agenti di **Reinforcement Learning** consente la risposta automatica e l'isolamento dinamico delle minacce:
\begin{itemize}
    \item \textbf{Modellazione dell'Ambiente di Difesa}: l'agente osserva lo stato di compromissione dei nodi e sceglie azioni quali l'isolamento di una rete tramite modifica dinamica dei \emph{NetworkPolicy} di Kubernetes, la revoca istantanea di token JWT/IAM, la migrazione del container o la limitazione di banda selettiva.
    \item \textbf{Algoritmi di Risposta}: vengono adottati approcci continui e discreti quali **Q-Learning**, **DQN**, **PPO** e **DDPG** (\emph{Deep Deterministic Policy Gradient}). Nei sistemi multi-cloud estesi si ricorre a formulazioni di **Multi-Agent Reinforcement Learning** (MARL), in cui agenti difensivi locali negoziano strategie cooperative globali.
\end{itemize}

L'ecosistema include modelli aperti pre-addestrati e framework dedicati:
\begin{itemize}
    \item \textbf{LogBERT} e \textbf{DeepLog}: modelli neurali basati rispettivamente su Transformer e LSTM per l'estrazione di anomalie strutturali dalle sequenze di log;
    \item \textbf{Drain3}: estrattore deterministico ad albero di template di log ad altissima velocità, che converte stringhe non strutturate di log in messaggi parametrici categorizzabili;
    \item \textbf{PyG} (\emph{PyTorch Geometric}) e \textbf{DGL} (\emph{Deep Graph Library}): librerie di riferimento per l'addestramento distribuito di reti GNN su topologie complesse;
    \item \textbf{OpenAI Whisper}: impiegato in contesti specialistici per la trascrizione e l'analisi automatica di log vocali provenienti da centralini e flussi operativi di supporto.
\end{itemize}

\newpage


\section{Accessibilità Dei Dati E Conversational Analytics}

L'integrazione di modelli di intelligenza artificiale all'interno del cloud trasforma la modalità di interazione con i grandi volumi di dati aziendali, democratizzando l'accesso alle informazioni. Gli utenti di qualsiasi reparto aziendale, privi di competenze avanzate di programmazione o di sintassi per basi di dati, possono interrogare archivi distribuiti complessi attraverso interfacce conversazionali avanzate e cruscotti interattivi.

% [IMMAGINE DEL FILE DI SLIDE CON NOME "CloudComputing_4"] Obiettivi di democratizzazione dei dati e abbattimento dei silos informativi — Slide 32 / Pagina 32 del PDF

Questo paradigma elimina la frammentazione informativa e i silos dipartimentali, permettendo decisioni operative tempestive basate su evidenze quantitative aggiornate in tempo reale.

% [IMMAGINE DEL FILE DI SLIDE CON NOME "CloudComputing_4"] Tabella riassuntiva: algoritmi di NLP, Text-to-SQL e modelli di Speech-to-Text — Slide 33 / Pagina 33 del PDF

\subsection{Elaborazione Del Linguaggio Naturale Per L'Accesso Ai Dati}

L'architettura alla base delle piattaforme di \emph{Conversational Analytics} si articola in una serie di moduli cooperanti:
\begin{itemize}
    \item \textbf{Classificazione dell'Intento}: categorizza la natura dell'interrogazione posta dall'operatore attraverso modelli Transformer compatti e veloci, quali **BERT**, **RoBERTa** e la sua controparte alleggerita **DistilBERT**;
    \item \textbf{Named Entity Recognition} (\textbf{NER}): estrae dall'interrogazione i parametri salienti (intervalli temporali, nomi di istanze, identificativi di pod, soglie di costo) utilizzando pipeline altamente ottimizzate come quelle fornite da **spaCy** o modelli specificamente calibrati come **BERT-NER**;
    \item \textbf{Parsing Semantico Avanzato}: trasforma la struttura dell'enunciato linguistico in una rappresentazione logico-formale intermedia o direttamente in un linguaggio di interrogazione strutturato, impiegando modelli sequenza-a-sequenza (\emph{seq2seq}), architetture basate su **T5** (\emph{Text-to-Text Transfer Transformer}) o modelli open di classe **LLaMA**.
\end{itemize}

\subsection{Modelli Text-To-SQL}

Il passaggio critico per estrarre informazioni da data warehouse relazionali (es. PostgreSQL, Snowflake, ClickHouse, Google BigQuery) risiede nella conversione sicura ed affidabile dell'interrogazione utente in sintassi SQL formalmente valida:
\begin{itemize}
    \item \textbf{Seq2SQL} e \textbf{SQLNet}: architetture neurali pionieristiche che hanno introdotto meccanismi di generazione modulare della clausola SQL per prevenire errori di ordine sintattico;
    \item \textbf{RAT-SQL} (\emph{Relation-Aware Transformer for Text-to-SQL}): codifica congiuntamente l'interrogazione formulata in linguaggio naturale e lo schema del database (tabelle, colonne, chiavi primarie e relazioni di chiave esterna), sfruttando una variante del meccanismo di auto-attenzione che apprende le correlazioni strutturali tra i campi;
    \item \textbf{PICARD} (\emph{Parsing Incrementally for Constrained Auto-Regressive Decoding}): algoritmo di decodifica incrementale vincolata che si interfaccia con modelli Transformer generativi; a ogni passo di generazione del token SQL successivo, PICARD esegue un parsing sintattico incrementale verificando la conformità grammaticale e la reale esistenza degli attributi rispetto allo schema del database, impedendo allucinazioni semantiche e la produzione di codice SQL non valido.
\end{itemize}

\subsection{Piattaforme Di Conversational Analytics Vocale}

A completamento dell'accessibilità conversazionale multimodale, le piattaforme cloud integrano stack di trascrizione e sintesi vocale operativi a bassa latenza:
\begin{itemize}
    \item \textbf{Coqui AI} e \textbf{Mozilla DeepSpeech}: implementazioni aperte e personalizzabili di algoritmi ASR (\emph{Automatic Speech Recognition}) basati su reti neurali end-to-end, idonee all'integrazione on-premise o su nodi dedicati per garantire piena sovranità sul dato;
    \item \textbf{Vosk}: toolkit leggero e altamente ottimizzato per il riconoscimento vocale continuo offline a bassissimo impatto computazionale, ideale per edge gateway o terminali operativi privi di connessione continuativa a banda larga.
\end{itemize}

L'unificazione di questi componenti all'interno del cloud computing trasforma l'accesso all'infrastruttura e all'informazione: l'interfaccia a riga di comando tradizionale (\emph{CLI}) e le configurazioni manuali cedono progressivamente il passo a un ecosistema reattivo, in cui l'infrastruttura comprende, ottimizza e protegge se stessa in modo completamente autonomo.

% DD --------------------------------------------------------------

\section{Accessibilità E Usabilità Avanzata Dei Dati}

L'accumulo di volumi eterogenei all'interno dei data lake aziendali pone una sfida cruciale: la maggior parte delle informazioni rimane non strutturata o inaccessibile al personale privo di competenze avanzate di programmazione e algebra relazionale. L'intersezione tra cloud computing ed elaborazione del linguaggio naturale (\emph{Natural Language Processing}, NLP) consente di superare questo divario mediante pipeline analitiche distribuite capaci di sintetizzare, interrogare e classificare i dati in modo autonomo.

% [IMMAGINE DEL FILE DI SLIDE CON NOME "CloudComputing_5"] Schematizzazione degli algoritmi di Data Summarization, framework Text-to-SQL e interfacce conversazionali — Slide 31 / Pagina 1 del PDF

\subsection{Sintesi Automatica, Spiegabilità Dei Modelli E Pipeline Text-To-SQL}

La generazione automatizzata di sommari narrativi ed evidenze quantitative si articola lungo due metodologie primarie:
\begin{itemize}
    \item \textbf{Sintesi Estrattiva} (\emph{Extractive Summarization}): seleziona e combina i frammenti più informativi estraendoli direttamente dal testo sorgente senza alterarne il lessico. L'approccio computazionale classico è rappresentato da \textbf{TextRank}, che modella il documento come un grafo non orientato $G = (V, E)$, dove i vertici $V$ corrispondono alle frasi e i pesi degli archi $w_{ij} \in E$ riflettono la similarità lessicale (es. tramite coefficiente di Jaccard o coseno tra embedding). Il punteggio di centralità $S(V_i)$ per ogni frase viene calcolato iterativamente applicando la formulazione di PageRank:
    \begin{equation}
        S(V_i) = (1 - d) + d \sum_{V_j \in \mathcal{N}(V_i)} \frac{w_{ji}}{\sum_{V_k \in \mathcal{N}(V_j)} w_{jk}} S(V_j)
    \end{equation}
    dove $d \in (0, 1)$ è il fattore di smorzamento (\emph{damping factor}, tipicamente $0.85$). In ambito deep learning, modelli come **BART-extractive** affiancano a questo principio classificatori binari sui singoli token o frasi.
    \item \textbf{Sintesi Astrattiva} (\emph{Abstractive Summarization}): genera frasi inedite sintetizzando il significato semantico complessivo mediante architetture encoder-decoder basate su Transformer, quali **T5** (\emph{Text-to-Text Transfer Transformer}), **BART** e **Pegasus** (ottimizzato tramite l'obiettivo di pre-addestramento basato su \emph{Gap Sentences Generation}).
\end{itemize}

L'adozione di decisioni automatizzate richiede garanzie di trasparenza, affrontate dai modelli di \textbf{Explainable AI} (XAI):
\begin{itemize}
    \item \textbf{SHAP} (\emph{SHapley Additive exPlanations}): quantifica il contributo marginale di ogni singola feature $i$ rispetto all'output predetto del modello, fondandosi sulla teoria dei giochi cooperativi e sui valori di Shapley:
    \begin{equation}
        \phi_i(x) = \sum_{S \subseteq F \setminus \{i\}} \frac{|S|!(|F| - |S| - 1)!}{|F|!} \left[ f_x(S \cup \{i\}) - f_x(S) \right]
    \end{equation}
    dove $F$ rappresenta l'insieme totale delle feature e $S$ un qualsiasi sottoinsieme che non include la feature $i$.
    \item \textbf{LIME} (\emph{Local Interpretable Model-agnostic Explanations}): approssima localmente il comportamento di un modello complesso non lineare addestrando un modello surrogato interpretabile $g \in G$ (es. un regressore lineare) nell'intorno del punto di predizione $x$:
    \begin{equation}
        \xi(x) = \arg\min_{g \in G} \mathcal{L}(f, g, \pi_x) + \Omega(g)
    \end{equation}
    dove $\pi_x(z)$ definisce la misura di prossimità del campione perturbato $z$ rispetto ad $x$, e $\Omega(g)$ penalizza l'eccessiva complessità del modello spiegatore.
\end{itemize}

Nel campo dell'estrazione delle caratteristiche, gli algoritmi di \textbf{AutoML} automatizzano la pulizia del dato, l'imputazione dei valori nulli, l'encoding categorico e la selezione delle combinazioni polinomiali ottimali mediante ricerca bayesiana o programmazione genetica.

Per quanto concerne l'interrogazione diretta dei database relazionali mediante linguaggio naturale (\textbf{Text-to-SQL}), l'infrastruttura sfrutta modelli a pesi aperti (\emph{Open-Source Software}, OSS) compatibili con le interfacce API OpenAI: **LLaMA**, **Mixtral** (modello basato su architettura \emph{Mixture of Experts}, MoE), **Falcon** e **StarCoder** (specializzato nell'elaborazione e traduzione di codice). Tali motori alimentano framework dedicati:
\begin{itemize}
    \item \textbf{SQLCoder}: modello open-source derivato da architetture LLM e sintonizzato specificamente per mappare interrogazioni arbitrarie su schemi complessi PostgreSQL o Snowflake;
    \item \textbf{RAT-SQL} (\emph{Relation-Aware Transformer for Text-to-SQL}): sfrutta matrici di attenzione arricchite con le relazioni semantiche tra colonne, tabelle e vincoli di chiave primaria ed esterna;
    \item \textbf{TypeSQL}: integra la tipizzazione statica dei dati direttamente nella rappresentazione dei token per prevenire filtri di confronto tra tipi incompatibili (es. stringhe su campi numerici);
    \item \textbf{Text2SQL}: famiglia di modelli e benchmark di riferimento della community open source.
\end{itemize}

A livello applicativo, l'orchestrazione delle pipeline semantiche impiega **Rasa** per la gestione deterministica dello stato della conversazione (\emph{dialogue management}), **Haystack** per pipeline di \emph{Question Answering} (QA) modulari, e orchestratori complessi come **LangChain** e **LlamaIndex** per strutturare agenti autonomi e flussi di interrogazione potenziati dal recupero contestuale.

% [IMMAGINE DEL FILE DI SLIDE CON NOME "CloudComputing_5"] Riconoscimento di pattern, riduzione della dimensionalità e cruscotti di Self-Service BI — Slide 32 / Pagina 2 del PDF

\subsection{Clustering, Riduzione Della Dimensionalità E Self-Service BI}

La scoperta non supervisionata di correlazioni all'interno di dataset multi-dimensionali supporta l'esplorazione analitica visiva:
\begin{itemize}
    \item \textbf{Clustering}: tecniche partizionali veloci come **K-Means** e la sua estensione per flussi massivi su larga scala **MiniBatch K-Means** (che aggiorna i centroidi campionando sottoinsiemi casuali), affiancate da algoritmi basati sulla densità come **DBSCAN** ed **HDBSCAN** (in grado di isolare cluster di forma arbitraria senza specificare a priori il parametro $k$), e modelli probabilistici come i **Gaussian Mixture Models** (GMM) risolti tramite l'algoritmo di \emph{Expectation-Maximization} (EM).
    \item \textbf{Riduzione della Dimensionalità}: per mappare spazi tensoriali ad alta dimensionalità $\mathbb{R}^D$ in coordinate bidimensionali $\mathbb{R}^2$ o tridimensionali $\mathbb{R}^3$ adatte all'ispezione visiva senza perdita di vicinanza locale:
    \begin{itemize}
        \item \textbf{t-SNE} (\emph{t-Distributed Stochastic Neighbor Embedding}): converte le distanze euclidee in probabilità condizionate gaussiane nello spazio originario e distribuzioni $t$ di Student a una coda pesante nello spazio a bassa dimensionalità, minimizzando la divergenza di Kullback-Leibler:
        \begin{equation}
            D_{\text{KL}}(P \,||\, Q) = \sum_{i \neq j} p_{ij} \log \frac{p_{ij}}{q_{ij}}
        \end{equation}
        \item \textbf{UMAP} (\emph{Uniform Manifold Approximation and Projection}): poggia su solide basi di geometria Riemanniana e topologia algebrica; modella i dati assumendo che siano campionati uniformemente da una varietà locale e ottimizza la corrispondenza tra i complessi simpliciali fuzzy mediante la minimizzazione dell'entropia incrociata, risultando scalabile computazionalmente e preservando le distanze globali meglio di t-SNE.
    \end{itemize}
    \item \textbf{Estrazione di Regole Associative}: algoritmi **Apriori** e **FP-Growth** (\emph{Frequent Pattern Growth}), che utilizzano alberi di pattern frequenti (\emph{FP-tree}) per calcolare indici di confidenza, supporto e \emph{lift} su miliardi di record transazionali senza generare candidati combinatori ridondanti.
    \item \textbf{Autoencoders}: reti neurali profonde addestrate ad apprendere codifiche latenti (\emph{embeddings}) compatte delle entità aziendali eliminando il rumore di fondo.
\end{itemize}

Questi modelli alimentano i moderni cruscotti di \textbf{Self-Service Business Intelligence} (BI): **Apache Superset** (piattaforma enterprise modulare), **Metabase** (interfaccia visuale orientata al linguaggio naturale), **Redash** (focalizzato sull'interrogazione SQL collaborativa) e **Grafana** (dedicato a metriche e serie temporali operative). Il rendering grafico interattivo sul browser client viene demandato a motori grafici scalabili quali **ECharts**, **Vega** (e la grammatica concisa **Vega-Lite**) e **D3.js**.

% [IMMAGINE DEL FILE DI SLIDE CON NOME "CloudComputing_5"] Motori OLAP ad altissime prestazioni e algoritmi di Insight Ranking — Slide 33 / Pagina 3 del PDF

\subsection{Motori Analitici Interattivi E Suggerimento Degli Insight}

L'esplorazione visiva interattiva richiede che le query analitiche eseguite su miliardi di tuple restituiscano risposte in frazioni di secondo. I database transazionali tradizionali (\emph{OLTP}) risultano strutturalmente inadeguati; lo stack cloud adotta pertanto motori \textbf{OLAP} (\emph{Online Analytical Processing}) orientati alla memorizzazione su disco e in RAM per colonne:
\begin{itemize}
    \item \textbf{ClickHouse}: database colonnare distribuito progettato per l'esecuzione parallela vettorializzata (istruzioni hardware SIMD), garantendo throughput analitici elevatissimi;
    \item \textbf{DuckDB}: motore analitico relazionale integrabile (\emph{in-process}), operante interamente in memoria locale, ideale per elaborazioni analitiche rapide senza l'infrastruttura di un demone server separato;
    \item \textbf{Apache Druid} e \textbf{Apache Pinot}: architetture distribuite real-time specializzate nell'ingestion diretta da code Kafka/Pulsar a bassa latenza, indicizzando i dati in segmenti immutabili per aggregazioni istantanee;
    \item \textbf{Trino} e **Presto**: motori di computazione SQL distribuita basati su architettura \emph{engine-only} (disaccoppiati dallo storage sottostante), in grado di federare interrogazioni parallele su fonti dati eterogenee (S3, Hive, PostgreSQL) in un unico schema logico.
\end{itemize}

Su questa base computazionale, algoritmi di \textbf{Insight Ranking} analizzano i cubi di dati per evidenziare anomalie rilevanti senza intervento umano (\emph{smart filters}): le regole associative (Apriori/FP-Growth) e il filtraggio collaborativo identificano schemi atipici; algoritmi di ranking su grafo derivati da **PageRank** ordinano le entità per influenza sistemica; infine, modelli linguistici avanzati sintetizzano il risultato traducendolo in descrizioni sintetiche per i decisori operativi.

\newpage

% ==============================================================================
% SEZIONE: POTENZIAMENTO DELLA PRODUTTIVITÀ OPERATIVA TRAMITE AI
% ==============================================================================

\section{Potenziamento Della Produttività Operativa Tramite AI}

L'integrazione di sistemi intelligenti all'interno dei processi operativi non si limita a snellire le attività ripetitive, ma introduce meccanismi di interazione avanzata (\emph{human-in-the-loop}) in grado di elaborare flussi documentali disomogenei e automatizzare la risoluzione di disservizi tecnici.

% [IMMAGINE DEL FILE DI SLIDE CON NOME "CloudComputing_5"] Panoramica sul potenziamento della produttività e automazione dei workflow — Slide 34 / Pagina 4 del PDF

% [IMMAGINE DEL FILE DI SLIDE CON NOME "CloudComputing_5"] Tabella riassuntiva: NLP, Conversational AI e modelli LLM open source — Slide 35 / Pagina 5 del PDF

\subsection{Elaborazione Del Linguaggio Naturale Ed Agenti Conversazionali}

Nei centri di supporto e nell'automazione del customer care, l'interazione testuale viene gestita da una catena modulare di modelli:
\begin{itemize}
    \item \textbf{Modelli Basati su Encoder} (**BERT**, **RoBERTa**, **DistilBERT**): pre-addestrati mediante l'obiettivo di modellazione linguistica mascherata (\emph{Masked Language Modeling}, MLM), ideali per compiti analitici quali l'estrazione di feature dense e la comprensione testuale profonda;
    \item \textbf{Modelli Generativi Autoregressivi} (\emph{Decoder-Only}): architetture come **LLaMA**, **Mixtral**, **Mistral**, **Falcon**, **MPT**, **GPT-J** e **GPT-NeoX**, ottimizzate per la produzione sequenziale di risposte coerenti mediante campionamento probabilistico della parola successiva condizionata al contesto:
    \begin{equation}
        P(w_1, w_2, \dots, w_T) = \prod_{t=1}^T P(w_t \mid w_1, \dots, w_{t-1})
    \end{equation}
    \item \textbf{Modelli Sequenza-a-Sequenza} (\emph{Seq2Seq}): modelli come **T5**, **BART** e la variante multilingue **mT5**, composti da un blocco encoder che processa l'intero input e un decoder autoregressivo che emette la traduzione o il riassunto strutturato.
\end{itemize}

A supporto di questi modelli operano sottomoduli specializzati:
\begin{itemize}
    \item \textbf{Classificazione dell'Intento}: mappatura delle espressioni dell'utente in classi semantiche di servizio mediante regressori logistici lineari (per compiti leggeri), reti CNN 1D, RNN o classificatori neurali sintonizzati (\emph{fine-tuned}) basati su BERT.
    \item \textbf{Named Entity Recognition} (\textbf{NER}): identificazione e taggatura di parametri critici (es. numeri di contratto, indirizzi IP, codici di errore). Storicamente formalizzata mediante campi casuali condizionali lineari (\textbf{CRF}, \emph{Conditional Random Fields}), i quali massimizzano la probabilità congiunta globale delle etichette della sequenza:
    \begin{equation}
        P(\bm{y}|\bm{x}) = \frac{1}{Z(\bm{x})} \exp\left( \sum_{t=1}^T \sum_{k} \lambda_k f_k(y_{t-1}, y_t, \bm{x}, t) \right)
    \end{equation}
    e oggi implementata tramite combinazioni **BiLSTM-CRF** (dove una rete neurale bidirezionale genera i vettori di emissione per lo strato CRF) o modelli dedicati **BERT-NER**.
    \item \textbf{Dialog Management}: la transizione tra stati della conversazione unisce macchine a stati a regole deterministiche con controllori di **Reinforcement Learning** (DQN, PPO), dove l'agente apprende a minimizzare il numero di passaggi per turno massimizzando il gradimento dell'utente.
\end{itemize}

% [IMMAGINE DEL FILE DI SLIDE CON NOME "CloudComputing_5"] Tabella riassuntiva: Machine Learning per Process Automation e Document Processing — Slide 36 / Pagina 6 del PDF

\subsection{Automazione Dei Processi E Document Processing Intelligente}

L'automazione dei processi aziendali (\emph{Process Automation}) sfrutta algoritmi predittivi per categorizzare ticket, smistarli dinamicamente alle unità competenti e intercettare anomalie prima che impattino gli utenti finali:
\begin{itemize}
    \item \textbf{Classificazione e Regressione}: modelli basati su alberi decisionali ad alte prestazioni (**Random Forest**, **XGBoost**, **LightGBM**) indirizzano i ticket di supporto in base a severità, metadati e storico di risoluzione;
    \item \textbf{Clustering}: **K-Means** e **DBSCAN** raggruppano automaticamente i disservizi simili emergenti, permettendo di identificare incidenti sistemici estesi a partire da molteplici segnalazioni isolate;
    \item \textbf{Previsione e Pianificazione della Capacità}: modelli di serie temporali (**Prophet**, **ARIMA**, **LSTM**) stimano l'afflusso di richieste attese per dimensionare dinamicamente i turni del personale di supporto o le quote di computing;
    \item \textbf{Gestione delle Eccezioni}: algoritmi di **Isolation Forest** e **LOF** rilevano deviazioni nei flussi operativi standard generate da bug procedurali o frodi transazionali.
\end{itemize}

Nell'elaborazione di documenti non strutturati (es. fatture scansionate, log di sistema compressi, moduli cartacei digitalizzati), il testo deve essere convertito in record strutturati interrogabili:
\begin{itemize}
    \item \textbf{OCR Classico ed Estrazione File}: **Tesseract OCR**, **PaddleOCR** (motore multilingue ad alta accuratezza basato su architetture deep learning compatte), **Apache Tika** (per l'estrazione agnostica di testo e metadati da oltre un migliaio di formati di file differenti) e **OCRmyPDF**;
    \item \textbf{Modelli Visuo-Linguistici Multimodali}:
    \begin{itemize}
        \item \textbf{LayoutLM}: modello che combina l'informazione lessicale dei token, le coordinate spaziali dei riquadri di delimitazione bidimensionali (\emph{2D bounding boxes}) e le caratteristiche visive dell'immagine sorgente, comprendendo la disposizione di tabelle, intestazioni e campi modulo;
        \item \textbf{Donut} (\emph{Document Understanding Transformer}): architettura avanzata completamente priva di motore OCR (\emph{OCR-free}), basata su un encoder visivo Swin Transformer e un decoder linguistico BART, che traduce direttamente i pixel del documento scansionato in alberi sintattici strutturati JSON;
        \item \textbf{TrOCR}: architettura sequenza-a-sequenza per il riconoscimento ottico dei caratteri end-to-end, basata su un Vision Transformer (ViT) come codificatore e un Transformer linguistico autoregressivo come decodificatore.
    \end{itemize}
\end{itemize}

% [IMMAGINE DEL FILE DI SLIDE CON NOME "CloudComputing_5"] Tabella riassuntiva: Computer Vision per automazione documentale e compiti fisici — Slide 37 / Pagina 7 del PDF

\subsection{Visione Artificiale Per Flussi Documentali E Compiti Fisici}

Nei casi in cui i compiti di automazione investono il mondo fisico (ispezione della componentistica fisica nei data center, controllo qualità nella logistica e monitoraggio visivo dei varchi di accesso), la visione artificiale adotta modelli computazionali dedicati:
\begin{itemize}
    \item \textbf{Reti Convoluzionali Profonde}: **ResNet** (basata su connessioni residue per abilitare l'addestramento stabile di centinaia di strati convoluzionali prevenendo la saturazione del gradiente), **EfficientNet** (che ottimizza congiuntamente profondità, larghezza e risoluzione spaziale mediante scalamento composto) e **MobileNet** (basata su convoluzioni separabili in profondità per dispositivi a risorse ristrette);
    \item \textbf{Vision Transformers}: modelli come **ViT** (che suddivide l'immagine in patch rettangolari proiettate come token sequenziali) e **DETR** (\emph{DEtection TRansformer}, che formula il rilevamento di oggetti come un problema di predizione insiemistica globale basata su accoppiamento bipartito);
    \item \textbf{Rilevamento Oggetti in Tempo Reale}: modelli della famiglia **YOLO** (\emph{You Only Look Once}, es. **YOLOv8**), capaci di inferenza a framerate elevati tramite una singola rete convoluzionale, affiancati da rilevatori a due stadi ad alta precisione come **Faster R-CNN**;
    \item \textbf{Riconoscimento di Testo Sequenziale}: modelli **CRNN** (\emph{Convolutional Recurrent Neural Network}) con funzioni di perdita basate su **CTC** (\emph{Connectionist Temporal Classification}), che allineano sequenze di lunghezze variabili senza richiedere annotazioni di coordinate per singolo carattere;
    \item \textbf{Librerie di Supporto}: framework di elaborazione grafica ad alta velocità come **OpenCV**, modelli per il tracciamento delle pose e dell'interazione corporea come **MediaPipe** e **OpenPose**, e framework di segmentazione semantica avanzati come **Detectron2**.
\end{itemize}

\newpage

% ==============================================================================
% SEZIONE: GESTIONE ROBUSTA DEI DATI SU SCALA CLOUD
% ==============================================================================

\section{Gestione Robusta Dei Dati Su Scala Cloud}

La convergenza tra la capacità di calcolo orizzontale del cloud e le tecniche di intelligenza artificiale permette di elaborare moli di dati eterogenee non strutturate (audio, flussi video, log applicativi, telemetria di rete) estraendone caratteristiche numeriche fruibili per decisioni strategiche. L'architettura dei data lake collaborativi garantisce che tali pipeline convergano verso un'unica sorgente di verità informativa (\emph{Single Source of Truth}, SSOT).

% [IMMAGINE DEL FILE DI SLIDE CON NOME "CloudComputing_5"] Panoramica sul Robust Data Management e creazione di collaborative data lake — Slide 38 / Pagina 8 del PDF

% [IMMAGINE DEL FILE DI SLIDE CON NOME "CloudComputing_5"] Tabella riassuntiva: estrazione di segnali da dati non strutturati e tecnologie NLP open source — Slide 39 / Pagina 9 del PDF

\subsection{Estrazione Di Segnali Da Dati Non Strutturati}

L'analisi di stringhe di testo non formattate (email di assistenza, dump di log, trascrizioni telefoniche) presuppone la conversione delle parole in rappresentazioni vettoriali dense:
\begin{itemize}
    \item \textbf{Modelli di Vettorizzazione Classici e Densità Lessicale}: **Word2Vec** (che apprende le rappresentazioni latenti massimizzando la probabilità delle parole di contesto tramite l'architettura Skip-Gram con campionamento negativo), **FastText** (che decompone le parole in n-grammi di caratteri, gestendo efficacemente vocaboli rari ed errori ortografici) e modelli avanzati basati su Transformer come **Sentence-Transformers** (SBERT), che impiegano reti siamesi per produrre embedding a livello di intera frase semanticamente confrontabili tramite distanza angolare.
    \item \textbf{Named Entity Recognition (NER)}: combinazioni ad alta fedeltà di modelli **BiLSTM-CRF** e modelli con pesi pre-addestrati sintonizzati su domini specifici (**BERT-NER**).
    \item \textbf{Topic Modeling}: estrazione automatizzata dei temi prevalenti all'interno di archivi documentali:
    \begin{itemize}
        \item \textbf{LDA} (\emph{Latent Dirichlet Allocation}): modello generativo probabilistico bayesiano in tre livelli, che assume che ogni documento sia modellato come una mistura di distribuzioni su argomenti latenti governata da una distribuzione di Dirichlet prioritaria;
        \item \textbf{HDP} (\emph{Hierarchical Dirichlet Process}): estensione non parametrica di LDA in cui il numero di argomenti non è fissato a priori ma viene inferito direttamente dalla complessità del corpus;
        \item \textbf{BERTopic}: approccio moderno che sfrutta embedding prodotti da modelli linguistici densi, riduce la dimensionalità dello spazio latente tramite UMAP, aggrega i vettori tramite HDBSCAN ed estrae le descrizioni degli argomenti calcolando varianti basate su TF-IDF della classe (\emph{c-TF-IDF}).
    \end{itemize}
    \item \textbf{Ricerca Semantica Densa} (\emph{Dense Retrieval}): individua i documenti rilevanti calcolando la similarità del coseno tra l'embedding dell'interrogazione e gli indici vettoriali dei passaggi archiviati, superando i limiti lessicali della corrispondenza esatta delle parole chiave tipica di indici invertiti come BM25.
\end{itemize}

Lo sviluppo di queste componenti poggia su framework consolidati: **spaCy** (ottimizzato per l'esecuzione industriale rapida di pipeline NLP), **NLTK** (orientato alla prototipazione accademica), **Hugging Face Transformers** (hub di modelli pre-addestrati allo stato dell'arte), **Gensim** (libreria storica specializzata in modellazione d'argomento su larga scala) e **Haystack**.

% [IMMAGINE DEL FILE DI SLIDE CON NOME "CloudComputing_5"] Tabella riassuntiva: Feature Engineering, Representation Learning e Pattern Discovery — Slide 40 / Pagina 10 del PDF

\subsection{Apprendimento Delle Rappresentazioni E Scoperta Di Pattern}

La modellazione di dati complessi richiede l'estrazione di rappresentazioni numeriche invarianti rispetto a trasformazioni non rilevanti dei dati d'ingresso:
\begin{itemize}
    \item \textbf{Apprendimento Contrastivo} (\emph{Contrastive Learning}): metodologie non supervisionate che addestrano una rete di estrazione a proiettare istanze generate dalla medesima sorgente (es. la stessa immagine sottoposta a ritaglio casuale e distorsione cromatica) vicine nello spazio latente, respingendo istanze provenienti da campioni differenti. Modelli di riferimento sono:
    \begin{itemize}
        \item \textbf{SimCLR}: minimizza la funzione di perdita InfoNCE normalizzata con temperatura:
        \begin{equation}
            \mathcal{L}_{i, j} = -\log \frac{\exp\left( \text{sim}(\bm{z}_i, \bm{z}_j) / \tau \right)}{\sum_{k=1}^{2N} \mathbb{I}_{[k \neq i]} \exp\left( \text{sim}(\bm{z}_i, \bm{z}_k) / \tau \right)}
        \end{equation}
        \item \textbf{MoCo} (\emph{Momentum Contrast}): affronta il problema dell'esaurimento della memoria GPU per grandi batch memorizzando le chiavi negative all'interno di una coda dinamica, aggiornando l'encoder delle chiavi tramite una media mobile esponenziale dei pesi (\emph{momentum update}) per garantire stabilità delle rappresentazioni.
    \end{itemize}
    \item \textbf{Analisi su Grafo e Community Detection}: estrazione di centralità informative mediante varianti dell'algoritmo **PageRank** e raggruppamento di entità correlate mediante l'algoritmo di **Louvain** (che massimizza la modularità del grafo partizionato).
\end{itemize}

L'infrastruttura software per queste elaborazioni combina **scikit-learn** (algoritmi canonici CPU-bound), **PyTorch** e **TensorFlow** (calcolo tensoriale accelerato su GPU), **NetworkX** (manipolazione di grafi in memoria locale), **Neo4j** (database a grafo nativo) e **RAPIDS cuML** (suite NVIDIA che reimplementa gli algoritmi classici di machine learning direttamente sui core CUDA e sulla memoria HBM delle schede acceleratrici).

% [IMMAGINE DEL FILE DI SLIDE CON NOME "CloudComputing_5"] Tabella riassuntiva: elaborazione vocale, audio e algoritmi avanzati di Computer Vision — Slide 41 / Pagina 11 del PDF

\subsection{Elaborazione Multimodale Di Flussi Audio E Video}

L'ingestion e l'analisi di dati multimediali generati dagli utenti impongono l'adozione di modelli specializzati:
\begin{itemize}
    \item \textbf{Elaborazione di Voce e Audio}:
    \begin{itemize}
        \item \emph{Automatic Speech Recognition} (ASR): modelli basati su decodifica probabilistica non allineata (**CTC**, \emph{Connectionist Temporal Classification}), architetture **RNN-Transducer** che disaccoppiano la trascrizione acustica dal modello linguistico predittivo, e architetture Transformer end-to-end massicciamente multilingue come **OpenAI Whisper**;
        \item \emph{Estrazione di Feature Acustiche}: estrazione dei coefficienti cepstrali su scala Mel (**MFCC**, \emph{Mel-Frequency Cepstral Coefficients}), che mappano le frequenze acustiche lineari su una scala logaritmica conforme alla percezione dell'orecchio umano;
        \item \emph{Diarizzazione dello Speaker} (\emph{Speaker Diarization}): segmentazione e raggruppamento temporale delle tracce vocali per rispondere alla domanda: \emph{``chi ha parlato e quando?''};
        \item \emph{Toolkit Specializzati}: **Kaldi** (framework storico ad alte prestazioni in C++ per ASR), **Vosk** (riconoscimento offline leggero per dispositivi edge), **Mozilla DeepSpeech** e **pyAudioAnalysis** (per analisi spettrale, classificazione timbrica e rilevamento delle emozioni).
    \end{itemize}
    \item \textbf{Visione Artificiale e Comprensione Video}:
    \begin{itemize}
        \item \emph{Comprensione dei Flussi Video}: modelli convoluzionali tridimensionali (**3D CNNs**), che integrano la dimensione temporale dei fotogrammi direttamente all'interno dei filtri di convoluzione spaziale, affiancati da architetture pure Transformer come **TimeSformer**, che disaccoppia l'attenzione temporale tra frame successivi dall'attenzione spaziale interna al singolo frame;
        \item \emph{Rilevamento e Tracciamento delle Pose}: librerie ad alte prestazioni come **OpenPose** e **MediaPipe** per estrarre in tempo reale i vettori scheletrici dei punti chiave corporei;
        \item \emph{Segmentazione di Oggetti}: framework modulari ad alta accuratezza come **Detectron2** di Meta.
    \end{itemize}
\end{itemize}

% [IMMAGINE DEL FILE DI SLIDE CON NOME "CloudComputing_5"] Tabella riassuntiva: algoritmi di Forecasting per Demand Planning e Trend Prediction — Slide 42 / Pagina 12 del PDF

% [IMMAGINE DEL FILE DI SLIDE CON NOME "CloudComputing_5"] Tabella riassuntiva: tecnologie per Cloud-Scale Data Lakes e piattaforme MLOps — Slide 43 / Pagina 13 del PDF

\subsection{Data Lakehouse Su Scala Cloud E Strumenti MLOps}

L'archiviazione e la gestione del ciclo di vita di moli di dati dell'ordine dei petabyte hanno portato alla convergenza tra i data lake (ad elevata scalabilità ma privi di garanzie transazionali) e i data warehouse relazionali (strutturati e prestazionali ma rigidi), dando origine all'architettura nota come \textbf{Data Lakehouse}:
\begin{itemize}
    \item \textbf{Formati File Colonnari Compressi}: standard aperti binari basati su frammentazione orizzontale in gruppi di righe (\emph{row groups}) e memorizzazione verticale per colonne compressa tramite codifica a dizionario e RLE (\emph{Run-Length Encoding}), quali **Apache Parquet** e **Apache ORC**;
    \item \textbf{Livello Metadati Transazionale per Lakehouse}:
    \begin{itemize}
        \item \textbf{Apache Iceberg}: struttura metadati ad albero gerarchico indipendente dal motore di calcolo, che supporta isolamento delle transazioni conforme al modello \emph{Snapshot Isolation}, evoluzione sicura dello schema e partizionamento nascosto (\emph{hidden partitioning});
        \item \textbf{Delta Lake}: estensione aperta che mantiene un log transazionale immutabile e sequenziale (\emph{DeltaLog}) basato su file JSON ordinati con checkpoint Parquet periodici, garantendo piene proprietà ACID;
        \item \textbf{Apache Hudi}: ottimizzato per pipeline di ingestion in tempo reale con aggiornamenti atomici frequenti (\emph{upsert}) tramite due tabelle operative distinte: \emph{Copy-on-Write} (massimizza le prestazioni in lettura) e \emph{Merge-on-Read} (minimizza la latenza in scrittura);
        \item \textbf{Apache Hive}: metastore di catalogo storico utilizzato per mappare file su tabelle relazionali logiche.
    \end{itemize}
    \item \textbf{Storage a Oggetti Cloud-Native}: protocolli di archiviazione distribuiti e disaccoppiati dalla computazione: **MinIO** (storage conforme alle API Amazon S3 a bassissima latenza per container), **Ceph** (file system e object store software-defined per cluster bare-metal) e lo storico **Apache Hadoop HDFS**.
\end{itemize}

A coronamento dell'infrastruttura operano i sistemi di orchestrazione del ciclo di vita dei modelli (**MLOps**):
\begin{itemize}
    \item \textbf{Kubeflow}: estensione nativa per Kubernetes deputata alla composizione, tracciamento e deployment scalabile di DAG di machine learning;
    \item \textbf{MLflow}: piattaforma per il logging sistematico dei parametri di addestramento, versionamento dei pesi dei modelli (\emph{Model Registry}) e riproducibilità degli esperimenti;
    \item \textbf{Orchestratori di Flussi}: motori di workflow basati su codice imperativo in Python come **Apache Airflow**, **Prefect** e **Dagster** (orientato nativamente al governo degli asset di dati anziché delle sole attività operative);
    \item \textbf{Federazione Dati e Monitoraggio}: motori computazionali come **Trino** e **Presto** che eseguono aggregazioni SQL inter-piattaforma, affiancati da **Grafana** per il monitoraggio continuativo del consumo di risorse e del degrado prestazionale dei modelli.
\end{itemize}

\newpage

% ==============================================================================
% SEZIONE: ACCESSIBILITÀ INCLUSIVA ED INTERFACCE CONVERSAZIONALI
% ==============================================================================

\section{Accessibilità Inclusiva Ed Interfacce Conversazionali}

La democratizzazione dell'accesso alle infrastrutture cloud richiede che i dati aziendali e i sistemi decisionali siano interrogabili da qualsiasi utente, a prescindere da disabilità fisiche o da specifiche competenze matematiche o ingegneristiche.

% [IMMAGINE DEL FILE DI SLIDE CON NOME "CloudComputing_5"] Principi base di accessibilità on-demand e strumenti di visual discovery — Slide 44 / Pagina 14 del PDF

% [IMMAGINE DEL FILE DI SLIDE CON NOME "CloudComputing_5"] Tabella riassuntiva: algoritmi di accessibilità visiva, XAI e speech-based accessibility — Slide 45 / Pagina 15 del PDF

\subsection{Accessibilità Visiva, Spiegabilità E Sintesi Vocale}

L'accessibilità visiva poggia su algoritmi capaci di tradurre informazioni numeriche complesse in rappresentazioni grafiche intuitive o in narrazioni testuali:
\begin{itemize}
    \item \textbf{Selezione Automatica dei Grafici} (\emph{Automated Chart Selection}): algoritmi euristici e classificatori supervisionati analizzano il tipo di dati degli attributi selezionati (categorico, ordinale, continuo temporale, geografico) e scelgono automaticamente la forma visuale più appropriata (es. istogrammi per distribuzioni, grafici a linee per serie storiche, mappe coroplete per coordinate spaziali);
    \item \textbf{Rilevamento di Tendenze e Spiegazione in Linguaggio Naturale}: moduli analitici identificano derive e anomalie all'interno dei dati, componendo brevi spiegazioni discorsive generate automaticamente;
    \item \textbf{Modelli di Spiegabilità} (**SHAP**, **LIME**): permettono agli operatori di comprendere quali variabili abbiano determinato una classificazione di rischio o una previsione d'insolvenza.
\end{itemize}

Per gli utenti con disabilità visiva o operanti in contesti a mani libere (\emph{hands-free}, come operatori di magazzino o chirurghi), l'interazione viene convertita interamente sul canale acustico:
\begin{itemize}
    \item \textbf{ASR} (\emph{Speech-to-Text}): trascrizione continua tramite **Whisper** o **Vosk**;
    \item \textbf{TTS} (\emph{Text-to-Speech}): sintesi vocale neurale ad alta fedeltà che converte il testo sintetizzato in forma d'onda acustica naturale. Piattaforme open di riferimento includono **Coqui TTS**, **Mozilla TTS** e **ESPnet** (toolkit end-to-end che unisce ASR, TTS, traduzione e arricchimento acustico all'interno di una medesima pipeline PyTorch).
\end{itemize}

% [IMMAGINE DEL FILE DI SLIDE CON NOME "CloudComputing_5"] Tabella riassuntiva: NLU per dati conversazionali, Answer Generation e RAG — Slide 46 / Pagina 16 del PDF

\subsection{Comprensione Del Linguaggio Naturale E Retrieval-Augmented Generation}

La costruzione di interfacce conversazionali avanzate si articola lungo due sottosistemi cooperanti:
\begin{enumerate}
    \item \textbf{Natural Language Understanding (NLU)}: decostruisce la frase utente determinandone l'intento tramite modelli encoder compatti (**BERT**, **RoBERTa**, **DistilBERT**), estrae i riferimenti a entità concrete mediante modelli **CRF**, **BiLSTM-CRF** o **BERT-NER**, ed esegue il parsing semantico formale (**Seq2Seq**, **T5**). Inoltre, proietta le interrogazioni in vettori latenti densi tramite **Sentence Transformers** per calcolare la correlazione semantica rispetto a basi di conoscenza indicizzate.
    \item \textbf{Conversational Reasoning e Answer Generation}: genera la risposta da presentare all'operatore sfruttando modelli linguistici generativi ad ampia scala (**LLaMA 3**, **Mixtral**, **Mistral**, **Falcon**, **GPT-J**, **GPT-NeoX**).
\end{enumerate}

Per garantire che le risposte fornite dagli LLM siano fattualmente accurate e non soffrano di allucinazioni sintattiche, l'architettura implementa il pattern \textbf{Retrieval-Augmented Generation} (RAG):
\begin{enumerate}
    \item L'interrogazione dell'utente viene codificata in un embedding denso;
    \item L'embedding interroga un database vettoriale distribuito estraendo i $K$ frammenti di documentazione interna aziendale più pertinenti;
    \item I frammenti recuperati vengono inseriti all'interno della finestra di contesto del prompt rivolto al modello linguistico, vincolando la generazione alla sola evidenza empirica fornita:
    \begin{equation}
        P(y \mid x) = \sum_{z \in \text{Top-}K} P(y \mid x, z) P(z \mid x)
    \end{equation}
\end{enumerate}

Sul versante delle interrogazioni a database relazionali, operano framework specializzati come **DB-GPT**, **SQLCoder**, **RAT-SQL** e i modelli della community **Text2SQL**, orchestrati da framework applicativi quali **LangChain**, **LlamaIndex**, **Haystack** e **Rasa**.

\newpage

% ==============================================================================
% SEZIONE: AUTOMAZIONE INTELLIGENTE ED AUTO-ADATTAMENTO NEL CLOUD
% ==============================================================================

\section{Automazione Intelligente Ed Auto-Adattamento Nel Cloud}

La gestione autonoma delle infrastrutture computazionali mira alla realizzazione di ambienti cloud \emph{self-healing} ed auto-ottimizzanti, capaci di correggere malfunzionamenti operativi, riallocare carichi degradati e aggiornare le proprie politiche decisionali sul campo senza richiedere la supervisione di ingegneri sistemisti.

% [IMMAGINE DEL FILE DI SLIDE CON NOME "CloudComputing_5"] Principi base di Intelligent Automation, auto-correzione e apprendimento continuo — Slide 47 / Pagina 17 del PDF

% [IMMAGINE DEL FILE DI SLIDE CON NOME "CloudComputing_5"] Tabella riassuntiva: algoritmi per Process Automation, Decision-Making e Self-Correcting RL — Slide 48 / Pagina 18 del PDF

\subsection{Processi Decisionali Di Processo Ed Apprendimento Per Rinforzo}

I motori di automazione dei processi integrano modelli probabilistici formali per governare transizioni decisionali complesse:
\begin{itemize}
    \item \textbf{Sistemi Basati su Regole}: motori inferenziali a catena diretta o inversa basati sull'algoritmo di Rete (quali **Drools** o piattaforme BPMN industriali come **Camunda**), affiancati da alberi di decisione statici;
    \item \textbf{Processi Decisionali Bayesiani}: aggiornano la distribuzione di probabilità a posteriori $P(\theta \mid D)$ relativa allo stato di salute dei microservizi incorporando progressivamente i dati di telemetria osservati;
    \item \textbf{Processi Decisionali di Markov} (\textbf{MDP}): definiscono la formalizzazione matematica fondamentale per gli agenti auto-correttivi, definiti dalla 5-upla $(\mathcal{S}, \mathcal{A}, \mathcal{P}, \mathcal{R}, \gamma)$, in cui:
    \begin{equation}
        \mathcal{P}_{ss'}^a = \mathbb{P}(S_{t+1} = s' \mid S_t = s, A_t = a)
    \end{equation}
    e l'obiettivo risiede nella ricerca della politica ottima $\pi^*(a \mid s)$ che massimizza il valore atteso della somma scontata dei premi futuri:
    \begin{equation}
        V^*(s) = \max_{a \in \mathcal{A}} \left[ \mathcal{R}(s, a) + \gamma \sum_{s' \in \mathcal{S}} \mathcal{P}_{ss'}^a V^*(s') \right]
    \end{equation}
    \item \textbf{Case-Based Reasoning} (CBR): risolve nuovi incidenti sistemici richiamando e adattando strategie adottate in contesti storici analoghi (\emph{retrieve, reuse, revise, retain}).
\end{itemize}

I framework per l'addestramento e il dispiegamento di agenti auto-adattivi includono:
\begin{itemize}
    \item Algoritmi a stima di valore e politica: **Q-Learning**, **DQN**, **PPO**, **Actor-Critic** sincroni e asincroni (**A2C / A3C**);
    \item **Multi-Agent Reinforcement Learning** (\textbf{MARL}): in cui molteplici agenti allocati su nodi distinti apprendono politiche cooperative o competitive;
    \item Ecosistemi di simulazione e sviluppo: **Ray RLlib** (esecuzione parallela su vasta scala), **Stable-Baselines3**, ambienti standardizzati **OpenAI Gymnasium** e framework specifici per la cooperazione multi-agente come **PettingZoo**;
    \item Motori industriali per l'orchestrazione deterministica di workflow a lunga durata: **Apache Airflow**, **Temporal** e **Netflix Conductor**.
\end{itemize}

% [IMMAGINE DEL FILE DI SLIDE CON NOME "CloudComputing_5"] Tabella riassuntiva: Continuous Learning, rilevamento del drift e AI per Robotic Process Automation — Slide 49 / Pagina 19 del PDF

\subsection{Apprendimento Continuo, Concept Drift E Robotic Process Automation}

In contesti cloud operativi, i dati del traffico e i comportamenti degli utenti non sono stazionari: le correlazioni statistiche mutano nel tempo. L'adattamento richiede algoritmi di **Continuous Learning**:
\begin{itemize}
    \item \textbf{Online and Incremental Learning}: modelli che aggiornano i propri parametri campione per campione senza dover riprocessare l'intero archivio storico, implementati in librerie come **River** (specializzata in machine learning per flussi di streaming in tempo reale), **Vowpal Wabbit** (motore di discesa del gradiente stocastica ottimizzato per vettori sparsi ad altissima efficienza) e **scikit-multiflow**;
    \item \textbf{Rilevamento del Concept Drift} (\emph{Concept Drift Detection}): formalmente, si ha una deriva quando la distribuzione congiunta al tempo $t$ e quella al tempo $t + \Delta$ differiscono:
    \begin{equation}
        P_t(\bm{x}, y) \neq P_{t+\Delta}(\bm{x}, y)
    \end{equation}
    Tale variazione può derivare da un cambiamento della sola distribuzione a priori delle feature $P(\bm{x})$ (\emph{covariate shift}), della distribuzione dei target $P(y)$ (\emph{prior probability shift}) o della relazione funzionale condizionata $P(y \mid \bm{x})$ (\emph{concept shift}). Strumenti aperti dedicati alla rilevazione continua del drift statistico sui dati di produzione includono **Alibi Detect**, **Evidently AI** e **WhyLogs**;
    \item \textbf{Meta-Learning e Transfer Learning}: consentono a modelli di apprendere nuovi task di classificazione su un numero ridotto di esempi (\emph{few-shot learning}) basandosi su basi neurali precedentemente regolarizzate.
\end{itemize}

A livello operativo di sistema, l'intelligenza artificiale arricchisce la **Robotic Process Automation** (RPA), che automatizzava in origine le sole interazioni deterministiche sulle interfacce grafiche utente (GUI). Con l'introduzione di OCR, NLP, visione artificiale e modelli di sequenza, il software bot apprende tramite \textbf{Imitation Learning} ad estrarre informazioni da documenti scansionati, compilare form complessi e prendere decisioni procedurali al variare del layout. Strumenti aperti di riferimento includono **Robocorp**, **OpenRPA**, **TagUI** e **Robot Framework**.

% [IMMAGINE DEL FILE DI SLIDE CON NOME "CloudComputing_5"] Tabella riassuntiva: Computer Vision per flussi documentali e strumenti di Workflow Orchestration — Slide 50 / Pagina 20 del PDF

\subsection{Visione Artificiale E Workflow Orchestration Nel Cloud}

La scansione visiva di flussi operativi automatizza la validazione di documenti di identità, la lettura di contratti cartacei o il monitoraggio continuo di impianti hardware nei data center. La combinazione di modelli convoluzionali (**CNNs**), **Vision Transformers**, modelli di rilevamento istantaneo (**YOLOv8**) e pipeline OCR complete (**OpenCV**, **Tesseract OCR**, **PaddleOCR**) viene inserita all'interno di motori di orchestrazione:
\begin{itemize}
    \item \textbf{Kubernetes}: fornisce il piano di orchestrazione nativo per allocare i container di visione e calcolo scalando le repliche in funzione del carico;
    \item \textbf{Argo Workflows}: motore di flussi di lavoro nativo Kubernetes basato su definizioni di DAG in formato dichiarativo YAML;
    \item \textbf{Apache Airflow}, **Prefect** e **Dagster**: piattaforme per il governo, il versionamento e l'allarmistica su pipeline di elaborazione batch ed asincrone;
    \item \textbf{n8n}: piattaforma modulare per l'automazione dei flussi a basso impiego di codice (\emph{low-code}).
\end{itemize}

I parametri di scheduling di questi workflow complessi vengono costantemente ottimizzati mediante solutori di programmazione matematica (lineare ed intera), algoritmi genetici, simulated annealing, ottimizzazione a sciame di particelle (\textbf{PSO}) e ricerca bayesiana.

\newpage

\section{Flessibilità Applicativa E Decision-Making Guidato Dai Dati}

L'elasticità del cloud offre alle organizzazioni la capacità di sviluppare, validare e rilasciare simultaneamente centinaia di modelli analitici eterogenei senza dover riconfigurare l'infrastruttura di base né investire in hardware dedicato per singola applicazione.

% [IMMAGINE DEL FILE DI SLIDE CON NOME "CloudComputing_5"] Introduzione alla flessibilità applicativa e scaling parallelo di modelli eterogenei — Slide 51 / Pagina 21 del PDF

% [IMMAGINE DEL FILE DI SLIDE CON NOME "CloudComputing_5"] Tabella riassuntiva: algoritmi e tecnologie per Backend Automation e Parallel Model Development — Slide 52 / Pagina 22 del PDF

\subsection{Sviluppo Ed Addestramento Parallelo Su Scala Cloud}

La sperimentazione concorrente su vasta scala si avvale di architetture distribuite che disaccoppiano le risorse fisiche dai framework di calcolo:
\begin{itemize}
    \item \textbf{Parallelismo Dati e Modello}: le sessioni di addestramento distribuito vengono gestite nativamente da librerie di calcolo ad alte prestazioni come **Horovod** (sviluppata da Uber per coordinare comunicazioni Ring-AllReduce via MPI) o **Ray Train**, scalando linearmente su decine di nodi GPU;
    \item \textbf{Orchestrazione su Kubernetes}: la gestione distribuita dei carichi di machine learning viene affidata al gestore di pacchetti **Helm** e ad operatori specializzati come **KubeRay** (per gestire cluster Ray nativi direttamente all'interno dell'ecosistema Kubernetes);
    \item \textbf{Ottimizzazione Concorrente degli Iperparametri e AutoML}: librerie come **Optuna** (che sfrutta l'algoritmo di Tree-structured Parzen Estimator, TPE, per la selezione bayesiana rapida) e **AutoGluon** (che costruisce ensemble multi-strato di modelli impilati massimizzando l'accuratezza senza sintonizzazione manuale);
    \item \textbf{Process Mining e Backend Automation}: analisi degli eventi registrati nei log dei sistemi informativi per ricostruire i grafi di flusso effettivi tramite algoritmi matematici come l'algoritmo Alpha e l'euristica di mining, confrontandoli con le specifiche di progetto per identificare colli di bottiglia o violazioni procedurali.
\end{itemize}

Gli strumenti a supporto spaziano da suite di ottimizzazione vincolata come **Google OR-Tools** e **Pyomo** a piattaforme di business rule come **Drools**, affiancate da **Ray RLlib**, **Stable-Baselines3**, **Temporal** e **Robot Framework**.

% [IMMAGINE DEL FILE DI SLIDE CON NOME "CloudComputing_5"] Concetti base del Data-Driven Decision Making e transizione da scelte empiriche ad evidenze oggettive — Slide 53 / Pagina 23 del PDF

% [IMMAGINE DEL FILE DI SLIDE CON NOME "CloudComputing_5"] Tabella riassuntiva: Data Ingestion, Streaming, Feature Intelligence e Feature Stores — Slide 54 / Pagina 24 del PDF

\subsection{Ingestion In Tempo Reale E Feature Store}

Il processo decisionale algoritmico fondato sui dati richiede la capacità di ingerire flussi continui di telemetria, trasformare i payload grezzi in feature numeriche pronte per l'inferenza e sincronizzare tali informazioni tra i sistemi di calcolo online e i repository di addestramento offline:
\begin{itemize}
    \item \textbf{Data Ingestion e Stream Processing}:
    \begin{itemize}
        \item **Apache Kafka** e **Apache Pulsar** (con architettura multi-tenant che disaccoppia la memorizzazione logica dei messaggi tramite Apache BookKeeper dalla gestione del broker di calcolo);
        \item **Apache Flink** e **Apache Spark Structured Streaming**: elaborano aggregazioni su finestre mobili (\emph{sliding windows}) e finestre temporali d'inattività (\emph{session windows}) gestendo eventi disordinati tramite meccanismi di \emph{watermarking};
        \item **Apache NiFi**: piattaforma visiva per il routing, la trasformazione e la mediazione automatizzata di flussi dati su larga scala.
    \end{itemize}
    \item \textbf{Feature Intelligence e Feature Store}:
    \begin{itemize}
        \item Il disallineamento tra le feature usate durante l'addestramento offline e quelle calcolate al volo durante l'inferenza in produzione (\emph{training-serving skew}) costituisce uno dei guasti architetturali più gravi nei sistemi ML. Tale rischio viene neutralizzato adottando un \textbf{Feature Store} centralizzato come **Feast**: esso gestisce un archivio a bassa latenza su database in-memory (es. Redis) per servire l'inferenza online in frazioni di millisecondo, affiancato da uno storage colonnare persistente (es. Parquet su S3) per consentire addestramenti storicamente coerenti (\emph{time-travel queries});
        \item Trasformazione delle feature dense tramite **scikit-learn**, pipeline tensoriali su **PyTorch / TensorFlow** e modelli di rappresentazione densa quali **Sentence-Transformers**.
    \end{itemize}
\end{itemize}

% [IMMAGINE DEL FILE DI SLIDE CON NOME "CloudComputing_5"] Tabella riassuntiva: Prescriptive Analytics, solutori matematici e Decision Refinement continuo — Slide 55 / Pagina 25 del PDF

\subsection{Analisi Prescrittiva Ed Apprendimento Ad Anello Chiuso}

Il ciclo di maturità analitica culmina nell'\textbf{Analisi Prescrittiva}: a differenza dell'analisi predittiva (che risponde a: \emph{``cosa accadrà?''}), l'analisi prescrittiva identifica la migliore azione operativa da intraprendere per massimizzare un obiettivo di business soggetto a vincoli di budget e capacità:
\begin{equation}
    \bm{a}^* = \arg\max_{\bm{a} \in \mathcal{A}_{\text{ammissibile}}} \mathbb{E}[U(\bm{x}, \bm{a})]
\end{equation}
dove l'insieme delle azioni ammissibili $\mathcal{A}_{\text{ammissibile}}$ è delimitato da disequazioni lineari o non lineari risolte tramite solutori matematici integrati (**OR-Tools**, **Pyomo**), agenti di apprendimento per rinforzo (**Ray RLlib**, **Stable-Baselines3**) o ottimizzatori di tipo black-box (**Optuna**).

La robustezza di queste raccomandazioni nel tempo viene preservata integrando meccanismi di \textbf{Feedback-Loop Learning}: l'esito reale delle decisioni intraprese viene continuamente reintrodotto all'interno del flusso d'addestramento, consentendo ad algoritmi incrementali (**River**, **Vowpal Wabbit**) e a sonde di monitoraggio statistico (**Alibi Detect**, **Evidently AI**, **WhyLogs**) di riallineare i pesi dei modelli al variare delle condizioni esterne.

\newpage

\section{Applicazioni Pratiche, Settori Industriali E Sfide Future}

L'impatto trasformativo dell'intelligenza artificiale integrata nel cloud attraversa tutti i comparti produttivi, ridefinendo il modello operativo dei settori industriali ad alta intensità di dati e ponendo complesse sfide di progettazione ingegneristica per il prossimo decennio.

% [IMMAGINE DEL FILE DI SLIDE CON NOME "CloudComputing_5"] Tassonomia delle 9 applicazioni primarie dell'AI nel Cloud Computing — Slide 56 / Pagina 26 del PDF

\subsection{Tassonomia Delle Applicazioni Cloud-AI}

Le principali categorie di applicazione si articolano lungo le seguenti direttrici operative:
\begin{enumerate}
    \item \textbf{Ottimizzazione delle Risorse}: allocazione predittiva ed elastica di capacità computazionale, memoria e connettività di rete per minimizzare il TCO garantendo gli SLA contrattuali.
    \item \textbf{Servizio Clienti e Supporto}: assistenti virtuali intelligenti operativi a ciclo continuo (24/7) capaci di gestire transazioni complesse e risolvere anomalie procedurali.
    \item \textbf{Analisi Predittiva}: modelli stocastici per stimare tendenze dei consumatori, flussi di cassa e necessità di manutenzione infrastrutturale.
    \item \textbf{Elaborazione del Linguaggio Naturale (NLP)}: pipeline multilingue per la traduzione istantanea, l'estrazione di sentiment da comunicazioni non strutturate e la sintesi di report.
    \item \textbf{Visione Artificiale}: sistemi di ispezione visiva automatica per il controllo qualità, la videosorveglianza dei data center e l'autenticazione biometrica.
    \item \textbf{Sistemi di Raccomandazione}: motori di ranking neurale per la personalizzazione istantanea delle interfacce e la massimizzazione delle conversioni commerciali.
    \item \textbf{Rilevamento delle Anomalie}: isolamento in tempo reale di violazioni di sicurezza, attacchi zero-day, frodi bancarie e derive dei parametri fisici dell'hardware.
    \item \textbf{Intelligenza Artificiale Generativa}: velocizzazione dello sviluppo software mediante modelli capaci di redigere script di configurazione, generare codice o comporre documentazione tecnica.
    \item \textbf{Business Intelligence Evoluta}: federazione di sorgenti dati interne ed esterne per alimentare cruscotti decisionali dinamici a supporto del management.
\end{enumerate}

% [IMMAGINE DEL FILE DI SLIDE CON NOME "CloudComputing_5"] Impatto dell'AI in ambito Healthcare: imaging medico, assistenti virtuali e medicina di precisione — Slide 57 / Pagina 27 del PDF

% [IMMAGINE DEL FILE DI SLIDE CON NOME "CloudComputing_5"] Impatto dell'AI nel settore Retail: chatbot 24/7, raccomandazioni e self-checkout visivo — Slide 58 / Pagina 28 del PDF

% [IMMAGINE DEL FILE DI SLIDE CON NOME "CloudComputing_5"] Impatto dell'AI nei Servizi Finanziari: sentiment analysis, robo-advisor e credit risk scoring — Slide 59 / Pagina 29 del PDF

% [IMMAGINE DEL FILE DI SLIDE CON NOME "CloudComputing_5"] Impatto dell'AI nel Manifatturiero: controllo qualità visivo, manutenzione predittiva e supply chain — Slide 60 / Pagina 30 del PDF

\subsection{Settori Industriali Di Riferimento}

La versatilità dell'infrastruttura si riflette in specifiche verticalizzazioni di settore:
\begin{itemize}
    \item \textbf{Sanità} (\emph{Healthcare}): elaborazione ad alte prestazioni di scansioni diagnostiche (tomografie, risonanze magnetiche) mediante reti convoluzionali profonde e Vision Transformers capaci di identificare patologie con accuratezza comparabile o superiore a quella specialistica; assistenti virtuali medici che incrociano lo storico anamnestico del paziente con basi di conoscenza cliniche per suggerire percorsi diagnostici; pipeline di elaborazione genomica su larga scala per la formulazione di terapie farmacologiche personalizzate (\emph{precision medicine}).
    \item \textbf{Commercio e Vendita al Dettaglio} (\emph{Retail}): motori di raccomandazione neurale a latenza ultra-bassa basati su database vettoriali per suggerimenti mirati; agenti conversazionali per la gestione di resi e tracciamento spedizioni; modelli di visione artificiale distribuiti per il monitoraggio delle scorte sugli scaffali, l'automazione dei pagamenti nei punti vendita senza casse (\emph{self-checkout}) e la rilevazione in tempo reale di furti o anomalie comportamentali.
    \item \textbf{Servizi Finanziari e Fintech}: elaborazione in tempo reale di notizie finanziarie e flussi social mediante sentiment analysis per generare segnali di trading quantitativo ad alta frequenza; identificazione istantanea di frodi creditizie analizzando le transazioni tramite modelli ensemble e grafi di correlazione; piattaforme di consulenza finanziaria automatizzata (\emph{robo-advisors}) per il bilanciamento dinamico dei portafogli; calcolo avanzato del merito creditizio (\emph{credit scoring}) integrando dati alternativi mediante regressori non lineari.
    \item \textbf{Settore Manifatturiero} (\emph{Manufacturing}): automazione del controllo qualità sui nastri trasportatori tramite telecamere industriali connesse a modelli YOLO o Faster R-CNN; abbattimento dei tempi di fermo macchina tramite modelli di vita utile residua (RUL) sui macchinari d'officina; ottimizzazione globale delle rotte di approvvigionamento e delle consistenze di magazzino mediante solutori di programmazione lineare intera e modelli previsionali su serie temporali.
\end{itemize}

% [IMMAGINE DEL FILE DI SLIDE CON NOME "CloudComputing_5"] Schema riassuntivo delle sfide architetturali aperte: costi, multi-cloud, sicurezza e portabilità — Slide 61 / Pagina 31 del PDF

\subsection{Sfide Architetturali E Tecnologiche Aperte}

Nonostante i vantaggi descritti, l'integrazione su larga scala dell'intelligenza artificiale nei sistemi cloud affronta cinque ordini di complessità strutturale:
\begin{enumerate}
    \item \textbf{Gestione dei Costi} (\emph{Cost Management}): l'addestramento e l'inferenza di modelli massivi con centinaia di miliardi di parametri generano costi operativi esponenziali in assenza di politiche rigorose di spegnimento e right-sizing.
    \item \textbf{Eterogeneità Multi-Cloud}: l'implementazione di pipeline analitiche distribuite tra molteplici fornitori infrastrutturali (\emph{multi-cloud}) incontra barriere insormontabili di interconnessione e costi di esportazione del traffico di rete (\emph{data egress fees}).
    \item \textbf{Prestazioni e Latenza Computazionale}: carichi complessi di inferenza visiva e linguistica richiedono un'allocazione hardware ottimizzata per prevenire colli di bottiglia nel trasferimento dei pesi tra bus PCIe, memorie HBM e interfacce di rete ad altissima velocità (es. InfiniBand o RoCE).
    \item \textbf{Portabilità dei Carichi di Lavoro} (\emph{Portability}): il rischio di dipendenza da un singolo fornitore (\emph{vendor lock-in}) imposto da API proprietarie e servizi AI chiavi in mano richiede la standardizzazione dello stack su contenitori aperti, runtime OCI e manifest dichiarativi Kubernetes.
    \item \textbf{Privacy e Sicurezza dei Dati}: il transito e l'elaborazione di dati confidenziali nei cluster condivisi (\emph{multi-tenancy}) impongono l'adozione di architetture di calcolo confidenziale (\emph{Confidential Computing}) basate su enclave hardware sicure a livello CPU (es. AMD SEV, Intel SGX) e metodologie di apprendimento federato.
\end{enumerate}

% [IMMAGINE DEL FILE DI SLIDE CON NOME "CloudComputing_5"] Prospettive evolutive: architetture AI-native, convergenza Edge-IoT, governance etica e calcolo quantistico — Slide 62 / Pagina 32 del PDF

\subsection{Prospettive Evolutive E Nuove Frontiere}

L'evoluzione della convergenza tra cloud ed intelligenza artificiale si delinea lungo sei vettori primari:
\begin{enumerate}
    \item \textbf{Architetture Cloud Native-AI}: i piani di controllo infrastrutturali non si limiteranno ad ospitare modelli intelligenti, ma integreranno nativamente algoritmi di machine learning per pilotare lo scheduling dei nodi, l'orchestrazione delle reti e la risoluzione automatizzata dei guasti senza supervisione umana.
    \item \textbf{Democratizzazione Totale delle Piattaforme}: la diffusione di framework \emph{low-code} e \emph{no-code} consentirà agli utenti operativi di configurare modelli analitici avanzati in modo autonomo, abbattendo la dipendenza specialistica dai team di data science.
    \item \textbf{Convergenza Tra Cloud, Edge Computing ed IoT}: per soddisfare requisiti di latenza nell'ordine del millisecondo o operare in contesti con connettività di rete intermittente (es. robotica autonoma, fabbriche intelligenti), i modelli verranno addestrati in modo scalabile nel cloud centrale ed eseguiti in locale sui nodi edge periferici.
    \item \textbf{Servizi Cloud Altamente Specializzati per Dominio}: transizione da modelli generici orizzontali verso piattaforme di settore sintonizzate su requisiti normativi e tassonomie verticali (es. sanità, aerospazio, finanza).
    \item \textbf{Governance, Trasparenza ed Etica dell'AI}: introduzione nativa nei control plane di strumenti per la conformità normativa, il tracciamento della provenienza dei dati (\emph{data lineage}), la certificazione dei dataset di addestramento e il rilevamento in tempo reale di distorsioni statistiche (\emph{bias}).
    \item \textbf{Carichi di Lavoro AI Predisposti per il Calcolo Quantistico}: l'integrazione di processori quantistici (\emph{Quantum Processing Units}, QPU) accessibili su piattaforma cloud amplierà le frontiere dell'apprendimento automatico potenziato da calcolo quantistico (\emph{Quantum Machine Learning}, QML), consentendo la risoluzione efficiente di problemi di ottimizzazione combinatoria e chimica computazionale intrattabili per le architetture binarie tradizionali.
\end{enumerate}

% DD --------------------------------------------------------------

\section{Scenari Di Rete Virtuale: Dal Livello 2 Al Livello 3}

\subsection{Richiamo Concettuale: Modello A Livelli E Reti Virtuali Private}
Nel modello di riferimento \textbf{OSI} (\emph{Open Systems Interconnection}) e nell'architettura \textbf{TCP/IP}, i livelli di collegamento (\emph{Data Link}, Livello 2) e di rete (\emph{Network}, Livello 3) assolvono a compiti profondamente distinti:
\begin{itemize}[noitemsep]
    \item Il \textbf{Livello 2} garantisce il recapito locale dei frame tra nodi fisicamente o logicamente adiacenti all'interno del medesimo \textbf{dominio di broadcast}, identificando le interfacce per mezzo di indirizzi fisici MAC (\emph{Media Access Control}).
    \item Il \textbf{Livello 3} realizza l'\textbf{instradamento gerarchico} (\emph{routing}) globale end-to-end su scala inter-rete attraverso indirizzi logici IP, disaccoppiando l'identità del terminale dalla topologia fisica immediata.
\end{itemize}

Quando un'azienda o un operatore necessita di connettere sedi geograficamente distanti, l'infrastruttura condivisa del fornitore di connettività (\emph{Service Provider}, SP) deve fornire canali logici isolati e riservati, noti come \textbf{Virtual Private Network (VPN)}. A seconda del livello architetturale in cui avviene la federazione dei siti remoti, si distinguono due paradigmi fondamentali: le reti di Livello 2 (\textbf{L2VPN}) e le reti di Livello 3 (\textbf{L3VPN}).

% [IMMAGINE DEL FILE DI SLIDE CON NOME "CloudComputing_5"] Copertina del corso "Virtual Networks and Cloud Networking - MPLS VxLAN EVPN" — Slide 1 / Pagina 1 del PDF

\subsection{Confronto Architetturale: L2VPN Contro L3VPN}
La scelta tra una soluzione L2VPN e una L3VPN definisce dove risiede l'intelligenza di instradamento e quale sia il confine di responsabilità tra cliente e fornitore del servizio.

\subsubsection{Reti Private Virtuali Di Livello 2 (L2VPN)}
In una \textbf{L2VPN}, il provider estende i segmenti di rete locale (VLAN) attraverso la propria rete di trasporto (\emph{metro} o \emph{backbone}), facendo apparire i dispositivi situati in siti geograficamente remoti come se fossero attestati sul medesimo segmento \emph{broadcast} locale.
Le caratteristiche distintive sono:
\begin{itemize}[noitemsep]
    \item \textbf{Trasparenza di routing}: il \emph{Service Provider} non partecipa né è a conoscenza delle tabelle o dei protocolli di instradamento del cliente (ad esempio OSPF, EIGRP o BGP interno). Il provider esamina e commuta unicamente i frame in base agli indirizzi MAC o a etichette logiche dedicate.
    \item \textbf{Casi d'uso ideali}: interconnessione trasparente di \emph{Data Center} (\emph{Data Center Interconnect}, DCI), reti metropolitane Ethernet (\emph{Metro Ethernet}) e federazione di cluster applicativi o nodi di elaborazione che impongono la condivisione della medesima sottorete IP (indirizzamento IP uniforme senza routing intermedio).
    \item \textbf{Esempi applicativi reali}: distribuzione di servizi a banda larga per utenti residenziali o aziendali (come Voice over IP - \emph{VoIP}, accesso Internet ad alta velocità - \emph{HSI}, e televisione broadcast - \emph{BTV}), attestati su gateway centralizzati quali i \textbf{BRAS} (\emph{Broadband Remote Access Server}).
\end{itemize}

\subsubsection{Reti Private Virtuali Di Livello 3 (L3VPN)}
In una \textbf{L3VPN}, l'interconnessione tra i siti avviene al livello IP. L'infrastruttura del fornitore partecipa direttamente al piano di instradamento del cliente:
\begin{itemize}[noitemsep]
    \item \textbf{Condivisione delle rotte}: i router di bordo del fornitore (\textbf{Provider Edge}, PE) instaurano adiacenze di routing (tipicamente tramite \textbf{BGP}) con i router periferici del cliente (\textbf{Customer Edge}, CE). I router PE gestiscono tabelle di routing dedicate e segregate per ciascun cliente per mezzo di istanze \textbf{VRF} (\emph{Virtual Routing and Forwarding}).
    \item \textbf{Scalabilità WAN}: solleva il cliente dall'onere di mantenere una magliatura complessa di collegamenti punto-punto; è la tecnologia cardine per le reti geografiche enterprise di grandi dimensioni con politiche di instradamento eterogenee.
\end{itemize}

\subsection{Tassonomia L2VPN: VPWS Contro VPLS}
All'interno del dominio L2VPN, le due principali implementazioni storiche su dorsale MPLS sono il servizio punto-punto \textbf{VPWS} e il servizio multipunto \textbf{VPLS}.

\subsubsection{VPWS (Virtual Private Wire Service)}
Il modello \textbf{VPWS} emula un circuito dedicato punto-punto (\emph{Point-to-Point}, P2P) ed è noto nello standard MEF (\emph{Metro Ethernet Forum}) come \textbf{E-Line}. Funziona come un "cavo virtuale" bidirezionale (\emph{pseudowire}) che collega esattamente due interfacce remote.
\begin{itemize}[noitemsep]
    \item \textbf{Gestione degli indirizzi MAC}: il dispositivo PE non effettua alcun apprendimento (\emph{MAC learning}). Poiché esiste una sola destinazione logica possibile all'altro capo dello pseudowire, qualsiasi frame Ethernet che entra nell'interfaccia locale viene incapsulato e trasmesso direttamente al PE remoto senza ispezionare la tabella dei MAC.
    \item \textbf{Topologia e complessità}: adotta una topologia a coppie (\emph{pairwise}), semplice da configurare, con un consumo computazionale e di memoria minimo sui router di bordo.
\end{itemize}

\subsubsection{VPLS (Virtual Private LAN Service)}
Il modello \textbf{VPLS} realizza una connettività punto-multipunto (\emph{P2MP}) o multipunto-multipunto (\emph{MP2MP}), nota come \textbf{E-LAN}. Emula a tutti gli effetti il comportamento di un commutatore Ethernet fisico distribuito su scala metropolitana o geografica.
\begin{itemize}[noitemsep]
    \item \textbf{Gestione degli indirizzi MAC}: i nodi PE devono apprendere dinamicamente gli indirizzi MAC sorgente associandoli alle rispettive pseudowire o interfacce locali, mantenendo una tabella di commutazione (\emph{forwarding database}) per decidere verso quale pseudowire inoltrare ciascun frame unicast, e replicando i frame \emph{Broadcast, Unknown Unicast e Multicast} (BUM).
    \item \textbf{Topologia e complessità}: per prevenire loop senza ricorrere a Spanning Tree sopra la dorsale, VPLS richiede una topologia completamente magliata (\emph{full-mesh}) di pseudowire tra tutti i router PE coinvolti, con regole di \emph{split horizon} (un frame ricevuto da uno pseudowire non viene mai re-inoltrato su un altro pseudowire). Per ovviare ai limiti di scalabilità del full-mesh, si ricorre all'architettura gerarchica \textbf{H-VPLS} (\emph{Hierarchical VPLS}), che introduce router di aggregazione intermedi a scapito di una maggiore complessità gestionale e di consumo di risorse hardware.
\end{itemize}

\subsection{Evoluzione Delle Architetture Di Protocollo}
Per superare la rigidità operativa e i limiti di scalabilità delle architetture native di commutazione e routing, l'ingegneria delle reti ha seguito una precisa traiettoria evolutiva:
\begin{enumerate}[noitemsep]
    \item \textbf{MultiProtocol Label Switching (MPLS)}: sostituzione dell'ispezione della tabella di routing IP con etichette corte a lunghezza fissa, integrando meccanismi di ingegneria del traffico e supporto a servizi VPN L2/L3.
    \item \textbf{Virtual eXtensible LAN (VxLAN)}: incapsulamento di livello 2 su reti IP di livello 3 (\emph{MAC-in-UDP}) per risolvere le criticità di scala nei data center virtualizzati.
    \item \textbf{BGP $\rightarrow$ MultiProtocol BGP (MP-BGP) $\rightarrow$ Ethernet VPN (EVPN)}: unificazione del piano di controllo attraverso BGP per distribuire informazioni di raggiungibilità sia MAC (L2) che IP (L3), superando definitivamente l'apprendimento basato sul \emph{data-plane flooding}.
\end{enumerate}

\newpage

\section{MultiProtocol Label Switching (MPLS): Architettura E Piani Di Funzionamento}

\subsection{La Genesi Di MPLS E Il Disaccoppiamento Dei Piani}
Nelle reti IP tradizionali, ogni router lungo il cammino analizza l'intestazione di Livello 3 di ciascun pacchetto in transito ed esegue una ricerca nella tabella di instradamento basata sul principio del \textbf{longest prefix match} (LPM). Storicamente, tale operazione risultava computazionalmente onerosa per le CPU dei router.

Sebbene l'evoluzione dei circuiti dedicati \textbf{ASIC} (\emph{Application-Specific Integrated Circuit}) abbia in seguito azzerato il collo di bottiglia della ricerca dei prefissi IP, MPLS ha mantenuto e ampliato il proprio valore strategico grazie alla convergenza tra la flessibilità del Livello 3 e l'efficienza della commutazione di Livello 2.

L'architettura MPLS poggia sulla rigida separazione tra due entità logiche:
\begin{itemize}[noitemsep]
    \item \textbf{Piano di Controllo (\emph{Control Plane})}: opera in modalità priva di connessione (\emph{connectionless}). Ha il compito di scambiare le informazioni di raggiungibilità IP tramite protocolli di routing interni (OSPF, IS-IS) o esterni (BGP) e di distribuire le etichette per mezzo di protocolli dedicati come \textbf{LDP} (\emph{Label Distribution Protocol}) o estensioni di RSVP/BGP.
    \item \textbf{Piano di Inoltro Dati (\emph{Data Plane})}: opera in modalità orientata alla connessione (\emph{connection-oriented}). Inoltra i pacchetti lungo percorsi virtuali prefissati denominati \textbf{LSP} (\emph{Label Switched Path}), basandosi unicamente su etichette numeriche anteposte al pacchetto, senza interrogare l'indirizzo IP di destinazione nei nodi interni di transito.
\end{itemize}

% [IMMAGINE DEL FILE DI SLIDE CON NOME "CloudComputing_5"] Diagramma del principio di inoltro MPLS (RTA, RTB, RTC, RTD con etichette 1024, 1029, 1039) — Slide 5 / Pagina 5 del PDF

\subsection{Elementi Della Topologia MPLS: LSR, LER Ed LSP}
All'interno di un dominio MPLS si distinguono ruoli differenziati per gli apparati di rete:
\begin{itemize}[noitemsep]
    \item \textbf{Label Edge Router (LER)}: router posizionato al perimetro del dominio MPLS. Si interfaccia da un lato con la rete IP convenzionale e dall'altro con la dorsale MPLS. Si distingue in:
    \begin{itemize}[noitemsep]
        \item \emph{Ingress LER}: riceve pacchetti IP puri, classifica il flusso in una specifica classe di equivalenza, appone l'etichetta MPLS (\emph{Push}) e immette il pacchetto nel dominio.
        \item \emph{Egress LER}: riceve pacchetti etichettati dall'interno del dominio, rimuove l'etichetta (\emph{Pop}) ed esegue il tradizionale instradamento IP verso la rete di destinazione.
    \end{itemize}
    \item \textbf{Label Switch Router (LSR)}: router o commutatore collocato nel nucleo (\emph{core}) del dominio MPLS. Non ispeziona l'header IP: riceve il pacchetto etichettato, consulta una tabella di commutazione rapida, sostituisce l'etichetta in ingresso con una nuova etichetta in uscita (\emph{Swap}) e lo trasmette al salto successivo.
    \item \textbf{Label Switched Path (LSP)}: la sequenza unidirezionale di nodi (LER di ingresso, uno o più LSR, LER di uscita) attraversata da pacchetti appartenenti alla medesima classe, stabilita preventivamente mediante i protocolli del piano di controllo.
\end{itemize}

% [IMMAGINE DEL FILE DI SLIDE CON NOME "CloudComputing_5"] Modello di rete MPLS: distinzione tra LER al bordo e LSR nel core con percorsi LSP — Slide 8 / Pagina 8 del PDF

\subsection{Ambiti Applicativi Di MPLS}
L'astrazione offerta dalla commutazione ad etichette ha reso MPLS il pilastro delle dorsali dei carrier di telecomunicazione per tre scopi principali:

\subsubsection{1. MPLS VPN}
Permette di costruire reti private su scala geografica aggregando sedi distribuite. Nella terminologia di rete:
\begin{itemize}[noitemsep]
    \item I router del cliente sono denominati \textbf{CE} (\emph{Customer Edge}) e non sono consapevoli dell'infrastruttura MPLS.
    \item I router del provider collegati ai CE sono i \textbf{PE} (\emph{Provider Edge}, equivalenti ai LER).
    \item I router di transito del provider sono denominati \textbf{P} (\emph{Provider}, equivalenti agli LSR), i quali commutano esclusivamente in base a etichette esterne senza mantenere alcuna rotta relativa alle VPN dei clienti.
\end{itemize}

% [IMMAGINE DEL FILE DI SLIDE CON NOME "CloudComputing_5"] Architettura MPLS VPN con router CE, PE e router di transito P — Slide 6 / Pagina 6 del PDF

\subsubsection{2. MPLS Traffic Engineering (MPLS TE)}
Nel routing IP ordinario basato su metrica IGP (ad esempio Shortest Path First), tutto il traffico converge lungo il percorso a costo minimo, portando a congestione sui collegamenti ottimali e sottoutilizzo dei cammini alternativi.

Con \textbf{MPLS TE}, la rete stabilisce tunnel LSP espliciti e vincolati (\emph{Constraint-based Routing}), riservando risorse di banda lungo percorsi specifici. Ad esempio, una frazione del traffico ($70\%$) può essere indirizzata lungo un percorso geografico dorsale ad alta capacità (es. RTB-RTC-RTD), mentre la quota rimanente ($30\%$) viene deviata su un cammino secondario (es. RTB-RTG-RTH-RTD), bilanciando dinamicamente il carico ed evitando nodi critici.

% [IMMAGINE DEL FILE DI SLIDE CON NOME "CloudComputing_5"] Applicazione di Traffic Engineering (MPLS TE) con ripartizione del traffico (70% e 30%) su percorsi differenziati — Slide 7 / Pagina 7 del PDF

\subsubsection{3. Integrazione Multiservizio E Qualità Del Servizio (QoS)}
MPLS consente di incapsulare nativamente frame provenienti da tecnologie eterogenee di Livello 2 (quali ATM, Frame Relay, Ethernet), mappando i campi di priorità del traffico direttamente nell'etichetta MPLS per garantire garanzie stringenti di latenza e banda passante.

\newpage

\section{Struttura Del Pacchetto E Meccanismi Di Forwarding MPLS}

\subsection{Incapsulamento In Modalità Frame E Cell}
Lo standard MPLS definisce due modalità di incapsulamento:
\begin{itemize}[noitemsep]
    \item \textbf{Cell Mode}: concepita per reti ATM (\emph{Asynchronous Transfer Mode}), dove l'etichetta viene mappata direttamente nei campi VPI/VCI della cella.
    \item \textbf{Frame Mode}: utilizzata nelle reti a commutazione di pacchetto contemporanee (Ethernet, PPP). In questa modalità, l'intestazione MPLS viene inserita come un \emph{shim header} direttamente tra l'intestazione di Livello 2 (MAC/Frame Header) e l'intestazione di Livello 3 (IP Header).
\end{itemize}

% [IMMAGINE DEL FILE DI SLIDE CON NOME "CloudComputing_5"] Schema di incapsulamento Frame Mode MPLS (shim header inserito tra Frame Header e IP Header) — Slide 10 / Pagina 10 del PDF

\subsection{Anatomia Dell'Intestazione MPLS}
L'intestazione MPLS ha una dimensione fissa di \textbf{4 byte} (ovvero \textbf{32 bit}) ed è strutturata nei seguenti quattro campi:

% [IMMAGINE DEL FILE DI SLIDE CON NOME "CloudComputing_5"] Struttura a 32 bit dell'intestazione MPLS (Label, EXP, S, TTL) — Slide 11 / Pagina 11 del PDF

\begin{table}[H]
\centering
\small
\begin{tabularx}{\textwidth}{l c X}
\toprule
\textbf{Campo} & \textbf{Lunghezza} & \textbf{Descrizione Funzionale} \\
\midrule
\textbf{LABEL} & 20 bit & Valore numerico dell'etichetta utilizzato per la decisione di commutazione. Permette $2^{20} \approx 10^6$ valori possibili. \\
\textbf{EXP}   & 3 bit  & Campo sperimentale (\emph{Experimental Use}), comunemente impiegato per la \textbf{Class of Service (CoS)} e la mappatura dei bit IP Precedence / DiffServ (QoS). \\
\textbf{S}     & 1 bit  & \textbf{Bottom of Stack}. Posto a $1$ se l'etichetta corrente è l'ultima dello stack (ossia il payload successivo è l'header IP o L2 nativo); posto a $0$ se seguono ulteriori etichette annidate. \\
\textbf{TTL}   & 8 bit  & \textbf{Time to Live}. Funzione analoga al campo TTL del protocollo IP; decrementato ad ogni transito di nodo per prevenire loop infiniti di pacchetti nel dominio. \\
\bottomrule
\end{tabularx}
\caption{Campi componenti l'header MPLS a 32 bit.}
\end{table}

\subsection{Annidamento Delle Etichette (Label Stacking) E Protocol Identifier}
MPLS supporta l'impilamento gerarchico di etichette multiple (\textbf{Label Nesting} o \emph{Stacking}). Il ricevente riconosce la natura del carico utile grazie al campo \textbf{PID} (\emph{Protocol Identifier} o \emph{EtherType}) situato nell'intestazione di Livello 2:
\begin{itemize}[noitemsep]
    \item Su reti \textbf{Ethernet}:
    \begin{itemize}[noitemsep]
        \item \texttt{0x0800}: Pacchetto IPv4 nativo non etichettato.
        \item \texttt{0x8847}: Pacchetto \textbf{MPLS Unicast}.
        \item \texttt{0x8848}: Pacchetto \textbf{MPLS Multicast}.
        \item \texttt{0x8100}: Frame Ethernet con tag VLAN IEEE 802.1Q.
    \end{itemize}
    \item Su collegamenti \textbf{Point-to-Point (PPP)}:
    \begin{itemize}[noitemsep]
        \item \texttt{0x8021}: IPv4.
        \item \texttt{0x8281}: MPLS Unicast.
        \item \texttt{0x8283}: MPLS Multicast.
    \end{itemize}
\end{itemize}

Quando sono presenti più etichette, il commutatore esamina e modifica unicamente l'etichetta in cima allo stack (\emph{outer label}).
\begin{itemize}[noitemsep]
    \item In \textbf{MPLS L3VPN}, si impiegano solitamente due livelli di etichette: l'etichetta esterna (\emph{IGP/Tunnel Label}) instrada il pacchetto attraverso la rete pubblica del provider fino al PE corretto, mentre l'etichetta interna (\emph{VPN/Service Label}) segnala al PE di destinazione a quale istanza VRF o interfaccia cliente appartiene il flusso.
    \item In \textbf{MPLS TE}, l'etichetta esterna definisce il tunnel con vincoli di traffico, mentre l'etichetta interna indirizza il payload verso il terminale o il servizio finale.
\end{itemize}

% [IMMAGINE DEL FILE DI SLIDE CON NOME "CloudComputing_5"] Schema di annidamento MPLS Label Nesting con S-bit ed EtherType di Livello 2 — Slide 12 / Pagina 12 del PDF

\subsection{Strutture Dati: FEC, NHLFE, FTN Ed ILM}
La commutazione MPLS poggia su tabelle e costrutti concettuali ben definiti:

\begin{theorem}
Una \textbf{Forwarding Equivalence Class (FEC)} rappresenta un insieme di pacchetti IP che condividono caratteristiche comuni (es. medesimo prefisso di destinazione, classe di servizio, tipo di traffico) e che, pertanto, vengono trattati e instradati in modo identico lungo il medesimo LSP con la medesima etichetta.
\end{theorem}

Nel piano di inoltro vengono utilizzate le seguenti tabelle e strutture informative:
\begin{itemize}[noitemsep]
    \item \textbf{NHLFE (Next Hop Label Forwarding Entry)}: contiene le istruzioni esecutive per il salto successivo. Specifica il router \emph{next hop}, l'interfaccia di uscita, l'operazione da eseguire sullo stack di etichette (\textbf{Push}, \textbf{Swap}, oppure \textbf{Pop}) ed eventuali parametri di incapsulamento di livello 2.
    \item \textbf{FTN (FEC-to-NHLFE)}: mappa una FEC direttamente a una voce NHLFE. Viene consultata dall'\textbf{Ingress LER} quando riceve un pacchetto IP privo di etichette che deve essere immesso nella rete MPLS.
    \item \textbf{ILM (Incoming Label Map)}: mappa ciascuna etichetta in arrivo (\emph{In-Label}) a una o più voci NHLFE. Viene consultata dai nodi intermedi (\textbf{LSR}) e di uscita (\textbf{Egress LER}).
    \item \textbf{Token (Forwarding Table Pointer)}: nelle implementazioni industriali (quali Huawei o H3C), la \emph{Forwarding Information Base} (FIB) IP non contiene direttamente i puntatori fisici alle etichette, ma associa a ciascuna FEC un identificativo numerico generato dal sistema denominato \textbf{Token}, che funge da puntatore rapido all'appropriata struttura NHLFE.
\end{itemize}

% [IMMAGINE DEL FILE DI SLIDE CON NOME "CloudComputing_5"] Diagramma del flusso di controllo e piano dati in MPLS (Routing Protocol, LDP, FIB, LFIB) — Slide 9 / Pagina 9 del PDF

\subsubsection{Ispezione Di Una Voce LSP Tramite CLI}
Un esempio tangibile di questa correlazione è osservabile nell'output diagnostico di un router di bordo:

\begin{lstlisting}[language=bash]
<RTA>display mpls lsp include 10.2.0.0 24 verbose
---------------------------------------------------
LSP Information: LDP LSP
---------------------------------------------------
No              : 1
VrfIndex        :
Fec             : 10.2.0.0/24
Nexthop         : 10.1.1.2
In-Label        : NULL
Out-Label       : 1030
In-Interface    : ----------
Out-Interface   : Serial0
LspIndex        : 10249
Token           : 0x22005
LsrType         : Ingress
Outgoing token  : 0x0
Label Operation : PUSH
Mpls-Mtu        : 1500
TimeStamp       : 822sec
\end{lstlisting}

In questo frammento, il prefisso di destinazione IP \texttt{10.2.0.0/24} costituisce la \textbf{FEC}. Poiché l'apparato agisce come \textbf{Ingress LER} (\texttt{LsrType: Ingress}), il campo \texttt{In-Label} è nullo (\texttt{NULL}). Il router associa la FEC all'etichetta di uscita \texttt{1030} da trasmettere sul salto successivo \texttt{10.1.1.2} tramite l'interfaccia \texttt{Serial0}, applicando l'operazione \texttt{PUSH}.

\subsection{Ciclo Di Vita Di Inoltro Di Un Pacchetto: Analisi Dettagliata}
Si consideri il percorso di un pacchetto destinato all'host \texttt{10.2.0.1} attraverso la catena di nodi $\text{RTA} \rightarrow \text{RTB} \rightarrow \text{RTC} \rightarrow \text{RTD}$:

\begin{enumerate}
    \item \textbf{Ingress LER (RTA)}:
    \begin{itemize}[noitemsep]
        \item Riceve un pacchetto IP convenzionale per \texttt{10.2.0.1}.
        \item Esegue la ricerca della FEC (\texttt{10.2.0.0/24}) nella tabella FTN / FIB e reperisce il puntatore NHLFE.
        \item Esegue un'operazione di \textbf{PUSH}, anteponendo l'etichetta MPLS \texttt{1030}.
        \item Invia il frame etichettato dall'interfaccia \texttt{Serial0} verso il next-hop \texttt{10.1.1.2} (RTB).
    \end{itemize}
    % [IMMAGINE DEL FILE DI SLIDE CON NOME "CloudComputing_5"] MPLS Forwarding - Ingress LER RTA: operazione di PUSH dell'etichetta 1030 — Slide 15 / Pagina 15 del PDF

    \item \textbf{Transit LSR (RTB)}:
    \begin{itemize}[noitemsep]
        \item Riceve il pacchetto recante l'etichetta di ingresso \texttt{1030}.
        \item Consulta la tabella ILM e individua la NHLFE associata.
        \item Esegue un'operazione di \textbf{SWAP}: in questo caso specifico, l'etichetta in uscita assegnata verso RTC rimane \texttt{1030}.
        \item Inoltra il pacchetto etichettato tramite l'interfaccia \texttt{Serial3} verso \texttt{10.1.1.6} (RTC).
    \end{itemize}
    % [IMMAGINE DEL FILE DI SLIDE CON NOME "CloudComputing_5"] MPLS Forwarding - Transit LSR RTB: operazione di SWAP verso RTC — Slide 16 / Pagina 16 del PDF

    \item \textbf{Transit LSR (RTC)}:
    \begin{itemize}[noitemsep]
        \item Riceve il pacchetto con etichetta \texttt{1030}.
        \item Interroga la propria ILM e rileva una nuova istruzione di \textbf{SWAP} verso la NHLFE con etichetta di uscita \texttt{1032}.
        \item Sostituisce \texttt{1030} con \texttt{1032} e trasmette il pacchetto via \texttt{Serial3} verso \texttt{10.1.1.10} (RTD).
    \end{itemize}
    % [IMMAGINE DEL FILE DI SLIDE CON NOME "CloudComputing_5"] MPLS Forwarding - Transit LSR RTC: operazione di SWAP da etichetta 1030 a 1032 — Slide 17 / Pagina 17 del PDF

    \item \textbf{Egress LER (RTD)}:
    \begin{itemize}[noitemsep]
        \item Riceve il pacchetto con etichetta \texttt{1032}.
        \item La propria tabella ILM indica l'operazione di \textbf{POP} (\texttt{Label Operation: POP}).
        \item Rimuove integralmente l'intestazione MPLS, ripristinando il pacchetto IP originario.
        \item Consulta la classica tabella di instradamento IP e recapita il pacchetto all'interfaccia locale verso il terminale \texttt{10.2.0.2}.
    \end{itemize}
    % [IMMAGINE DEL FILE DI SLIDE CON NOME "CloudComputing_5"] MPLS Forwarding - Egress LER RTD: operazione di POP dell'etichetta 1032 e routing IP — Slide 18 / Pagina 18 del PDF
\end{enumerate}

\subsection{Distribuzione Delle Etichette Tramite LDP}
Affinché il piano di inoltro possa commutare i pacchetti, le etichette devono essere concordate tra nodi adiacenti. Il protocollo \textbf{LDP} assegna e distribuisce i binding tra FEC ed etichette lungo l'albero di instradamento:
\begin{itemize}[noitemsep]
    \item Un router a valle (\emph{Downstream}) notifica al router a monte (\emph{Upstream}) il valore dell'etichetta con cui desidera ricevere il traffico destinato a un dato prefisso (\emph{Label Mapping}).
    \item Il processo può operare in modalità su richiesta (\emph{Downstream-on-Demand}) oppure non sollecitata (\emph{Unsolicited Downstream}), costruendo in modo deterministico la catena coerente di etichette da monte a valle per ogni rotta appresa dal protocollo IGP sottostante.
\end{itemize}

% [IMMAGINE DEL FILE DI SLIDE CON NOME "CloudComputing_5"] Diagramma del protocollo LDP per la distribuzione delle etichette (Mapping e Request) — Slide 19 / Pagina 19 del PDF

\newpage

\section{Reti Nei Data Center Moderni E Limiti Delle Architetture Classiche}

\subsection{Fondamenti Del Data Center E Topologia Spine-Leaf}
Un moderno \textbf{Data Center (DC)} accentra potenza di calcolo, storage e connettività di rete per erogare servizi e applicazioni critiche con quattro requisiti essenziali: affidabilità, flessibilità, efficienza energetica ed elevata produttività.

% [IMMAGINE DEL FILE DI SLIDE CON NOME "CloudComputing_5"] Concetti base del Data Center e topologia Spine-and-Leaf per Network, Computing e Storage — Slide 21 / Pagina 21 del PDF

Nelle infrastrutture tradizionali, la rete era concepita su una gerarchia rigida a tre livelli: \textbf{Accesso}, \textbf{Aggregazione} e \textbf{Core}. Con la virtualizzazione massiva dei server e la nascita delle architetture cloud a microservizi:
\begin{itemize}[noitemsep]
    \item Il modello di traffico è mutato radicalmente da prevalentemente \textbf{Nord-Sud} (traffico client-server tra l'esterno del DC e l'infrastruttura interna) a prevalentemente \textbf{Est-Ovest} (traffico server-server tra macchine virtuali e microservizi distribuiti su nodi fisici diversi).
    \item L'architettura convenzionale a tre livelli introduce colli di bottiglia e latenze non deterministiche dovute ai molteplici stadi di commutazione attraversati dal traffico Est-Ovest.
\end{itemize}

Per rispondere a questa criticità, i data center hanno adottato la topologia a matrice di commutazione \textbf{Spine-and-Leaf} (derivata dalle reti Clos bi-stadio):
\begin{itemize}[noitemsep]
    \item Ogni commutatore di accesso (\textbf{Leaf}, tipicamente integrato come switch \textbf{ToR} - \emph{Top-of-Rack}) è direttamente collegato a tutti i commutatori dorsali (\textbf{Spine}).
    \item I commutatori \textbf{Spine} non sono interconnessi tra loro, così come non esistono interconnessioni dirette tra leaf switch.
    \item La distanza tra qualsiasi coppia di server collegati al DC è rigidamente costante e pari a due salti (\emph{hops}), garantendo una latenza uniforme e predicibile.
\end{itemize}

\subsection{Bilanciamento Del Carico Con ECMP}
Per sfruttare appieno la molteplicità dei collegamenti fisici tra leaf e spine senza incappare nei blocchi topologici del protocollo Spanning Tree (STP), la rete di interconnessione viene instradata al Livello 3 sfruttando l'\textbf{Equal-Cost Multi-Path (ECMP)}.

% [IMMAGINE DEL FILE DI SLIDE CON NOME "CloudComputing_5"] Criticità della latenza per nodi di calcolo e bilanciamento del carico tramite 5-tuple hash ECMP — Slide 22 / Pagina 22 del PDF

Il meccanismo ECMP ripartisce i flussi sui link paralleli calcolando un valore di hash basato sulla \textbf{5-tupla}:
\[
\text{Hash} = h\Big(\text{IP}_{\text{src}},\, \text{IP}_{\text{dst}},\, \text{Protocol},\, \text{Port}_{\text{src}},\, \text{Port}_{\text{dst}}\Big)
\]
In questo modo, pacchetti appartenenti a flussi applicativi distinti presentano valori di hash differenti e vengono instradati in modo deterministico su cammini fisici alternativi a pari costo, massimizzando il \emph{throughput} complessivo ed evitando la congestione dei singoli link.

\subsection{Il Paradigma Del Multitenancy}
La coesistenza sicura di più clienti (\textbf{tenant}) sulla medesima infrastruttura condivisa rappresenta il fondamento operativo del Cloud Computing. In ambito \emph{Software-as-a-Service} (SaaS) e \emph{Infrastructure-as-a-Service} (IaaS):
\begin{itemize}[noitemsep]
    \item \textbf{Single-Tenant}: ciascun cliente possiede un'istanza software dedicata e risorse infrastrutturali isolate fisicamente. Non vi è condivisione di codice, storage o elaborazione con terzi.
    \item \textbf{Multi-Tenant}: una singola istanza infrastrutturale o software serve molteplici clienti concorrenti. I dati, le macchine virtuali e i segmenti di rete dei diversi tenant sono segregati logicamente tramite barriere software e crittografiche, garantendo la totale confidenzialità a fronte di un'efficienza economica superiore.
\end{itemize}

% [IMMAGINE DEL FILE DI SLIDE CON NOME "CloudComputing_5"] Confronto tra architetture applicative Single-Tenant e Multi-Tenant — Slide 23 / Pagina 23 del PDF

\subsection{Analisi Quantitativa: Il Crollo Dello Standard IEEE 802.1Q}
Nelle reti Ethernet tradizionali, l'isolamento logico dei tenant al Livello 2 è delegato allo standard \textbf{IEEE 802.1Q}. L'intestazione 802.1Q interpone un tag di 4 byte tra l'indirizzo MAC sorgente e il campo Type/Length del frame Ethernet:
\begin{itemize}[noitemsep]
    \item \textbf{TPID} (16 bit): \emph{Tag Protocol Identifier}, fissato al valore \texttt{0x8100}.
    \item \textbf{PRI / PCP} (3 bit): \emph{Priority Code Point} per la differenziazione del traffico.
    \item \textbf{CFI / DEI} (1 bit): \emph{Canonical Format Indicator / Drop Eligible Indicator}.
    \item \textbf{VLAN ID} (12 bit): identificativo univoco della VLAN.
\end{itemize}

% [IMMAGINE DEL FILE DI SLIDE CON NOME "CloudComputing_5"] Formato del frame Ethernet con tag IEEE 802.1Q evidenziando il campo a 12 bit VLAN ID — Slide 27 / Pagina 27 del PDF

Il vincolo strutturale dei \textbf{12 bit} riserva uno spazio di indirizzamento massimo teorico pari a:
\[
2^{12} = 4096 \text{ VLAN}
\]

% [IMMAGINE DEL FILE DI SLIDE CON NOME "CloudComputing_5"] Diagramma architetturale Top-of-Rack con 20 rack e switch ToR — Slide 25 / Pagina 25 del PDF

Si valuti l'inadeguatezza di tale limite attraverso uno scenario dimensionale quantitativo relativo a un data center di dimensioni medio-piccole:
\begin{itemize}[noitemsep]
    \item Numero di rack fisici: $20$.
    \item Server per rack collegati allo switch ToR: $48$.
    \item Totale server fisici: $20 \times 48 = 960$ server.
    \item Densità di virtualizzazione: $50$ Macchine Virtuali (\textbf{VM}) ospitate per ciascun server fisico.
    \item Totale VM distribuibili nel cluster: $960 \times 50 = \mathbf{48\,000 \text{ VM}}$.
\end{itemize}

% [IMMAGINE DEL FILE DI SLIDE CON NOME "CloudComputing_5"] Tabella dimensionale quantitativa: Large Enterprise vs Small Business e limiti 802.1Q — Slide 26 / Pagina 26 del PDF

Si analizzino i requisiti di segmentazione di rete per due profili di clientela:
\begin{table}[H]
\centering
\small
\begin{tabularx}{\textwidth}{l X X}
\toprule
\textbf{Parametro Dimensionale} & \textbf{Profilo Large Enterprise} & \textbf{Profilo Small Business} \\
\midrule
VM per singolo Tenant & $10$ & $1$ \\
Sottoreti (VLAN) per Tenant & $20$ (es. Web, DB, App, Backup, DMZ) & $5$ \\
\textbf{Numero di Tenant Attivabili} & $\frac{48\,000}{10} = \mathbf{4\,800}$ & $\frac{48\,000}{1} = \mathbf{48\,000}$ \\
\textbf{Sottoreti Totali Richieste} & $4\,800 \times 20 = \mathbf{96\,000}$ & $48\,000 \times 5 = \mathbf{240\,000}$ \\
\midrule
\textbf{Capacità Massima IEEE 802.1Q} & \multicolumn{2}{c}{$\mathbf{4\,096 \text{ VLAN}}$} \\
\bottomrule
\end{tabularx}
\caption{Confronto quantitativo tra i requisiti di isolamento e il limite di IEEE 802.1Q.}
\end{table}

Il divario tra le $4096$ VLAN fornite dallo standard IEEE 802.1Q e le oltre $96\,000$ o $240\,000$ sottoreti richieste rende impraticabile la segmentazione a livello 2 fisico. Inoltre, l'esplosione delle tabelle MAC degli switch di accesso (dovuta all'esposizione diretta dei MAC virtuali di decine di migliaia di VM) e la necessità di consentire la \textbf{migrazione a caldo} (\emph{vMotion/live migration}) delle VM mantenendo inalterato l'indirizzo IP richiedevano un cambio radicale di paradigma.

\newpage

\section{Virtual eXtensible Local Area Network (VxLAN): La Soluzione Overlay}

\subsection{Principio Di Funzionamento: Rete Overlay Su Underlay IP}
Per scardinare i vincoli del Livello 2 fisico, l'organismo IETF ha standardizzato le tecnologie \textbf{NVO3} (\emph{Network Virtualization over Layer 3}), tra cui spiccano NVGRE, STT e, universalmente adottata, \textbf{VxLAN} (RFC 7348).

L'architettura separa nettamente l'infrastruttura di rete in due strati sovrapposti:
\begin{itemize}[noitemsep]
    \item \textbf{Underlay Network}: la rete fisica di trasporto di Livello 3 (instradata tramite IP/OSPF/BGP), robusta, scalabile e dotata di bilanciamento di carico ECMP nativo.
    \item \textbf{Overlay Network}: la rete virtuale di Livello 2 stabilita logicamente al di sopra dell'underlay per mezzo di tunnel punto-multipunto.
\end{itemize}

% [IMMAGINE DEL FILE DI SLIDE CON NOME "CloudComputing_5"] Modello di rete VxLAN: disaccoppiamento tra Overlay (VxLAN) e Underlay (IP) tramite VTEP — Slide 30 / Pagina 30 del PDF

La tecnica impiegata da VxLAN è l'incapsulamento \textbf{MAC-in-UDP}: l'intero frame Ethernet generato dalla macchina virtuale (compreso di header L2 originale e payload) viene racchiuso come carico utile all'interno di un pacchetto UDP/IP convenzionale. I router del nucleo non vedono né apprendono i MAC delle VM, ma instradano esclusivamente i pacchetti UDP in base agli indirizzi IP degli endpoint del tunnel.

% [IMMAGINE DEL FILE DI SLIDE CON NOME "CloudComputing_5"] Evoluzione dell'interconnessione Layer 2 tra Data Center mediante tunnel VxLAN — Slide 29 / Pagina 29 del PDF

\subsection{Superamento Dei Limiti: Il VNI A 24 Bit}
Al posto del tag VLAN a 12 bit, VxLAN introduce il \textbf{VXLAN Network Identifier (VNI)}:
\begin{itemize}[noitemsep]
    \item Lo spazio del VNI ha una lunghezza di \textbf{24 bit}.
    \item Il numero totale di segmenti logici isolati supportabili contemporaneamente ammonta a:
    \[
    \boxed{2^{24} = 16\,777\,216 \text{ domini di isolamento}}
    \]
\end{itemize}
Questo volume di oltre 16 milioni di reti virtuali consente a qualsiasi operatore cloud di allocare decine di domini di broadcast dedicati a centinaia di migliaia di tenant distinti, eliminando definitivamente le collisioni di indirizzamento.

\subsection{Formato Dettagliato Del Pacchetto VxLAN}
L'incapsulamento completo di un frame all'interno di una trama VxLAN comporta un incremento della dimensione totale (\emph{overhead}) pari ad almeno 50 byte (o 54 byte in caso di tag 802.1Q preservato):

% [IMMAGINE DEL FILE DI SLIDE CON NOME "CloudComputing_5"] Struttura dettagliata dei campi di un pacchetto incapsulato VxLAN (Outer MAC, Outer IP, UDP, VxLAN Header) — Slide 31 / Pagina 31 del PDF

\begin{figure}[H]
\centering
\begin{tikzpicture}[x=0.13cm, y=0.8cm, every node/.style={font=\scriptsize}]
    % Outer Frame
    \draw[fill=gray!20] (0,0) rectangle (14,1.2) node[midway, align=center]{Outer MAC\\(14 B)};
    \draw[fill=blue!15] (14,0) rectangle (34,1.2) node[midway, align=center]{Outer IP\\(20 B)};
    \draw[fill=cyan!20] (34,0) rectangle (42,1.2) node[midway, align=center]{UDP\\(8 B)};
    \draw[fill=green!20] (42,0) rectangle (50,1.2) node[midway, align=center]{VxLAN\\(8 B)};
    \draw[fill=yellow!20] (50,0) rectangle (90,1.2) node[midway, align=center]{Original L2 Frame\\(MAC + IP + Payload)};
    \draw[fill=gray!30] (90,0) rectangle (96,1.2) node[midway, align=center]{FCS\\(4 B)};
\end{tikzpicture}
\caption{Rappresentazione schematica dell'incapsulamento MAC-in-UDP in VxLAN.}
\end{figure}

L'articolazione interna dei singoli campi comprende:
\begin{enumerate}
    \item \textbf{Outer MAC Header} (14 byte):
    \begin{itemize}[noitemsep]
        \item \emph{Source MAC}: indirizzo MAC dell'interfaccia fisica del VTEP sorgente.
        \item \emph{Destination MAC}: indirizzo MAC del dispositivo di salto successivo (\emph{next-hop router}) verso il VTEP remoto.
    \end{itemize}
    \item \textbf{Outer IP Header} (20 byte):
    \begin{itemize}[noitemsep]
        \item \emph{Source IP}: indirizzo IP univoco assegnato al VTEP sorgente.
        \item \emph{Destination IP}: indirizzo IP univoco del VTEP di destinazione (oppure un indirizzo IP Multicast in caso di flooding BUM).
    \end{itemize}
    \item \textbf{UDP Header} (8 byte):
    \begin{itemize}[noitemsep]
        \item \emph{Destination Port}: rigidamente fissata al valore standardizzato \textbf{4789} assegnato da IANA per VxLAN.
        \item \emph{Source Port}: valore calcolato dinamicamente mediante una funzione di hash applicata all'intestazione interna del frame nativo della VM (es. hash della 5-tupla originale).
        \item \emph{UDP Length} e \emph{Checksum}: include la lunghezza complessiva e un checksum opzionale per la verifica di integrità.
    \end{itemize}
    \item \textbf{VXLAN Header} (8 byte / 64 bit):
    \begin{itemize}[noitemsep]
        \item \emph{Flags} (8 bit): il flag $I$ (bit in posizione 3) deve essere asserito a $1$ per certificare la validità del VNI; i restanti bit sono riservati.
        \item \emph{Reserved} (24 bit): riservati e impostati a zero in trasmissione.
        \item \textbf{VNI} (24 bit): l'identificativo univoco della rete virtuale overlay.
        \item \emph{Reserved} (8 bit): riservati e impostati a zero.
    \end{itemize}
    \item \textbf{Payload}: il frame Ethernet originario integro, seguito dal relativo campo di controllo FCS (\emph{Frame Check Sequence}).
\end{enumerate}

\subsection{Motivazioni Ingegneristiche Della Scelta Di UDP}
La decisione di veicolare il traffico incapsulato sopra lo strato \textbf{UDP} (rispetto a un incapsulamento IP grezzo, come in GRE, o sopra TCP) risponde a quattro precisi imperativi architetturali:
\begin{enumerate}
    \item \textbf{Sfruttamento nativo del bilanciamento ECMP nel core}: i router di dorsale ispezionano la 5-tupla standard IP/UDP. Se l'incapsulamento avvenisse direttamente su IP (senza porte L4), tutti i pacchetti tra due VTEP avrebbero la medesima tupla (stesso IP sorgente e destinazione, stesso protocollo), venendo confinati dal calcolo di hash sul medesimo link fisico.
    Calcolando la \textbf{UDP Source Port} come funzione hash del traffico interno della VM:
    \[
    \text{Port}_{\text{src}}^{\text{UDP}} = \text{hash}\Big(\text{Inner-5-Tuple}\Big)
    \]
    flussi applicativi diversi tra i medesimi due host ottengono porte UDP sorgente differenti. I commutatori underlay instradano automaticamente i pacchetti del medesimo tunnel su link fisici paralleli mediante ECMP, eliminando i punti caldi di congestione.
    \item \textbf{Elaborazione ad alta velocità via ASIC}: i componenti hardware commerciali nei commutatori moderni sono ottimizzati per analizzare e commutare frame UDP/IP a \emph{line rate}. L'introduzione di un nuovo protocollo a livello IP avrebbe richiesto la riprogettazione delle pipeline hardware dei router.
    \item \textbf{Attraversamento trasparente di Firewall e NAT}: UDP è universalmente gestito da tutti gli apparati di sicurezza perimetrale, bilanciatori di carico e dispositivi middlebox, consentendo al tunnel di attraversare domini eterogenei senza richiedere configurazioni speciali nel core.
    \item \textbf{Integrità opzionale}: il campo checksum UDP offre una verifica aggiuntiva di coerenza del pacchetto attraverso la rete underlay, se richiesto dall'applicazione.
\end{enumerate}

\subsection{Implementazione Ed Architettura Dei VTEP}
L'entità preposta a instaurare, mantenere ed eseguire l'incapsulamento e la decapsulazione dei tunnel VxLAN è il \textbf{VXLAN Tunnel Endpoint (VTEP)}. Un tunnel VxLAN è definito formalmente da una coppia di VTEP.

% [IMMAGINE DEL FILE DI SLIDE CON NOME "CloudComputing_5"] Modalità di implementazione del VTEP: Software nell'Hypervisor vs Hardware nel Top-of-Rack Switch — Slide 32 / Pagina 32 del PDF

A seconda del punto della rete in cui risiede la terminazione del tunnel, si distinguono due paradigmi realizzativi:
\begin{itemize}[noitemsep]
    \item \textbf{VTEP Software (Host-based)}: l'endpoint è implementato all'interno del software di virtualizzazione (\textbf{Hypervisor}, quali KVM, Microsoft Hyper-V, VMware NSX o switch virtuali come Open vSwitch). Il frame viene incapsulato prima di lasciare l'interfaccia di rete fisica del server. Questo approccio garantisce la massima programmabilità e visibilità end-to-end, ma grava parzialmente sul consumo di cicli CPU del nodo di calcolo.
    \item \textbf{VTEP Hardware (ToR Switch-based)}: l'incapsulamento e la decapsulazione vengono delegati agli ASIC dello switch di sommità del rack (\textbf{ToR Switch}). Ciò consente di integrare nel fabric VxLAN anche server dedicati \emph{bare-metal} o appliance fisiche non virtualizzate prive di supporto nativo VxLAN.
\end{itemize}

\subsubsection{Ruoli Operativi Dei VTEP}
\begin{itemize}[noitemsep]
    \item \textbf{VTEP come Gateway di Accesso}: effettua il bridging trasparente tra i frame Ethernet classici (basati su tag VLAN 802.1Q) e i segmenti logici VxLAN (identificati dal VNI corrispondente).
    \item \textbf{VTEP in Ambiente Routed (Layer 3 Gateway)}: funge da default gateway per i carichi di lavoro attestati sui diversi VNI, eseguendo l'instradamento tra sottoreti distinte o connettendo l'overlay con l'esterno (verso reti WAN aziendali, la rete Internet o altre infrastrutture cloud).
    \item \textbf{Integrazione con Protocolli di Routing}: l'apprendimento delle corrispondenze tra MAC e VTEP remoto (inizialmente demandato a tecniche di flooding multicast) viene oggi orchestrato in modo scalabile e deterministico mediante l'integrazione con \textbf{MP-BGP EVPN} (Ethernet VPN), che standardizza il piano di controllo per distribuire i binding tramite annunci di rotte BGP dedicate.
\end{itemize}

% DD -----------------------------------------------------------

% ==============================================================================
% CAPITOLO: ARCHITETTURE VXLAN, FORWARDING BUM, PRATICA LINUX ED EVPN/BGP
% ==============================================================================

\section{Architettura VxLAN Nei Data Center Moderni}

Nelle moderne infrastrutture di data center, la crescente densità computazionale introdotta dalla virtualizzazione e dai container ha reso obsoleti i tradizionali progetti di rete basati esclusivamente su switching di Livello~2 (Ethernet standard e protocollo \emph{Spanning Tree}, STP). Il protocollo \textbf{VxLAN} (\emph{Virtual Extensible LAN}, definito nella RFC~7432 e formalizzato nella specifica RFC~7348) risponde a questi limiti implementando una rete logica sovrapposta (\emph{overlay}) al di sopra di un'infrastruttura fisica di instradamento di Livello~3 (\emph{underlay}). Prima di analizzare i dettagli implementativi, è opportuno richiamare il principio cardine dell'incapsulamento \emph{packet-over-packet}: un'intera trama Ethernet di Livello~2 generata da una macchina virtuale sorgente viene racchiusa all'interno di un pacchetto UDP/IP instradabile, consentendo la trasparenza di Livello~2 attraverso domini di routing puramente IP.

\subsection{Collocazione Dei VTEP: Topologia Spine-and-Leaf E Gateway Di Bordo}

Il punto terminale del tunnel VxLAN è identificato dal componente concettuale denominato \textbf{VTEP} (\emph{VxLAN Tunnel Endpoint}). Il VTEP svolge le operazioni fondamentali di incapsulamento (aggiunta degli header esterni su trame uscenti dal carico computazionale) e decapsulamento (rimozione degli header esterni e inoltro della trama originale al carico computazionale di destinazione).

Nelle architetture di rete contemporanee, l'underlay adotta prevalentemente un modello topologico a due livelli di tipo \textbf{Spine-and-Leaf} (noto anche come architettura Clos a due stadi):
\begin{itemize}[noitemsep]
    \item \textbf{Livello Spine}: gli switch spine costituiscono il nucleo centrale (\emph{core fabric}) dell'infrastruttura. Il loro unico compito consiste nell'instradare velocemente il traffico IP tra i diversi switch di accesso tramite protocolli di routing ad alte prestazioni e percorsi equi (\emph{Equal-Cost Multi-Pathing}, ECMP). Gli switch spine sono completamente agnostici rispetto alla presenza della rete VxLAN: operano a puro Livello~3 e non eseguono alcuna funzione di decapsulamento o gestione dei tunnel overlay.
    \item \textbf{Livello Leaf}: gli switch leaf fungono da nodi di accesso (\emph{Top-of-Rack}, ToR) e interconnettono direttamente i sistemi terminali fisici (\emph{bare metal}), i server hypervisor o gli chassis computazionali. È esattamente sul livello leaf che risiede la funzione di VTEP hardware/software: tali dispositivi associano le VLAN locali o le interfacce fisiche/logiche al corrispondente \textbf{VNI} (\emph{VxLAN Network Identifier}), gestiscono l'incapsulamento dei pacchetti verso l'underlay e mantengono le tabelle di corrispondenza tra indirizzi MAC e VTEP remoti.
\end{itemize}

% [IMMAGINE DEL FILE DI SLIDE CON NOME "CloudComputing_5"] Topologia Spine-and-Leaf con indicazione dei VTEP implementati a livello Leaf e trasparenza degli switch Spine al traffico VxLAN — Slide 31 / Pagina 2 del PDF

Accanto alla collocazione sui nodi leaf ToR, una funzione cruciale è svolta dai \textbf{VTEP at the Edge} (o \emph{VxLAN Gateways}). I gateway di rete operano al perimetro del dominio overlay, ovvero nel punto in cui il traffico deve abbandonare la rete virtuale incapsulata per raggiungere la rete fisica tradizionale (\emph{legacy bare-metal}), data center remoti o la rete pubblica Internet. Il VTEP di edge ha il compito di rimuovere l'incapsulamento VxLAN e instradare i pacchetti a Livello~3 (routing inter-VxLAN o routing verso l'esterno) o tradurre il VNI in una classica VLAN 802.1Q (bridging L2 edge). In contesti avanzati, il gateway VTEP si integra direttamente con protocolli di routing dinamico convenzionali (come \textbf{BGP} od \textbf{OSPF}) per pubblicare la raggiungibilità dei prefissi overlay verso i domini esterni.

\subsection{Modello Di Rete Ed Esigenza Di Connettività Full-Mesh}

La connessione tra macchine virtuali collocate su host differenti richiede che ciascun VTEP sia in grado di comunicare con tutti gli altri VTEP che condividono lo stesso identificatore di segmento virtuale.

% [IMMAGINE DEL FILE DI SLIDE CON NOME "CloudComputing_5"] Modello di rete VxLAN con VTEP sorgente 1.1.1.1 e VTEP destinazione 2.2.2.2 attestati su rete IP underlay — Slide 32 / Pagina 3 del PDF

Ogni VTEP possiede un indirizzo IP di Livello~3 univoco sull'infrastruttura di underlay e può avere associati uno o più identificatori VNI. Quando un frame di Livello~2 generato da un nodo terminale perviene all'interfaccia di accesso del VTEP di ingresso (\emph{ingress VTEP}), quest'ultimo mappa la porta o il tag VLAN di origine nel rispettivo VNI ed esegue l'incapsulamento all'interno di un pacchetto UDP/IP, indirizzato verso l'IP underlay del VTEP di uscita (\emph{egress VTEP}).

Qualora una rete logica estesa debba interconnettere macchine virtuali attestate su un insieme distribuito di nodi (ad esempio $\text{VTEP}_1$, $\text{VTEP}_2$ e $\text{VTEP}_3$), ciascuna coppia di VTEP deve instaurare un tunnel punto-punto bidirezionale, dando luogo a una topologia overlay \emph{full-mesh}.

% [IMMAGINE DEL FILE DI SLIDE CON NOME "CloudComputing_5"] Rete overlay full-mesh con tunnel VxLAN punto-punto stabiliti tra nodi VTEP_1, VTEP_2 e VTEP_3 — Slide 33 / Pagina 4 del PDF

Il problema cardine di questo modello architetturale risiede nella gestione, apprendimento e distribuzione coerente delle associazioni tra gli indirizzi di Livello~2 (MAC address delle VM), gli indirizzi di Livello~3 (IP delle VM) e gli indirizzi IP dei VTEP ospitanti (\emph{MAC-to-VTEP binding}).

\subsection{Formato Del Pacchetto VxLAN Ed Analisi Dell'Overhead (RFC 7348)}

L'incapsulamento VxLAN prevede l'aggiunta di una sequenza prefissata di header al frame Ethernet originale di Livello~2. La struttura complessiva del pacchetto incapsulato è mostrata formalmente di seguito.

% [IMMAGINE DEL FILE DI SLIDE CON NOME "CloudComputing_5"] Formato dettagliato del pacchetto VxLAN: Outer MAC, Outer IP, Outer UDP, VxLAN Header e payload originale — Slide 34 / Pagina 5 del PDF

L'architettura del pacchetto comprende le seguenti sezioni:
\begin{enumerate}
    \item \textbf{Outer Ethernet Header (14 byte, opzionalmente 18 con 802.1Q)}: contiene il MAC sorgente del router/switch di hop corrente nell'underlay e il MAC di destinazione del salto successivo (\emph{next-hop}).
    \item \textbf{Outer IP Header (20 byte per IPv4)}: contiene come indirizzo IP sorgente l'IP underlay del VTEP locale e come destinazione l'IP underlay del VTEP remoto (in modalità unicast) oppure l'indirizzo di un gruppo multicast.
    \item \textbf{Outer UDP Header (8 byte)}:
    \begin{itemize}[noitemsep]
        \item \textbf{Destination Port}: l'ente IANA ha assegnato a VxLAN la porta UDP standard \texttt{4789} (RFC~7348). Le implementazioni devono supportare la configurabilità di tale porta per garantire compatibilità con bozze pre-standard che impiegavano porte differenti (ad esempio la porta 8472).
        \item \textbf{Source Port}: per massimizzare l'entropia del traffico, la specifica stabilisce che il valore della porta UDP sorgente sia calcolato applicando una funzione di hash ai campi dell'header interno (ad esempio una tupla basata su IP sorgente/destinazione e protocollo interni, oppure sui campi dell'header Ethernet interno). Ciò permette agli switch dell'underlay, privi di conoscenza su VxLAN, di ripartire il traffico dei tunnel lungo percorsi paralleli tramite algoritmi di bilanciamento di carico ECMP (\emph{Equal-Cost Multi-Pathing}). È raccomandato l'impiego dell'intervallo di porte effimere/private $[49152, 65535]$.
        \item \textbf{UDP Checksum}: da specifica standard, il checksum UDP deve essere trasmesso con valore impostato a zero (\texttt{0x0000}) per evitare il ricalcolo computazionale su frame incapsulati ad alta velocità. Se ricevuto con valore pari a zero, il pacchetto deve essere accettato per la decapsulazione. Qualora un VTEP decida di generare un checksum valido (non nullo), il ricevente può opzionalmente verificarlo e, qualora l'esito risulti corrotto, il pacchetto deve essere scartato.
    \end{itemize}
    \item \textbf{VxLAN Header (8 byte)}:
    \begin{itemize}[noitemsep]
        \item \textbf{Flags (8 bit)}: il bit in posizione 3 (indicato convenzionalmente come flag $I$, \emph{VNI validity flag}) deve essere tassativamente impostato a $1$ affinché il VNI trasportato sia considerato valido. I restanti 7 bit sono etichettati come riservati ($R$), devono essere trasmessi come zero e ignorati in ricezione.
        \item \textbf{Campi Riservati (24 bit e 8 bit)}: riservati per future estensioni dello standard, obbligatoriamente valorizzati a zero all'atto della trasmissione.
        \item \textbf{VxLAN Network Identifier (VNI, 24 bit)}: consente di identificare in modo univoco fino a $2^{24} = 16{,}777{,}216$ segmenti logici sovrapposti, superando radicalmente la barriera dimensionale imposta dallo spazio a 12 bit delle tradizionali VLAN IEEE 802.1Q (limitate a sole 4096 istanze).
    \end{itemize}
    \item \textbf{Original Layer 2 Frame}: la trama Ethernet originaria emessa dalla macchina sorgente, comprensiva di indirizzi MAC sorgente e destinazione, eventuale etichetta interna e payload IP.
\end{enumerate}

\paragraph{Bilancio Computazionale dell'Overhead}
L'overhead totale introdotto dall'incapsulamento varia a seconda del conteggio della sequenza di controllo di trama (\emph{Frame Check Sequence}, FCS):
\begin{align*}
    \text{Overhead}_{\text{senza FCS}} &= \underbrace{14\,\text{B}}_{\text{Outer MAC}} + \underbrace{20\,\text{B}}_{\text{Outer IP}} + \underbrace{8\,\text{B}}_{\text{Outer UDP}} + \underbrace{8\,\text{B}}_{\text{VxLAN Hdr}} = \mathbf{50\,\text{byte}} \\
    \text{Overhead}_{\text{con FCS}} &= 50\,\text{B} + \underbrace{4\,\text{B}}_{\text{FCS}} = \mathbf{54\,\text{byte}}
\end{align*}
Nelle reti reali, il campo FCS di 4 byte viene solitamente rigenerato a livello hardware dai controller di rete (NIC) e non è visibile al software di sistema; di conseguenza, l'overhead effettivo è computato comunemente in \textbf{50 byte}. Ne deriva che, per trasportare trame Ethernet standard con MTU pari a 1500 byte senza incorrere in frammentazione a livello IP nell'underlay, le interfacce fisiche dell'infrastruttura di transito devono essere riconfigurate per supportare trame giganti (\emph{Jumbo Frame}), con una MTU minima consigliata pari a:
\[
    \text{MTU}_{\text{underlay}} \ge 1500\,\text{B} + 50\,\text{B} = 1550\,\text{byte} \quad (\text{tipicamente impostata a } 1600\,\text{B o } 9000\,\text{B}).
\]

\subsection{Combinazioni Di Incapsulamento Underlay Ed Overlay}

L'architettura definita da VxLAN disaccoppia completamente lo stack protocollare del servizio utente dalla tecnologia di instradamento della rete fisica. Sono pertanto possibili e standardizzate quattro combinazioni architetturali distinte tra overlay ed underlay:

% [IMMAGINE DEL FILE DI SLIDE CON NOME "CloudComputing_5"] Combinazioni di incapsulamento VxLAN tra underlay e overlay (IPv4/IPv6 over IPv4/IPv6) — Slide 35 / Pagina 6 del PDF

\begin{table}[H]
    \centering
    \small
    \begin{tabularx}{\textwidth}{lXX}
        \toprule
        \textbf{Tipologia Incapsulamento} & \textbf{Descrizione Strutturale} & \textbf{Caso d'Uso Tipico} \\
        \midrule
        \textbf{IPv4 over IPv4} & Payload IPv4 incapsulato in header esterno IPv4. & Standard consolidato nei data center di generazione attuale. \\
        \textbf{IPv6 over IPv4} & Carichi di lavoro IPv6 veicolati su fabric underlay legacy IPv4. & Transizione verso IPv6 senza richiedere riconfigurazione dell'underlay. \\
        \textbf{IPv4 over IPv6} & Reti legacy IPv4 erogate al di sopra di un moderno fabric IPv6-only. & Data center di nuova generazione con underlay nativo IPv6. \\
        \textbf{IPv6 over IPv6} & Stack completamente IPv6 sia a livello overlay che underlay. & Ambienti \emph{greenfield} privi di dipendenze con indirizzamento legacy. \\
        \bottomrule
    \end{tabularx}
    \caption{Combinazioni ammesse per il trasporto dei tunnel VxLAN.}
\end{table}

\newpage

% ==============================================================================
% SEZIONE 2: FORWARDING DEL TRAFFICO E GESTIONE BUM
% ==============================================================================

\section{Meccanismi Di Inoltro E Gestione Del Traffico BUM}

Dal punto di vista dell'inoltro sul piano dati, il traffico che attraversa una rete VxLAN si suddivide tassativamente in due categorie fondamentali:
\begin{enumerate}[noitemsep]
    \item \textbf{Known Unicast}: traffico unicast per il quale il VTEP di ingresso conosce preventivamente la corrispondenza tra l'indirizzo MAC di destinazione e l'indirizzo IP del VTEP remoto che serve quel terminale.
    \item \textbf{BUM Traffic} (\emph{Broadcast, Unknown Unicast, Multicast}): insieme di trame che non possono essere trasmesse a un singolo VTEP tramite un tunnel unicast diretto. L'\textbf{unknown unicast} si manifesta specificamente quando uno switch o VTEP riceve un frame con MAC destinatario unicast ma privo di riscontro nella propria tabella di inoltro (\emph{Forwarding Information Base}, FIB).
\end{enumerate}

\subsection{Il Principio Del Flood and Learn Sul Piano Dati}

In assenza di un piano di controllo centralizzato o distribuito (quale MP-BGP EVPN), il meccanismo nativo per gestire la dispersione delle informazioni di raggiungibilità è il \textbf{Flood and Learn} (mutuato dal funzionamento classico dei bridge trasparenti IEEE 802.1D):
\begin{itemize}[noitemsep]
    \item Quando un host sorgente invia un frame verso una destinazione sconosciuta, il VTEP di ingresso inietta la trama in modalità \emph{flooding} verso tutti gli altri VTEP appartenenti al medesimo VNI.
    \item Ciascun VTEP ricevente analizza la trama decapsulata: esaminando l'indirizzo MAC sorgente e associandolo all'indirizzo IP del VTEP mittente (appreso dall'header IP esterno), il nodo popola dinamicamente la propria tabella locale di Livello~2 con la tupla:
    \[
        \langle \text{MAC}_{\text{sorgente}},\, \text{VNI},\, \text{IP}_{\text{VTEP\_remoto}} \rangle.
    \]
\end{itemize}

\subsection{Modalità Di Replica Del Traffico BUM}

La replica fisica delle copie dei frame BUM attraverso la rete IP di underlay può essere implementata secondo tre paradigmi architetturali differenti:

\subsubsection{1. Ingress Replication (Head-End Replication)}
Nel paradigma \textbf{Ingress Replication} (o replica alla sorgente), il compito di duplicare il traffico grava interamente sul VTEP di ingresso:
\begin{itemize}
    \item Il VTEP mittente consulta una lista locale preconfigurata di tutti i VTEP remoti che partecipano allo stesso VNI (\emph{ingress replication list}).
    \item Per ciascun peer remoto presente nella lista, il VTEP crea una copia distinta del pacchetto originale, applica l'header esterno VxLAN e invia la trama tramite un pacchetto unicast IP punto-punto.
\end{itemize}
\textbf{Vantaggi}: Non richiede alcun supporto multicast nella rete di underlay; la configurazione fisica della rete di transito resta pura, semplice ed economica. \\
\textbf{Svantaggi}: Scalabilità limitata. Se a un VNI appartengono $N$ VTEP, il VTEP di ingresso deve generare $N-1$ copie dello stesso frame, provocando un picco nell'occupazione della banda del link di uplink locale e un elevato consumo di risorse di calcolo.

% [IMMAGINE DEL FILE DI SLIDE CON NOME "CloudComputing_5"] Sequenza di inoltro BUM tramite Ingress Replication con richiesta ARP di VM_A replicata a VTEP_2 e VTEP_3 — Slide 40 / Pagina 10 del PDF

\paragraph{Flusso Operativo con Richiesta ARP}:
Considerando lo scenario in cui la macchina $\text{VM}_A$ debba inviare una richiesta ARP in broadcast per apprendere il MAC di $\text{VM}_C$:
\begin{enumerate}[noitemsep]
    \item $\text{VM}_A$ trasmette un frame con $\text{DMAC} = \texttt{FF:FF:FF:FF:FF:FF}$.
    \item $\text{VTEP}_1$ mappa il frame nel VNI 5000 e consulta la propria lista di replica, contenente $\text{VTEP}_2$ e $\text{VTEP}_3$.
    \item $\text{VTEP}_1$ crea due copie distinte con destinazioni IP esterne pari rispettivamente a $\text{IP}_{\text{VTEP\_2}}$ e $\text{IP}_{\text{VTEP\_3}}$.
    \item All'arrivo, $\text{VTEP}_2$ e $\text{VTEP}_3$ decapsulano il pacchetto, memorizzano la corrispondenza $\langle \text{MAC}_A, 5000, \text{IP}_{\text{VTEP\_1}} \rangle$ e iniettano il frame nei domini di broadcast locali.
    \item $\text{VM}_B$ scarta il pacchetto, mentre $\text{VM}_C$ genera una risposta ARP diretta verso $\text{MAC}_A$.
\end{enumerate}

\subsubsection{2. Multicast Replication}
In questa modalità, la complessità della replicazione viene delegata interamente all'infrastruttura di routing dell'underlay:
\begin{itemize}
    \item Ciascun VNI viene associato a un indirizzo di gruppo multicast IP (ad esempio \texttt{225.0.0.1}).
    \item Tutti i VTEP che ospitano carichi di lavoro attestati su tale VNI inviano segnalazioni IGMP/PIM per unirsi all'albero di distribuzione multicast corrispondente.
    \item Il VTEP di ingresso incapsula il pacchetto BUM una sola volta, impostando l'indirizzo del gruppo multicast nell'header IP esterno. Gli apparati di rete dell'underlay duplicano il frame esclusivamente in corrispondenza dei nodi di diramazione dell'albero.
\end{itemize}

% [IMMAGINE DEL FILE DI SLIDE CON NOME "CloudComputing_5"] Topologia di Multicast Replication con gruppo multicast nell'underlay IP che unisce i VTEP — Slide 41 / Pagina 11 del PDF

La costruzione dell'albero di distribuzione si appoggia a protocolli di multicast routing quali \textbf{PIM-SM} (\emph{Sparse Mode}), \textbf{PIM-SSM} (\emph{Source-Specific Multicast}) o più frequentemente \textbf{PIM-BiDir} (\emph{Bidirectional PIM}).

% [IMMAGINE DEL FILE DI SLIDE CON NOME "CloudComputing_5"] Dettaglio di inoltro BUM su Spine-Leaf con albero multicast e ruolo del Bud Node su Leaf 4 — Slide 43 / Pagina 13 del PDF

\paragraph{Concetti Architetturali Avanzati nel Modello Multicast}:
\begin{itemize}
    \item \textbf{Bridge Domain (BD)}: Dominio logico di trasmissione interno allo switch che definisce il perimetro di broadcast. Sostituisce l'accoppiamento rigido basato su IEEE 802.1Q e permette di mappare diversi ID VLAN locali sullo stesso identificatore globale VNI.
    \item \textbf{NVE (Network Virtual Interface)}: Interfaccia logica virtuale in cui risiedono l'indirizzo IP del VTEP e le funzioni hardware di incapsulamento e decapsulamento.
    \item \textbf{Bud Node}: Un nodo leaf (come Leaf~4 nell'esempio) che opera simultaneamente sia da switch di transito multicast per altri leaf sia da terminazione di servizio: esso deve inoltrare la copia del pacchetto lungo l'albero di distribuzione e contemporaneamente decapsularla per i carichi di lavoro locali.
    \item \textbf{Condivisione degli Alberi Multicast}: Data la discrepanza tra il numero massimo di gruppi multicast gestibili dall'hardware (poche migliaia) e i 16 milioni di VNI teorici, più VNI possono condividere il medesimo gruppo multicast nell'underlay. La segregazione del traffico è garantita poiché il ricevente valida sempre il campo VNI dell'header VxLAN; tuttavia, un albero condiviso comporta uno spreco di banda qualora un VTEP riceva traffico per un VNI che non possiede localmente.
\end{itemize}

\subsubsection{3. Centralized Replication (Flood Proxy)}
Per superare i limiti computazionali della replica alla sorgente e l'onere gestionale del multicast nell'underlay, si adotta talvolta un'architettura basata su un \textbf{Flood Proxy} (o \emph{Flood Gateway}):
\begin{itemize}
    \item Il VTEP di ingresso invia una singola copia unicast del frame BUM verso un nodo centrale dedicato (il replicatore centralizzato).
    \item Il replicatore centrale riceve il pacchetto, valida l'integrità dei parametri, lo decapsula e lo incapsula nuovamente replicandolo verso tutti i VTEP periferici registrati.
    \item Per preservare l'apprendimento del MAC sul piano dati, il replicatore reimposta come indirizzo IP sorgente esterno l'indirizzo IP del VTEP di ingresso originale ($\text{IP}_{\text{Leaf\_1}}$), garantendo la corretta convergenza della FIB sui nodi riceventi.
\end{itemize}

% [IMMAGINE DEL FILE DI SLIDE CON NOME "CloudComputing_5"] Forwarding BUM centralizzato con flood proxy configurato su Leaf 4 — Slide 44 / Pagina 14 del PDF

\textbf{Vantaggi}: Minimizza il consumo di banda sui VTEP di accesso e annulla la necessità di protocolli multicast nell'underlay. \\
\textbf{Svantaggi}: Il nodo replicatore centrale rappresenta un potenziale collo di bottiglia prestazionale (\emph{bottleneck}) e un punto di potenziale vulnerabilità sistemica (\emph{single point of failure}).

\subsection{Inoltro Del Traffico Unicast Noto (Known Unicast)}

Una volta che il processo di apprendimento sul piano dati o sul piano di controllo ha popolato le tabelle di inoltro dei VTEP, il traffico applicativo viaggia esclusivamente come unicast punto-punto.

% [IMMAGINE DEL FILE DI SLIDE CON NOME "CloudComputing_5"] Flusso di risposta unicast diretto (Known Unicast) tra VTEP_3 e VTEP_1 con incapsulamento UDP punto-punto — Slide 46 / Pagina 16 del PDF

Nel momento in cui la macchina $\text{VM}_C$ risponde a $\text{VM}_A$, il frame ha come $\text{DMAC} = \text{MAC}_A$. Il dispositivo $\text{VTEP}_3$, avendo memorizzato che $\text{MAC}_A$ è raggiungibile tramite l'indirizzo underlay $\text{IP}_{\text{VTEP\_1}}$, genera un pacchetto con:
\begin{itemize}[noitemsep]
    \item $\text{Outer Source IP} = \text{IP}_{\text{VTEP\_3}}$, $\text{Outer Destination IP} = \text{IP}_{\text{VTEP\_1}}$;
    \item $\text{Outer Destination UDP Port} = 4789$;
    \item $\text{VNI} = 5000$.
\end{itemize}
Il pacchetto attraversa l'underlay come puro traffico IP unicast; all'arrivo, $\text{VTEP}_1$ rimuove gli header esterni e recapita il payload originale alla porta di $\text{VM}_A$.

\newpage

% ==============================================================================
% SEZIONE 3: PRATICA LINUX ED IMPLEMENTAZIONE DEI VTEP
% ==============================================================================

\section{Implementazione Pratica Di VxLAN In Ambiente Linux}

Il sottosistema di rete del kernel Linux offre un supporto nativo di altissimo livello per le interfacce VxLAN tramite la suite di strumenti standard \texttt{iproute2}. In un ambiente di laboratorio virtualizzato (quale VirtualBox con rete interna dedicata su subnet underlay \texttt{192.168.56.0/24}), un'interfaccia VTEP può essere configurata direttamente all'interno delle macchine virtuali.

\subsection{Configurazione Punto-a-Punto Senza Bridge}

Nello scenario più elementare, si intende connettere due host tramite un tunnel punto-punto bidirezionale assegnando l'indirizzo IP overlay direttamente all'interfaccia virtuale VxLAN. In questo caso non è richiesta la presenza di un bridge di sistema: il tunnel opera a tutti gli effetti come un'interfaccia di rete virtuale dotata di indirizzo di Livello~3.

% [IMMAGINE DEL FILE DI SLIDE CON NOME "CloudComputing_5"] Topologia di laboratorio Linux con due nodi VTEP diretti (vxlan10) su rete underlay 192.168.56.0/24 — Slide 47 / Pagina 17 del PDF

\paragraph{Comandi Operativi per la Creazione}:
Su ciascun nodo viene creata l'interfaccia specificando il VNI identificativo (\texttt{id 10}), l'indirizzo underlay remoto e l'interfaccia fisica di ancoraggio:

\begin{lstlisting}[language=bash, caption={Configurazione VTEP point-to-point su Node 1 e Node 2.}]
# --- CONFIGURAZIONE SU NODO 1 (Underlay: 192.168.56.10) ---
# Creazione dell'interfaccia VxLAN associata al VNI 10 verso Node 2
ip link add vxlan10 type vxlan id 10 remote 192.168.56.20 \
  local 192.168.56.10 dstport 4789 dev enp0s3

# Assegnazione dell'indirizzo IP del segmento overlay
ip addr add 10.10.10.1/24 dev vxlan10

# Attivazione dell'interfaccia
ip link set vxlan10 up

# --- CONFIGURAZIONE SU NODO 2 (Underlay: 192.168.56.20) ---
# Creazione dell'interfaccia VxLAN speculare verso Node 1
ip link add vxlan10 type vxlan id 10 remote 192.168.56.10 \
  local 192.168.56.20 dstport 4789 dev enp0s3

# Assegnazione dell'indirizzo IP overlay su Node 2
ip addr add 10.10.10.2/24 dev vxlan10

# Attivazione dell'interfaccia
ip link set vxlan10 up
\end{lstlisting}

Al termine della procedura, un comando di verifica della connettività overlay (\texttt{ping 10.10.10.2} eseguito da Node~1) genererà pacchetti ICMP incapsulati in trame Ethernet, a loro volta racchiuse nei pacchetti UDP diretti alla porta 4789 dell'underlay.

\subsection{Configurazione Con Linux Bridge Per Reti Multi-Host}

Qualora l'host Linux debba collegare al medesimo segmento overlay una pluralità di macchine virtuali locali, container (ad esempio veth pair) o interfacce fisiche aggiuntive, la semplice interfaccia punto-punto non è più sufficiente: è necessario introdurre un'istanza di **Linux Bridge** (\texttt{type bridge}). L'interfaccia VxLAN assume in tal caso il ruolo di porta di uplink del bridge software.

% [IMMAGINE DEL FILE DI SLIDE CON NOME "CloudComputing_5"] Topologia Linux con interfaccia vxlan10 associata come porta subordinata a un bridge software br-vxlan10 — Slide 49 / Pagina 19 del PDF

In tale assetto strutturale, l'indirizzo IP di gestione dell'overlay non viene assegnato all'interfaccia VxLAN ma direttamente al bridge, che funge da punto di convergenza per il forwarding locale di Livello~2.

\begin{lstlisting}[language=bash, caption={Integrazione tra interfaccia VxLAN e Linux Bridge su Node 1 e Node 2.}]
# --- CONFIGURAZIONE SU NODO 1 ---
# 1. Creazione dell'interfaccia VxLAN
ip link add vxlan10 type vxlan id 10 remote 192.168.56.20 \
  local 192.168.56.10 dstport 4789 dev enp0s3

# 2. Creazione dell'istanza Bridge software
ip link add br-vxlan10 type bridge

# 3. Asservimento dell'interfaccia VxLAN al bridge
ip link set vxlan10 master br-vxlan10

# 4. Attivazione delle interfacce
ip link set vxlan10 up
ip link set br-vxlan10 up

# 5. Assegnazione dell'indirizzo IP overlay al bridge
ip addr add 10.10.10.1/24 dev br-vxlan10

# --- CONFIGURAZIONE SU NODO 2 ---
# 1. Creazione dell'interfaccia VxLAN
ip link add vxlan10 type vxlan id 10 remote 192.168.56.10 \
  local 192.168.56.20 dstport 4789 dev enp0s3

# 2. Creazione dell'istanza Bridge software
ip link add br-vxlan10 type bridge

# 3. Asservimento al bridge
ip link set vxlan10 master br-vxlan10

# 4. Attivazione delle interfacce
ip link set vxlan10 up
ip link set br-vxlan10 up

# 5. Assegnazione dell'indirizzo IP overlay al bridge
ip addr add 10.10.10.2/24 dev br-vxlan10
\end{lstlisting}

\newpage

% ==============================================================================
% SEZIONE 4: EVOLUZIONE DEL CONTROL PLANE: EVPN IN VXLAN
% ==============================================================================

\section{Evoluzione Del Control Plane: Ethernet VPN (EVPN) In VxLAN}

I meccanismi nativi basati su tunneling manuale e data-plane \emph{Flood and Learn} presentano criticità strutturali insormontabili nelle architetture cloud hyperscale:
\begin{itemize}[noitemsep]
    \item \textbf{Carico di configurazione manuale}: richiede la definizione statica di elenchi di peer VTEP su ciascun dispositivo, operazione insostenibile in ambienti dinamici con continue creazioni e migrazioni di VM.
    \item \textbf{Saturazione da traffico BUM}: l'apprendimento basato sul flooding genera tempeste di broadcast e un'occupazione massiccia della banda di underlay, vanificando i benefici dell'architettura a tunnel.
\end{itemize}

Per ovviare a tali inefficienze, lo standard introduce \textbf{EVPN} (\emph{Ethernet Virtual Private Network}, formalizzato nella specifica RFC~7432). EVPN dota VxLAN di un \textbf{piano di controllo distribuito} standardizzato basato sull'estensione multiprotocollo del Border Gateway Protocol (\textbf{MP-BGP}). In questo modello integrato, BGP non si limita a propagare la raggiungibilità degli indirizzi IP, ma trasporta esplicitamente gli indirizzi MAC di Livello~2 e i parametri di binding dei VTEP.

\subsection{Comportamento Del Traffico BUM In Presenza Di EVPN}

L'introduzione di EVPN trasforma l'apprendimento delle informazioni di raggiungibilità da un paradigma reattivo basato sui dati a un paradigma proattivo governato dal piano di controllo. Nonostante la segnalazione sia accentrata in BGP, la necessità fisica di trasportare pacchetti destinati a più stazioni (\emph{one-to-all} o \emph{one-to-many}) permane. La gestione dei singoli flussi BUM evolve come segue:

\begin{enumerate}
    \item \textbf{Unknown Unicast (Quasi Completamente Eliminato)}:
    Grazie alla diffusione tempestiva delle rotte EVPN di tipo 2 (\textbf{Route Type 2 - MAC/IP Advertisement Route}), ciascun VTEP riceve in tempo reale la collocazione esatta di ogni nuovo indirizzo MAC attivato nella rete. Di conseguenza, il VTEP non deve ricorrere al flooding per localizzare la destinazione: se un MAC non compare nella tabella di inoltro, esso viene trattato come non esistente (salvo rari casi di transitorio), azzerando l'unknown unicast nel fabric.

    \item \textbf{Broadcast e Richieste ARP (Fortemente Ridotto)}:
    Viene abilitata la funzionalità di \textbf{ARP Suppression} (o \emph{ARP Proxy}). Quando un host locale genera una richiesta broadcast ARP per un determinato IP, il VTEP ToR locale intercetta la richiesta a livello di accesso. Se l'associazione MAC/IP è già nota al VTEP grazie alle tabelle popolate da BGP EVPN, il VTEP stesso sintetizza la risposta ARP e la invia all'host mittente. La richiesta broadcast originale non viene mai incapsulata e non attraversa i tunnel di underlay. Solo nel caso in cui la corrispondenza risulti sconosciuta nel database EVPN, il pacchetto viene propagato.

    \item \textbf{Multicast e Broadcast Residuo (Replicazione Controllata)}:
    I pacchetti intrinsecamente multi-destinatario che non possono essere risolti localmente (ad esempio streaming multicast o discovery DHCP) devono essere comunque inoltrati attraverso la rete. EVPN demanda il trasporto fisico a una delle due modalità sotto-stanti:
    \begin{itemize}[noitemsep]
        \item \textbf{Ingress Replication}: il VTEP genera $N$ repliche unicast; approccio ottimale per reti di dimensioni medio-piccole poiché non vincola la configurazione dell'underlay, pur consumando banda di trasmissione locale.
        \item \textbf{Multicast Replication}: la copia singola viene veicolata su un gruppo multicast di underlay governato da protocolli PIM, garantendo la massima efficienza trasmissiva nelle reti su larghissima scala al prezzo di una maggiore complessità gestionale.
    \end{itemize}
\end{enumerate}

\newpage

% ==============================================================================
% SEZIONE 5: FONDAMENTI DEL PROTOCOLLO BGP COME PIANO DI CONTROLLO
% ==============================================================================

\section{Fondamenti Del Protocollo BGP Come Piano Di Controllo}

Per comprendere appieno il funzionamento di EVPN, è imprescindibile dominare i meccanismi operativi del protocollo sottostante: il **Border Gateway Protocol** (BGP). Prima di esaminarne la configurazione e i formati di messaggio, richiamiamo la differenza concettuale tra instradamento intra-dominio e inter-dominio:
\begin{itemize}[noitemsep]
    \item I protocolli **IGP** (\emph{Interior Gateway Protocol}, come OSPF, IS-IS e RIP) operano all'interno di un singolo dominio amministrativo e si concentrano sulla topologia e sul calcolo del percorso a costo minimo (fornendo la \emph{raggiungibilità interna}).
    \item I protocolli **EGP** (\emph{Exterior Gateway Protocol}), di cui BGP rappresenta lo standard universale de facto, operano al confine tra diversi domini amministrativi (denominati **Sistemi Autonomi**, AS) e si occupano di applicare policy di instradamento complesse e scalabili (fornendo la \emph{conoscenza delle destinazioni mondiali}).
\end{itemize}

% [IMMAGINE DEL FILE DI SLIDE CON NOME "CloudComputing_5"] Architettura a Sistemi Autonomi con interconnessione BGP tra AS 65000, 65001, 65002, 65003 ed esecuzione di IGP interno — Slide 50 / Pagina 23 del PDF

\subsection{Evoluzione Da EGP A BGP E Numerazione Dei Sistemi Autonomi}

Lo storico protocollo EGP presentava una struttura rudimentale: si limitava a trasferire vettori di raggiungibilità senza disporre di logiche atte a prevenire la formazione di anelli di instradamento (\emph{routing loops}) o a selezionare cammini ottimali. BGP ha superato questi limiti introducendo:
\begin{itemize}[noitemsep]
    \item Instaurazione esplicita e affidabile delle relazioni di vicinato (\emph{peering});
    \item Riconoscimento ed eliminazione intrinseca dei cicli mediante l'ispezione della sequenza dei Sistemi Autonomi attraversati (\emph{AS-Path});
    \item Ricco assortimento di attributi metrici associati a ciascun prefisso per un controllo granulare delle politiche di traffico (\emph{policy-based routing}).
\end{itemize}

Un **Sistema Autonomo (AS)** è formalmente definito come un insieme coeso di router e reti sotto la guida di un'unica autorità amministrativa che adotta una politica di routing omogenea. L'autorità \textbf{IANA} (\emph{Internet Assigned Numbers Authority}) assegna i codici identificativi degli AS:
\begin{itemize}[noitemsep]
    \item \textbf{Formato a 2 byte (16 bit)}: utilizzato storicamente (fino a gennaio 2009), con intervallo $[1, 65535]$. Gli identificatori da $1$ a $64511$ sono AS pubblici globalmente unici; i numeri da $64512$ a $65534$ sono riservati per utilizzi privati (analogamente agli indirizzi RFC 1918).
    \item \textbf{Formato a 4 byte (32 bit)}: introdotto per far fronte all'esaurimento dello spazio di indirizzamento a 16 bit, estende l'intervallo da $65536$ fino a $4\,294\,967\,295$.
\end{itemize}

\subsection{Caratteristiche Meccaniche Di BGP: Trasporto Su Canale TCP}

A differenza dei protocolli IGP (quali OSPF o RIP) che utilizzano indirizzi multicast e pacchetti raw IP o frame diretti sul link fisico locale per la scoperta automatica dei vicini, BGP basa interamente la propria affidabilità sul protocollo **TCP**, attestandosi sulla porta ben nota \textbf{179}.

% [IMMAGINE DEL FILE DI SLIDE CON NOME "CloudComputing_5"] Proprietà strutturali di BGP e sessioni di adiacenza stabilite tramite canali TCP su porte 179 — Slide 51 / Pagina 24 del PDF

L'adozione di un canale affidabile orientato alla connessione comporta conseguenze architetturali decisive:
\begin{itemize}[noitemsep]
    \item \textbf{Assenza di Multicast Discovery}: BGP non invia pacchetti "Hello" periodici in multicast. Tutti i peer devono essere configurati staticamente per indirizzo IP.
    \item \textbf{Adiacenze Multi-Hop}: Poiché TCP supporta l'instradamento attraverso nodi intermedi, due router BGP possono negoziare una sessione di adiacenza anche senza essere direttamente connessi dallo stesso cavo o switch fisico, purché l'infrastruttura di routing sottostante garantisca la connettività IP tra i loro endpoint.
\end{itemize}

% [IMMAGINE DEL FILE DI SLIDE CON NOME "CloudComputing_5"] Handshake a tre vie TCP tra RTA (1.1.1.1) e RTB (2.2.2.2) per l'apertura del canale di peering BGP — Slide 52 / Pagina 25 del PDF

Il router che avvia la procedura di connessione agisce da client TCP, impiegando una porta locale casuale per aprire la sessione verso la porta standard 179 dell'interfaccia remota del vicino.

\subsection{Tipologie Di Vicinato: EBGP Vs IBGP}

A seconda della collocazione amministrativa dei peer, si distinguono due modalità di peering con requisiti operativi profondamente differenti:

\begin{figure}[H]
    \centering
    \begin{subfigure}[b]{0.48\textwidth}
        \centering
        % [IMMAGINE DEL FILE DI SLIDE CON NOME "CloudComputing_5"] Topologia EBGP con router RTB e RTC che dialogano verso router esterni RTD ed RTE — Slide 53 / Pagina 26 del PDF
        \caption{Relazione EBGP tra domini AS distinti.}
    \end{subfigure}
    \hfill
    \begin{subfigure}[b]{0.48\textwidth}
        \centering
        % [IMMAGINE DEL FILE DI SLIDE CON NOME "CloudComputing_5"] Topologia IBGP interna con peering tra RTB e RTC attraverso l'infrastruttura OSPF di RTA — Slide 54 / Pagina 27 del PDF
        \caption{Relazione IBGP all'interno del medesimo AS.}
    \end{subfigure}
    \caption{Confronto topologico tra sessioni BGP esterne (EBGP) e interne (IBGP).}
\end{figure}

\begin{enumerate}
    \item \textbf{EBGP (External BGP)}: sessione stabilita tra router appartenenti a Sistemi Autonomi differenti (ad esempio tra RTB in AS~100 ed RTD in AS~200). Solitamente le sessioni EBGP vengono instaurate tra indirizzi IP appartenenti a interfacce fisiche direttamente connesse (\emph{single-hop}), con valore del Time-To-Live (TTL) IP impostato di default a 1.
    \item \textbf{IBGP (Internal BGP)}: sessione stabilita tra router appartenenti al medesimo Sistema Autonomo (ad esempio tra RTB ed RTC all'interno dell'AS~100).
\end{enumerate}

\paragraph{Perché Utilizzare IBGP invece di Ridistribuire in IGP?}
Una soluzione ingenua per propagare all'interno dell'AS~100 le rotte apprese esternamente da AS~200 e AS~300 consisterebbe nella ridistribuzione (\emph{route redistribution}) delle rotte BGP all'interno del protocollo IGP locale (es. OSPF) eseguito sui router RTB, RTA ed RTC. Tale approccio è categoricamente sconsigliato in ingegneria di rete per i seguenti motivi:
\begin{itemize}[noitemsep]
    \item \textbf{Scalabilità della memoria e della CPU}: la tabella di routing BGP pubblica globale contiene centinaia di migliaia di prefissi (Full Internet Routing Table). Gli algoritmi basati sullo stato dei collegamenti come OSPF (Dijkstra SPF) non sono dimensionati per elaborare e ricalcolare grafi topologici contenenti tali volumi di dati, rischiando il blocco sistemico dei router interni.
    \item \textbf{Perdita irreversibile degli attributi}: i protocolli IGP supportano unicamente metriche scalari di costo (bandwidth/ritardo). Ridistribuendo una rotta BGP in OSPF, tutti i ricchi attributi di policy (AS-Path, Local Preference, Multi-Exit Discriminator, Community) andrebbero irrimediabilmente perduti.
    \item \textbf{Complessità gestionale}: richiede complesse mappe di instradamento e filtri su ogni nodo di confine per prevenire anelli di ridistribuzione bidirezionale.
\end{itemize}

La strategia corretta consiste nel mantenere disaccoppiati i due livelli:
\begin{itemize}[noitemsep]
    \item \textbf{L'IGP (OSPF)} è eseguito su tutti i nodi interni (incluso RTA, che non necessita di parlare BGP) per fornire la pura \emph{reachability} di transito tra gli indirizzi IP dei router interni.
    \item \textbf{BGP (IBGP)} viene instaurato logicamente tra i router di confine (RTB ed RTC) per propagare la \emph{knowledge} delle reti esterne.
\end{itemize}

\paragraph{Pipeline di Inoltro BGP e Tabelle di Routing}:
Quando un router riceve molteplici annunci per la medesima destinazione, la selezione della rotta avviene attraverso stadi sequenziali ben definiti:
\begin{enumerate}[noitemsep]
    \item \textbf{BGP RIB (Loc-RIB)}: archivio locale in cui vengono memorizzate tutte le rotte valide apprese da ciascun vicino BGP.
    \item \textbf{BGP Decision Process}: algoritmo deterministico a cascata che valuta gli attributi BGP (Next-Hop reachability, Weight, Local Preference, AS-Path length, Origin, MED, preferenza EBGP rispetto a IBGP, costo IGP verso il Next-Hop) per identificare il percorso migliore (\emph{Best Path}).
    \item \textbf{Global Routing Table (Global RIB / FIB)}: solo la rotta proclamata come "best path" viene effettivamente inviata al gestore della tabella di routing di sistema ed installata nella FIB hardware per l'inoltro dei pacchetti.
\end{enumerate}

\subsection{Configurazione Delle Adiacenze E Best Practice Con Loopback}

La configurazione delle relazioni di peering richiede la dichiarazione esplicita del BGP Identifier univoco del router (\textbf{Router ID}) e la specifica dei parametri di sessione.

\paragraph{Regole di Elezione Automatica del Router ID}:
Qualora il parametro non venga assegnato manualmente dall'amministratore tramite l'istruzione \texttt{router id X.X.X.X}, il software seleziona l'identificatore applicando il seguente ordine tassativo:
\begin{enumerate}[noitemsep]
    \item Viene selezionato l'indirizzo IPv4 numericamente più alto configurato tra tutte le interfacce logiche di \textbf{Loopback} attive.
    \item Se non è presente alcuna interfaccia di loopback, viene selezionato l'indirizzo IPv4 numericamente più alto configurato tra le interfacce \textbf{fisiche} operative.
\end{enumerate}

% [IMMAGINE DEL FILE DI SLIDE CON NOME "CloudComputing_5"] Topologia di configurazione di adiacenza BGP tra nodi RTB, RTC, RTD, RTE — Slide 55 / Pagina 28 del PDF

La natura della sessione è determinata dal confronto tra l'AS locale e l'AS dichiarato nella direttiva \texttt{peer}: se i valori coincidono, la sessione viene classificata e gestita come IBGP; se differiscono, viene negoziata come EBGP.

% [IMMAGINE DEL FILE DI SLIDE CON NOME "CloudComputing_5"] Ottimizzazione dell'adiacenza IBGP tra RTB e RTC mediante interfacce di Loopback 0 e connect-interface — Slide 56 / Pagina 29 del PDF

\paragraph{Ottimizzazione delle Sessioni IBGP tramite Interfacce di Loopback}:
Nelle configurazioni di base tra router adiacenti, le sessioni IBGP possono essere attestate sugli indirizzi fisici dei link punto-punto. Tuttavia, se il collegamento fisico di connessione diretta subisce un guasto, la sessione TCP sottostante decade immediatamente, provocando il collasso di tutti gli annunci di rotta associati.

La best practice universale prescrive di stabilire le sessioni IBGP impiegando gli indirizzi delle interfacce virtuali di \textbf{Loopback} (ad esempio \texttt{2.2.2.2} per RTB e \texttt{3.3.3.3} per RTC) e forzando l'indirizzo di bind del socket tramite il comando dedicato:
\begin{lstlisting}[language=bash, caption={Configurazione ottimizzata di sessione BGP con interfaccia di Loopback.}]
# --- CONFIGURAZIONE SU ROUTER RTC (AS 100) ---
router id 3.3.3.3
bgp 100
 # Peering EBGP verso RTE su interfaccia fisica diretta
 peer 10.1.35.5 as-number 300
 # Peering IBGP verso RTB attestato su indirizzo di Loopback
 peer 2.2.2.2 as-number 100
 peer 2.2.2.2 connect-interface loopback 0
\end{lstlisting}
Poiché l'interfaccia di loopback è puramente logica e non decade mai a meno di uno spegnimento manuale del dispositivo, finché la rete interna governata da OSPF possiede un cammino fisico alternativo per raggiungere l'IP \texttt{2.2.2.2}, la sessione TCP di BGP rimane perfettamente stabile ed operativa, prevenendo oscillazioni distruttive nella propagazione delle rotte (\emph{route flapping}).

\subsection{Macchina A Stati Finiti (FSM) Per L'Instaurazione Del Peering}

La creazione di una relazione operativa di vicinato BGP è regolata da una rigida macchina a stati finiti (FSM), le cui transizioni sequenziali sono scandite da eventi TCP e messaggi di protocollo:

% [IMMAGINE DEL FILE DI SLIDE CON NOME "CloudComputing_5"] Diagramma temporale della macchina a stati BGP: Idle, Connect, OpenSent, OpenConfirm, Established — Slide 57 / Pagina 30 del PDF

La dinamica degli stati procede come descritto di seguito:
\begin{enumerate}
    \item \textbf{Idle}: stato iniziale di riposo. Il router alloca le risorse e resta in attesa di un evento di avvio configurato (ad esempio l'abilitazione del comando di peer).
    \item \textbf{Connect}: il router tenta l'apertura della connessione a tre vie TCP (\emph{Three-Way Handshake}: SYN, SYN-ACK, ACK) verso la porta 179 dell'interfaccia del vicino.
    \item \textbf{OpenSent}: completato l'handshake TCP, il router trasmette immediatamente e in modo indipendente il proprio messaggio di benvenuto \textbf{OPEN} contenente i parametri di negoziazione, posizionandosi in ascolto del messaggio speculare emesso dall'omologo.
    \item \textbf{OpenConfirm}: ricevuto l'OPEN del vicino, il nodo ne valida la conformità (correttezza dell'AS number, versione di protocollo, valore del timer di hold e unicità del BGP Identifier). In caso di verifica positiva, risponde inviando un messaggio di conferma \textbf{KEEPALIVE}.
    \item \textbf{Established}: ricevuto il Keepalive di conferma dal peer, la sessione entra stabilmente nello stato operativo finale. Solo a questo punto i router iniziano la disseminazione e l'aggiornamento massivo dei cammini di routing scambiando messaggi \textbf{UPDATE}.
\end{enumerate}

\subsection{Tassonomia Dettagliata Dei Messaggi BGP}

La specifica BGP definisce tassativamente cinque tipologie di messaggi per regolare l'intero ciclo di vita della connessione:

\begin{table}[H]
    \centering
    \small
    \begin{tabularx}{\textwidth}{lp{4cm}X}
        \toprule
        \textbf{Tipo Messaggio} & \textbf{Finalità Operativa} & \textbf{Parametri e Meccanismi Chiave} \\
        \midrule
        \textbf{OPEN} & Negoziazione iniziale dei parametri di sessione. & Versione BGP (standard v4), Numero di AS locale, Hold Time negoziato, BGP Router Identifier e Capacità estese (Multiprotocol extensions per EVPN, IPv6). \\
        \addlinespace
        \textbf{UPDATE} & Trasmissione proattiva o revoca dei prefissi di routing. & - Annuncio di rotte raggiungibili congiunte al medesimo set di Path Attributes (Local Pref, AS-Path, Next Hop) per i prefissi NLRI (\emph{Network Layer Reachability Information}). \newline - Ritiro esplicito di rotte non più raggiungibili (\emph{Withdrawn Routes}) specificando unicamente i prefissi IP. \\
        \addlinespace
        \textbf{KEEPALIVE} & Manutenzione dello stato di adiacenza e prevenzione del timeout. & Pacchetto privo di payload dati (soli 19 byte di header BGP standard). Trasmesso periodicamente a cadenza pari a un terzo del valore di Hold Time (default tipico: invio ogni 60s con Hold Time a 180s). \\
        \addlinespace
        \textbf{NOTIFICATION} & Segnalazione di eccezioni o errori fatali e chiusura sessione. & Inviato prima della terminazione unilaterale del socket TCP. Trasporta un Error Code (es. errore di sincronizzazione, formato OPEN errato, scadenza Hold Timer) e relativi dati diagnostici. \\
        \addlinespace
        \textbf{ROUTE-REFRESH} & Richiesta di ritrasmissione completa degli annunci di rotta. & Definito nella RFC 2918; consente a un router di richiedere al vicino la ri-pubblicazione delle rotte a seguito di modifiche alle politiche di filtro locali, evitando la disconnessione e ricostruzione traumatica del canale TCP. \\
        \bottomrule
    \end{tabularx}
    \caption{Tassonomia e compiti funzionali dei messaggi del protocollo BGP.}
\end{table}

Grazie alla flessibilità intrinseca del messaggio di tipo \textbf{UPDATE} e alla sua capacità di veicolare strutture NLRI generalizzate tramite estensioni multi-protocollo (BGP-4 MP Extensions, RFC 4760), BGP si configura come il veicolo architetturale ideale per il controllo delle reti VxLAN distribuite: sostituendo la semplice reachability di indirizzi IP univariati con le definizioni complesse delle rotte EVPN (Route Type 2 per indirizzi MAC/IP e Route Type 3 per l'auto-discovery dei tunnel di Ingress Replication), il fabric di data center converte integralmente un sistema di bridging storicamente caotico e orientato al flooding in una rete instradata deterministica, altamente scalabile e pienamente integrata.

% DD ----------------------------------------------

% =============================================================================
% SEZIONE: BGP IMPLEMENTAZIONE, DATABASE E FLUSSO ELABORATIVO
% =============================================================================

\section{Architettura Di Un Router BGP: Strutture Dati E Basi Dati Di Instradamento}

Prima di analizzare i meccanismi di selezione e manipolazione del percorso all'interno del \emph{Border Gateway Protocol} (BGP), è fondamentale comprendere come un router (definito formalmente \textbf{BGP speaker}) organizza le informazioni ricevute dall'esterno e quelle destinate alla rete locale.

Nello stack architetturale di un dispositivo di rete coesistono protocolli di \emph{routing} interno (IGP quali OSPF o IS-IS) e protocolli di frontiera (BGP). Per operare in maniera deterministica e scalabile, BGP non scrive direttamente gli aggiornamenti di percorso nella tabella di inoltro del kernel (\emph{Forwarding Information Base}, FIB), né sovrascrive istantaneamente la tabella di instradamento globale del sistema operativo. Al contrario, implementa un'architettura modulare a database stratificati, nota come \textbf{Routing Information Base} (RIB), specificata dallo standard RFC~4271.

% [IMMAGINE DEL FILE DI SLIDE CON NOME "CloudComputing_5"] Struttura dei database BGP (IP-RIB, Loc-RIB, BGP Peer Database, Adj-RIB-In, Adj-RIB-Out) — Slide 59 / Pagina 1 del PDF

La memoria e la logica di memorizzazione di un nodo BGP si articolano in cinque entità primarie:
\begin{itemize}[noitemsep]
    \item \textbf{Adj-RIB-In} (\emph{Adjacency Routing Information Base, Incoming}): memorizza le informazioni di instradamento grezze, non elaborate, ricevute tramite messaggi di \texttt{UPDATE} da ciascun peer BGP connesso. Esiste una struttura Adj-RIB-In logica (o fisica) per ogni adiacenza BGP attiva. Le rotte qui contenute non sono state ancora sottoposte ad alcun filtro locale di ingresso.
    \item \textbf{Input Policy Engine}: rappresenta il motore di filtraggio e manipolazione in ingresso. Valuta le rotte presenti nell'Adj-RIB-In attraverso \emph{route-policy}, \emph{route-map} o filtri di prefisso/comunità. È in questa fase che gli attributi possono essere alterati (ad esempio assegnando una preferenza locale) o interi prefissi possono essere scartati.
    \item \textbf{Loc-RIB} (\emph{Local Routing Information Base}): è la tabella di instradamento BGP locale del router. Contiene esclusivamente le rotte ottimali selezionate dal motore decisionale (\emph{BGP path selection algorithm}) applicato all'insieme delle rotte ammesse dall'Input Policy Engine, integrate con le rotte originate localmente dal router stesso (\emph{locally originated routes}).
    \item \textbf{IP-RIB} (\emph{IP Routing Information Base}): è la tabella di instradamento globale di livello 3 del dispositivo di rete. Raccoglie i migliori candidati proposti da tutti i protocolli attivi (BGP, OSPF, rotte statiche e connesse). Sulla base della distanza amministrativa o preferenza di protocollo, l'IP-RIB seleziona la rotta definitiva da installare nella FIB hardware per il \emph{forwarding} dei pacchetti dati.
    \item \textbf{Output Policy Engine}: motore di policy che processa le rotte presenti nella Loc-RIB prima che vengano inviate all'esterno. Decide se annunciare o sopprimere un prefisso verso un dato vicino e può riscrivere gli attributi (ad esempio aggiungendo comunità, modificando il \texttt{MED} o eseguendo l'\emph{AS-Path prepending}).
    \item \textbf{Adj-RIB-Out} (\emph{Adjacency Routing Information Base, Outgoing}): contiene i percorsi che il router ha validato per l'annuncio a uno specifico peer adiacente, già processati e pronti per essere serializzati all'interno di messaggi \texttt{UPDATE}.
    \item \textbf{BGP Peer Database}: archivio interno che conserva lo stato della macchina a stati finiti (FSM), le capacità negoziate, i contatori di keepalive/hold time e i parametri di trasporto TCP di ciascun vicino con cui è instaurata una sessione di \emph{peering}.
\end{itemize}

\subsection{Flusso Elaborativo Delle Informazioni Di Instradamento}

Il ciclo di vita di un annuncio di rete attraversa una sequenza deterministica a stadi, schematizzata nel flusso operativo di elaborazione BGP:

% [IMMAGINE DEL FILE DI SLIDE CON NOME "CloudComputing_5"] Diagramma del flusso di elaborazione delle rotte BGP da Adj-RIB-In attraverso le policy fino a Loc-RIB, IP-RIB e Adj-RIB-Out — Slide 60 / Pagina 2 del PDF

Il trattamento dei percorsi si articola nei seguenti passi:
\begin{enumerate}
    \item \textbf{Ricezione ed archiviazione iniziale} (\ding{182}): Quando il BGP speaker riceve un messaggio di \texttt{UPDATE} da un peer adiacente, estrae i prefissi e i relativi attributi di percorso, registrandoli nella struttura \textbf{Adj-RIB-In} associata a quello specifico peer.
    \item \textbf{Filtraggio in ingresso} (\ding{183}): Il motore delle policy in ingresso (\textbf{Input Policy Engine}) esamina le rotte dell'Adj-RIB-In. Se le regole locali (ad esempio ACL, filtri su prefissi o espressioni regolari su \texttt{AS\_Path}) bocciano il prefisso, questo viene scartato; se invece la policy ne ammette l'accettazione, possono essere modificate metriche interne come la \texttt{Local\_Pref}.
    \item \textbf{Esecuzione dell'algoritmo di selezione del percorso} (\ding{184}): Tra tutti i cammini validi pervenuti da peer diversi e le rotte iniettate localmente per lo stesso prefisso di destinazione, il router invoca l'algoritmo di \emph{BGP path selection} per eleggere un'unica rotta ottima (\emph{best path}). Tale rotta viene memorizzata nella \textbf{Loc-RIB}.
    \item \textbf{Installazione nell'IP-RIB} (\ding{185}): Il \emph{best path} estratto dalla Loc-RIB viene sottoposto all'IP-RIB globale. Se non esistono rotte con distanza amministrativa migliore per lo stesso prefisso, la rotta viene definitivamente installata nella tabella di instradamento del sistema e inviata alla FIB. Nel caso in cui sul nodo sia attivo il supporto \textbf{BGP Multipath}, il router sottomette alla tabella IP-RIB non solo la rotta ottima principale, ma anche tutte le rotte di pari costo (\emph{equal-cost}), permettendo il bilanciamento del traffico.
    \item \textbf{Filtraggio in uscita e propagazione} (\ding{186}): Dalla Loc-RIB, le rotte ottimali e le rotte originate localmente vengono inoltrate all'\textbf{Output Policy Engine}. Soltanto i prefissi che superano indenni i criteri stabiliti dalle regole di esportazione vengono scritti nell'\textbf{Adj-RIB-Out} e successivamente inviati tramite messaggi \texttt{UPDATE} ai rispettivi vicini BGP.
\end{enumerate}

% [IMMAGINE DEL FILE DI SLIDE CON NOME "CloudComputing_5"] Schema sintetico dell'elaborazione di instradamento: Update in ingresso, Route Selection verso Local_RIB e IP_RIB, e generazione Update in uscita — Slide 62 / Pagina 4 del PDF

Da questa rigorosa gerarchia discende una regola cardine dell'architettura BGP: \emph{un router non annuncia mai a un vicino percorsi che non abbia preventivamente selezionato e installato come ottimali nella propria Loc-RIB}.

% =============================================================================
% SEZIONE: IL MESSAGGIO UPDATE E GLI ATTRIBUTI BGP
% =============================================================================

\section{Il Messaggio UPDATE E La Classificazione Degli Attributi Di Percorso}

A differenza dei protocolli IGP di tipo link-state, che propagano la topologia della rete tramite flooded \emph{Link State Advertisements}, BGP opera come protocollo di tipo \textbf{Path-Vector}. L'intero scambio delle informazioni di raggiungibilità e delle modifiche topologiche è affidato al messaggio di \texttt{UPDATE}.

% [IMMAGINE DEL FILE DI SLIDE CON NOME "CloudComputing_5"] Struttura logica e campi del messaggio BGP UPDATE (NLRI, Path Attributes, Withdrawn Routes) — Slide 61 / Pagina 3 del PDF

Un messaggio di \texttt{UPDATE} trasporta tre componenti informative primarie:
\begin{enumerate}[noitemsep]
    \item \textbf{Network Layer Reachability Information (NLRI)}: elenca una o più coppie formate da prefisso IP e lunghezza della maschera (\emph{prefix length}) che condividono la medesima traiettoria e i medesimi attributi.
    \item \textbf{Path Attributes}: vettori e parametri che descrivono puntualmente la storia del cammino, forniscono meccanismi deterministici di prevenzione dei loop e indirizzano il processo decisionale multi-criterio. In conformità a RFC~4271, gli attributi di percorso sono veicolati \textbf{esclusivamente} all'interno dei messaggi \texttt{UPDATE}.
    \item \textbf{Withdrawn Routes}: insieme di prefissi e relative lunghezze di maschera che il mittente non è più in grado di raggiungere, consentendo la cancellazione immediata delle rotte non più attive senza attendere la scadenza di timer di hold.
\end{enumerate}

\subsection{Tassonomia Degli Attributi Di Percorso (Path Attributes)}

Gli attributi di percorso BGP non presentano tutti il medesimo grado di obbligatorietà o trasparenza nel transito tra sistemi autonomi differenti. La specifica li suddivide in quattro categorie formali:

% [IMMAGINE DEL FILE DI SLIDE CON NOME "CloudComputing_5"] Matrice tassonomica dei BGP Path Attributes: Well-known vs Optional, Mandatory vs Discretionary, Transitive vs Non-transitive — Slide 63 / Pagina 5 del PDF

\begin{table}[H]
\centering
\small
\begin{tabularx}{\textwidth}{l p{5.5cm} X}
\toprule
\textbf{Categoria} & \textbf{Definizione e Requisiti} & \textbf{Esempi Noti} \\
\midrule
\textbf{Well-known Mandatory} &
Devono essere necessariamente riconosciuti e supportati da qualsiasi implementazione BGP conforme. \textbf{Devono essere obbligatoriamente inclusi} in ogni messaggio di \texttt{UPDATE} che contenga NLRI raggiungibili. &
\texttt{ORIGIN}, \texttt{AS\_PATH}, \texttt{NEXT\_HOP} \\
\addlinespace
\textbf{Well-known Discretionary} &
Devono essere universalmente riconosciuti da ogni router BGP, ma la loro presenza nel messaggio \texttt{UPDATE} è opzionale e subordinata a precise scelte di configurazione. &
\texttt{LOCAL\_PREF}, \texttt{ATOMIC\_AGGREGATE} \\
\addlinespace
\textbf{Optional Transitive} &
Non è richiesto che tutti i router BGP li comprendano. Se un router intermedio non riconosce l'attributo, ha comunque l'\textbf{obbligo} di accettarlo, mantenerlo inalterato e propagarlo ai peer successivi. &
\texttt{COMMUNITY}, \texttt{AGGREGATOR} \\
\addlinespace
\textbf{Optional Non-transitive} &
Non richiedono un supporto universale. Se un router ricevente non riconosce l'attributo, può semplicemente ignorarlo e \textbf{non deve} propagarlo ulteriormente ai propri vicini. &
\texttt{MED} (\emph{Multi-Exit Discriminator}) \\
\bottomrule
\end{tabularx}
\caption{Classificazione formale dei BGP Path Attributes secondo RFC~4271.}
\end{table}

\begin{theorem}
\textbf{Distinzione Concettuale}: È imperativo non confondere i parametri opzionali negoziati nella fase di apertura della sessione (\emph{Optional Parameters} del messaggio \texttt{OPEN}, utilizzati ad esempio per negoziare il supporto multiprotocollo o ASN a 4 byte) con gli attributi opzionali di percorso (\emph{Optional Path Attributes}), che risiedono unicamente nel corpo dei messaggi \texttt{UPDATE}.
\end{theorem}

% =============================================================================
% SEZIONE: ANALISI DETTAGLIATA DEI PRINCIPALI ATTRIBUTI
% =============================================================================

\section{Analisi Dettagliata Degli Attributi Di Instradamento}

\subsection{Attributo Origin}

L'attributo \textbf{Origin} è un attributo \emph{well-known mandatory} che specifica l'origine della rotta e le modalità formali con cui essa è stata iniettata all'interno del processo BGP.

% [IMMAGINE DEL FILE DI SLIDE CON NOME "CloudComputing_5"] Attributo BGP Origin: definizione, valori possibili (IGP, EGP, Incomplete) e ordine di priorità — Slide 64 / Pagina 6 del PDF

L'attributo assume uno tra tre valori canonici:
\begin{enumerate}
    \item \textbf{IGP} (codificato con \texttt{i} nella tabella BGP): indica che la rotta ha avuto origine all'interno del sistema autonomo di provenienza tramite un protocollo di instradamento interno. È tipicamente generata tramite il comando manuale \texttt{network}. Rappresenta un controllo selettivo ed esplicito (\emph{``desidero annunciare esattamente questa subnet''}).
    \item \textbf{EGP} (codificato con \texttt{e}): la rotta è stata appresa tramite l'\emph{Exterior Gateway Protocol} storico (il predecessore di BGP-4). Nella moderna rete Internet questo protocollo è deprecato; tale valore può presentarsi solo a fronte di manipolazioni esplicite mediante policy.
    \item \textbf{Incomplete} (codificato con \texttt{?}): il prefisso è stato appreso mediante meccanismi eterogenei o redistribuzione massiva da altri protocolli (ad esempio tramite comandi del tipo \texttt{import-route} o \texttt{redistribute}). Indica un'origine non rigorosamente strutturata dal punto di vista BGP (\emph{``annuncia indistintamente tutto ciò che trovi''}) ed è considerata la tipologia meno affidabile.
\end{enumerate}

Nell'algoritmo di selezione del percorso BGP, a parità di altri fattori decisionali più prioritari, la preferenza segue un ordine gerarchico decrescente basato sull'affidabilità:
\begin{equation}
    \text{IGP (\texttt{i})} \;>\; \text{EGP (\texttt{e})} \;>\; \text{Incomplete (\texttt{?})}
\end{equation}

% [IMMAGINE DEL FILE DI SLIDE CON NOME "CloudComputing_5"] Topologia applicativa dell'attributo Origin tra AS 100 e AS 200 con output del comando display bgp routing-table — Slide 65 / Pagina 7 del PDF

\paragraph{Caso di studio pratico (Slide 65--66).}
Si consideri lo scenario operativo in cui nel dominio AS~200 sia attivo il protocollo OSPF, che annuncia il prefisso \texttt{200.0.0.0/24}. Due router di frontiera, RTB e RTC, esportano tale subnet verso il router RTA (situato in AS~100) attraverso collegamenti EBGP:
\begin{itemize}[noitemsep]
    \item \textbf{RTB} converte il prefisso \texttt{200.0.0.0/24} in rotta BGP utilizzando il comando \texttt{network}. In virtù di ciò, assegna alla rotta l'attributo Origin \textbf{IGP} (\texttt{i}).
    \item \textbf{RTC} inietta la stessa rotta mediante il comando di redistribuzione \texttt{import-route}, contrassegnandola con l'attributo Origin \textbf{Incomplete} (\texttt{?}).
\end{itemize}
Esaminando la tabella di instradamento su RTA tramite il comando \texttt{display bgp routing-table}:
\begin{lstlisting}[language=bash]
<RTA>display bgp routing-table
   Network          NextHop     MED  LocPrf  PrefVal  Path/Ogn
*> 200.0.0.0/24     10.1.12.2   1             0        200i
*                   10.1.13.3   1             0        200?
\end{lstlisting}
Il router RTA rileva identici valori di preferenza locale, assenza di lunghezze disuguali di \texttt{AS\_Path} e identico \texttt{MED}. In conformità alla gerarchia, RTA seleziona come rotta attiva e ottima (\texttt{* >}) quella appresa da RTB (\texttt{NextHop 10.1.12.2}) poiché il flag di origine \texttt{i} possiede priorità strettamente superiore rispetto a \texttt{?}.

\subsection{Attributo AS\_Path}

L'attributo \textbf{AS\_Path} è un attributo \emph{well-known mandatory} che tiene traccia, in ordine sequenziale vettoriale (\emph{vector order}), dei numeri di AS attraversati da un annuncio a partire dal sistema autonomo originario fino al router corrente.

% [IMMAGINE DEL FILE DI SLIDE CON NOME "CloudComputing_5"] Topologia multi-AS con anello topologico e rotta 100.0.0.0/24 per illustrare AS_Path e prevenzione loop — Slide 66 / Pagina 9 del PDF

L'attributo adempie a due compiti architetturali indispensabili:
\begin{enumerate}
    \item \textbf{Prevenzione dei loop inter-dominio}: quando un router riceve un annuncio da un peer EBGP, esamina l'elenco degli AS memorizzati nell'\texttt{AS\_Path}. Se il proprio numero di AS compare all'interno della lista, il router riconosce l'innesco di un anello di instradamento e \textbf{scarta immediatamente} il pacchetto di \texttt{UPDATE}. Nello scenario in Slide~66, se la rotta \texttt{100.0.0.0/24} originata da RTE (AS~4) compisse il circuito $\text{RTE} \to \text{RTB} \to \text{RTC} \to \text{RTD} \to \text{RTE}$, RTE scarterebbe il pacchetto ricevuto da RTD rilevando la presenza del proprio AS~4 nell'elenco.
    \item \textbf{Metrica di selezione del miglior cammino}: un percorso avente un numero inferiore di sistemi autonomi attraversati (lunghezza dell'\texttt{AS\_Path} minore) viene preferito rispetto a percorsi più lunghi.
\end{enumerate}

% [IMMAGINE DEL FILE DI SLIDE CON NOME "CloudComputing_5"] Rappresentazione vettoriale della composizione dell'AS_Path tra sessioni IBGP ed EBGP — Slide 67 / Pagina 10 del PDF

\paragraph{Regole formali di manipolazione dell'AS\_Path:}
\begin{itemize}
    \item \textbf{Origine locale}: se un router annuncia un prefisso originato localmente verso un peer \textbf{EBGP}, vi appone il proprio AS; se lo propaga all'interno del proprio dominio tramite sessione \textbf{IBGP}, genera un elenco \texttt{AS\_Path} vuoto (\emph{empty list}).
    \item \textbf{Transito EBGP}: quando un BGP speaker rilancia una rotta a un vicino EBGP appartenente a un altro dominio, aggiunge il proprio AS locale \textbf{in testa} alla lista (\emph{prepending}), ponendolo come primo elemento del vettore.
    \item \textbf{Transito IBGP}: quando una rotta viene riflessa o inoltrata tra peer IBGP interni allo stesso AS, l'\texttt{AS\_Path} \textbf{non viene modificato in alcun modo}.
\end{itemize}

\paragraph{Sottotipi di AS\_Path.}
Esistono quattro tipologie di segmenti:
\begin{itemize}[noitemsep]
    \item \textbf{AS\_Sequence}: lista ordinata e rigorosa di numeri di AS attraversati (rappresentazione standard).
    \item \textbf{AS\_Set}: insieme disordinato di AS generato nei processi di aggregazione di rotte per preservare la cronologia degli AS originali senza specificarne la sequenza sequenziale, evitando loop.
    \item \textbf{AS\_Confed\_Sequence} e \textbf{AS\_Confed\_Set}: utilizzati nell'ambito delle Confederazioni BGP (RFC~5065) per segmentare grandi domini interni.
\end{itemize}

\subsection{Attributo Next\_Hop}

L'attributo \textbf{Next\_Hop} è un attributo \emph{well-known mandatory} che descrive l'indirizzo IP del salto successivo necessario per instradare i pacchetti verso i prefissi descritti nell'NLRI. A differenza degli IGP convenzionali (in cui l'indirizzo di next hop coincide sempre con l'interfaccia fisica del router adiacente direttamente connesso), in BGP il \texttt{Next\_Hop} possiede una semantica peculiare dettata dalle tipologie di sessione.

% [IMMAGINE DEL FILE DI SLIDE CON NOME "CloudComputing_5"] Scenari di gestione dell'attributo Next_Hop (originate su IBGP, annuncio EBGP, transito da EBGP a IBGP) — Slide 68 / Pagina 12 del PDF

Le tre regole standard per la determinazione dell'indirizzo sono:
\begin{enumerate}
    \item \textbf{Rotte originate localmente e trasmesse a peer IBGP}: il \texttt{Next\_Hop} è impostato all'indirizzo dell'interfaccia locale attraverso cui è stata aperta la sessione BGP (tipicamente l'indirizzo di \emph{Loopback}).
    \item \textbf{Annuncio verso un peer EBGP}: il router sostituisce il \texttt{Next\_Hop} con l'indirizzo IP della propria interfaccia tramite cui è stabilita l'adiacenza EBGP verso il vicino.
    \item \textbf{Propagazione da EBGP a IBGP}: quando un router di frontiera riceve una rotta da un peer esterno EBGP e la propaga internamente ai propri peer IBGP, di norma \textbf{non modifica} il \texttt{Next\_Hop}. Questo valore continua a riflettere l'indirizzo del peer EBGP esterno. È compito dell'IGP interno garantire la risoluzione di raggiungibilità di tale indirizzo (oppure si ricorre al comando operativo \texttt{peer next-hop-local}).
\end{enumerate}

\subsection{Attributo Local\_Pref (Local Preference)}

L'attributo \textbf{Local\_Pref} è un attributo \emph{well-known discretionary} impiegato per stabilire la preferenza interna di uscita da un sistema autonomo per il raggiungimento di una determinata destinazione.

% [IMMAGINE DEL FILE DI SLIDE CON NOME "CloudComputing_5"] Uso di Local_Pref per indirizzare il traffico in uscita da AS 100 attraverso R2 anziché R3 — Slide 69 / Pagina 13 del PDF

\paragraph{Caratteristiche distintive:}
\begin{itemize}[noitemsep]
    \item \textbf{Portata di diffusione}: è scambiato \textbf{esclusivamente} tra router appartenenti al medesimo AS tramite sessioni IBGP. Non viene \textbf{mai} propagato attraverso collegamenti EBGP a sistemi autonomi limitrofi.
    \item \textbf{Metrica di decisione}: assume un valore intero senza segno a 32 bit (valore di default tipico: $100$). A un valore numericamente \textbf{più elevato} corrisponde una priorità di instradamento \textbf{maggiore}.
\end{itemize}

Nello scenario descritto alla Slide~69, AS~100 è interconnesso verso la rete esterna mediante due router di frontiera, R2 (collegato ad AS~200) e R3 (collegato ad AS~300), che apprendono entrambi la subnet \texttt{10.0.0.0/24}. Se l'amministratore di AS~100 desidera che l'intero traffico uscente diretto verso tale rete transiti sistematicamente da R2, è sufficiente configurare su R2 un valore di \texttt{Local\_Pref} più alto (es. $300$) rispetto a quello configurato su R3 (es. $200$). Tutti i nodi interni di AS~100 (incluso R1) eleggeranno la rotta via R2 come preferita.

\subsection{Attributo Multi-Exit Discriminator (MED)}

L'attributo \textbf{MED} (noto anche come metrica esterna) è un attributo \textbf{optional non-transitive}. Esso consente a un sistema autonomo di suggerire o influenzare il punto di ingresso privilegiato per il traffico proveniente da un AS adiacente, nel caso in cui sussistano interconnessioni multiple tra i due domini.

% [IMMAGINE DEL FILE DI SLIDE CON NOME "CloudComputing_5"] Applicazione dell'attributo MED da AS 100 per guidare i punti di ingresso da AS 200 — Slide 70 / Pagina 14 del PDF

\paragraph{Regole di comportamento:}
\begin{itemize}[noitemsep]
    \item \textbf{Natura non transitiva limitata}: il \texttt{MED} viene scambiato esclusivamente tra due sistemi autonomi confinanti diretti. L'AS ricevente utilizza il valore per prendere decisioni locali, ma \textbf{non deve trasmetterlo} ad alcun AS terzo.
    \item \textbf{Metrica decrescente}: opera analogamente a un costo metrica interno (IGP metric); di conseguenza, a parità di altri attributi, un router seleziona la rotta con il valore di \texttt{MED} \textbf{più basso}. Il valore di default in assenza di configurazione è pari a $0$.
\end{itemize}

Nello schema di Slide~70, i router R1 ed R2 in AS~100 annunciano la rete \texttt{10.0.0.0/24} verso i router R3 ed R4 in AS~200. Configurando un \texttt{MED} pari a $0$ su R1 e pari a $300$ su R2, i nodi dell'AS~200 preferiranno sistematicamente l'interconnessione via R1 come via di accesso verso la subnet remota.

\subsection{Attributo Community}

L'attributo \textbf{Community} è un attributo \textbf{optional transitive} (RFC~1997) concepito per semplificare drasticamente la gestione e l'applicazione di criteri di routing su insiemi aggregati di prefissi, svincolando la logica di controllo dal match puntuale dei singoli indirizzi IP.

% [IMMAGINE DEL FILE DI SLIDE CON NOME "CloudComputing_5"] Topologia con Service Provider AS 12 e clienti AS 10 e AS 11 che impiegano BGP Community — Slide 71 / Pagina 15 del PDF

Una \textbf{Community} è un'etichetta numerica a 32 bit (generalmente rappresentata nel formato \texttt{ASN:valore}, dove la prima metà identifica il numero di AS che ha definito la semantica e la seconda il codice di policy). Essa consente di raggruppare prefissi aventi proprietà comuni al fine di applicare policy unificate (quali preferenza locale, scarto di rotte o percorsi di aggregazione).

\paragraph{Comunità Standard Well-Known:}
\begin{itemize}[noitemsep]
    \item \texttt{Internet}: valore predefinito; la rotta può essere pubblicizzata a qualsiasi peer.
    \item \texttt{No-Export} (\texttt{0xFFFFFF01}): la rotta non deve essere annunciata al di fuori dell'AS locale (o confederazione).
    \item \texttt{No-Advertise} (\texttt{0xFFFFFF02}): la rotta non deve essere annunciata ad alcun altro peer BGP (né IBGP né EBGP).
    \item \texttt{No-Export-Subconfed}: impedisce la propagazione anche verso peer esterni all'interno di una sotto-confederazione BGP.
\end{itemize}

\paragraph{Utilità Ingegneristica.}
Senza comunità, la modifica delle preferenze di centinaia di subnet richiederebbe complesse liste di controllo accessi (ACL) o \emph{prefix-list} replicate su tutti i router di un provider. Assegnando invece all'origine un'etichetta di comunità (ad esempio \texttt{65000:100} per il traffico voce VoIP), i nodi di transito possono implementare una singola regola:
\begin{center}
\emph{``Se la rotta presenta la comunità 65000:100, imposta Local\_Pref = 200.''}
\end{center}

Nello scenario di Slide~71--72, il provider AS~12 gestisce i clienti AS~10 (subnet \texttt{10.1.10.0/24}) e AS~11 (subnet \texttt{10.1.11.0/24}). Per aggregare le subnet e comunicare unicamente il blocco cumulativo \texttt{10.1.10.0/23} verso AS~13 imponendo che quest'ultimo non lo rilanci a livello globale:
\begin{enumerate}[noitemsep]
    \item AS~10 marca la propria rotta con la community \texttt{10:12}; AS~11 la marca con \texttt{11:12}.
    \item Il provider RTC (AS~12) riconosce tali etichette mediante un \texttt{community-filter}, unifica le rotte nel blocco aggregato supernet \texttt{10.1.10.0/23} e vi associa la comunità well-known \texttt{no-export} in uscita verso AS~13.
    \item AS~13 riceve la rotta aggregata, ma la direttiva \texttt{no-export} gli vieta di riannunciarla oltre i propri confini.

\end{enumerate}

\newpage

% =============================================================================
% SEZIONE: MULTIPROTOCOL BGP (MP-BGP / BGP4+)
% =============================================================================

\section{Multiprotocol BGP (MP-BGP / BGP4+) E Supporto Multi-Famiglia}

Il protocollo BGP-4 tradizionale era intrinsecamente vincolato all'infrastruttura di rete IPv4: le informazioni NLRI e il campo \texttt{NEXT\_HOP} all'interno del messaggio \texttt{UPDATE} assumevano in modo fisso indirizzi a 32 bit. Con l'evoluzione verso IPv6 e l'emergere di tecnologie di virtualizzazione di rete (MPLS L3VPN, EVPN), l'IETF ha formalizzato le estensioni \textbf{Multiprotocol BGP} (MP-BGP, RFC~4760), designate anche come \textbf{BGP4+}.

MP-BGP converte il protocollo in un piano di controllo agnostico rispetto alla famiglia di indirizzi trasportata. La flessibilità è ottenuta disaccoppiando l'identificativo del protocollo di rete mediante due costrutti gerarchici:
\begin{itemize}[noitemsep]
    \item \textbf{AFI} (\emph{Address Family Identifier}, 2 ottetti): definisce il protocollo di rete a cui le informazioni appartengono (es. $\text{AFI}=1$ per IPv4, $\text{AFI}=2$ per IPv6, $\text{AFI}=25$ per L2VPN).
    \item \textbf{SAFI} (\emph{Subsequent Address Family Identifier}, 1 ottetto): specifica la semantica e il tipo di informazione veicolata per quella specifica famiglia (es. $\text{SAFI}=1$ per Unicast standard, $\text{SAFI}=128$ per etichette MPLS L3VPN, $\text{SAFI}=70$ per Ethernet VPN - EVPN).
\end{itemize}

% [IMMAGINE DEL FILE DI SLIDE CON NOME "CloudComputing_5"] Struttura dei campi di MP_REACH_NLRI e mapping AFI/SAFI — Slide 73 / Pagina 18 del PDF

\subsection{I Nuovi Attributi Estesi Di MP-BGP}

Per garantire la piena retrocompatibilità senza compromettere la struttura base dei pacchetti BGP-4, MP-BGP ha introdotto due nuovi attributi di percorso, trasportati \textbf{esclusivamente} all'interno del corpo dei messaggi \texttt{UPDATE}:

\subsubsection{Attributo MP\_REACH\_NLRI (Type Code = 14)}
Utilizzato per annunciare nuove rotte raggiungibili multiprotocollo e per notificare l'indirizzo del \emph{Next-Hop} corrispondente. La sua composizione binaria include:
\begin{itemize}[noitemsep]
    \item \textbf{AFI} (2 byte) e \textbf{SAFI} (1 byte): identificano la specifica famiglia e sotto-famiglia di indirizzi.
    \item \textbf{Length of Next Hop Network Address} (1 byte): specifica l'estensione in byte dell'indirizzo del salto successivo (ad esempio $16$ byte nel caso standard di un indirizzo IPv6 globale unicast, o $32$ byte qualora sia presente anche il link-local).
    \item \textbf{Network Address of Next Hop} (lunghezza variabile): indirizzo IP effettivo del salto successivo.
    \item \textbf{Reserved} (1 byte): campo riservato, tipicamente impostato a zero.
    \item \textbf{Network Layer Reachability Information (NLRI)} (lunghezza variabile): sequenza di prefissi di rete preceduti dalla rispettiva lunghezza di maschera (es. \texttt{2000:0:0:1::/64}). Se la lunghezza è $0$, il messaggio descrive una rotta di default.
\end{itemize}

\subsubsection{Attributo MP\_UNREACH\_NLRI (Type Code = 15)}
Utilizzato per revocare (\emph{withdraw}) percorsi multiprotocollo non più raggiungibili. È composto dai campi identificativi \textbf{AFI} (2 byte), \textbf{SAFI} (1 byte) e dalla sezione \textbf{Withdrawn Routes}, contenente l'elenco dei prefissi da cancellare nella forma $\langle\text{lunghezza maschera}, \text{prefisso}\rangle$.

% [IMMAGINE DEL FILE DI SLIDE CON NOME "CloudComputing_5"] Struttura del messaggio MP_UNREACH_NLRI per la revoca dei prefissi — Slide 74 / Pagina 20 del PDF

\subsection{Negoziazione Delle Capacità (Capabilities Advertisement)}

Affinché due router possano scambiarsi prefissi appartenenti a famiglie di indirizzi eterogenee, essi devono verificare preliminarmente il mutuo supporto alle estensioni attraverso il meccanismo di negoziazione delle capacità (\emph{Capabilities Advertisement}, RFC~5492) eseguito durante la fase di \emph{handshake} nel messaggio di \texttt{OPEN}.

% [IMMAGINE DEL FILE DI SLIDE CON NOME "CloudComputing_5"] Sessione BGP4+ tra R1 ed R2 su link IPv6 con scambio messaggi OPEN — Slide 75 / Pagina 21 del PDF

% [IMMAGINE DEL FILE DI SLIDE CON NOME "CloudComputing_5"] Pacchetto Wireshark di un messaggio BGP OPEN che negozia l'estensione Multiprotocol IPv6 unicast — Slide 76 / Pagina 22 del PDF

All'interno del campo dei parametri opzionali del messaggio \texttt{OPEN}, i nodi presentano un parametro di tipo $2$ (\emph{Capabilities}), specificando:
\begin{itemize}[noitemsep]
    \item \textbf{Multiprotocol Extensions Capability} (Codice di capacità 1): indica la tupla $\langle\text{AFI}, \text{SAFI}\rangle$ supportata (es. $\text{AFI}=2$ IPv6, $\text{SAFI}=1$ Unicast).
    \item \textbf{Support for 4-octet AS number capability} (RFC~6793): permette l'impiego di ASN a 32 bit.
    \item \textbf{Route-refresh capability} (RFC~2918): consente la ri-sincronizzazione dinamica delle tabelle senza interruzione della sessione TCP.
    \item \textbf{Multi-layer label capability}: necessaria nei contesti di etichettatura MPLS stratificata.
\end{itemize}

% [IMMAGINE DEL FILE DI SLIDE CON NOME "CloudComputing_5"] Cattura Wireshark del messaggio UPDATE con attributo MP_REACH_NLRI per il prefisso 2000:0:0:1::/64 — Slide 77 / Pagina 23 del PDF

% =============================================================================
% SEZIONE: ARCHITETTURE VPN: L2VPN, L3VPN ED EVPN
% =============================================================================

\section{Tecnologie VPN Su Backbone MPLS: Dal Livello 2 Al Livello 3}

Nelle reti di telecomunicazione di livello geografico e nei moderni data center, il problema cardine consiste nel connettere siti cliente dislocati geograficamente su un'infrastruttura condivisa garantendo isolamento di traffico, supporto a spazi di indirizzamento sovrapposti (\emph{overlapping subnets}) e scalabilità.

\subsection{Infrastrutture Di Livello 2: Pseudowire (PW) E VPLS}

La virtualizzazione di rete a livello di collegamento (L2VPN) mira a trasportare frame Ethernet attraverso una rete core MPLS (\emph{Multi-Protocol Label Switching}).

% [IMMAGINE DEL FILE DI SLIDE CON NOME "CloudComputing_5"] Differenze architetturali tra Pseudowire (PW) e VPLS sopra i tunnel LSP — Slide 78 / Pagina 24 del PDF

I due paradigmi di base si differenziano per la topologia del servizio erogato:
\begin{itemize}
    \item \textbf{Pseudowire (PW)}: realizza un circuito virtuale punto-a-punto (\emph{point-to-point}) trasparente che interconnette esattamente due router di bordo (\textbf{PE}, \emph{Provider Edge}). Dal punto di vista del cliente (\textbf{CE}, \emph{Customer Edge}), il circuito appare come un cavo Ethernet virtuale dedicato.
    \item \textbf{VPLS} (\emph{Virtual Private LAN Service}): realizza un servizio multipunto-a-multipunto (\emph{multipoint-to-multipoint}). L'infrastruttura del provider emula il comportamento di un gigantesco \textbf{switch Ethernet virtuale distribuito}. Ciascun sito cliente può comunicare direttamente con qualsiasi altro sito appartenente al medesimo dominio di broadcast.
\end{itemize}

\subsubsection{Architettura Di Incapsulamento: LSP E Pseudowire (Slide 78--80)}
Il trasporto di trame Layer 2 sopra una rete MPLS richiede una chiara \textbf{separazione dei ruoli} strutturata su due livelli di etichettatura sovrapposti (\emph{label stack}):

% [IMMAGINE DEL FILE DI SLIDE CON NOME "CloudComputing_5"] Creazione e segnalazione di un Pseudowire con protocollo LDP sopra un tunnel LSP — Slide 80 / Pagina 26 del PDF

\begin{enumerate}
    \item \textbf{Piano di Trasporto (\emph{Transport Plane}) - LSP}: il trasporto fisico-logico tra router PE è delegato a cammini commutati ad etichette (\textbf{LSP}, \emph{Label Switched Path}), istituiti tramite protocolli di segnalazione come LDP (\emph{Label Distribution Protocol}) o RSVP-TE. Un LSP è per sua natura un tunnel unidirezionale; per implementare un servizio bidirezionale sono richiesti due LSP simmetrici indipendenti oppure un LSP bidirezionale tramite RSVP-TE.
    \item \textbf{Piano di Servizio (\emph{Service Plane}) - Pseudowire (PW)}: sopra gli LSP di trasporto viene istanziato il circuito virtuale PW associato allo specifico cliente. Il legame gerarchico è noto con il nome di \textbf{Layered Over}. Due router PE concordano la creazione del PW scambiandosi un messaggio LDP \emph{Label Mapping} contenente:
    \begin{itemize}[noitemsep]
        \item Un identificatore numerico di circuito virtuale: \textbf{VC ID} (\emph{Virtual Circuit Identifier}).
        \item Una \textbf{MPLS Service Label} (etichetta interna di servizio) assegnata dal PE ricevente.
    \end{itemize}
\end{enumerate}

All'atto dell'inoltro, il router PE sorgente associa l'etichetta di servizio del PW (interna) con l'etichetta di trasporto dell'LSP che punta al PE remoto (esterna). Nel caso di VPLS, per abilitare la connettività any-to-any tra $N$ nodi PE appartenenti allo stesso dominio, è mandatorio istanziare una magliatura completa (\emph{full-mesh}) di Pseudowire.

\begin{theorem}
\textbf{Complessità della Topologia Full-Mesh}:
Dato un dominio VPLS comprendente $N$ router PE, il numero totale di Pseudowire bidirezionali che devono essere configurati e mantenuti attivi nella rete ammonta a:
\begin{equation}
    P_{\text{PW}} = \frac{N(N - 1)}{2}
\end{equation}
Per un'infrastruttura costituita da $N = 10$ router PE, la magliatura esige l'istituzione e il monitoraggio di esattamente $\frac{10 \times 9}{2} = 45$ Pseudowire, a prescindere dal numero effettivo di LSP fisici sottostanti.
\end{theorem}

% =============================================================================
% SEZIONE: TRADITIONAL L3VPN (BGP/MPLS IP VPN - RFC 4364)
% =============================================================================

\section{Traditional L3VPN (RFC 4364): Control Plane E Forwarding Data Plane}

Mentre i sistemi L2VPN emulano collegamenti a livello di collegamento dati, la tecnologia \textbf{Traditional L3VPN} (formalizzata nello standard RFC~4364 e storicamente nota come BGP/MPLS IP VPN) opera al livello 3 del modello di riferimento. Essa consente al fornitore di connettività di agire come router di transito per il traffico IP dei clienti, isolando le istanze di routing mediante tabelle di inoltro virtuali.

% [IMMAGINE DEL FILE DI SLIDE CON NOME "CloudComputing_5"] Architettura di un'MPLS L3VPN tra CE Site A e CE Site B attraverso rete di router P e PE — Slide 82 / Pagina 28 del PDF

L'architettura poggia su quattro pilastri concettuali:
\begin{itemize}
    \item \textbf{VRF} (\emph{Virtual Routing and Forwarding}): ciascun router PE ospita tabelle di instradamento e inoltro virtuali e completamente indipendenti per ciascun cliente (\emph{tenant}). In questo modo, clienti diversi possono utilizzare blocchi di indirizzi privati sovrapposti (ad es. RFC~1918, come la rete \texttt{10.0.0.0/8}) senza che si verifichi alcuna collisione nel routing.
    \item \textbf{Route Distinguisher (RD)}: è un prefisso a 64 bit aggiunto a monte di un indirizzo IPv4 a 32 bit per renderlo globalmente univoco all'interno del piano di controllo del provider. L'unione forma un indirizzo \textbf{VPNv4} di 96 bit (12 byte):
    \begin{equation}
        \text{VPNv4 Address} = \underbrace{\text{RD}}_{\text{64 bit}} \;+\; \underbrace{\text{IPv4 Prefix}}_{\text{32 bit}}
    \end{equation}
    Grazie a tale conversione, se due clienti utilizzano entrambi il prefisso \texttt{10.1.1.0/24}, essi vengono mappati su due rotte VPNv4 distinte (ad es. \texttt{1:100:10.1.1.0/24} e \texttt{2:100:10.1.1.0/24}).
    \item \textbf{Route Target (RT)}: è un attributo BGP esteso (formattato come \emph{BGP Extended Community}) che disciplina le politiche di importazione ed esportazione delle rotte all'interno delle rispettive VRF. Quando un router PE esporta una rotta VPNv4 in MP-BGP, vi applica un \textbf{Export RT}; all'estremità opposta, il PE ricevente installerà la rotta all'interno delle sole VRF che hanno un \textbf{Import RT} corrispondente.
    \item \textbf{Doppia Etichettatura MPLS (Dual MPLS Label)}:
    \begin{itemize}[noitemsep]
        \item \textbf{Internal Label (Service Label)}: etichetta assegnata dal PE remoto tramite MP-BGP; identifica univocamente la specifica istanza VRF di destinazione sul PE ricevente.
        \item \textbf{External Label (Transport Label)}: etichetta assegnata dai protocolli LDP o RSVP-TE per instradare il pacchetto lungo l'LSP fino a raggiungere il PE opposto.
    \end{itemize}
\end{itemize}

\subsection{Dinamica Del Piano Di Controllo E Del Piano Dei Dati}

Il funzionamento integrato di una rete L3VPN (illustrato in Slide~81--82 e descritto a pagina~29) si scompone rigorosamente in due fasi:

\paragraph{1. Control Plane (MP-BGP VPNv4):}
\begin{enumerate}[noitemsep]
    \item Il router CE del Sito B annuncia la subnet locale \texttt{10.2.2.0/24} al router PE2 via routing convenzionale (statico, OSPF o EBGP).
    \item PE2 riceve il prefisso e lo trasforma in una rotta VPNv4 aggiungendo il proprio RD (es. \texttt{1:100}). Vi associa il Route Target stabilito (es. \texttt{100:1}), alloca localmente una \textbf{Service Label} univoca (es. etichetta $100$) e imposta se stesso come \texttt{Next\_Hop} (\texttt{Next\_Hop = PE2}).
    \item PE2 diffonde la rotta verso gli altri PE (incluso PE1) impiegando una sessione MP-BGP VPNv4 ($\text{AFI}=1, \text{SAFI}=128$).
    \item PE1 riceve l'annuncio, verifica che l'Export RT corrisponda all'Import RT della VRF configurata per il cliente e installa la rotta \texttt{10.2.2.0/24} nella VRF del Cliente A, memorizzando l'etichetta di servizio $100$ e l'indirizzo di uscita PE2.
\end{enumerate}

\paragraph{2. Data Plane (Inoltro del Traffico Utente):}
\begin{enumerate}
    \item Il nodo CE del Sito A genera un normale pacchetto IP diretto a \texttt{10.2.2.0/24} e lo invia a PE1. Il CE non ha alcuna nozione di MPLS: per lui si tratta di puro instradamento IPv4 nativo.
    \item PE1 riceve il pacchetto sulla sotto-interfaccia legata alla VRF del Cliente A e cerca la destinazione nella tabella di routing della VRF.
    \item PE1 incapsula il pacchetto applicando uno \textbf{stack di due etichette MPLS}:
    \begin{itemize}[noitemsep]
        \item Inserisce sul fondo dello stack l'etichetta interna (\textbf{Service Label 100}) per identificare la VRF su PE2.
        \item Inserisce in cima l'etichetta esterna (\textbf{Transport Label 200}) relativa all'LSP che conduce a PE2.
    \end{itemize}
    \item Il pacchetto viene instradato nella rete di core del provider, popolata da router intermedi (\textbf{P Router}).
    \item I router P operano esclusivamente la commutazione rapida dell'etichetta esterna (\emph{label swapping}), ignorando la VRF, i percorsi del cliente e l'etichetta interna:
    \begin{center}
        \emph{Il core è totalmente agnostico rispetto al traffico cliente (\emph{stateless rispetto al cliente}), garantendo massima scalabilità e throughput.}
    \end{center}
    \item Quando il pacchetto raggiunge PE2, l'etichetta di trasporto $200$ è stata rimossa (tramite \emph{Penultimate Hop Popping} o sul PE stesso). PE2 esamina la Service Label interna $100$, identifica la VRF del Cliente A, decapsula integralmente il payload IP nativo e lo recapita al router CE del Sito B.
\end{enumerate}

% =============================================================================
% SEZIONE: CRITICITÀ DI VPLS E INTRODUZIONE AD EVPN
% =============================================================================

\section{Limiti Architetturali Di VPLS E Motivazioni Dell'Adozione Di EVPN}

Sebbene ampiamente implementata, la tecnologia VPLS manifesta gravi carenze strutturali all'aumentare delle dimensioni dell'infrastruttura, in particolare nelle interconnessioni tra data center (\emph{Data Center Interconnect}, DCI) e nelle reti di campus virtualizzate.

% [IMMAGINE DEL FILE DI SLIDE CON NOME "CloudComputing_5"] Schematizzazione delle problematiche e dei colli di bottiglia di VPLS che motivano la nascita di EVPN — Slide 83 / Pagina 30 del PDF

La letteratura ingegneristica e le slide (pagine 30--32) evidenziano sei colli di bottiglia critici di VPLS:

\begin{enumerate}
    \item \textbf{Carenza nell'Active/Active Multihoming}:
    Se un sito cliente si collega a due router PE distinti per ridondanza (\emph{multihoming}), si instaura un anello di livello 2. In ambiente VPLS convenzionale, l'unico rimedio standard per prevenire loop consiste nell'adozione dello \emph{Spanning Tree Protocol} (STP). STP opera bloccando completamente uno dei due collegamenti, che rimane confinato a uno stato inattivo di puro standby. VPLS manca di un meccanismo nativo standardizzato per consentire l'inoltro \emph{active/active} con bilanciamento del carico (\emph{per-flow multipathing}), portando a uno spreco sistematico di ampiezza di banda.

    \item \textbf{Trattamento inefficiente del traffico Multicast}:
    In VPLS, il comportamento di default tratta i flussi multicast esattamente come il traffico di broadcast. Le trame vengono replicate e inviate in \emph{flooding} a tutti i nodi PE appartenenti al dominio, costringendoli a riversarle sui rispettivi CE anche qualora nessun host collegato abbia richiesto quel flusso video o dati, saturando la banda del backbone.

    \item \textbf{Complessità di configurazione e proliferazione di protocolli}:
    VPLS non possiede un piano di controllo unificato. Ciascun PE richiede la configurazione indipendente dell'istanza di commutazione virtuale (\emph{VSI}), del protocollo di segnalazione dei label (Targeted LDP), dei protocolli di trasporto (LDP/RSVP-TE) e dei meccanismi di loop prevention (STP). L'elevato carico protocollare moltiplica la probabilità di errori di configurazione umana.

    \item \textbf{Scalabilità limitata in scenari Multi-tenant DCI}:
    In un data center cloud che ospita centinaia o migliaia di tenant distinti, ciascuno dei quali richiede un dominio Layer 2 isolato esteso su più sedi, VPLS impone la creazione di magliature complete di centinaia di Pseudowire dedicati per singolo tenant. Il piano di controllo basato su LDP e l'apprendimento massivo di decine di migliaia di MAC address nel piano dati generano un sovraccarico insostenibile di memoria e segnalazione.

    \item \textbf{Convergenza lenta a fronte di guasti}:
    Mentre la rete di sottofondo (\emph{underlay}) IP/MPLS può beneficiare di algoritmi veloci di ripristino (\emph{Fast Reroute}), i guasti che colpiscono i servizi VPLS dipendono dai lenti timer di riconvergenza di LDP e dallo svuotamento/re-imparamento delle tabelle MAC, con tempi che superano abbondantemente la soglia critica dei $\SI{50}{\milli\second}$, provocando cadute di connessione in servizi real-time come VoIP o video streaming.

    \item \textbf{Inondazione del traffico BUM e Broadcast Storms}:
    VPLS è progettato per emulare fedelmente un dominio di trasmissione broadcast. Di conseguenza, il traffico \textbf{BUM} (\emph{Broadcast, Unknown unicast, Multicast}) viene propagato a tappeto sull'intera rete metropolitana/geografica. Un host malconfigurato o un'anomalia locale su una porta possono scatenare tempeste di broadcast che saturano i collegamenti inter-sede dell'intero dominio di servizio.
\end{enumerate}

Tali marcate limitazioni hanno imposto la transizione verso una tecnologia moderna, nota come \textbf{EVPN} (\emph{Ethernet VPN}), concepita per trasferire l'apprendimento dei MAC address dal piano dati al piano di controllo attraverso estensioni MP-BGP dedicate ($\text{AFI}=25, \text{SAFI}=70$).

% DD -----------------------------------------------------------

% =============================================================================
% CAPITOLO / SEZIONE: EVPN (ETHERNET VPN) E INTEGRAZIONE VXLAN
% =============================================================================

\section{Evoluzione Dell'Architettura EVPN: Dai Backbone MPLS Alle Rete Overlay VXLAN}

Nelle moderne infrastrutture di telecomunicazione e nei data center cloud, la necessità di interconnettere carichi di lavoro virtualizzati su scala geografica ha reso evidenti i limiti delle tecnologie tradizionali. La tecnologia \textbf{EVPN} (\emph{Ethernet Virtual Private Network}), standardizzata originariamente in RFC~7432 sotto la denominazione di \emph{BGP MPLS-Based Ethernet VPN}, nasce proprio per superare le rigidità strutturali dei protocolli L2VPN convenzionali (in primo luogo VPLS, \emph{Virtual Private LAN Service}). La motivazione progettuale primaria risiedeva nell'esigenza di supportare ampiezze di banda estremamente elevate e politiche complesse di \emph{Quality of Service} (QoS) all'interno di dorsali MPLS.

La vera rivoluzione architetturale introdotta da EVPN consiste nella netta astrazione e separazione tra il \textbf{Piano di Controllo} (\emph{Control Plane}) e il \textbf{Piano dei Dati} (\emph{Data Plane}):
\begin{itemize}[noitemsep]
    \item \textbf{Control Plane Unificato}: le informazioni di raggiungibilità topologica non vengono apprese tramite flooding di pacchetti nel piano dati, ma sono distribuite in modo deterministico attraverso il protocollo \textbf{MP-BGP} (\emph{Multiprotocol BGP}, RFC~4760), sfruttando una famiglia di indirizzi dedicata ($\text{AFI} = 25$ per L2VPN, $\text{SAFI} = 70$ per EVPN). Tramite MP-BGP vengono scambiate rotte contenenti indirizzi MAC, indirizzi IP e informazioni di segmentazione.
    \item \textbf{Data Plane Agnostico e Flessibile}: il piano di controllo BGP-EVPN è completamente disaccoppiato dall'incapsulamento sottostante. Se nelle implementazioni originarie il piano dati era vincolato alla commutazione di etichette MPLS (\emph{Label Switching}) o a tunnel IP/GRE, la sua flessibilità ha consentito una naturale convergenza verso le reti overlay incapsulate su protocolli standard di livello 4, prime fra tutte \textbf{VXLAN} (\emph{Virtual eXtensible Local Area Network}).
\end{itemize}

% [IMMAGINE DEL FILE DI SLIDE CON NOME "CloudComputing_5"] Evoluzione dell'architettura EVPN da MPLS/GRE a VXLAN con separazione del piano di controllo MP-BGP e piano dati — Slide 84 / Pagina 1 del PDF

\subsection{Integrazione Di EVPN Con VXLAN: Il Control Plane Mancante}

L'architettura originaria di VXLAN (formalizzata in RFC~7348) era concepita come un mero meccanismo di incapsulamento nel data plane (\emph{MAC-in-UDP}) privo di un piano di controllo standardizzato. Nelle implementazioni VXLAN native, la scoperta dei tunnel endpoint (\textbf{VTEP}, \emph{VXLAN Tunnel End Point}) e l'apprendimento delle informazioni sugli host collegati (indirizzi MAC, IP associati, mapping con i \textbf{VNI}, \emph{VXLAN Network Identifier}) avvenivano mediante il paradigma del \emph{flood-and-learn}. I frame di broadcast, sconosciuti unicast e multicast (traffico \textbf{BUM}) venivano replicati inondando la rete IP di sottofondo (\emph{underlay}), spesso gravando su configurazioni multicast IP complesse e poco scalabili.

L'adozione di EVPN quale piano di controllo per VXLAN risolve alla radice questa inefficienza. Attraverso lo scambio di rotte BGP EVPN tra i VTEP:
\begin{enumerate}[noitemsep]
    \item I VTEP eseguono una scoperta automatica reciproca (\emph{automatic discovery}).
    \item Le associazioni tra indirizzi MAC, indirizzi IP e identificativi VNI vengono notificate tramite messaggi di controllo BGP prima che il traffico utente attraversi la rete.
    \item Il fenomeno del traffico di broadcast non necessario all'interno dell'infrastruttura del data center viene drasticamente ridotto o completamente soppresso.
\end{enumerate}

Questo percorso evolutivo ha visto l'estensione del framework originario RFC~7432 attraverso una serie di specifiche e draft IETF di rilievo:
\begin{itemize}[noitemsep]
    \item Il documento \texttt{draft-ietf-bess-evpn-overlay} si è formalmente evoluto nello standard \textbf{RFC~8365} (\emph{A Network Virtualization Overlay Solution Using Ethernet VPN}), definendo l'interoperabilità di EVPN su overlay VXLAN, NVGRE e MPLS-over-GRE.
    \item I draft complementari \texttt{draft-ietf-bess-evpn-inter-subnet-forwarding} e \texttt{draft-ietf-bess-evpn-prefix-advertisement} (successivamente confluiti in RFC~9135 e RFC~9136) hanno formalizzato l'inoltro integrato inter-subnet e l'annuncio dei prefissi IP (Rotte di Tipo~5).
\end{itemize}

% =============================================================================
% SEZIONE: RISOLUZIONE DEI LIMITI STORICI DI VPLS
% =============================================================================

\section{Risoluzione Delle Criticità Di VPLS Tramite EVPN}

Nelle architetture L2VPN basate su VPLS, l'assenza di un piano di controllo intelligente per i MAC address e l'impiego esclusivo dell'apprendimento nel piano di inoltro causavano gravi anomalie operative, in particolare negli scenari ad alta disponibilità.

% [IMMAGINE DEL FILE DI SLIDE CON NOME "CloudComputing_5"] Criticità storiche di VPLS (loopback, repliche broadcast, MAC flapping) a confronto con le soluzioni fornite da EVPN RFC 7432 — Slide 85 / Pagina 2 del PDF

Quando un sito cliente (\textbf{CE}, \emph{Customer Edge}) veniva collegato a due o più router di bordo del provider (\textbf{PE}, \emph{Provider Edge}) tramite collegamenti indipendenti per garantire ridondanza (\emph{multihoming}), la rete incorreva in tre problematiche distruttive:
\begin{enumerate}
    \item \textbf{Traffic Loopback (Eco)}: un frame trasmesso dal CE verso uno dei router PE veniva riflesso dall'infrastruttura di core e re-iniettato verso lo stesso CE tramite il secondo PE connesso al medesimo segmento.
    \item \textbf{Ricezione di copie multiple di broadcast}: un pacchetto broadcast generato dal CE veniva propagato su entrambi i collegamenti attivi; ciascun PE lo duplicava verso la rete di transito, recapitandolo in duplice copia a tutti i siti remoti.
    \item \textbf{Continuous MAC Address Flapping}: l'arrivo alternato delle repliche da collegamenti fisici distinti costringeva i router di core e i PE remoti ad aggiornare continuamente la porta di associazione della tabella MAC, destabilizzando la commutazione locale e degradando le prestazioni hardware.
\end{enumerate}

Nelle reti VPLS, l'unica soluzione percorribile consisteva nell'eseguire lo \emph{Spanning Tree Protocol} (STP) per bloccare a livello logico uno dei collegamenti, vincolando l'infrastruttura a un modello \emph{Active/Standby} che lasciava inutilizzata la capacità trasmissiva del link secondario.

\subsection{I Meccanismi Risolutivi Di EVPN (RFC 7432)}

EVPN introduce un insieme coordinato di costrutti protocollari volti a eliminare alla radice tali problematiche:
\begin{itemize}
    \item \textbf{Ethernet Segment Identifier (ESI)}: identifica in modo univoco il segmento fisico multi-homed che interconnette il CE ai router PE. L'identità condivisa del collegamento permette ai nodi della rete di riconoscere che i link afferiscono alla medesima entità fisica.
    \item \textbf{Etichette ESI per Split Horizon}: nelle reti EVPN-MPLS, a ciascun segmento multi-homed viene associata un'etichetta ESI. Quando un PE inoltra traffico BUM nel core verso i peer collegati allo stesso segmento, appone tale etichetta; i PE gemelli identificano l'etichetta e scartano il pacchetto, prevenendo la re-iniezione verso il CE originario (\emph{Split Horizon}).
    \item \textbf{Elezione del Designated Forwarder (DF)}: per evitare la duplicazione dei pacchetti broadcast e multicast verso il cliente, i router PE collegati al medesimo ESI negoziano un unico nodo responsabile dell'inoltro del traffico BUM proveniente dal core.
    \item \textbf{Apprendimento MAC via BGP ed ECMP}: la distribuzione delle tabelle di inoltro tramite rotte BGP sostituisce l'apprendimento reattivo nel data plane. Poiché i MAC address sono trattati come prefissi di instradamento, i flussi unicast diretti al sito multi-homed possono essere distribuiti uniformemente su tutti i link attivi sfruttando algoritmi \textbf{ECMP} (\emph{Equal-Cost Multi-Path}), abilitando un reale funzionamento \textbf{All-Active Multihoming}.
\end{itemize}

\subsection{Vantaggi Operativi Nel Moderno Data Center}

L'adozione di un piano di controllo centralizzato su BGP conferisce a EVPN vantaggi decisivi rispetto al modello VPLS:
\begin{itemize}
    \item \textbf{Ottimizzazione del Traffico Multicast}: integrandosi con l'\emph{IGMP Snooping}, i router PE monitorano le richieste degli host e le propagano via BGP; in tal modo si costruiscono alberi di distribuzione ottimali, evitando l'inondazione indiscriminata del traffico video e multicast verso PE privi di ricevitori interessati.
    \item \textbf{Semplificazione e Consistenza Configurale}: scompare la proliferazione di protocolli disparati (LDP per la segnalazione dei label, RSVP-TE per i tunnel, STP per i loop). BGP gestisce l'intero spettro di segnalazione (MAC, IP e topologia).
    \item \textbf{Scalabilità per Data Center Interconnect (DCI)}: la collaudata capacità di BGP di scalare fino a milioni di rotte consente di gestire agevolmente decine di migliaia di tenant isolati e relative tabelle MAC.
    \item \textbf{Convergenza Ultrarapida (\emph{Fast Convergence})}: a fronte del guasto di un link o di un intero router PE, EVPN non richiede la cancellazione puntuale di ciascun MAC address memorizzato, ma permette il ritiro massivo in batch (\emph{Mass Withdraw}) tramite una singola rotta BGP, ripristinando il traffico su percorsi alternativi in tempi inferiori al secondo.
    \item \textbf{Broadcast Suppression e ARP Suppression}: poiché i VTEP conoscono preventivamente le corrispondenze IP-MAC annunciate via BGP, un PE può intercettare una richiesta ARP broadcast generata da un host locale e rispondere direttamente per conto dell'host remoto (\emph{ARP Proxy/Suppression}), azzerando il flooding delle richieste ARP sul backbone WAN o sulla fabric IP.
\end{itemize}

% =============================================================================
% SEZIONE: IL MECCANISMO DEL DESIGNATED FORWARDER (DF)
% =============================================================================

\section{L'Elezione Del Designated Forwarder (DF) E L'Active/Active Multihoming}

Per comprendere a fondo la dinamica del traffico BUM in un contesto multi-homed, si consideri lo scenario operativo illustrato in Slide~85 e dettagliato nelle pagine 3 e 4.

% [IMMAGINE DEL FILE DI SLIDE CON NOME "CloudComputing_5"] Topologia di All-Active Multihoming: differenziazione tra traffico da CE a Core e da Core a CE mediato da DF e Non-DF — Slide 85 / Pagina 2 del PDF

Supponiamo che un dispositivo cliente (CE1) sia fisicamente connesso a due nodi di bordo del provider, denominati PE1 e PE2, che condividono lo stesso \textbf{Ethernet Segment (ES)}. In assenza di vincoli sul data plane, qualora CE1 emetta una trama broadcast (ad esempio una richiesta ARP), tale frame verrebbe acquisito simultaneamente sia da PE1 che da PE2:
\begin{itemize}[noitemsep]
    \item Entrambi i PE, operando come bridge per l'estensione del dominio di livello 2, incapsulerebbero la trama inviandola attraverso il core verso i restanti nodi (es. PE3 e PE4).
    \item I router PE remoti riceverebbero due copie identiche dello stesso pacchetto broadcast, recapitandole ai CE destinatari e scatenando tempeste di rete (\emph{broadcast storm}).
\end{itemize}

Per risolvere l'anomalia senza inibire l'uso simultaneo dei collegamenti, EVPN introduce il meccanismo di elezione del \textbf{Designated Forwarder (DF)}. I router PE attestati sullo stesso Ethernet Segment eleggono automaticamente un singolo nodo DF per ciascun dominio di broadcast (o per VLAN).

\subsection{Ruoli Operativi: DF Vs Non-DF}

La suddivisione dei compiti tra i PE adiacenti è asimmetrica e dipende rigorosamente dalla direzione del vettore di traffico:
\begin{itemize}
    \item \textbf{Designated Forwarder (DF)}: è l'unico router PE abilitato a inoltrare il traffico BUM proveniente dal core MPLS/IP in direzione discendente verso il dispositivo CE.
    \item \textbf{Non-Designated Forwarders (Non-DF)}: sono tutti gli altri router PE afferenti al medesimo segmento Ethernet. Il loro compito consiste nel filtrare e scartare silenziosamente qualsiasi pacchetto BUM pervenuto dal core e indirizzato a quel segmento CE.
\end{itemize}

\begin{theorem}
\textbf{Distinzione Cardine Sulla Direzionalità Del Flusso}:
L'elezione del Designated Forwarder disciplina unicamente la traiettoria del traffico \textbf{dal Core verso il CE}. Esso non applica alcun blocco al traffico emesso \textbf{dal CE in direzione del Core}. Ne consegue che:
\begin{itemize}[noitemsep]
    \item \textbf{Traffico Upstream (CE $\to$ Core)}: sia PE1 che PE2 possono acquisire, commutare e inoltrare le trame inviate dal cliente, ripartendo i carichi trasmissivi su base per-flusso (\emph{Load Balancing} / ECMP).
    \item \textbf{Traffico Downstream (Core $\to$ CE)}: i pacchetti unicast noti vengono recapitati da entrambi i router in base alle rotte BGP, mentre il traffico BUM viene riversato sul CE esclusivamente dal nodo eletto come DF (es. PE1), mentre il non-DF (es. PE2) provvede allo scarto immediato.
\end{itemize}
\end{theorem}

% =============================================================================
% SEZIONE: VOCABOLARIO FONDAMENTALE E STRUTTURE IDENTIFICATIVE EVPN
% =============================================================================

\section{Vocabolario Fondamentale E Parametri Strutturali Di EVPN}

La corretta implementazione di una rete EVPN poggia su una tassonomia formale che mappa le entità logiche e fisiche all'interno degli attributi di instradamento BGP.

\subsection{Route Distinguisher (RD)}

Il \textbf{Route Distinguisher (RD)} è un prefisso a 64 bit (8 ottetti) impiegato per garantire l'univocità globale di prefissi IP o indirizzi MAC che risultino identici e sovrapposti all'interno di tenant differenti.

% [IMMAGINE DEL FILE DI SLIDE CON NOME "CloudComputing_5"] Esemplificazione della coesistenza di indirizzi IP sovrapposti grazie all'assegnazione di RD distinti — Slide 87 / Pagina 6 del PDF

L'RD è un concetto puramente locale al router che origina la rotta BGP; non prescrive verso quale tabella debba confluire il traffico, ma impedisce che l'algoritmo di BGP path selection consideri in conflitto annunci coincidenti generati da clienti distinti. Come documentato a pagina~7 e 8, l'RD di 8 ottetti è composto da un campo Type a 2 byte seguito da un payload di 6 byte, strutturato secondo tre formati standard:
\begin{enumerate}
    \item \textbf{Formato $\text{ASN}:\text{NN}$ con ASN a 2 byte}: i primi 2 byte rappresentano il numero di sistema autonomo ($\text{ASN} \in [0, 65535]$) e i restanti 4 byte un numero assegnato dall'amministratore ($\text{NN} \in [0, 4294967295]$), ad esempio \texttt{65001:100} o \texttt{1:3}. I due campi non possono essere simultaneamente nulli (la combinazione \texttt{0:0} è illegittima).
    \item \textbf{Formato $\text{ASN}:\text{NN}$ con ASN a 4 byte}: l'ASN a 4 byte occupa un intero senza segno a 32 bit ($\text{ASN} \in [65536, 4294967295]$), mentre il valore definito dall'utente è confinato a 2 byte ($\text{NN} \in [0, 65535]$), ad esempio \texttt{65537:3}. In alternativa, la medesima struttura a 4 byte può essere codificata in notazione decimale puntata $x.y:\text{NN}$ (con $x, y \in [0, 65535]$), ad esempio \texttt{0.1:0} o \texttt{0.0:3} (con il vincolo che \texttt{0.0:0} non è ammesso).
    \item \textbf{Formato $\text{IP}:\text{NN}$ con indirizzo IPv4}: i primi 4 byte ospitano un indirizzo IP a 32 bit (es. l'indirizzo di loopback del VTEP sorgente, compreso tra \texttt{0.0.0.0} e \texttt{255.255.255.255}), seguiti da un identificatore numerico a 2 byte ($\text{NN} \in [0, 65535]$), ad esempio \texttt{192.168.122.15:1} o \texttt{192.168.1.1:100}.
\end{enumerate}

\subsection{Route Target (RT) E Politiche Di Import/Export Sincrone}

Il \textbf{Route Target (RT)} è un attributo BGP esteso (\emph{BGP Extended Community} a 8 byte) che determina la portata di distribuzione delle rotte, governando formalmente \emph{quali} entità di rete siano autorizzate a ricevere e installare un annuncio.

% [IMMAGINE DEL FILE DI SLIDE CON NOME "CloudComputing_5"] Meccanismo simmetrico di Export RT e Import RT tra VRF collocate su nodi distinti — Slide 88 / Pagina 7 del PDF

Mentre l'RD rende il prefisso univoco nel control plane, l'RT agisce da filtro di appartenenza:
\begin{itemize}[noitemsep]
    \item \textbf{Export RT}: etichetta applicata alla rotta nel momento in cui il PE locale la pubblica verso il piano di controllo MP-BGP.
    \item \textbf{Import RT}: regola di filtraggio associata a una specifica istanza VRF o EVI. Il dispositivo importa e installa esclusivamente gli annunci il cui Export RT coincide con l'elenco dei propri Import RT configurati.
\end{itemize}

Nelle architetture standard tra due nodi (Nodo A e Nodo B), per garantire la bidirezionalità dello scambio si configura una coppia di RT secondo un modello simmetrico:
\begin{lstlisting}[language=bash]
Nodo A (VRF A):
  Export RT: 200:1   (esporta verso il piano BGP con tag 200:1)
  Import RT: 200:2   (importa nella VRF A i prefissi recanti tag 200:2)

Nodo B (VRF B):
  Export RT: 200:2   (esporta verso il piano BGP con tag 200:2)
  Import RT: 200:1   (importa nella VRF B i prefissi recanti tag 200:1)
\end{lstlisting}
In questo modo, il canale logico di routing viene stabilito in maniera rigorosa senza che i router intermedi debbano conoscere le singole sottoreti dei clienti.

\subsection{Ethernet Segment Identifier (ESI)}

L'\textbf{ESI} è un identificatore numerico globale univoco a \textbf{10 byte (80 bit)} che denota l'interfaccia fisica o il fascio aggregato di collegamenti (LAG) che connette un sito CE a uno o più router PE. Rappresenta l'identità del segmento físico per tutti i nodi della fabric. Se un host o un apparato cliente è monoconnesso a un singolo switch (\emph{single-homed}), il valore dell'ESI è convenzionalmente impostato a $0$ ($00:00:00:00:00:00:00:00:00:00$). In contesti multi-homing assume valori univoci esadecimali (ad esempio \texttt{aa:bb:cc:dd:ee:ff:00:00:01}).

\subsection{EVPN Instance (EVI)}

L'\textbf{EVI} rappresenta l'istanza globale della VPN per un determinato cliente o tenant estesa sull'intera infrastruttura del provider. È definita dall'insieme coordinato di Route Targets e raggruppa tutti i domini di broadcast (\emph{Bridge Domains}) associati a quel tenant. Quando un router rileva un nuovo indirizzo MAC su una porta assegnata a una specifica EVI, genera una rotta EVPN di Tipo~2, vi allega l'Export RT proprio dell'EVI e la propaga su MP-BGP. Tutti i PE che condividono la medesima EVI (ossia configurati con il corrispondente Import RT) acquisiranno la rotta.

\subsection{Bridge Domain (BD) Ed Ethernet Tag ID}

Il \textbf{Bridge Domain (BD)} è l'entità logica di livello 2 istanziata sul router PE che delimita un dominio di broadcast isolato. Funge da container logico che aggrega interfacce fisiche, sotto-interfacce VLAN o tunnel virtuali VXLAN consentendo alle entità collegate di comunicare come se fossero attestate sul medesimo switch fisico locale. Ciascun BD mantiene la propria tabella di commutazione MAC isolata (\textbf{MAC-VRF}).

Nel piano di controllo EVPN, ogni Bridge Domain è formalmente contraddistinto da un \textbf{Ethernet Tag ID}, un identificatore numerico a \textbf{24 bit}. Nelle architetture che adottano l'incapsulamento VXLAN sul data plane:
\begin{itemize}[noitemsep]
    \item Il Bridge Domain viene mappato in rapporto $1:1$ con un \textbf{Layer 2 VNI} (\emph{L2 VNI}).
    \item Se più VLAN/domini condividono la stessa istanza EVI e lo stesso VNI, l'Ethernet Tag ID viene valorizzato con il tag VLAN originario per discriminare i singoli domini all'interno dello stesso flusso logico.
\end{itemize}

\subsection{VRF E Interfaccia IRB (Integrated Routing And Bridging)}

A livello 3, la gestione del tenant è demandata all'istanza \textbf{VRF} (\emph{Virtual Routing and Forwarding}), che ospita la tabella di routing IP isolata (\textbf{IP-RIB}). Il collegamento architetturale tra il mondo a commutazione di frame (Livello 2 - Bridge Domain) e il mondo a instradamento di pacchetti (Livello 3 - VRF) è operato dall'interfaccia \textbf{IRB} (\emph{Integrated Routing and Bridging}, denominata anche BDIF):
\begin{center}
\emph{L'interfaccia IRB incarna il Default Gateway distribuito del tenant: è il punto di giunzione logico in cui una trama Ethernet viene decapsulata per trasformarsi in pacchetto IP instradabile inter-subnet, e viceversa.}
\end{center}

\newpage

% =============================================================================
% SEZIONE: GERARCHIA STRATIFICATA DEI LIVELLI IN EVPN
% =============================================================================

\section{Gerarchia Dei Livelli In EVPN E Confronto Con MPLS L3VPN}

La convergenza tra funzionalità Layer 2 e Layer 3 in un'unica piattaforma impone una rigorosa organizzazione ad albero dei parametri operativi. La gerarchia strutturale illustrata alle pagine 14--16 mappa le dipendenze dal livello fisico fino all'instradamento di prefisso:

% [IMMAGINE DEL FILE DI SLIDE CON NOME "CloudComputing_5"] Struttura ad albero gerarchica di EVPN: Level 1 Fisico (ESI), Level 2 Tenant (EVI, VNI, Ethernet Tag), Level 3 Tenant (VRF, IRB, Type-5) — Slide 95 / Pagina 16 del PDF

La gerarchia si compone di tre livelli logici fondamentali:
\begin{enumerate}
    \item \textbf{Level 1 -- Physical (Segmento Fisico)}:
    Definito dall'\textbf{ESI} (es. \texttt{aa:bb:cc:dd:ee:ff:00:00:01}). Ha validità locale-segmento e identifica \emph{chi è collegato a cosa} a livello hardware. Risulta indispensabile per l'elezione del Designated Forwarder e per le funzioni di Split Horizon.
    \item \textbf{Level 2 -- L2 Tenant (Servizio di Livello 2)}:
    Rappresentato dall'\textbf{EVI} (es. EVI 10) mappata sul piano dati con il \textbf{Layer 2 VNI} (es. VNI 10010). Al di sotto dell'EVI si attestano i diversi \textbf{Ethernet Tag ID} che identificano i broadcast domain del cliente (es. Tag 10 per VLAN 10, Tag 20 per VLAN 20). Ciascun tag L2 alimenta la rispettiva sotto-interfaccia logica di gateway \textbf{IRB 10} e \textbf{IRB 20}.
    \item \textbf{Level 3 -- L3 Tenant (Instradamento Inter-Subnet)}:
    Le interfacce IRB convergono nella specifica \textbf{VRF-A} del Tenant A. A questo livello risiede il \textbf{Layer 3 VNI} (es. L3 VNI 50010), impiegato per l'incapsulamento del traffico instradato tra subnet differenti dello stesso tenant. Dalla VRF vengono originate e importate le rotte di rete \textbf{Type-5 IP Prefix} (ad es. \texttt{192.168.10.0/24} e \texttt{192.168.20.0/24}).
\end{enumerate}

\subsection{Transizione Da Traditional MPLS L3VPN A VXLAN EVPN L3VPN}

Come illustrato nelle Slide~96--97 (pagine 17 e 18), la convergenza di EVPN determina una profonda mutazione dell'infrastruttura di trasporto e del piano di controllo rispetto alle tradizionali L3VPN (RFC~4364):

% [IMMAGINE DEL FILE DI SLIDE CON NOME "CloudComputing_5"] Schema architetturale di EVPN L3VPN con incapsulamento VXLAN over IP su topologia Spine-Leaf — Slide 96 / Pagina 17 del PDF

\begin{table}[H]
\centering
\small
\begin{tabularx}{\textwidth}{l p{6.5cm} X}
\toprule
\textbf{Caratteristica} & \textbf{Traditional MPLS L3VPN (RFC 4364)} & \textbf{VXLAN EVPN L3VPN} \\
\midrule
\textbf{Underlay Core} &
Richiede una dorsale etichettata MPLS con protocolli LDP/RSVP-TE su tutti i router P di transito. &
\textbf{Pure IP Fabric}: il core è una semplice rete IP routed (es. Leaf-Spine con eBGP o OSPF/IS-IS); non richiede MPLS nel core. \\
\addlinespace
\textbf{Incapsulamento} &
Doppio stack di etichette MPLS: Transport Label (esterna) + Service Label/VRF (interna). &
\textbf{VXLAN over IP/UDP}: Outer IP $\to$ UDP (porta 4789) $\to$ VXLAN Header contenente il VNI $\to$ Original Payload. \\
\addlinespace
\textbf{Piano di Controllo} &
MP-BGP con famiglia VPNv4/VPNv6 ($\text{AFI} = 1/2, \text{SAFI} = 128$). &
MP-BGP EVPN ($\text{AFI} = 25, \text{SAFI} = 70$), che unifica la gestione sia di livello 2 che di livello 3. \\
\addlinespace
\textbf{Ruolo dei Nodi Core} &
I router P eseguono la commutazione delle sole etichette di trasporto MPLS (\emph{label swapping}). &
Gli switch Spine operano come meri router IP di transito e, sul piano di controllo, come \textbf{BGP Route Reflector} (RR). \\
\addlinespace
\textbf{Flessibilità Servizio} &
Supporta unicamente l'instradamento di prefissi IP a livello 3 (VPN isolate per clienti WAN). &
\textbf{Multi-service nativo}: supporta sia bridging MAC (L2) che instradamento IP (L3) all'interno dello stesso control plane. \\
\bottomrule
\end{tabularx}
\caption{Confronto architetturale tra MPLS L3VPN classica e VXLAN EVPN.}
\end{table}

% =============================================================================
% SEZIONE: LE CINQUE ROTTE FONDAMENTALI DI EVPN NLRI
% =============================================================================

\section{Le Tipologie Di Rotte EVPN NLRI (RFC 7432 E RFC 9136)}

Nel formato binario definito da RFC~7432, le informazioni di raggiungibilità viaggiano all'interno di un'apposita struttura NLRI inserita nell'attributo esteso \texttt{MP\_REACH\_NLRI} o \texttt{MP\_UNREACH\_NLRI}. La testata comune presenta il seguente formato a ottetti:

% [IMMAGINE DEL FILE DI SLIDE CON NOME "CloudComputing_5"] Struttura binaria dell'EVPN NLRI con campo Route Type a 1 ottetto, Length e payload specifico — Slide 97 / Pagina 19 del PDF

\begin{lstlisting}[language=bash]
+---------------------------------------+
| Route Type (1 octet)                  |
+---------------------------------------+
| Length (1 octet)                      |
+---------------------------------------+
| Route Type Specific (lunghezza variabile)|
+---------------------------------------+
\end{lstlisting}

Il valore assunto dal campo \textbf{Route Type} nel primo byte determina l'interpretazione semantica e il formato dei dati contenuti nel payload:
\begin{itemize}[noitemsep]
    \item \textbf{Type 1}: \emph{Ethernet Auto-Discovery (A-D) Route} (RFC~7432).
    \item \textbf{Type 2}: \emph{MAC/IP Advertisement Route} (RFC~7432).
    \item \textbf{Type 3}: \emph{Inclusive Multicast Ethernet Tag Route} (RFC~7432).
    \item \textbf{Type 4}: \emph{Ethernet Segment Route} (RFC~7432).
    \item \textbf{Type 5}: \emph{IP Prefix Route} (RFC~9136).
\end{itemize}

% =============================================================================
% SEZIONE: ANALISI DETTAGLIATA DELLA ROTTA DI TIPO 1 (ETHERNET A-D)
% =============================================================================

\section{Rotta EVPN Type 1: Ethernet Auto-Discovery Route}

La rotta \textbf{Type 1 (Ethernet Auto-Discovery)} è un messaggio specializzato scambiato tra router PE necessario esclusivamente quando l'interconnessione verso il cliente impiega l'isolamento multi-homing identificato da un ESI valido. Il suo scopo principale consiste nella scoperta degli endpoint e nella segnalazione avanzata delle capacità del segmento.

% [IMMAGINE DEL FILE DI SLIDE CON NOME "CloudComputing_5"] Formato del messaggio NLRI Type 1 Ethernet Auto-Discovery Route — Slide 98 / Pagina 20 del PDF

La composizione dei campi del payload della rotta Type 1 presenta una lunghezza totale di $25$ byte:
\begin{lstlisting}[language=bash]
+--------------------------------------------------+
| Route Distinguisher (RD) (8 octets)              |
+--------------------------------------------------+
| Ethernet Segment Identifier (ESI) (10 octets)    |
+--------------------------------------------------+
| Ethernet Tag ID (4 octets)                       |
+--------------------------------------------------+
| MPLS Label / VNI (3 octets)                      |
+--------------------------------------------------+
\end{lstlisting}

Nel contesto operativo di **EVPN su VXLAN**, il campo \emph{MPLS Label} a 3 ottetti perde la sua valenza originaria di etichetta di commutazione: il valore viene impostato a zero oppure ignorato dal ricevente, in quanto la discriminazione del contesto di servizio avviene tramite il **L2 VNI** (per i Bridge Domain) o tramite il **L3 VNI** (per le interfacce IRB).

La rotta Type 1 assolve a tre funzioni cardine:

\subsection{1. Split Horizon Ed ESI Label Assignment (Slide 99--100)}

Per scongiurare il loopback del traffico broadcast multi-homed illustrato alle pagine 21 e 22, la rotta di Tipo 1 viene specializzata nella forma \textbf{Ethernet A-D per ES Route}.

% [IMMAGINE DEL FILE DI SLIDE CON NOME "CloudComputing_5"] Dinamica di Split Horizon tra PE1, PE2 e CE2 con segnalazione ESI per evitare il reinoltro del traffico BUM — Slide 99 / Pagina 21 del PDF

Si consideri un sito CE2 dual-homed attestato su PE1 e PE2 (segmento \texttt{ES2}):
\begin{enumerate}
    \item CE2 emette un frame broadcast; assumendo che PE1 sia il Designated Forwarder, PE2 scarta la copia pervenuta dal link locale.
    \item PE1 acquisisce il frame e, per garantire la connettività LAN estesa, ha l'obbligo di propagarlo a tutti i PE della rete, compreso PE2.
    \item Se PE2 trattasse il frame ricevuto dal core come un ordinario pacchetto BUM, lo ritrasmetterebbe a CE2 sul proprio link fisico, provocando una ricezione duplicata o innescando un anello su larga scala.
\end{enumerate}

La soluzione a questo scenario critico si articola a seconda della tecnologia del piano dati:
\begin{itemize}
    \item \textbf{In ambiente EVPN-MPLS (RFC 7432)}: PE2 alloca una specifica etichetta denominata \textbf{ESI Label} per identificare l'interfaccia collegata a \texttt{ES2}, e la annuncia a PE1 all'interno di una rotta Type 1. Quando PE1 invia la replica BUM a PE2 attraverso il core MPLS, vi applica in testa l'etichetta ESI. PE2 decapsula il pacchetto, rileva la presenza dell'etichetta ESI corrispondente al segmento \texttt{ES2} e blocca l'uscita della trama verso CE2.
    \item \textbf{In ambiente EVPN-VXLAN (RFC 8365)}: l'header standard VXLAN non possiede un campo per trasportare etichette ESI. Di conseguenza, il problema di Split Horizon viene risolto tramite il meccanismo noto come \textbf{Local Bias}:
    \begin{enumerate}[noitemsep]
        \item Tramite il piano di controllo BGP (grazie alle rotte di Tipo 4, \emph{Ethernet Segment Route}), PE1 e PE2 si scoprono reciprocamente connessi allo stesso segmento \texttt{ES2} e memorizzano i rispettivi indirizzi IP di loopback (\emph{VTEP IP}).
        \item Quando PE2 riceve un pacchetto incapsulato VXLAN proveniente da PE1, esamina l'indirizzo IP sorgente nell'intestazione IP esterna (\emph{Outer IP Source Address}).
        \item PE2 riscontra che la sorgente è PE1, suo peer associato al segmento \texttt{ES2}.
        \item Per il principio dello Split Horizon, PE2 conclude che il traffico è già stato consegnato (o lo sarà a breve) alle interfacce locali di \texttt{ES2} ad opera di PE1.
        \item PE2 procede dunque a replicare il frame verso tutti i propri Bridge Domain e porte locali, \textbf{escludendo tassativamente tutte le interfacce fisiche facenti parte di ES2}.
    \end{enumerate}
\end{itemize}

\subsection{2. Fast Convergence E Ritiro Massivo Delle Rotte (Mass Withdraw)}

Nelle reti tradizionali, se un collegamento tra un apparato CE e un router PE si interrompe, il router deve cancellare ed emettere messaggi di \texttt{WITHDRAW} per ciascuno dei singoli indirizzi MAC precedentemente appresi su quel link. Se il CE ospita centinaia di macchine virtuali o container, il flooding di migliaia di messaggi BGP satura il control plane e dilata i tempi di riconvergenza ben oltre le soglie tollerabili.

% [IMMAGINE DEL FILE DI SLIDE CON NOME "CloudComputing_5"] Meccanismo di Fast Convergence (Mass Withdraw) a seguito del guasto del link CE1-PE1 tramite revoca Type 1 — Slide 100 / Pagina 23 del PDF

Come schematizzato in Slide~100, EVPN implementa la convergenza accelerata (\textbf{Mass Withdraw}):
\begin{center}
\emph{A fronte del guasto del link fisico tra CE1 e PE1, il router PE1 non invia una revoca per ciascun MAC (Type 2). Invia un'unica rotta Type 1 Ethernet A-D per ES recante l'attributo MP\_UNREACH\_NLRI riferito all'intero segmento ES1.}
\end{center}
I PE remoti (es. PE3), alla ricezione di questa singola notifica, invalidano all'istante l'associazione di PE1 come next-hop per tutte le decine di migliaia di indirizzi MAC legati a quel segmento, commutando l'intero volume di traffico in direzione di PE2 in frazioni di secondo (\emph{sub-second fast switchover}).

\subsection{3. Aliasing Tramite Ethernet A-D Per EVI Routes}

Negli scenari multi-homing in modalità All-Active, può verificarsi una situazione asimmetrica per cui un dispositivo cliente CE1 trasmetta traffico unidirezionale verso la rete attraverso unicamente uno dei due PE (ad esempio inviando frame solo a PE1). In tal caso, PE1 impara l'indirizzo MAC sorgente e lo annuncia via BGP a PE3; al contrario, il router gemello PE2 non riceve traffico dal CE e la sua tabella MAC locale ignora la presenza dell'host.

% [IMMAGINE DEL FILE DI SLIDE CON NOME "CloudComputing_5"] Principio di funzionamento dell'Aliasing per instradare traffico in ECMP anche verso il PE che non ha ancora appreso il MAC — Slide 101 / Pagina 24 del PDF

In condizioni ordinarie, il router remoto PE3 invierebbe il traffico di risposta esclusivamente a PE1, vanificando la ridondanza e il bilanciamento del carico. Per risolvere questa asimmetria, EVPN introduce il meccanismo di \textbf{Aliasing} tramite rotte \textbf{Ethernet A-D per EVI}:
\begin{itemize}
    \item Sebbene PE2 non abbia ancora registrato il singolo indirizzo MAC di quell'host, esso sa di essere connesso al medesimo segmento fisico (\texttt{ES1}) e pubblica una rotta di Tipo 1 specifica per l'EVI del cliente (\emph{Ethernet A-D per EVI Route}) indicando il medesimo ESI.
    \item Il messaggio è associato all'RD e all'Export RT propri dell'EVI; pertanto PE3, importando la rotta nella medesima tabella, deduce che \emph{qualsiasi indirizzo MAC associato a quell'ESI è raggiungibile indifferentemente sia via PE1 che via PE2}.
    \item PE3 può così instaurare percorsi multipath bilanciati (\textbf{ECMP}) ripartendo il traffico tra entrambi i router PE.
\end{itemize}

\subsection{L'Attributo ESI Label Extended Community E Le Modalità Di Ridondanza}

La rotta Type 1 è accompagnata da un attributo transitorio obbligatorio denominato \textbf{ESI Label Extended Community} (Tipo \texttt{0x06}, Sotto-tipo \texttt{0x01}, estensione 8 byte), la cui formattazione specifica è mostrata a pagina~25:

% [IMMAGINE DEL FILE DI SLIDE CON NOME "CloudComputing_5"] Struttura dell'ESI Label Extended Community e dettaglio del flag di ridondanza (All-Active vs Single-Active) — Slide 102 / Pagina 25 del PDF

\begin{lstlisting}[language=bash]
 0                   1                   2                   3
 0 1 2 3 4 5 6 7 8 9 0 1 2 3 4 5 6 7 8 9 0 1 2 3 4 5 6 7 8 9 0 1
+-+-+-+-+-+-+-+-+-+-+-+-+-+-+-+-+-+-+-+-+-+-+-+-+-+-+-+-+-+-+-+-+
| Type = 0x06   |Sub-Type = 0x01| Flags (1 Byte)|  Reserved = 0 |
+-+-+-+-+-+-+-+-+-+-+-+-+-+-+-+-+-+-+-+-+-+-+-+-+-+-+-+-+-+-+-+-+
|  Reserved = 0 |               ESI Label (3 Bytes)             |
+-+-+-+-+-+-+-+-+-+-+-+-+-+-+-+-+-+-+-+-+-+-+-+-+-+-+-+-+-+-+-+-+
\end{lstlisting}

Il campo \textbf{Flags} (1 ottetto) riveste un ruolo discriminante mediante il suo bit meno significativo (\emph{Least Significant Bit}, bit 0):
\begin{itemize}
    \item \textbf{Valore 0 -- Modalità All-Active Redundancy}:
    Attesta che il segmento opera in modalità attiva/attiva simultanea. Tutti i router PE collegati possono inoltrare traffico unicast proveniente dal o diretto al CE. Il router remoto (es. PE3) riconosce il flag $0$ ed è autorizzato a eseguire il bilanciamento di carico ECMP su tutti i nodi attestati sul segmento.
    \item \textbf{Valore 1 -- Modalità Single-Active Redundancy}:
    Attesta che il segmento opera secondo uno schema logico primario/secondario (\emph{Active/Standby}). Un solo router (quello eletto come Designated Forwarder) è abilitato all'inoltro del traffico per quel segmento, mentre gli altri fungono da mere adiacenze di backup in attesa di subentro a seguito di guasti. Questa modalità è indispensabile nei protocolli (come PPPoE o implementazioni storage Fibre Channel over Ethernet) che non tollerano la ricezione disordinata di pacchetti appartenenti allo stesso flusso su collegamenti distinti. In questa modalità, il campo \emph{ESI Label} è valorizzato a zero.
\end{itemize}

% =============================================================================
% SEZIONE: ANALISI DETTAGLIATA DELLA ROTTA DI TIPO 2 (MAC/IP ADVERTISEMENT)
% =============================================================================

\section{Rotta EVPN Type 2: MAC/IP Advertisement Route}

La rotta \textbf{Type 2 (MAC/IP Advertisement Route)} costituisce l'elemento fondante e più versatile dell'architettura EVPN. Essa incarna il meccanismo di annuncio puntuale degli endpoint (host, macchine virtuali o container), consentendo la registrazione nel control plane degli indirizzi MAC e, opzionalmente, dei corrispettivi indirizzi IP.

% [IMMAGINE DEL FILE DI SLIDE CON NOME "CloudComputing_5"] Struttura binaria dettagliata dei campi dell'NLRI Type 2 MAC/IP Advertisement Route — Slide 104 / Pagina 29 del PDF

\subsection{Struttura Binaria Dei Campi NLRI}

Come dettagliato alle pagine 29 e 30, il layout del payload per una rotta di Tipo 2 presenta una lunghezza variabile tra $33$ e $49$ byte in funzione della tipologia di indirizzamento IP impiegato:

\begin{lstlisting}[language=bash]
+--------------------------------------------------+
| Route Distinguisher (RD) (8 octets)              |
+--------------------------------------------------+
| Ethernet Segment Identifier (ESI) (10 octets)    |
+--------------------------------------------------+
| Ethernet Tag ID (4 octets)                       |
+--------------------------------------------------+
| MAC Address Length (1 octet) [tipicamente 48]    |
+--------------------------------------------------+
| MAC Address (6 octets)                           |
+--------------------------------------------------+
| IP Address Length (1 octet) [0, 32 o 128 bit]    |
+--------------------------------------------------+
| IP Address (0, 4 o 16 octets)                    |
+--------------------------------------------------+
| MPLS Label 1 (3 octets)                          |
+--------------------------------------------------+
| MPLS Label 2 (0 o 3 octets)                      |
+--------------------------------------------------+
\end{lstlisting}

\paragraph{Indici di Prefisso e Attributi di Percorso.}
All'interno dell'algoritmo di BGP path selection, i campi \textbf{Ethernet Tag ID}, \textbf{MAC Address} e \textbf{IP Address} agiscono congiuntamente da \emph{chiave di prefisso} (\emph{prefix indexes}) per discriminare la singola rotta, mentre i campi \textbf{ESI} e \textbf{MPLS Label} svolgono il ruolo di attributi associati al percorso.

\subsection{I Tre Profili Funzionali Della Rotta Type 2}

La rotta Type 2 può essere impiegata in tre scenari applicativi distinti modulando i campi valorizzati:

\begin{enumerate}
    \item \textbf{MAC Advertisement Route (Puro Bridging L2)}:
    È la funzione basilare di commutazione. Quando un host si attiva, l'ingress PE ne apprende l'indirizzo MAC sull'interfaccia di accesso e genera una rotta Type 2 contenente unicamente il campo MAC, impostando il campo \emph{IP Address Length} a $0$. I nodi PE remoti popolano le rispettive tabelle di inoltro MAC (\emph{MAC-VRF}) associando quel MAC all'indirizzo del PE mittente, rimpiazzando il tradizionale flood-and-learn nel data plane.

    \item \textbf{ARP/ND Advertisement Route (ARP Suppression)}:
    Qualora il router PE sia in grado di ispezionare le comunicazioni ARP (IPv4) o \emph{Neighbor Discovery} (IPv6) dell'host, valorizza contemporaneamente sia il campo MAC sia il campo IP Address all'interno della stessa rotta Type 2.
    Questa associazione consente la funzionalità di \textbf{ARP Suppression}: quando un host remoto emette una richiesta broadcast chiedendo \emph{``Chi possiede questo indirizzo IP?''}, il VTEP locale intercetta la richiesta e formula direttamente la risposta ARP di replica prelevando le informazioni dalla propria tabella di controllo EVPN, azzerando l'inoltro di pacchetti broadcast lungo la fabric IP.

    \item \textbf{IRB Host Route (Instradamento Distribuito L3)}:
    Nelle architetture a gateway distribuito (\emph{Distributed Anycast Gateway}), dove ciascuno switch Leaf funge da default gateway per le rispettive subnet, la rotta Type 2 propaga contestualmente l'indirizzo IP host (con maschera \texttt{/32} o \texttt{/128}) unitamente a un'etichetta di livello 3 (\textbf{Layer 3 VNI}). Questo meccanismo abilita il modello di \textbf{routing simmetrico}: l'ingress VTEP esegue l'instradamento nella VRF locale, incapsula il pacchetto con il VNI L3 e il nodo egress VTEP ricevente commuta il pacchetto direttamente all'interno della medesima VRF di destinazione senza transitare per gateway centralizzati.
\end{enumerate}

\paragraph{Mobilità Degli Host.}
In caso di migrazione a caldo di una macchina virtuale (\emph{VM Live Migration}), la rotta Type 2 viene accompagnata dalla \textbf{MAC Mobility Extended Community}, che reca un contatore di sequenza (\emph{Sequence Number}). Ogni qualvolta l'host si sposta su un nuovo hypervisor/VTEP, il nuovo leaf switch emette una rotta Type 2 incrementando il numero di sequenza; i nodi della rete eleggono come autoritativo l'annuncio con sequenza maggiore e rilasciano le vecchie associazioni, prevenendo la formazione di anelli o disservizi durante la transizione.

\subsection{Mappatura Delle Etichette MPLS Label 1 E Label 2}

La tabella comparativa analizzata alle pagine 30 e 31 formalizza come i campi \textbf{MPLS Label 1} e \textbf{MPLS Label 2} vengano semanticamente interpretati a seconda della tecnologia di incapsulamento impiegata nel piano dati:

\begin{table}[H]
\centering
\small
\begin{tabularx}{\textwidth}{l p{3.5cm} X X}
\toprule
\textbf{Scenario di Servizio} & \textbf{Incapsulamento Data Plane} & \textbf{Significato di ``MPLS Label 1''} & \textbf{Significato di ``MPLS Label 2''} \\
\midrule
\textbf{L2 Service (Bridging)} &
EVPN-MPLS &
Etichetta MPLS che identifica l'istanza di servizio L2 (\textbf{EVI}). &
Non impiegata (campo assente o valorizzato a $0$). \\
\addlinespace
\textbf{L2 Service (Bridging)} &
EVPN-VXLAN &
\textbf{Layer 2 VNI} (il VNI che identifica il dominio di broadcast L2). &
Non impiegata (campo assente o valorizzato a $0$). \\
\addlinespace
\textbf{L3 Service (IRB Routing)} &
EVPN-MPLS &
Etichetta MPLS associata all'istanza di bridging locale (EVI). &
\textbf{Etichetta MPLS di VRF} (impiegata per l'instradamento inter-subnet). \\
\addlinespace
\textbf{L3 Service (IRB Routing)} &
EVPN-VXLAN &
\textbf{Layer 2 VNI} associato alla subnet del mittente. &
\textbf{Layer 3 VNI} (il VNI che identifica univocamente la VRF di routing). \\
\bottomrule
\end{tabularx}
\caption{Interpretazione e mappatura delle etichette Label 1 e Label 2 in rotte Type 2.}
\end{table}

\subsection{Dinamica Del Flusso Informativo In Un Dominio EVPN Type 2}

La Slide~105 (pagina~32) illustra visivamente la convergenza dei database locali di due nodi di bordo denominati NVE1 (\emph{Network Virtualization Edge 1}) ed NVE2:

% [IMMAGINE DEL FILE DI SLIDE CON NOME "CloudComputing_5"] Scambio di rotta Type 2 tra NVE1 e NVE2 e popolamento delle tabelle MAC, ARP e IP Route Table — Slide 105 / Pagina 32 del PDF

Il processo si sviluppa attraverso i seguenti stadi:
\begin{enumerate}
    \item L'host \texttt{HOST1} (caratterizzato dalla coppia \texttt{MAC1} e \texttt{IP1}) invia una trama di rete che approda sulla porta \texttt{Port1} del dispositivo NVE1.
    \item NVE1 apprende localmente le informazioni e aggiorna i propri database di accesso:
    \begin{itemize}[noitemsep]
        \item Nella \textbf{MAC Table} locale registra: $\text{Destination} = \texttt{MAC1}$, $\text{Nexthop} = \texttt{Port1}$, $\text{Learned from} = \texttt{Port1}$.
        \item Nella \textbf{ARP Table} associa: $\text{IP} = \texttt{IP1}$, $\text{MAC} = \texttt{MAC1}$, $\text{Learned from} = \texttt{Port1}$.
    \end{itemize}
    \item Il processo BGP-EVPN attivo su NVE1 converte immediatamente tali associazioni in una rotta \textbf{Type 2 MAC/IP Advertisement} e la serializza in un messaggio MP-BGP verso i peer della fabric, impostando il proprio indirizzo come salto successivo (\texttt{Nexthop = NVE1}).
    \item Il dispositivo remoto NVE2 riceve la rotta Type 2 e, dopo aver validato la corrispondenza del Route Target, popola in modo deterministico le proprie strutture dati interne senza che si sia verificato alcun flooding nel data plane:
    \begin{itemize}[noitemsep]
        \item Nella propria \textbf{MAC Table} inserisce una voce remota: $\text{Destination} = \texttt{MAC1}$, $\text{Nexthop} = \texttt{NVE1}$, $\text{Learned from} = \texttt{BGP}$.
        \item Nella propria \textbf{Route Table} (se abilitato il routing distribuito) installa la rotta host a 32 bit: $\text{Destination} = \texttt{IP1/32}$, $\text{Nexthop} = \texttt{NVE1}$, $\text{Learned from} = \texttt{BGP}$.
    \end{itemize}
\end{enumerate}

Qualora in seguito un secondo host (\texttt{HOST2}) collegato su NVE2 necessiti di trasmettere dati verso \texttt{IP1}, NVE2 non dovrà ricorrere a broadcast ARP generici: risolverà l'indirizzo internamente e instraderà il traffico direttamente in unicast incapsulato verso l'indirizzo IP di NVE1, garantendo un'infrastruttura di data center deterministica, scalabile e priva di broadcast storm.

% DD -----------------------------------------------------------

% ==============================================================================
% CONTROLLO DEL PIANO DI CONTROLLO E INOLTRO IN RETI EVPN-VXLAN
% ==============================================================================

\section{Control Plane BGP EVPN Nelle Reti Overlay VXLAN}

\subsection{Richiamo Concettuale: Dal Flood-and-Learn Nel Data Plane Al Control Plane Strutturato}

Nelle architetture di rete tradizionali basate su protocollo Ethernet (IEEE 802.1D/802.1Q), l'apprendimento della topologia e la localizzazione dei nodi terminali avvengono esclusivamente a livello di \textbf{piano dati} (\emph{data plane}). Questo meccanismo, universalmente noto come \textbf{Flood-and-Learn}, prevede che uno switch, a fronte di una trama con indirizzo MAC di destinazione sconosciuto (\emph{Unknown Unicast}), di tipo broadcast (es. richieste ARP) o multicast (traffico collettivamente denominato \textbf{BUM}, \emph{Broadcast, Unknown unicast, Multicast}), replichi la trama su tutte le porte appartenenti allo stesso dominio di broadcast eccetto quella di ricezione. L'apprendimento dell'associazione tra porta fisica e indirizzo MAC sorgente avviene per auto-apprendimento passivo (\emph{backward learning}) osservando i pacchetti in transito.

Quando la tecnologia \textbf{VXLAN} (\emph{Virtual Extensible LAN}, formalizzata in \textbf{RFC 7348}) è stata introdotta per superare il limite canonico dei \num{4096} identificatori VLAN e consentire l'estensione di segmenti di Livello~2 al di sopra di un'infrastruttura di Livello~3 routed (\emph{underlay network}), il meccanismo iniziale di funzionamento replicava fedelmente il modello flood-and-learn. In assenza di un piano di controllo dedicato, i punti terminali del tunnel (\textbf{VTEP}, \emph{VXLAN Tunnel Endpoint}) erano costretti a ricorrere al multicast nell'underlay per distribuire le trame BUM. Tale approccio presentava evidenti criticità architetturali:
\begin{itemize}[noitemsep]
    \item Complessità gestionale dell'underlay, costretto a implementare protocolli di instradamento multicast quali \textbf{PIM} (\emph{Protocol Independent Multicast}) in modalità Sparse Mode o Source-Specific Multicast.
    \item Elevato spreco di banda e consumo di risorse elaborative a causa della replicazione massiva del traffico non noto.
    \item Assoluta mancanza di sincronizzazione deterministica tra i nodi d'accesso e difficoltà nella gestione rapida della mobilità delle macchine virtuali (\emph{VM mobility}).
\end{itemize}

Per risolvere radicalmente tali limitazioni, l'IETF ha standardizzato l'impiego del protocollo \textbf{Multiprotocol BGP} (MP-BGP) esteso con la famiglia d'indirizzi \textbf{EVPN} (\emph{Ethernet Virtual Private Network}, originariamente definita in \textbf{RFC 7432}) come \textbf{piano di controllo} (\emph{control plane}) unificato per overlay VXLAN. In tale contesto:
\begin{enumerate}[noitemsep]
    \item BGP EVPN sostituisce il flood-and-learn sul data plane con una propagazione deterministica e scalabile delle informazioni di raggiungibilità (indirizzi MAC e IP degli host, binding con i VTEP remoti).
    \item I VTEP apprendono gli host locali tramite le interfacce d'accesso (tramite snooping ARP/NDP o traffico dati) e incapsulano tali informazioni all'interno di rotte BGP, distribuendole agli altri VTEP tramite \textbf{Route Reflector} (RR).
    \item Il piano dati VXLAN si occupa esclusivamente del trasporto incapsulato \emph{point-to-point} incapsulando il frame originale in pacchetti \texttt{MAC-in-UDP}, preservando l'astrazione di un unico switch logico distribuito.
\end{enumerate}

% [IMMAGINE DEL FILE DI SLIDE CON NOME "cloudComputing_5"] Topologia generale dell'annuncio di rotte MAC/IP tra host H1 e H4 tramite BGP EVPN e VTEP NVE1 e NVE2 — Slide 106 / Pagina 1 del PDF

---

\section{Annuncio Delle Rotte MAC/IP: EVPN Route Type 2}

\subsection{Struttura Del Network Layer Reachability Information (NLRI) Type 2}

La rotta EVPN di Tipo~2, denominata \textbf{MAC/IP Advertisement Route}, è il pilastro del piano di controllo per l'inoltro di Livello~2 e per l'instradamento distribuito host-based (\emph{host route}). La sua funzione principale è propagare l'associazione tra un indirizzo MAC, un eventuale indirizzo IP e il VTEP di ancoraggio attraverso il messaggio \texttt{MP\_REACH\_NLRI} di BGP.

La struttura dell'EVPN Route Type~2 comprende i seguenti campi fondamentali:
\begin{itemize}
    \item \textbf{Route Distinguisher (RD)} ($8$ ottetti): garantisce l'univocità globale del prefisso all'interno delle tabelle BGP, consentendo la coesistenza di indirizzi sovrapposti in \emph{Virtual Routing and Forwarding} (VRF) o istanze EVPN (\textbf{EVI}) distinte.
    \item \textbf{Ethernet Segment Identifier (ESI)} ($10$ ottetti): identifica il segmento fisico o logico a cui l'host è collegato (pari a zero nel caso di host \emph{single-homed}).
    \item \textbf{Ethernet Tag ID} ($4$ ottetti): identifica il contesto di broadcast domain specifico (spesso impostato a \num{0} nelle mappature VLAN-based standard).
    \item \textbf{MAC Address Length} ($1$ ottetto) e \textbf{MAC Address} ($6$ ottetti): specifica l'indirizzo fisico dell'host di destinazione, fondamentale per popolare la tabella di inoltro di Livello~2 (\emph{MAC Table}).
    \item \textbf{IP Address Length} ($1$ ottetto) e \textbf{IP Address} ($0$, $4$ o $16$ ottetti): campo opzionale. Se omesso, la rotta distribuisce unicamente connettività L2 (\emph{pure MAC route}); se valorizzato (con indirizzo IPv4 a \SI{32}{\bit} o IPv6 a \SI{128}{\bit}), abilita l'\textbf{Integrated Routing and Bridging} (IRB), consentendo l'apprendimento della rotta di host direttamente nel piano di controllo.
    \item \textbf{MPLS Label 1 (L2 VNI)} ($3$ ottetti / \SI{24}{\bit}): nelle implementazioni EVPN-VXLAN, il campo etichetta MPLS è reinterpretato per trasportare il \textbf{Layer 2 VNI} (\emph{VXLAN Network Identifier}), identificando univocamente il dominio di bridging di appartenenza.
    \item \textbf{MPLS Label 2 (L3 VNI)} ($3$ ottetti, opzionale): utilizzato nelle architetture \emph{Symmetric IRB} per indicare il VNI associato alla VRF di routing del tenant.
    \item \textbf{BGP Next-Hop}: corrisponde all'indirizzo IP del VTEP sorgente che ha originato la rotta (es. \texttt{1.1.1.9}), indicando all'overlay verso quale destinazione fisica nell'underlay instradare i pacchetti VXLAN.
    \item \textbf{BGP Tunnel Encapsulation Attribute}: attributo esteso (RFC 9012) che specifica la tipologia del tunnel, codificato con il valore numerico intero $8$ per indicare l'incapsulamento VXLAN.
\end{itemize}

\subsection{Procedura Di Annuncio E Inoltro Locale/Remoto}

Il processo operativo con cui un pacchetto generato da un host attiva la sincronizzazione del piano di controllo e l'aggiornamento del piano di inoltro si articola nelle seguenti fasi:

% [IMMAGINE DEL FILE DI SLIDE CON NOME "cloudComputing_5"] Schema del meccanismo di apprendimento locale del MAC e annuncio Type 2 al peer con popolamento della MAC Table e Tunnel Table — Slide 107 / Pagina 2 del PDF

\begin{enumerate}
    \item \textbf{Rilevamento locale dell'host}: l'host locale (es. \texttt{HOST1} con indirizzo \texttt{H1-MAC}) accede alla rete trasmettendo una trama (es. una richiesta o una risposta ARP). Il VTEP locale d'accesso (nodo \texttt{NVE1} con indirizzo IP di underlay \texttt{1.1.1.9}) intercetta la trama sulla propria porta di accesso associata al bridge domain (\textbf{BD}) con identificatore \texttt{VNI 100}.
    \item \textbf{Generazione della rotta BGP EVPN}: \texttt{NVE1} registra l'indirizzo \texttt{H1-MAC} nella propria tabella MAC locale e costruisce un aggiornamento BGP contenente un NLRI di Tipo~2:
    \[
    \text{RT2} = \Big\{ \text{RD}=1:1, \; \text{MAC}=\texttt{H1-MAC}, \; \text{Label}=\text{L2 VNI}:100, \; \text{Tunnel}=\text{VXLAN}, \; \text{Next-Hop}=1.1.1.9 \Big\}
    \]
    Tale annuncio viene inviato ai nodi Route Reflector (\textbf{RR}), i quali lo riflettono a tutti i VTEP peer (es. nodo \texttt{NVE2} con indirizzo \texttt{2.2.2.9}).
    \item \textbf{Ricezione ed elaborazione remota}: il VTEP remoto \texttt{NVE2} riceve l'annuncio e ne analizza gli attributi:
    \begin{itemize}[noitemsep]
        \item Estrae l'identificatore VNI ($100$) e verifica se localmente è istanziata la medesima EVI. In caso affermativo, la rotta viene accettata.
        \item Ispeziona il campo BGP Next-Hop (\texttt{1.1.1.9}) ed effettua una ricerca nella tabella di instradamento dell'underlay (\emph{Routing Table} globale) per accertarsi che tale indirizzo sia raggiungibile.
        \item Se il next-hop è valido, \texttt{NVE2} programma il piano di inoltro hardware:
        \begin{itemize}
            \item Crea una riga nella \textbf{MAC Table}: destinazione \texttt{H1-MAC}, associata a \texttt{VNI 100} e con \emph{Destination IP} (DIP) impostato a \texttt{1.1.1.9}.
            \item Aggiorna la \textbf{Tunnel Table}, mappando il tunnel verso \texttt{1.1.1.9} con sorgente (\emph{Source IP}, SIP) \texttt{2.2.2.9}.
        \end{itemize}
    \end{itemize}
\end{enumerate}

Grazie a questo flusso, quando un secondo terminale (\texttt{HOST4}) collegato a \texttt{NVE2} deve comunicare con \texttt{HOST1}, la trama Ethernet non subisce alcun flooding: \texttt{NVE2} consulta direttamente la MAC table programmata dal control plane, incapsula la trama nel pacchetto VXLAN unicast diretto a \texttt{1.1.1.9} e la instrada sull'underlay.

\subsection{Mobilità Delle Macchine Virtuali (VM Migration) E Sequence Number}

Uno dei maggiori vantaggi operativi derivanti dall'adozione del piano di controllo BGP EVPN risiede nella gestione trasparente, tempestiva e priva di anomalie della migrazione a caldo (\emph{live migration}) delle macchine virtuali da un server fisico attestato su un leaf switch a un altro nodo d'accesso.

% [IMMAGINE DEL FILE DI SLIDE CON NOME "cloudComputing_5"] Flusso temporale e aggiornamento del Sequence Number durante la migrazione di VM1 da NVE1 a NVE2 — Slide 108 / Pagina 3 del PDF

Quando una macchina virtuale \texttt{VM1} (con indirizzo fisso \texttt{MAC1-IP1}) migra fisicamente dal nodo sorgente \texttt{NVE1} al nodo destinazione \texttt{NVE2}, si attiva la procedura standardizzata basata sull'attributo esteso \textbf{MAC Mobility Extended Community}:

\begin{enumerate}
    \item \textbf{Stato iniziale}: prima della migrazione, tutti i nodi della fabric (\texttt{NVE1}, \texttt{NVE2}, \texttt{NVE3}) hanno registrato la locazione di \texttt{VM1} su \texttt{NVE1}. La voce nella tabella MAC di \texttt{NVE1} indica come nexthop la porta fisica locale (\texttt{Port1}), con \textbf{Sequence Number} impostato a un valore base $N$ (tipicamente $N=0$ all'attivazione iniziale). Nei VTEP remoti (\texttt{NVE2}, \texttt{NVE3}), la voce è appresa via BGP con next-hop \texttt{NVE1} e sequenza $N$.
    \item \textbf{Migrazione e rilevamento locale}: \texttt{VM1} completa il trasferimento e trasmette una trama di rete (ad esempio un pacchetto RARP o un gratuitous ARP) sulla nuova porta fisica di \texttt{NVE2}.
    \item \textbf{Incremento del Sequence Number}: \texttt{NVE2} rileva la presenza locale di \texttt{MAC1-IP1}. Riscontrando nella propria tabella di routing BGP una voce preesistente per il medesimo MAC associata al numero di sequenza $N$, \texttt{NVE2} riconosce l'evento di mobilità:
    \[
    \text{Sequence Number}_{\text{nuovo}} = N + 1
    \]
    \texttt{NVE2} aggiorna la propria tabella di inoltro locale associando \texttt{MAC1} alla propria interfaccia fisica con sequenza $N+1$, e genera istantaneamente una rotta EVPN Type~2 contenente la community \emph{MAC Mobility} valorizzata con $\text{SEQ} = N + 1$.
    \item \textbf{Arbitraggio e convergenza della fabric}: la nuova rotta viene annunciata via BGP a tutti i nodi:
    \begin{itemize}[noitemsep]
        \item I nodi terzi (come \texttt{NVE3}) confrontano il valore di sequenza della nuova rotta ($N+1$) con quello precedentemente registrato ($N$). Poiché:
        \[
        N + 1 > N
        \]
        la rotta con numero di sequenza superiore prevale in modo deterministico. La MAC table viene aggiornata sostituendo il nexthop \texttt{NVE1} con \texttt{NVE2}.
        \item Il VTEP d'origine \texttt{NVE1} riceve la rotta Type~2 con sequenza $N+1$. Riconoscendo che la localizzazione remota è più recente rispetto alla propria registrazione locale (ferma a $N$), \texttt{NVE1} ritira la propria rotta precedente inviando un messaggio di \textbf{BGP Route Withdrawal} e cancella o disattiva l'entry associata alla porta locale.
    \end{itemize}
\end{enumerate}

Questo meccanismo elimina il rischio di loop di inoltro temporanei o di traffico \emph{black-holed} (inviato verso il vecchio host), garantendo una convergenza sub-secondaria senza dover ricorrere a complessi timer di invecchiamento (\emph{aging timer}) sul data plane.

\subsection{Host Route Learning In Ambito Integrated Routing And Bridging (IRB)}

Quando l'annuncio Type~2 viene impiegato congiuntamente per il bridging e per l'instradamento inter-sottorete (IRB), la rotta trasporta contestualmente le coordinate di Livello~2 e di Livello~3 dell'host.

% [IMMAGINE DEL FILE DI SLIDE CON NOME "cloudComputing_5"] Schema architetturale e tabelle di inoltro per Host Route Learning con Type 2 (MAC + IP) e dual-label (L2 VNI e L3 VNI) — Slide 109 / Pagina 4 del PDF

Come schematizzato nella documentazione tecnica, quando un host (\texttt{HOST1} con indirizzo \texttt{10.1.1.2} e \texttt{H1-MAC}) entra in funzione sul bridge domain associato a \texttt{VNI 100}:
\begin{itemize}
    \item Il VTEP locale intercetta il traffico ARP e apprende la corrispondenza univoca tra \texttt{10.1.1.2} e \texttt{H1-MAC}.
    \item Viene generata una rotta Type~2 contenente due etichette:
    \[
    \text{RT2} = \Big\{ \text{RD}=1:1, \; \text{MAC}=\texttt{H1-MAC}, \; \text{IP}=10.1.1.2, \; \text{Label1}=\text{L2VNI } 100, \; \text{Label2}=\text{L3VNI } 1000 \Big\}
    \]
    La rotta include l'attributo esteso BGP \textbf{Route Target} (\textbf{RT}) di esportazione (\textbf{ERT}, \emph{Export Route Target}), impostato a \texttt{1:1}.
    \item Il peer remoto \texttt{NVE2} riceve la rotta. L'associazione con la VRF locale del tenant non si basa su un generico identificatore locale di VRF (che ha valenza prettamente circoscritta al singolo apparato), bensì sulla corrispondenza dell'\textbf{Import Route Target} (\textbf{IRT}) configurato nella VRF con l'ERT trasportato dal messaggio BGP:
    \[
    \text{Se } \text{IRT}_{\text{locale}} == \text{ERT}_{\text{rotta}} \implies \text{Importazione nella VRF di pertinenza}
    \]
    \item All'atto dell'importazione, \texttt{NVE2} compila simultaneamente tre strutture dati fondamentali:
    \begin{enumerate}[noitemsep]
        \item \textbf{MAC Table}: associa \texttt{H1-MAC} al dominio di L2 (\texttt{VNI 100}) e al DIP \texttt{1.1.1.9}.
        \item \textbf{Tunnel Table}: assicura l'esistenza del tunnel VXLAN bidirezionale con \texttt{1.1.1.9}.
        \item \textbf{FIB Table (L3 Forwarding Information Base)}: installa un prefisso di host \texttt{10.1.1.2/32} vincolato a \texttt{VRF 1}, specificando come incapsulamento l'identificatore \texttt{L3 VNI 1000}, il MAC del router di destinazione (\textbf{Router MAC} di \texttt{NVE1}) e il nexthop IP di underlay \texttt{1.1.1.9}.
    \end{enumerate}
\end{itemize}

---

\newpage

\section{Gestione Del Traffico BUM E Scoperta Automatica Dei Tunnel: EVPN Route Type 3}

\subsection{Definizione E Formato Dell'Inclusive Multicast Ethernet Tag Route}

Nelle reti prive di multicast nell'infrastruttura fisica di trasporto, l'inoltro del traffico BUM richiede che ciascun VTEP mittente mantenga una lista accurata di tutti gli altri VTEP appartenenti al medesimo VNI, in modo da generare copie unicast del pacchetto solo verso i destinatari interessati (modalità nota come \textbf{Ingress Replication} o \emph{Head-End Replication}).

Il meccanismo con cui un VTEP annuncia proattivamente la propria presenza e il proprio interesse a ricevere traffico per uno specifico dominio broadcast è demandato alla rotta EVPN di Tipo~3, formalmente denominata \textbf{Inclusive Multicast Ethernet Tag Route}.

% [IMMAGINE DEL FILE DI SLIDE CON NOME "cloudComputing_5"] Formato del messaggio NLRI Type 3 ed estrazione del VNI tramite BGP Update Extended Community — Slide 110 / Pagina 5 del PDF

Il formato standard dell'NLRI di Tipo~3 è strutturato come segue:
\begin{center}
\begin{tabular}{|l|c|p{8cm}|}
\hline
\textbf{Campo NLRI} & \textbf{Lunghezza} & \textbf{Descrizione e Funzione Operativa} \\ \hline
\texttt{Route Distinguisher (RD)} & $8$ ottetti & Identifica l'univocità dell'annuncio BGP associato alla EVI. \\ \hline
\texttt{Ethernet Tag ID} & $4$ ottetti & Identificatore del tag di trasmissione broadcast (tipicamente \num{0}). \\ \hline
\texttt{IP Address Length} & $1$ ottetto & Lunghezza dell'indirizzo IP del VTEP originante ($32$ per IPv4, $128$ per IPv6). \\ \hline
\texttt{Originating Router's IP} & $4$ o $16$ ottetti & Indirizzo IP del VTEP sorgente che annuncia l'adesione al gruppo. \\ \hline
\end{tabular}
\end{center}

\begin{theorem}[Trasporto Del VNI Nelle Rotte Type 3]
Un aspetto fondamentale delle reti EVPN-VXLAN è che \textbf{il VNI non è codificato all'interno del campo NLRI primario della rotta Type~3}.
Il valore del \textbf{Layer 2 VNI} viene inserito e trasportato negli attributi di percorso (\emph{Path Attributes}) dell'aggiornamento BGP (\texttt{BGP UPDATE}), specificamente all'interno della \textbf{BGP Extended Community} (o nel PMSI Tunnel Attribute, riutilizzando il campo a \SI{24}{\bit} originariamente riservato all'etichetta MPLS).
\end{theorem}

Alla ricezione del messaggio \texttt{BGP UPDATE}, il router ricevente decodifica gli attributi estesi, estrae il valore dal campo etichetta, lo interpreta quale identificatore VNI e aggiorna la corrispondente lista di replicazione per quel dominio di broadcast.

\subsection{Scoperta Dinamica Dei VTEP E Stabilimento Dei Tunnel VXLAN}

La combinazione della rotta Type~3 con l'indirizzamento di underlay consente la completa automazione del ciclo di vita dei tunnel overlay:

% [IMMAGINE DEL FILE DI SLIDE CON NOME "cloudComputing_5"] Topologia Spine-Leaf con annuncio Type 3 tra NVE1 e NVE2 e compilazione dell'Ingress Replication List — Slide 111-112 / Pagine 6-7 del PDF

\begin{enumerate}
    \item Non appena su un VTEP (\texttt{NVE1}) viene configurato un bridge domain locale associato a un \texttt{L2 VNI} (es. \texttt{VNI 100}) e un'interfaccia di loopback come sorgente VTEP (es. \texttt{1.1.1.9}), il demone BGP genera automaticamente un annuncio Type~3:
    \[
    \text{RT3} = \Big\{ \text{RD}=1:1, \; \text{Originating IP}=1.1.1.9, \; \text{Label}=\text{L2 VNI } 100, \; \text{Tunnel ID}=\texttt{1.1.1.9} \Big\}
    \]
    \item Il messaggio raggiunge il nodo remoto \texttt{NVE2} attraverso l'infrastruttura BGP.
    \item \texttt{NVE2} confronta i parametri ricevuti con la propria configurazione locale:
    \begin{itemize}[noitemsep]
        \item Se \texttt{NVE2} possiede a sua volta istanziato il medesimo \texttt{L2 VNI 100}, riconosce che entrambi i nodi condividono lo stesso dominio broadcast.
        \item Verifica nella tabella di instradamento IP globale che l'\texttt{Originating IP} (\texttt{1.1.1.9}) sia raggiungibile nell'underlay.
        \item Se la verifica di raggiungibilità ha esito positivo, \texttt{NVE2} crea o valida il tunnel VXLAN verso \texttt{1.1.1.9} e aggiunge l'indirizzo \texttt{1.1.1.9} alla propria \textbf{Ingress Replication List} associata a \texttt{VNI 100}.
    \end{itemize}
\end{enumerate}

\subsection{Ruolo Dei Nodi Spine Nell'Underlay E Assenza Di Decapsulamento}

Nelle topologie moderne di tipo \textbf{Spine-Leaf}, l'interazione tra i piani di instradamento rispetta una rigorosa gerarchia architetturale:
\begin{itemize}
    \item \textbf{I nodi Leaf (NVE)} agiscono come punti terminali dell'incapsulamento VXLAN. Essi gestiscono le interfacce verso i server, eseguono il bridging L2, mantengono le VRF del tenant ed effettuano l'incapsulamento e il decapsulamento dell'intestazione VXLAN/UDP.
    \item \textbf{I nodi Spine non decapsulano mai il traffico VXLAN}. Per la dorsale Spine, i pacchetti overlay sono normali pacchetti unicast IP/UDP trasportati attraverso l'underlay. I nodi Spine eseguono esclusivamente l'inoltro di Livello~3 (\emph{IP routing}) basato sull'intestazione esterna (\textbf{outer header}), instradando il pacchetto verso l'indirizzo IP del VTEP di destinazione.
    \item Dal punto di vista del piano di controllo, i nodi Spine ospitano frequentemente la funzione di \textbf{BGP Route Reflector} (RR) per scambiare le rotte EVPN tra tutti i Leaf, pur rimanendo completamente disconnessi dai domini di bridging di Livello~2 sottostanti.
\end{itemize}

% [IMMAGINE DEL FILE DI SLIDE CON NOME "cloudComputing_5"] Riepilogo della gestione del traffico BUM tramite unicast replication e BGP Type 3 senza multicast nell'underlay — Slide 113 / Pagina 8 del PDF

---

\section{Multihoming E Designated Forwarder Election: EVPN Route Type 4}

\subsection{La Sfida Del Multihoming E Prevenzione Dei Loop Di Traffico BUM}

Nei datacenter ad alta affidabilità, i server o gli apparati di bordo (\textbf{Customer Edge}, \textbf{CE}) non vengono attestati su un singolo switch leaf, ma sono collegati simultaneamente a due o più leaf distinti (\textbf{Provider Edge}, \textbf{PE}) tramite aggregazioni di link ridondate (\emph{Link Aggregation Group}, LAG). In terminologia EVPN, tale configurazione definisce un \textbf{Ethernet Segment} multi-homed, contraddistinto da un identificatore univoco denominato \textbf{Ethernet Segment Identifier} (\textbf{ESI}).

La connessione contemporanea di un CE a molteplici PE introduce una severa criticità potenziale nell'instradamento del traffico BUM: se una trama BUM originata da un host remoto giunge simultaneamente a entrambi i PE collegati allo stesso segmento ESI, ed entrambi la inoltrassero verso il CE, il CE riceverebbe copie duplicate dei pacchetti. Inoltre, si innescherebbero pericolosi anelli di commutazione (\emph{forwarding loop}) tra i link del LAG.

Per prevenire tale condizione, l'architettura EVPN impone che per ogni segmento Ethernet condiviso venga eletto un unico apparato abilitato a inoltrare il traffico broadcast, sconosciuto e multicast verso il dispositivo d'accesso. Tale apparato assume il ruolo di \textbf{Designated Forwarder} (\textbf{DF}), mentre gli altri apparati collegati al medesimo ESI operano come \textbf{Non-DF} (o nodi di backup) bloccando la trasmissione di traffico BUM sul link d'accesso.

% [IMMAGINE DEL FILE DI SLIDE CON NOME "cloudComputing_5"] Topologia multihoming con CE2 dual-homed verso PE2 e PE3, rappresentazione del traffico BUM e segnalazione BGP — Slide 115 / Pagina 10 del PDF

\subsection{Formato Dell'Ethernet Segment Route (NLRI Type 4)}

L'elezione del DF richiede che tutti i router attestati sullo stesso ESI si riconoscano reciprocamente in modo deterministico. Questa scoperta avviene mediante lo scambio della rotta EVPN di Tipo~4, formalmente nota come \textbf{Ethernet Segment Route}.

% [IMMAGINE DEL FILE DI SLIDE CON NOME "cloudComputing_5"] Formato dei campi NLRI della rotta Type 4 per Ethernet Segment — Slide 114 / Pagina 9 del PDF

Il pacchetto NLRI di Tipo~4 presenta i seguenti campi:
\begin{itemize}[noitemsep]
    \item \textbf{Route Distinguisher (RD)} ($8$ ottetti).
    \item \textbf{Ethernet Segment Identifier (ESI)} ($10$ ottetti): valore a \SI{80}{\bit} condiviso tra i nodi che collegano il CE multi-homed.
    \item \textbf{IP Address Length} ($1$ ottetto): lunghezza dell'indirizzo del PE originante.
    \item \textbf{Originating Router's IP Address} ($4$ o $16$ ottetti): indirizzo IP di underlay del PE candidato.
\end{itemize}

\subsection{Algoritmo Di Elezione Deterministica Del DF (RFC 7432)}

L'elezione del DF non viene condotta in modo monolitico per l'intero segmento, ma è calibrata con granularità \textbf{per-segmento, per-VLAN} (o per-VNI). Ciò consente di bilanciare equamente il carico di traffico BUM tra i router multi-homed, assegnando la funzione di DF a PE differenti per insiemi disgiunti di VLAN.

L'algoritmo standardizzato opera sequenzialmente come segue:

\begin{enumerate}
    \item \textbf{Scambio delle rotte Type 4}: all'attivazione dell'interfaccia ESI, ciascun PE pubblica la propria rotta Type~4 annunciando il proprio \texttt{ESI} e il proprio \texttt{IP Address} (o Router-ID). Un timer di attesa (\emph{DF Election Timer}, impostato convenzionalmente a \SI{3}{\second}) previene instabilità transitorie.
    \item \textbf{Costruzione della lista ordinata dei candidati}: ciascun PE raccoglie tutte le rotte Type~4 ricevute che presentano il medesimo valore di ESI. Viene così costruita localmente una lista di tutti i PE partecipanti, ordinata in senso strettamente crescente in base al loro indirizzo IP numerico (o Router-ID):
    \[
    \mathcal{C} = \Big[ PE_0, \; PE_1, \; \dots, \; PE_{N-1} \Big] \quad \text{con } IP(PE_0) < IP(PE_1) < \dots < IP(PE_{N-1})
    \]
    dove $N$ rappresenta la cardinalità complessiva dei PE attestati sul segmento.
    \item \textbf{Calcolo dell'indice vincente tramite operatore modulo}: per ciascuna VLAN (o Ethernet Tag ID) $V$ configurata sul segmento, l'indice $i$ del PE eletto Designated Forwarder è calcolato mediante la relazione:
    \begin{equation}
        \boxed{i = V \pmod N}
    \end{equation}
    Il router la cui posizione ordinale nella lista $\mathcal{C}$ coincide con l'indice calcolato $i$ assume lo stato di \textbf{DF} per la specifica VLAN $V$; tutti gli altri router assumono lo stato di \textbf{Non-DF}.
\end{enumerate}

\subsubsection{Esempio Pratico Di Bilanciamento Del Traffico}

Si consideri uno scenario con due nodi Provider Edge, $PE_2$ e $PE_3$, attestati sul medesimo segmento ESI ($N=2$), con indirizzi ordinati in modo tale che $PE_2$ occupi l'indice ordinale $0$ e $PE_3$ l'indice ordinale $1$:
\begin{align*}
    \text{Per la VLAN } 100: \quad & i = 100 \pmod 2 = 0 \implies \mathbf{PE_2} \text{ è eletto DF}. \\
    \text{Per la VLAN } 101: \quad & i = 101 \pmod 2 = 1 \implies \mathbf{PE_3} \text{ è eletto DF}. \\
    \text{Per la VLAN } 102: \quad & i = 102 \pmod 2 = 0 \implies \mathbf{PE_2} \text{ è eletto DF}. \\
    \text{Per la VLAN } 103: \quad & i = 103 \pmod 2 = 1 \implies \mathbf{PE_3} \text{ è eletto DF}.
\end{align*}

Questo meccanismo assicura che il traffico BUM per la VLAN~$100$ venga inoltrato al CE esclusivamente da $PE_2$, mentre il traffico per la VLAN~$101$ venga erogato unicamente da $PE_3$, eliminando i loop e garantendo una suddivisione simmetrica del throughput disponibile.

---

\newpage

\section{Instradamento Inter-Subnet E Reti Esterne: EVPN Route Type 5}

\subsection{Formato Dell'IP Prefix Route (NLRI Type 5)}

Nelle architetture \textbf{Network Virtualization Overlay} (NVO) basate su VXLAN (RFC 7348), NVGRE (RFC 7637) o GENEVE (RFC 8926), l'interconnessione tra sottoreti distinte e l'integrazione con infrastrutture esterne (reti WAN, Internet o datacenter legacy) non possono basarsi unicamente sulla conoscenza degli host host-per-host (/32 o /128). È indispensabile poter iniettare e propagare prefissi di sottorete aggregati (\emph{network segment routes}).

A questo scopo è deputata la rotta EVPN di Tipo~5, definita formalmente come \textbf{IP Prefix Route} (standardizzata nell'ambito della specifica RFC 9136, ex \texttt{draft-ietf-bess-evpn-prefix-advertisement}).

% [IMMAGINE DEL FILE DI SLIDE CON NOME "cloudComputing_5"] Formato dettagliato dell'NLRI Type 5 per la propagazione di prefissi IP interni ed esterni — Slide 118 / Pagina 14 del PDF

Il blocco informativo dell'NLRI Type~5 è codificato nel seguente modo:
\begin{center}
\begin{tabular}{|l|c|p{8.5cm}|}
\hline
\textbf{Campo NLRI Type 5} & \textbf{Dimensione} & \textbf{Significato Operativo} \\ \hline
\texttt{Route Distinguisher (RD)} & $8$ ottetti & Identifica univocamente l'istanza VRF di provenienza della rotta. \\ \hline
\texttt{Ethernet Segment Identifier} & $10$ ottetti & ESI del gateway o del nodo originante (pari a $0$ nei prefissi globali). \\ \hline
\texttt{Ethernet Tag ID} & $4$ ottetti & Identificativo di contesto di commutazione. \\ \hline
\texttt{IP Prefix Length} & $1$ ottetto & Lunghezza della maschera di rete (da $0$ a $32$ per IPv4, da $0$ a $128$ per IPv6). \\ \hline
\texttt{IP Prefix} & Fino a $16$ ottetti & Prefisso IP della sottorete annunciata (es. \texttt{10.20.20.0}). \\ \hline
\texttt{GW IP Address} & $4$ o $16$ ottetti & Indirizzo IP del gateway di inoltro (utilizzato in scenari non-overlay). \\ \hline
\texttt{MPLS Label (L3 VNI)} & $3$ ottetti & Etichetta che specifica il \textbf{Layer 3 VNI} associato alla VRF di destinazione. \\ \hline
\end{tabular}
\end{center}

\subsection{Distinzione Concettuale Tra Route Type 2 E Route Type 5}

È essenziale comprendere la profonda differenza operativa e architetturale che intercorre tra le rotte di Tipo~2 e quelle di Tipo~5:
\begin{itemize}
    \item \textbf{Route Type 2 (MAC/IP Route)}: modella uno stato di \textbf{host puro}. Ogni indirizzo IP eventualmente trasportato è rigorosamente \emph{vincolato a un indirizzo MAC} ben definito. Ha maschera implicita fissa (/32 per IPv4, /128 per IPv6) e richiede che il VTEP crei un record simultaneo nella MAC table di Livello~2 e nella FIB di Livello~3.
    \item \textbf{Route Type 5 (IP Prefix Route)}: modella un'informazione di instradamento di \textbf{puro Livello~3}. Il prefisso IP \emph{non è vincolato ad alcun indirizzo MAC} di host. La maschera può assumere qualsiasi estensione (/24, /16, rotta di default 0.0.0.0/0). Viene utilizzata per importare sottoreti esterne apprese tramite protocolli di routing dinamico tradizionali (eBGP, OSPF) o rotte statiche all'interno della VRF del tenant.
\end{itemize}

% [IMMAGINE DEL FILE DI SLIDE CON NOME "cloudComputing_5"] Topologia di interconnessione con routing esterno: NVE3 apprende una rotta esterna e la distribuisce come Type 5 a NVE1 e NVE2 — Slide 119 / Pagina 15 del PDF

Si distinguono due macro-categorie di prefissi Type~5:
\begin{itemize}[noitemsep]
    \item \textbf{Internal Prefixes}: sottoreti appartenenti all'architettura interna della fabric VXLAN/EVPN, generate localmente mediante sintesi delle interfacce d'accesso.
    \item \textbf{External Prefixes}: prefissi provenienti dall'esterno del dominio overlay (es. connessione a un border router verso la WAN aziendale o il provider cloud), re-iniettati nel piano di controllo EVPN.
\end{itemize}

\subsection{Apprendimento Delle Rotte Di Segmento E Isolamento Multi-Tenant (L3 VNI)}

In un ambiente con gateway distribuiti, ciascun tenant definisce un contesto di instradamento privato (\textbf{VRF}). Poiché le intestazioni dei pacchetti di underlay non contengono metadati relativi all'identità della VRF, il disaccoppiamento multi-tenant lungo l'infrastruttura condivisa viene realizzato assegnando a ciascuna VRF un identificatore univoco globale: il \textbf{Layer 3 VNI} (\textbf{L3 VNI}).

% [IMMAGINE DEL FILE DI SLIDE CON NOME "cloudComputing_5"] Flusso completo di annuncio del prefisso 10.20.20.1/24 con Type 5, binding di VRF tramite ERT/IRT e L3 VNI 1000 — Slide 120 / Pagina 16 del PDF

Il meccanismo di propagazione e binding di una rotta di segmento opera secondo i seguenti passaggi:
\begin{enumerate}
    \item Il router di bordo o leaf d'accesso (\texttt{PE1} o \texttt{VTEP1} con IP \texttt{1.1.1.9}) importa o configura localmente la sottorete \texttt{10.20.20.0/24} all'interno della propria istanza \texttt{VRF1}.
    \item \texttt{VTEP1} emette un aggiornamento BGP contenente la rotta Type~5:
    \[
    \text{RT5} = \left\{
    \begin{aligned}
    &\text{RD} = 1:1, \quad \text{Prefix} = 10.20.20.1/24, \quad \text{Label} = \text{L3 VNI: } 1000, \\
    &\text{Next-Hop} = 1.1.1.9, \quad \text{Router MAC} = \texttt{VTEP1-MAC}, \quad \text{ERT} = 1:1
    \end{aligned}
    \right\}
    \]
    \item L'attributo BGP \textbf{Router MAC Extended Community} trasporta il MAC del router sorgente, parametro critico che il VTEP ricevente dovrà inserire come MAC di destinazione interno nella trama decapsulata.
    \item Il VTEP remoto (\texttt{VTEP2}) riceve l'annuncio. L'importazione avviene secondo la logica delle \textbf{Route Target}:
    \begin{itemize}[noitemsep]
        \item \textbf{Export Route Target (ERT)}: etichetta associata alla rotta al momento della generazione sul nodo sorgente.
        \item \textbf{Import Route Target (IRT)}: filtro configurato localmente su ciascuna VRF per selezionare quali rotte recepire.
    \end{itemize}
    Se sul nodo ricevente esiste una VRF configurata con $\text{IRT} = 1:1$, la rotta viene validata.
    \item \texttt{VTEP2} compila la propria tabella di routing IP (FIB) all'interno di tale VRF:
    \begin{center}
    \begin{tabular}{|c|c|c|c|c|}
    \hline
    \textbf{VRF ID} & \textbf{IP-Prefix} & \textbf{L3 VNI} & \textbf{Router MAC} & \textbf{Next Hop} \\ \hline
    $1$ & \texttt{10.20.20.1/24} & \num{1000} & \texttt{VTEP1 MAC} & \texttt{1.1.1.9} \\ \hline
    \end{tabular}
    \end{center}
    Contestualmente, viene programmato il tunnel VXLAN verso il DIP \texttt{1.1.1.9}. Quando qualsiasi host appartenente alla VRF locale deve comunicare con la sottorete \texttt{10.20.20.0/24}, il traffico viene direttamente instradato all'interno del \texttt{L3 VNI 1000}.
\end{enumerate}

\newpage

\section{Architetture Di Gateway VXLAN: Centralizzato Vs Distribuito}

\subsection{Tassonomia Dei Gateway: Layer 2 Vs Layer 3}

Un \textbf{VXLAN Gateway} è un'entità logica o fisica incaricata di garantire la continuità operativa tra il dominio incapsulato VXLAN e le reti non-VXLAN o tra sottoreti logiche distinte:

% [IMMAGINE DEL FILE DI SLIDE CON NOME "cloudComputing_5"] Schema generale di un gateway VXLAN per la traduzione e l'interconnessione tra domini VXLAN e reti esterne IP/VLAN — Slide 121 / Pagina 18 del PDF

\begin{itemize}
    \item \textbf{Layer 2 VXLAN Gateway (VXLAN Bridge)}:
    \begin{itemize}[noitemsep]
        \item Connette segmenti overlay VXLAN a domini di broadcast tradizionali basati su tag VLAN (IEEE 802.1Q).
        \item Esegue una pura operazione di terminazione: riceve il pacchetto incapsulato in UDP, elimina l'intestazione esterna e inserisce la trama Ethernet nel circuito VLAN corrispondente (e viceversa).
        \item Non esegue alcuna operazione di routing né altera l'indirizzamento IP.
    \end{itemize}
    \item \textbf{Layer 3 VXLAN Gateway (VXLAN Router)}:
    \begin{itemize}[noitemsep]
        \item Realizza la terminazione di Livello~3, instradando pacchetti tra differenti domini VNI (\emph{inter-VXLAN routing}) o tra un VNI e la rete IP esterna (WAN/Internet).
        \item Decapsula il pacchetto VXLAN, decrementa il TTL del pacchetto IP, riscrive le intestazioni Ethernet e reincapsula il pacchetto verso un nuovo VNI o lo immette nativamente su una subnet fisica.
    \end{itemize}
\end{itemize}

% [IMMAGINE DEL FILE DI SLIDE CON NOME "cloudComputing_5"] Classificazione funzionale tra Layer 2 Gateway (bridging) e Layer 3 Gateway (routing) con casi d'uso applicativi — Slide 122 / Pagina 19 del PDF

Nelle implementazioni pratiche, le interfacce virtuali di Livello~3 associate a un dominio di bridge prendono il nome di \textbf{VBDIF} (\emph{Virtual Bridge Domain Interface}, corrispondenti alle interfacce \textbf{SVI}, \emph{Switch Virtual Interface}, della letteratura di rete). Un'interfaccia VBDIF possiede un proprio indirizzo IP e MAC e funge da default gateway per i terminali afferenti a quel bridge domain.

\subsection{Architettura A Gateway Centralizzato}

Nel paradigma a \textbf{Gateway Centralizzato}, le funzioni di instradamento di Livello~3 sono confinate su uno o due nodi centrali della fabric, tipicamente gli switch di livello \textbf{Spine} o apparati dedicati definiti \textbf{Border Leaf}:

% [IMMAGINE DEL FILE DI SLIDE CON NOME "cloudComputing_5"] Topologia del Gateway Centralizzato su nodo Spine con evidenza del percorso di traffico tromboning (hairpinning) — Slide 124 / Pagina 23 del PDF

\begin{itemize}
    \item \textbf{Vantaggi}:
    \begin{itemize}[noitemsep]
        \item Semplicità di gestione iniziale: tutte le interfacce VBDIF/SVI e le relative politiche di routing risiedono su un unico apparato.
        \item Ideale per il \emph{service insertion}: facilita l'instradamento forzato di tutto il traffico inter-subnet attraverso apparati di sicurezza centralizzati (firewall perimetrali, IPS).
    \end{itemize}
    \item \textbf{Criticità Architetturali}:
    \begin{enumerate}
        \item \textbf{Traffic Hairpinning (o Tromboning)}: quando due macchine virtuali collegate allo \emph{stesso leaf switch} d'accesso appartengono a due sottoreti diverse (es. \texttt{VM1} su \texttt{VBDIF10} e \texttt{VM3} su \texttt{VBDIF20}), il traffico generato da \texttt{VM1} non può essere commutato localmente. Il leaf deve incapsulare il pacchetto e inviarlo allo Spine centrale; lo Spine esegue il routing L3 tra le due VBDIF e lo reincapsula rispedendolo verso il medesimo leaf d'accesso. Questo percorso a "forcina" satura inutilmente i collegamenti di uplink e incrementa significativamente la latenza.
        \item \textbf{Collo di bottiglia per le tabelle ARP}: lo Spine centrale deve apprendere e mantenere in memoria lo stato ARP di tutti i terminali del datacenter. La capacità limitata delle tabelle di inoltro hardware dei nodi centrali diviene rapidamente il fattore bloccante per la scalabilità complessiva dell'infrastruttura.
    \end{enumerate}
\end{itemize}

\subsection{Architettura A Gateway Distribuito (Anycast Gateway)}

Per ovviare a questi limiti strutturali, le reti moderne adottano l'architettura a \textbf{Gateway Distribuito}, in cui la funzione di routing di Livello~3 viene distribuita capillarmente su ogni singolo switch di accesso (\textbf{Leaf}).

% [IMMAGINE DEL FILE DI SLIDE CON NOME "cloudComputing_5"] Architettura del Distributed VXLAN Gateway su topologia Spine-Leaf con instradamento locale al primo hop — Slide 128 / Pagina 27 del PDF

Il meccanismo si fonda sul paradigma dell'\textbf{Anycast Gateway}:
\begin{itemize}
    \item Ciascun leaf switch della fabric su cui è attivo un dato dominio VXLAN viene configurato con la **medesima interfaccia logica di gateway**, avente esattamente lo **stesso indirizzo IP** e lo **stesso indirizzo MAC virtuale** (es. \texttt{10.1.1.1} e \texttt{AA:BB:CC:DD:EE:FF}).
    \item Qualsiasi server o macchina virtuale, indipendentemente dal rack in cui si trova, individua il proprio gateway locale sul leaf a cui è direttamente collegato. In caso di migrazione in un altro punto della rete, la VM non deve aggiornare né il proprio default gateway né la propria tabella ARP.
    \item L'instradamento inter-subnet avviene al **primo salto di rete** (\emph{ingress routing}): il leaf di ingresso esegue la decisione di routing non appena riceve la trama dall'host. Se la destinazione risiede su un leaf remoto, il pacchetto viene incapsulato e spedito direttamente al leaf target. Non si verifica alcun fenomeno di hairpinning.
    \item Ogni leaf deve apprendere le voci ARP unicamente per i terminali fisicamente connessi alle proprie porte di accesso, abbattendo la dimensione delle tabelle di inoltro.
\end{itemize}

\begin{theorem}[Principio Della Doppia Identità Del Leaf Switch]
Per consentire il corretto funzionamento dell'Anycast Gateway senza creare conflitti di instradamento nell'infrastruttura fisica di trasporto, ogni switch leaf implementa una rigorosa separazione funzionale su \textbf{due identità logiche distinte}:
\begin{enumerate}
    \item \textbf{Identità Pubblica (Anycast Overlay)}: costituita dalla coppia IP/MAC virtuale condivisa da tutti i nodi della fabric. Questa identità dialoga \emph{esclusivamente con i server terminali} collegati in locale, rispondendo alle loro richieste ARP e fungendo da default gateway. Tale identità \textbf{non viene mai utilizzata per la comunicazione inter-switch tra i Leaf}.
    \item \textbf{Identità Privata (Unicast Underlay)}: costituita dall'indirizzo IP univoco di loopback assegnato al VTEP fisico. Questa identità opera nell'underlay e viene utilizzata come indirizzo sorgente/destinazione per stabilire le sessioni BGP EVPN e per incapsulare le intestazioni esterne dei pacchetti VXLAN trasmessi attraverso la fabric.
\end{enumerate}
\end{theorem}

\newpage

\section{Modelli Di Integrated Routing And Bridging: Asymmetric Vs Symmetric}

La coesistenza del bridging L2 e del routing L3 sul medesimo apparato distribuito definisce il paradigma \textbf{Integrated Routing and Bridging} (IRB). In ambiente EVPN-VXLAN, la gestione dell'inoltro tra sottoreti distinte si articola in due modelli fondamentali: \textbf{Asymmetric IRB} e \textbf{Symmetric IRB}.

\subsection{Asymmetric IRB: Meccanismo E Limiti Di Scalabilità}

Nel modello \textbf{Asymmetric IRB}, l'intera complessità decisionale di instradamento è concentrata esclusivamente sul VTEP di ingresso (\emph{Ingress Leaf}), mentre il VTEP di uscita (\emph{Egress Leaf}) opera come un semplice ponte di Livello~2:

% [IMMAGINE DEL FILE DI SLIDE CON NOME "cloudComputing_5"] Diagramma operativo e flusso del pacchetto nel modello Asymmetric IRB — Slide 133 / Pagina 32 del PDF

\begin{enumerate}
    \item \textbf{Fase Ingress}: un host collegato al dominio \texttt{L2VNI 1000} su \texttt{VTEP1} invia un pacchetto a un terminale situato nel dominio \texttt{L2VNI 2000}.
    \item \texttt{VTEP1} riceve la trama, effettua la ricerca di Livello~3 nella tabella di instradamento della VRF e determina che la destinazione appartiene al dominio \texttt{L2VNI 2000}.
    \item \texttt{VTEP1} esegue la risoluzione dell'indirizzo MAC di destinazione, individua l'indirizzo IP del VTEP remoto (\texttt{VTEP2}) su cui tale host è attestato e **incapsula il pacchetto direttamente nell'identificatore del dominio di arrivo (\texttt{L2VNI 2000})**.
    \item \textbf{Attraversamento della Fabric}: il pacchetto attraversa l'underlay con intestazione VXLAN valorizzata con $\text{VNI} = 2000$.
    \item \textbf{Fase Egress}: \texttt{VTEP2} riceve il pacchetto, decapsula l'intestazione VXLAN e rileva il tag \texttt{L2VNI 2000}. \texttt{VTEP2} **non compie alcuna operazione di routing**: si limita a eseguire una classica commutazione L2 all'interno della propria MAC table e consegna la trama Ethernet al destinatario.
\end{enumerate}

\paragraph{Perché il modello Asimmetrico è inefficiente?}
Il limite strutturale di questo schema risiede nella parola "asimmetria": per poter instradare il pacchetto verso \texttt{L2VNI 2000}, l'ingress switch \texttt{VTEP1} deve **obbligatoriamente possedere la configurazione locale e lo stato di inoltro di \texttt{L2VNI 2000}**, anche se nessun server del proprio rack appartiene a quel segmento. Ciò implica che ogni leaf switch della fabric deve configurare preventivamente \textbf{tutti i VNI di Livello~2 esistenti nel datacenter}, vanificando la scalabilità e saturando rapidamente lo spazio di memoria per le tabelle di inoltro.

\subsection{Symmetric IRB: Architettura, Instradamento A Due Stadi E L3 VNI}

Nelle reti EVPN moderne, lo standard di riferimento assoluto è rappresentato dal modello \textbf{Symmetric IRB}. In questa architettura, la decisione di instradamento è ripartita in modo perfettamente simmetrico tra il leaf di ingresso e quello di uscita, disaccoppiando completamente i domini di Livello~2 grazie all'introduzione di una dorsale logica di transito basata su \textbf{L3 VNI}:

% [IMMAGINE DEL FILE DI SLIDE CON NOME "cloudComputing_5"] Architettura di instradamento a due stadi in Symmetric IRB con evidenza dei ruoli di L2-VNI locale e L3-VNI condiviso — Slide 131-132 / Pagine 30-31 del PDF

\begin{itemize}
    \item I \textbf{Layer 2 VNI} (es. Orange VNI per VLAN~A, Green VNI per VLAN~B) mantengono una portata puramente locale: estendono il dominio di broadcast unicamente tra i leaf switches in cui sono fisicamente attestati server appartenenti a tali VLAN.
    \item Il \textbf{Layer 3 VNI} (es. Purple VNI) è associato univocamente alla VRF del tenant ed è configurato su tutti i leaf switches che ospitano interfacce di quella VRF.
    \item **Qualsiasi traffico che deve essere instradato tra sottoreti diverse viaggia nella fabric incapsulato all'interno del Layer 3 VNI**.
\end{itemize}

\subsection{Analisi Dettagliata Del Flusso Di Inoltro In Symmetric IRB}

Si analizzi analiticamente il percorso compiuto da un pacchetto originato da un host attestato su \texttt{Leaf01} (\texttt{L2VNI 1000}, sottorete \texttt{10.1.1.0/24}) e destinato a un nodo attestato su \texttt{Leaf02} (\texttt{L2VNI 2000}, sottorete \texttt{10.1.2.0/24}), appartenenti alla medesima VRF di tenant:

\begin{enumerate}
    \item \textbf{Ingresso e Routing al Primo Salto (Ingress Leaf)}:
    \begin{itemize}[noitemsep]
        \item L'host sorgente trasmette una trama Ethernet avente come MAC di destinazione l'indirizzo dell'Anycast Gateway (VBDIF/SVI locale).
        \item \texttt{Leaf01} riceve la trama, riconosce il proprio Anycast MAC ed esegue il routing di Livello~3 all'interno della VRF.
        \item La consultazione della FIB identifica che il prefisso di destinazione è raggiungibile tramite l'indirizzo underlay di \texttt{Leaf02} attraverso il \textbf{L3 VNI 5000}.
        \item La trama Ethernet interna viene riscritta: il MAC sorgente diventa il \textbf{Router MAC} di \texttt{Leaf01}, mentre il MAC di destinazione viene impostato sul \textbf{Router MAC} di \texttt{Leaf02} (appreso preliminarmente via BGP EVPN).
    \end{itemize}
    \item \textbf{Incapsulamento e Attraversamento Della Fabric}:
    \begin{itemize}[noitemsep]
        \item \texttt{Leaf01} incapsula il pacchetto in un'intestazione VXLAN valorizzando il campo VNI con il valore del dominio L3:
        \[
        \text{VNI}_{\text{outer}} = \mathbf{5000} \quad (\text{L3 VNI})
        \]
        \item L'intestazione IP esterna presenta sorgente \texttt{IP(Leaf01)} e destinazione \texttt{IP(Leaf02)}.
        \item I router intermedi (Spine) inoltrano il pacchetto basandosi puramente sull'indirizzo IP di destinazione esterno, ignorando totalmente il contenuto VXLAN.
    \end{itemize}
    \item \textbf{Ricezione e Routing Al Secondo Salto (Egress Leaf)}:
    \begin{itemize}[noitemsep]
        \item \texttt{Leaf02} riceve il pacchetto e decapsula l'intestazione VXLAN.
        \item Rilevando che il VNI corrisponde a \texttt{L3 VNI 5000}, associa immediatamente il flusso alla corrispondente istanza VRF locale.
        \item Ispeziona il MAC di destinazione interno (che corrisponde al proprio \texttt{Router MAC}) ed esegue un **secondo passaggio di routing di Livello~3** all'interno della VRF.
        \item Dalla consultazione della tabella di instradamento locale, determina che l'IP appartiene alla sottorete direttamente connessa su \texttt{L2VNI 2000}.
        \item Esegue una ricerca nella propria tabella ARP/FIB locale per determinare l'interfaccia fisica e il MAC del server di destinazione, riscrive l'intestazione Ethernet di Livello~2 e inoltra la trama non incapsulata al destinatario finale.
    \end{itemize}
\end{enumerate}

\subsection{Quadro Comparativo: Asymmetric IRB Vs Symmetric IRB}

La tabella seguente riassume in modo esaustivo le divergenze ingegneristiche tra i due approcci di instradamento integrato:

\begin{table}[H]
\centering
\small
\renewcommand{\arraystretch}{1.4}
\begin{tabularx}{\textwidth}{|p{3.5cm}|X|X|}
\hline
\textbf{Caratteristica} & \textbf{Asymmetric IRB} & \textbf{Symmetric IRB} \\ \hline
\textbf{Stadio di Routing} & Instradamento unicamente sul leaf di ingresso; il leaf di uscita esegue solo bridging L2. & Instradamento distribuito su entrambi i lati (routing su ingress e routing su egress). \\ \hline
\textbf{VNI di Transito nella Fabric} & Il pacchetto attraversa l'underlay incapsulato nel \textbf{L2 VNI} del bridge domain di destinazione. & Il pacchetto attraversa l'underlay incapsulato nel \textbf{L3 VNI} associato alla VRF di tenant. \\ \hline
\textbf{Impiego di L3 VNI} & Non presente nel data path. & Fondamentale per disaccoppiare l'inoltro di L3 dai domini L2. \\ \hline
\textbf{Requisiti di Configurazione sui Leaf} & Ciascun leaf deve configurare \textbf{tutti} i L2 VNI per poter instradare verso di essi (esplosione delle VLAN/VNI). & Ciascun leaf configura unicamente i L2 VNI \textbf{locali} al proprio rack, oltre al L3 VNI condiviso della VRF. \\ \hline
\textbf{Consumo di Risorse (Tabelle HW)} & Elevato: MAC table e ARP table di tutti i segmenti replicate ovunque. & Ottimale e scalabile: tabelle L2 confinate localmente, tabelle L3 filtrate per VRF. \\ \hline
\textbf{Classificazione Standard} & Modello obsoleto, scarsamente scalabile. & **Standard de facto moderno** per datacenter e cloud infrastructure. \\ \hline
\end{tabularx}
\caption{Confronto sistematico tra i modelli di instradamento Asymmetric IRB e Symmetric IRB in EVPN-VXLAN.}
\end{table}

\newpage

\section{Inoltro Inter-Subnet in Overlay VXLAN: Architettura Asymmetric EVPN IRB}

\subsection{Richiamo Concettuale: Bridging, Routing e la Necessità dell'IRB}
Nelle reti tradizionali basate su stack TCP/IP e modello ISO/OSI, la separazione tra il livello di collegamento (Data Link, L2) e il livello di rete (Network, L3) impone una netta differenziazione operativa:
\begin{itemize}[noitemsep]
    \item Il \textbf{bridging} (o \emph{switching}) opera intra-subnet inoltrando trame Ethernet in base all'indirizzo MAC di destinazione, mantenendo confinato il dominio di broadcast all'interno di una singola VLAN (\emph{Virtual Local Area Network}).
    \item Il \textbf{routing} opera inter-subnet instradando pacchetti IP tra prefissi di rete distinti mediante l'ispezione dell'indirizzo IP di destinazione e la consultazione della tabella di instradamento (\textbf{FIB}, \emph{Forwarding Information Base}). Tale operazione comporta l'interruzione del dominio L2, il decremento del campo TTL (\emph{Time To Live}) e la completa riscrittura dell'intestazione di livello 2 (sostituzione del MAC mittente con il MAC dell'interfaccia di uscita del router e del MAC destinatario con quello del \emph{next-hop} o dell'host finale).
\end{itemize}

Nelle moderne infrastrutture di data center che adottano un'architettura a \emph{overlay} basata su incapsulamento \textbf{VXLAN} (\emph{Virtual Extensible LAN}) controllato da \textbf{MP-BGP EVPN} (\emph{Ethernet VPN}), il traffico viene trasportato su una rete \emph{underlay} puramente L3 (tipicamente con topologia Clos \emph{Spine-Leaf}). Per consentire la comunicazione sia tra nodi appartenenti alla medesima sottorete (intra-subnet, associata a un identificatore di segmento virtuale noto come \textbf{L2VNI}, \emph{Layer 2 Virtual Network Identifier}), sia tra nodi posizionati su sottoreti logiche distinte (inter-subnet), gli apparati di bordo o \textbf{VTEP} (\emph{VXLAN Tunnel Endpoint}) posizionati sui commutatori di tipo \textbf{Leaf} devono integrare entrambe le funzioni.

Questo paradigma prende il nome di \textbf{IRB} (\emph{Integrated Routing and Bridging}). Esistono due principali modelli architetturali per implementare l'IRB in una fabric EVPN: il modello asimmetrico (\textbf{Asymmetric IRB}) e il modello simmetrico (\textbf{Symmetric IRB}).

---

\subsection{Principio di Funzionamento dell'Asymmetric EVPN IRB}

Nel modello \textbf{Asymmetric EVPN IRB}, l'infrastruttura di rete \emph{overlay} non definisce alcun identificatore VXLAN dedicato all'instradamento di terzo livello; in altri termini, in questo schema progettuale:
\[
\text{\textbf{Non esiste alcun L3VNI (Layer 3 VNI)}}.
\]
Tutti i tunnel VXLAN attivi nel piano dati trasportano esclusivamente traffico associato a identificatori di livello 2 (\textbf{L2VNI}), corrispondenti ciascuno a una specifica VLAN della fabric.

Il principio fondamentale su cui poggia l'intera logica di inoltro dell'Asymmetric IRB è riassunto dalla regola operativa:
\begin{center}
    \textbf{Routing on Ingress, Bridging on Egress}
\end{center}

Ciò significa che la decisione di instradamento inter-subnet viene interamente presa e completata dal commutatore Leaf di ingresso (\emph{ingress VTEP}). Il commutatore Leaf di destinazione (\emph{egress VTEP}), al contrario, si limita a decapsulare il pacchetto VXLAN e a eseguire una pura commutazione di livello 2 verso la porta di accesso su cui è attestato l'host destinatario.

% [IMMAGINE DEL FILE DI SLIDE CON NOME "CloudComputing_5"] Topologia Leaf-Spine con architettura Asymmetric EVPN IRB: instradamento inter-subnet eseguito sull'ingress Leaf con incapsulamento nel VNI di destinazione e bridging sull'egress Leaf — Slide 1 / Pagina 1 del PDF

\subsubsection{Dinamica Passo-Passo del Flusso Dati Unicast}
Si consideri lo scenario operativo illustrato in figura, caratterizzato da due commutatori di bordo, $\text{Leaf01}$ e $\text{Leaf02}$, collegati attraverso una rete di trasporto L3 (\emph{Layer 3 Network}):
\begin{itemize}[noitemsep]
    \item $\text{Host A}$ è attestato su $\text{Leaf01}$ ed appartiene alla $\text{VLAN A}$, mappata nell'overlay sull'\textbf{Orange VNI} (L2VNI per la sottorete A).
    \item $\text{Host B}$ è attestato su $\text{Leaf02}$ ed appartiene alla $\text{VLAN B}$, mappata nell'overlay sul \textbf{Green VNI} (L2VNI per la sottorete B).
    \item $\text{Host C}$ è attestato anch'esso su $\text{Leaf02}$ ma appartiene alla $\text{VLAN A}$ (Orange VNI).
\end{itemize}

Quando l'$\text{Host A}$ deve trasmettere un pacchetto IP verso l'$\text{Host B}$ (comunicazione inter-subnet tra $\text{VLAN A}$ e $\text{VLAN B}$), la sequenza di elaborazione si articola come segue:

\begin{enumerate}
    \item \textbf{Invio del frame locale}: Poiché $\text{Host B}$ si trova su una sottorete IP differente, $\text{Host A}$ risolve l'indirizzo MAC del proprio \emph{Default Gateway} (configurato localmente su $\text{Leaf01}$ come interfaccia IRB/SVI) e invia a esso il frame Ethernet contenente il pacchetto IP.

    \item \textbf{Routing on Ingress (su Leaf01)}:
    \begin{itemize}[noitemsep]
        \item $\text{Leaf01}$ riceve il frame e riscontra che il MAC di destinazione corrisponde alla propria interfaccia di gateway per la $\text{VLAN A}$.
        \item L'intestazione Ethernet originaria viene rimossa e viene effettuata una consultazione della tabella di instradamento IP (\emph{FIB lookup}) per l'indirizzo IP di destinazione di $\text{Host B}$.
        \item La tabella di routing locale indica che l'indirizzo appartiene al prefisso di $\text{VLAN B}$. Poiché $\text{Leaf01}$ possiede un'interfaccia di gateway instanziata anche per $\text{VLAN B}$, esso esegue localmente l'operazione di instradamento di terzo livello, passando il pacchetto dal contesto di $\text{VLAN A}$ a quello di $\text{VLAN B}$.
        \item Mediante le informazioni preventivamente apprese via MP-BGP EVPN (tramite route di \emph{Tipo 2 - MAC/IP Advertisement}), $\text{Leaf01}$ conosce sia l'indirizzo MAC di $\text{Host B}$ sia l'indirizzo IP del VTEP remoto su cui tale host risiede ($\text{Leaf02}$).
    \end{itemize}

    \item \textbf{Incapsulamento VXLAN nel VNI di destinazione}:
    \begin{itemize}[noitemsep]
        \item $\text{Leaf01}$ riscrive l'intestazione Ethernet di livello 2: il MAC sorgente diventa il MAC dell'interfaccia di gateway di $\text{Leaf01}$ su $\text{VLAN B}$, mentre il MAC di destinazione diventa l'indirizzo fisico di $\text{Host B}$.
        \item Il pacchetto così formattato viene incapsulato in un pacchetto VXLAN specificando come identificatore di rete il \textbf{Green VNI} (il VNI di livello 2 della \emph{destinazione}, associato a $\text{VLAN B}$).
        \item L'intestazione esterna IP/UDP del tunnel VXLAN specifica come IP sorgente il VTEP di $\text{Leaf01}$ e come IP di destinazione il VTEP di $\text{Leaf02}$.
    \end{itemize}

    \item \textbf{Trasporto nell'underlay}: La rete di trasporto instaura un normale inoltro L3 via shortest-path (ECMP) recapitandone il pacchetto all'indirizzo IP di $\text{Leaf02}$.

    \item \textbf{Bridging on Egress (su Leaf02)}:
    \begin{itemize}[noitemsep]
        \item $\text{Leaf02}$ riceve il pacchetto, rileva che è indirizzato al proprio IP VTEP ed elimina l'incapsulamento VXLAN.
        \item Ispezionando il VNI estratto (\textbf{Green VNI}), associa immediatamente il frame decapsulato al proprio dominio di broadcast locale per la $\text{VLAN B}$.
        \item $\text{Leaf02}$ \emph{non esegue alcuna decisione di routing}: si limita a effettuare una pura ricerca L2 nella propria tabella dei MAC address associata a $\text{VLAN B}$.
        \item Il frame viene infine inoltrato in chiaro sulla porta fisica locale a cui è connesso $\text{Host B}$.
    \end{itemize}
\end{enumerate}

Nel tragitto inverso (risposta da $\text{Host B}$ a $\text{Host A}$), il meccanismo si ripete in modo speculare:
\begin{itemize}[noitemsep]
    \item $\text{Leaf02}$ esegue il \textbf{Routing on Ingress}, instradando il pacchetto da $\text{VLAN B}$ (Green) a $\text{VLAN A}$ (Orange) e incapsulandolo nell'\textbf{Orange VNI}.
    \item $\text{Leaf01}$ riceve il pacchetto con Orange VNI ed esegue unicamente il \textbf{Bridging on Egress} verso $\text{Host A}$.
\end{itemize}

L'aggettivo \emph{asimmetrico} deriva precisamente da questa divergenza operativa tra i due capi del tunnel: l'ingress switch compie sia il routing L3 sia la predisposizione del framing L2 per la destinazione, mentre l'egress switch opera passivamente da semplice ponte di livello 2.

---

\subsection{Analisi Critica: Limiti Architetturali e Problemi di Scalabilità}

Nonostante l'Asymmetric IRB presenti il vantaggio di una logica concettuale lineare (eliminando la necessità di configurare L3VNI e tabelle VRF dedicate al transito overlay), esso presenta gravi vincoli di scalabilità che ne limitano fortemente l'adozione nelle grandi fabric di data center:

\begin{enumerate}
    \item \textbf{Obbligo di istanziazione universale di tutte le VLAN/L2VNI}:
    Poiché qualsiasi Leaf della rete può dover fungere da nodo di ingresso e instradare direttamente verso qualsiasi sottorete di destinazione presente nella fabric, ogni singolo nodo Leaf deve obbligatoriamente possedere la configurazione attiva di \emph{tutti i domini di broadcast} (tutte le VLAN e tutti i rispettivi L2VNI), a prescindere dal fatto che abbia o meno host locali attestati su tali VLAN.

    \item \textbf{Esplosione delle tabelle MAC e delle adiacenze ARP/ND}:
    Affinché un ingress Leaf possa incapsulare direttamente il pacchetto nel VNI di destinazione, deve conoscere preventivamente l'indirizzo MAC dell'host ricevente. Di conseguenza, il piano di controllo BGP EVPN deve distribuire le tabelle di corrispondenza MAC-IP di \emph{tutti gli host dell'intera fabric} a \emph{tutti i nodi Leaf}. Ciascun apparato di bordo deve memorizzare l'intero stato L2/L3 della rete nelle proprie memorie hardware ad alta velocità (\textbf{TCAM}, \emph{Ternary Content-Addressable Memory}).

    \item \textbf{Saturazione delle risorse hardware nei commutatori di bordo}:
    Se in una fabric sono definite ad esempio $500$ VLAN con migliaia di macchine virtuali o container distribuiti:
    \[
    \text{Ogni singolo Leaf deve supportare ed emulare tutte le } 500 \text{ VLAN e conservare tutti i relativi record ARP/MAC.}
    \]
    I Leaf switch di fascia d'accesso esauriscono rapidamente lo spazio nelle tabelle di inoltro (MAC table e neighbor table), impedendo la crescita modulare (\emph{scale-out}) dell'infrastruttura.
\end{enumerate}

Questi colli di bottiglia hanno portato l'industria del networking a privilegiare, nelle implementazioni EVPN di moderna concezione ad alta densità, il modello \textbf{Symmetric IRB}, in cui l'instradamento inter-subnet viene distribuito simmetricamente su entrambi i nodi (routing su ingress verso un \textbf{L3VNI} condiviso e routing su egress dall'L3VNI alla subnet locale), liberando i singoli Leaf dalla necessità di conoscere e istanziare le VLAN non localmente attestate.

\end{document}
